% English translation, made from our own transcription (original.tex), not from any existing
% translation. Every math expression, symbol, \tag, \origpage{N} marker and \ednote is reproduced
% unchanged; only the prose is translated (HOUSESTYLE R21 for prose set via \text{...} inside
% formulas: et/seu/sive/unde and the like are rendered and the insert stays in place). The work
% has no figures. The existing \ednote notes are OUR notes, already in English, and are carried
% over verbatim in the same positions, even where the misprint they describe is Latin orthography
% that cannot survive translation.
%
% Ordinal suffixes written as Latin superscripts inside \mathrm (n^{\mathrm{ti}}, p^{\mathrm{tum}},
% Cl$^{\mathrm{i}}$, \S$^{\mathrm{o}}$) are math tokens, not \text inserts, so they are left as
% printed; the surrounding English prose carries the ordinal ("of the n^{ti} order").
%
% Translated in 14 chunks by subagents of the same model, then assembled and harmonised. Period
% terminology, applied consistently (see glossary.yaml): Modulus, Amplitude, Argument, Multiplier,
% Complement; "Elliptic Functions" and similar terms are capitalised only where the print
% capitalises them; transformatio prima/secunda = first/second transformation; species of
% elliptic integrals kept as "species" (not "kind"); evolutio = expansion; functio integra
% rationalis = rational integral function; Legendre's "functiones integrae" = complete functions.
% The honorific "Cl." (Clarissimus) is kept in the author's own abbreviated form, with one
% translator's note at its first occurrence (Preface); "Regiom." is rendered Konigsberg.
\documentclass{article}
\usepackage{readmasters}
\begin{document}


\origpage{i}

\section*{NEW FOUNDATIONS OF THE THEORY OF ELLIPTIC FUNCTIONS}

BY

D. CARL GUSTAV JACOB JACOBI,

ORDINARY PROFESSOR IN THE UNIVERSITY OF KÖNIGSBERG

KÖNIGSBERG

AT THE EXPENSE OF THE BORNTRÆGER BROTHERS

1829.

PARIS at Ponthieu \& Co. $\quad$ Treuttel \& Wuerz.

LONDON at Treuttel, Wuerz \& Richter. $\quad$ H. W. Koller. $\quad$ Black, Young \& Young.

AMSTERDAM at Mueller \& Co. $\quad$ C. G. Suelpke.

ST. PETERSBURG at Graeff.

\origpage{iii}

\section*{PREFACE.}

Almost two years ago, when I resolved to examine the theory of elliptic functions more closely, I came upon certain questions of the greatest weight, which seemed both to give that theory a new aspect and to advance the whole analytic art remarkably. When these had been brought to a happy conclusion, hardly expected on account of the difficulty of the matter, I communicated their first moments to the Geometers briefly and without demonstration, and soon afterwards, when it was desired more urgently and credence seemed hardly to be given to the new discovery, with the demonstration added. At the same time I was urged to publish the complete system of the questions I had undertaken. To satisfy this wish at least in part, I have resolved to publish the foundations on which my questions are built. These new foundations of the theory of elliptic functions we now commend to the indulgence of the Geometers.

\origpage{iv}

So that the book might come out as free from the printers' faults as could be managed, Cl.\ Scherk\ednote{translator's note: ``Cl.'' is the honorific Clarissimus, ``most distinguished'', which Jacobi prefixes to the name of a living colleague throughout; it is kept here in his own abbreviated form.} was willing to take charge, to whom I profess myself greatly obliged in this matter. What remains to be corrected has been added at the end.

I wrote this in the month of Feb. of the year 1829

at the Univ. of Königsberg.

\origpage{v}

\section*{INDEX OF CONTENTS.}

\subsection*{On the Transformation of Elliptic Functions. §§.\ 1---34. p.\ 1---83}

Statement of the general problem of transformation. §§.\ 1.\ 2. p.\ 1

Principles of transformation. §§.\ 3.\ 4. --- 8

The expression $\dfrac{dy}{\sqrt{\pm(y-\alpha)(y-\beta)(y-\gamma)(y-\delta)}}$ is proposed, to be reduced to the simpler form $\dfrac{dx}{M\sqrt{(1-x^{2})(1-k^{2}x^{2})}}$. §§.\ 5---9. --- 6

On the transformation of the expression $\dfrac{dy}{\sqrt{(1-y^{2})(1-\lambda^{2}y^{2})}}$ into another similar to it, $\dfrac{dx}{M\sqrt{(1-x^{2})(1-k^{2}x^{2})}}$. §§.\ 10---12. --- 17

The transformation of the third order is proposed. §§.\ 13.\ 14. --- 23

The transformation of the fifth order is proposed. §.\ 15. --- 26

How multiplication is arrived at by applying the transformation twice. §.\ 16. --- 29

On the new notation of elliptic functions. §.\ 17. --- 30

Fundamental formulae in the analysis of elliptic functions. §.\ 18. --- 32

On the imaginary values of elliptic functions. The principle of the double period. §.\ 19. --- 34

Analytic theory of the transformation of elliptic functions. §.\ 20. --- 36

Demonstration of the analytic formulae for the transformation. §§.\ 21---23. --- 39

On the various transformations of the same order. Two real transformations, of a greater modulus into a lesser and of a lesser into a greater. §.\ 24. --- 49

On complementary transformations, or how from a transformation of a modulus into a modulus another of the complement into the complement is derived. §.\ 25. --- 57

On supplementary transformations to multiplication. §§.\ 26.\ 27. --- 60

General analytic formulae for the multiplication of elliptic functions. §.\ 28. --- 64

On the affections of modular equations. §§.\ 29---34. --- 66

\origpage{vi}

\subsection*{Theory of the Expansion of Elliptic Functions. §§.\ 35---66. p.\ 84---188}

On the expansion of elliptic functions into infinite products. §.\ 35---38. p.\ 84

Expansion of elliptic functions into series proceeding according to the sines or cosines of multiples of the argument. §§.\ 39---42. --- 99

General formulae for expanding the functions $\sin^{n}\operatorname{am}\left(\dfrac{2Kx}{\pi}\right)$, $\dfrac{1}{\sin^{n}\operatorname{am}\left(\dfrac{2Kx}{\pi}\right)}$ into series proceeding according to the sines or cosines of multiples of $x$ itself. §§.\ 43---46. --- 115

The second species of elliptic integrals is expanded into series. §§.\ 47.\ 48. --- 133

The indefinite elliptic integrals of the third species are brought back to the definite case, in which the amplitude equals the parameter. §§.\ 49.\ 50. --- 137

The elliptic integrals of the third species are expanded into series. How they are conveniently expressed by means of the new transcendent $\Theta$. §§.\ 51.\ 52. --- 142

On the addition of the arguments, both of the amplitude and of the parameter, in the third species of elliptic integrals. §§.\ 53---55. --- 151

Reductions of the expressions $Z(iu)$, $\Theta(iu)$ to a real argument. General reduction of the third species of elliptic integrals, in which the arguments both of the amplitude and of the parameter are imaginary. §§.\ 56---60. --- 161

Elliptic functions are fractional functions. On the functions $H$, $\Theta$, which hold the place of numerator and denominator. §.\ 61. --- 172

On the expansion of the functions $H$, $\Theta$ into series. Third expansion of elliptic functions. Analytic demonstration of the Fermatian theorem, that every\ednote{The print has ununquemque, a misprint for unumquemque.} number is the sum of four squares. §§.\ 62---66. --- 176

\origpage{1}

\section*{ON THE TRANSFORMATION OF ELLIPTIC FUNCTIONS.}

\subsection*{STATEMENT OF THE GENERAL PROBLEM OF TRANSFORMATION.}

\subsection*{1.}

The most memorable integrals, which are exhibited by the Formula $\displaystyle\int\frac{d\Phi}{\sqrt{1-k^{2}\operatorname{Sin}\Phi^{2}}}$, and which constitute the first species of the so-called Elliptic Functions, depend on a twofold Argument, both on the Amplitude $\phi$ and on the Modulus $k$. By comparing with one another the values of such a function, which it takes for different Amplitudes while the Modulus remains the same, Analysts had discovered many excellent things, which concern the Addition and Multiplication of these. We have recently seen this question carried forward in a marvellous manner by Cl.\ \emph{Abel} in a Memoir, greater than our praise (See \emph{Crelle} Journal für reine und angewandte Mathematik V. II.).

There is another question, of no less moment --- indeed, taken in the widest sense, involving that one --- concerning the comparison of Elliptic Functions to be set up for different Moduli. This question, after the splendid discoveries of Cl$^{\mathrm{i}}$ \emph{Legendre} --- the Founder of the Theory of Elliptic Functions --- we were the first to recall to certain principles, and we have given its general solution (See \emph{Astronomische Nachrichten} A. 1827. No. 123. 127). This theory of ours on Transformation, and whatever else flows from it into the Analysis of Elliptic Functions, we shall now set forth more fully.

\origpage{2}

\subsection*{2.}

The Problem which we propose to ourselves is this, in general:

„\emph{A rational Function $y$ of the element $x$ is sought, such that there be:}
\[
\frac{dy}{\sqrt{A'+B'y+C'y^{2}+D'y^{3}+E'y^{4}}}=\frac{dx}{\sqrt{A+Bx+Cx^{2}+Dx^{3}+Ex^{4}}}\ \text{"}
\]
We see that this Problem embraces both Multiplication and Transformation.

Innumerable examples had long been established of such rational functions $y$, which satisfy the proposed problem. It was first known that, whatever whole odd number $n$ be given, a rational function $y$ of this kind can be exhibited, such that there be:
\[
\frac{dy}{\sqrt{A+By+Cy^{2}+Dy^{3}+Ey^{4}}}=\frac{n\,dx}{\sqrt{A+Bx+Cx^{2}+Dx^{3}+Ex^{4}}};
\]
which is the theorem on Multiplication. To this end the form must be employed:
\[
y=\frac{a+a'x+a''x^{2}+a'''x^{3}+\ldots+a^{(nn)}x^{nn}}{b+b'x+b''x^{2}+b'''x^{3}+\ldots+b^{(nn)}x^{nn}},
\]
the Coefficients $a, a', a'', .\,.\,.\,.;\ b, b', b'', \ldots$ being duly determined. It has also been sufficiently long explored that this form:
\[
y=\frac{a+a'x+a''x^{2}}{b+b'x+b''x^{2}},
\]
or this more general one:
\[
y=\frac{a+a'x+a''x^{2}+a'''x^{3}+\ldots+a^{(2^{m})}x^{2^{m}}}{b+b'x+b''x^{2}+b'''x^{3}+\ldots+b^{(2^{m})}x^{2^{m}}},
\]
which takes its origin from repetition of that substitution, can be so determined as to solve the problem. Quite recently it has also been proved by Cl$^{\mathrm{o}}$ Legendre that this form, duly determined, can be employed to that end:
\[
y=\frac{a+a'x+a''x^{2}+a'''x^{3}}{b+b'x+b''x^{2}+b'''x^{3}},
\]
or again, the same substitution being repeated, this more general one:
\[
y=\frac{a+a'x+a''x^{2}+a'''x^{3}+\ldots+a^{(3^{m})}x^{3^{m}}}{b+b'x+b''x^{2}+b'''x^{3}+\ldots+b^{(3^{m})}x^{3^{m}}}.
\]
From these forms joined together it is clear that the problem can be satisfied, with a suitable choice of the Coefficients, by putting:
\[
y=\frac{a+a'x+a''x^{2}+a'''x^{3}+\ldots+a^{(p)}x^{p}}{b+b'x+b''x^{2}+b'''x^{3}+\ldots+b^{(p)}x^{p}},
\]

\origpage{3}
provided $p$ be a number of the form $2^{\alpha}3^{\beta}(2m+1)^{2}$. In what follows it will now be proved that the same holds \emph{whatever number} $p$ \emph{may be.}

\subsection*{PRINCIPLES OF TRANSFORMATION.}

\subsection*{3.}

Let $U$, $V$ denote rational integral functions of the element $x$; let further $y=\dfrac{U}{V}$; there results:
\[
\frac{dy}{\sqrt{A'+B'y+C'y^{2}+D'y^{3}+E'y^{4}}}=\frac{V\,dU-U\,dV}{\sqrt{Y}},
\]
where for brevity's sake we have put:
\[
Y=A'V^{4}+B'V^{3}U+C'V^{2}U^{2}+D'VU^{3}+E'U^{4}.
\]
The fraction $\dfrac{V\,dU-U\,dV}{\sqrt{Y}}$ may be reduced to a simpler form whenever $Y$ has double factors; indeed, when besides four linear factors different from one another there exist among the remaining ones pairs equal to each other, that fraction returns of its own accord to the Differential of an Elliptic Function $\dfrac{dx}{U\sqrt{A+Bx+Cx^{2}+Dx^{3}+Ex^{4}}}$, $U$ denoting a rational function of the element $x$. Let us examine this case more accurately and see how many and what Conditions it demands for itself.

Let the functions $U$, $V$ be, the one of the $p^{\mathrm{ti}}$ order, the other of the $m^{\mathrm{ti}}$, so that $m\leqq p$: then $Y$ will be of the $(4p)^{\mathrm{ti}}$ order. Now in order that, four linear factors excepted, of the remaining factors of the function $Y$, whose number is $4p-4$, pairs may come out equal to each other, $(2p-2)$ Conditions will have to be satisfied. For as many double linear factors as the proposed function must have, so many Conditional Equations must hold among its Coefficients.

But in the functions $U$, $V$ there are present $m+p+2$ Constant Indeterminate Quantities, or rather $m+p+1$, since one of their number may be put $=1$. Their number is therefore either equal to the number of Conditions $2p-2$ or exceeds it, provided it be supposed that $m$ is one of the numbers $p-3$, $p-2$, $p-1$, $p$, in which cases the number of Indeterminates becomes respectively $2p-2$, $2p-1$, $2p$, $2p+1$. That the two former cases are to be rejected will be demonstrated below, and it is clear in this way as well. For when functions $U$, $V$ have been found which procure for the function $Y$ that prescribed form, if in place of $x$ there is substituted $\alpha+\beta x$, neither is the order of the functions $U$, $V$, $Y$ changed, nor the number of double factors of the function $Y$: whence we may at once introduce two Arbitrary Quantities into the solution found.

\origpage{4}
Therefore the number of Indeterminates must exceed the number of Conditions by at least two units, whence the cases $m=p-3$, $m=p-2$ are to be rejected. Further we see that, with $\dfrac{\alpha+\beta x}{1+\gamma x}$ put in place of $x$, the third case is reduced to the fourth and the fourth is not changed at all, so that in this case three Indeterminates both remain arbitrary and must remain so.

It is now therefore established, so far as it is permitted to conclude from comparing the number of Indeterminates and the number of Conditions with one another, that \emph{whatever number $p$ may be, the form:}
\[
y=\frac{a+a'x+a''x^{2}+\ldots+a^{(p)}x^{p}}{1+b'x+b''x^{2}+\ldots+b^{(p)}x^{p}};
\]
\emph{can be so determined that there be:}
\[
\frac{dy}{\sqrt{A'+B'y+C'y^{2}+D'y^{3}+E'y^{4}}}=\frac{dx}{M\sqrt{A+Bx+Cx^{2}+Dx^{3}+Ex^{4}}}
\]
\emph{$M$ denoting a rational function of $x$; indeed that the solution can involve three Arbitrary Quantities.}

\subsection*{4.}

In order that that function $M$ may be determined, let $Y=(A+Bx+Cx^{2}+Dx^{3}+Ex^{4})TT$, $T$ denoting a rational integral function of the element $x$: there will be
\[
M=\frac{T}{V\dfrac{dU}{dx}-U.\dfrac{dV}{dx}}.
\]
$T$ itself will be of the order $(2p-2)^{\mathrm{ti}}$; nor can $V\dfrac{dU}{dx}-U\dfrac{dV}{dx}$ be of greater. Now in certain cases it is established, namely where the number $p$ has that form $2^{\alpha}3^{\beta}(2n+1)^{2}$, that $M$ even becomes Constant. The same will be proved in general in what follows, whatever number $p$ may be.

We may suppose that the functions $U$, $V$ do not have a common factor; for with a common factor added, the fraction $\dfrac{U}{V}=y$ is not changed. Let us resolve the expression
\[
A'+B'y+C'y^{2}+D'y^{3}+E'y^{4}
\]
into linear factors, so that there be:
\[
A'+B'y+C'y^{2}+D'y^{3}+E'y^{4}=A'.(1-\alpha'y).(1-\beta'y)(1-\gamma'y)(1-\delta'y),
\]
whence also:
\[
Y=A'V^{4}+B'V^{3}U+C'V^{2}U^{2}+D'VU^{3}+E'U^{4}=A'(V-\alpha'U)(V-\beta'U)(V-\gamma'U)(V-\delta'U).
\]

\origpage{5}

Now there cannot exist a factor which is common to the quantities $V-\alpha'U$, $V-\beta'U$, $V-\gamma'U$, $V-\delta'U$, either to all or even to only two of their number; for the same would at once measure both $V$ and $U$, which we have supposed not to have a common factor. Therefore where some linear factor measures the function $Y$ twice, it is necessary that the same also measure twice some one of the quantities $V-\alpha'U$, $V-\beta'U$, $V-\gamma'U$, $V-\delta'U$.

Now let the following equations be noted:
\[
\begin{gathered}
(V-\alpha'U)\frac{dU}{dx}-\frac{d(V-\alpha'U)}{dx}.U=V\frac{dU}{dx}-U\frac{dV}{dx}\\
(V-\beta'U)\frac{dU}{dx}-\frac{d(V-\beta'U)}{dx}.U=V.\frac{dU}{dx}-U\frac{dV}{dx}\\
(V-\gamma'U)\frac{dU}{dx}-\frac{d(V-\gamma'U)}{dx}.U=V\frac{dU}{dx}-U\frac{dV}{dx}\\
(V-\delta'U)\frac{dU}{dx}-\frac{d(V-\delta'U)}{dx}.U=V.\frac{dU}{dx}-U\frac{dV}{dx},
\end{gathered}
\]
from which it follows that a factor which measures some one of the quantities $V-\alpha'U$, $V-\beta'U$, $V-\gamma'U$, $V-\delta'U$ twice, and therefore also its differential, measures the expression $V\dfrac{dU}{dx}-U\dfrac{dV}{dx}$ as well. But the product of all those factors, which also measure $Y$ itself twice, we have put $=T$, whence $T$ will measure $V\dfrac{dU}{dx}-U\dfrac{dV}{dx}$ itself. But $T$ is not of lower order than $V\dfrac{dU}{dx}-U\dfrac{dV}{dx}$ itself, whence we see that
\[
M=\frac{T}{V\dfrac{dU}{dx}-U\dfrac{dV}{dx}}
\]
passes over into a Constant.

Let us further remark that, where one of the functions $U$, $V$ had been of lower order than the $(p-1)^{\mathrm{ti}}$, $V\dfrac{dU}{dx}-U\dfrac{dV}{dx}$ itself would also have been of lower order than $T$, which nevertheless ought to measure it; since this is absurd, the cases $m=p-2$, $m=p-3$ had to be rejected.

It is now therefore demonstrated that \emph{the form:}
\[
y=\frac{a+a'x+a''^{2}+\ldots+a^{(p)}x^{p}}{b+b'x+b''^{2}+\ldots+b^{(p)}x^{p}},
\]\ednote{In the print the letter $x$ is missing before the exponent in the terms $a''x^{2}$ and $b''x^{2}$ (printed as $a''^{2}$ and $b''^{2}$), and the $x$ of the last terms is set low beside the $p$.}

\origpage{6}
\emph{whatever number $p$ may be, can be so determined that there result:}
\[
\frac{dy}{\sqrt{A'+B'y+C'y^{2}+D'y^{3}+E'y^{4}}}=\frac{dx}{\sqrt{A+Bx+Cx^{2}+Dx^{3}+Ex^{4}}}.
\]
This is the \emph{Fundamental Principle in the Theory of Transformations of Elliptic Functions.}

\subsection*{THE EXPRESSION $\dfrac{dy}{\sqrt{\pm(y-\alpha)(y-\beta)(y-\gamma)(y-\delta)}}$ IS PROPOSED, TO BE REDUCED TO THE SIMPLER FORM $\dfrac{dx}{M\sqrt{(1-x)(1-k^{2}x^{2})}}$.}

\subsection*{5.}

\ednote{In the heading above the factor is printed as $(1-x)$ instead of $(1-x^{2})$ (compare the index and the text of this section); an apparent misprint.}By the aid of the three Arbitrary Constants, which we have seen the solution of our Problem to admit, the expression $A+Bx+Cx^{2}+Dx^{3}+Ex^{4}$ can be reduced to this simpler one: $A(1-x^{2})(1-k^{2}x^{2})$. In order that this and the rest which have just been demonstrated may also be shown by examples, let it be proposed to transform the given expression:
\[
\frac{dy}{\sqrt{\pm(y-\alpha)(y-\beta)(y-\gamma)(y-\delta)}}
\]
by making the substitution:
\[
y=\frac{a+a'x+a''x^{2}}{b+b'x+b''x^{2}}
\]
into this simpler one:
\[
\frac{dx}{M\sqrt{(1-x^{2})(1-k^{2}x^{2})}}.
\]
The question is how the substitution to be employed, the Modulus $k$ and the Constant factor $M$ are to be determined from the given quantities $\alpha$, $\beta$, $\gamma$, $\delta$.

Let $a+a'x+a''x^{2}=U$, $b+b'x+b''x^{2}=V$, $y=\dfrac{U}{V}$: by the principles just set forth there must be:
\[
(U-\alpha V)(U-\beta V)(U-\gamma V)(U-\delta V)=K(1-x^{2})(1-k^{2}x^{2})(1+mx)^{2}(1+nx)^{2},
\]
$K$ denoting some arbitrary Constant. Hence we see that two of the number of factors $U-\alpha V$, $U-\beta V$, $U-\gamma V$, $U-\delta V$, which will be of the second order, even become squares. Let us therefore put:
\[
\begin{gathered}
U-\gamma V=C(1+mx)^{2}\\
U-\delta V=D(1+nx)^{2}.
\end{gathered}
\]

\origpage{7}

Now as regards the remaining factors $U-\alpha V$, $U-\beta V$, there can be put either:
\[
\begin{gathered}
U-\alpha V=A(1-x^{2}),\ U-\beta V=B(1-k^{2}x^{2}),\ \text{or:}\\
U-\alpha V=A.(1-x)(1-kx),\ U-\beta V=B(1+x)(1+kx),
\end{gathered}
\]
$A$, $B$, $C$, $D$ denoting Constant quantities. The former is to be rejected. For there would result
\[
\frac{U-\alpha V}{U-\beta V}=\frac{y-\alpha}{y-\beta}=\frac{A}{B}.\,\frac{1-x^{2}}{1-k^{2}x^{2}},
\]
whence it would follow that, the element $x$ being changed into $-x$, $y$ remains unchanged, which is clearly absurd from the equations:
\[
\begin{gathered}
\frac{U-\alpha V}{U-\gamma V}=\frac{y-\alpha}{y-\gamma}=\frac{A}{B}.\,\frac{1-k^{2}}{(1+mx)^{2}}\\
\frac{U-\alpha V}{U-\delta V}=\frac{y-\alpha}{y-\delta}=\frac{A}{D}.\,\frac{1-k^{2}}{(1+nx)^{2}}.
\end{gathered}
\]\ednote{In these two displays the print has $\frac{A}{B}.\frac{1-k^{2}}{(1+mx)^{2}}$ and $\frac{A}{D}.\frac{1-k^{2}}{(1+nx)^{2}}$, evidently a misprint for the correct factors $\frac{A}{C}.\frac{(1-x)(1-kx)}{(1+mx)^{2}}$ and $\frac{A}{D}.\frac{(1-x)(1-kx)}{(1+nx)^{2}}$ (the author's Corrigenda list these places). The copy also carries faint hand corrections in ink (a $C$ over the $B$ and an added $x^{2}$ in the second display), which are not transcribed.}
There must therefore be put:
\[
\begin{gathered}
\text{1)}\quad U-\alpha V=A(1-x).(1-kx)\\
\text{2)}\quad U-\beta V=B(1+x).(1+kx)\\
\text{3)}\quad U-\gamma V=C(1+mx)^{2}\\
\text{4)}\quad U-\delta V=D(1+nx)^{2}.
\end{gathered}
\]
It is fitting to remark that of the Constants $A$, $B$, $C$, $D$ one can be determined at will.

\subsection*{6.}

We see from equation 1), both on putting $x=1$ and on putting $x=\dfrac{1}{k}$, that $U=\alpha V$. Hence from the equation:
\[
\frac{U-\gamma V}{U-\beta V}=\frac{C}{B}.\,\frac{(1+mx)^{2}}{(1+x)(1+kx)}
\]
with $x=1$ put, there results:
\[
\frac{\alpha-\gamma}{\alpha-\beta}=\frac{C}{B}.\,\frac{(1+m)^{2}}{2(1+k)};
\]
with $x=\dfrac{1}{k}$ put:
\[
\frac{\alpha-\gamma}{\alpha-\beta}=\frac{C}{B}.\,\frac{\left(1+\dfrac{m}{k}\right)^{2}}{2\left(1+\dfrac{1}{k}\right)},
\]
whence:
\[
(1+m)^{2}=k\left(1+\frac{m}{k}\right)^{2}.
\]

\origpage{8}

In an entirely similar manner is found:
\[
(1+n)^{2}=k\left(1+\frac{n}{k}\right)^{2},
\]
whence $m=\sqrt{k}$, $n=-\sqrt{k}$. For it is not permitted to put $m$ and $n$ equal; for then the expression $\dfrac{U-\gamma V}{U-\delta V}=\dfrac{y-\gamma}{y-\delta}$, and therefore $y$ itself, would pass over into a Constant.

Now in the equation:
\[
\frac{U-\gamma V}{U-\delta V}=\frac{y-\gamma}{y-\delta}=\frac{C}{D}.\left\{\frac{1+\sqrt{k.x}}{1-\sqrt{k.x}}\right\}^{2}
\]\ednote{The print has the radicand $k.x$ under one bar in both roots; the author's Corrigenda ask for $\sqrt{k}.x$ instead.}
let there be put first $x=+1$, in which case $U=\alpha V$; then $x=-1$, in which case $U=\beta V$: the two following equations result:
\[
\begin{gathered}
\frac{\alpha-\gamma}{\alpha-\delta}=\frac{C}{D}\left\{\frac{1+\sqrt{k}}{1-\sqrt{k}}\right\}^{2}\\
\frac{\beta-\gamma}{\beta-\delta}=\frac{C}{D}\left\{\frac{1-\sqrt{k}}{1+\sqrt{k}}\right\}^{2}.
\end{gathered}
\]
These equations being multiplied into one another, there comes:
\[
\frac{C}{D}=\sqrt{\frac{(\alpha-\gamma)(\beta-\gamma)}{(\alpha-\delta)(\beta-\delta)}},
\]
whence it is permitted to put:
\[
\begin{gathered}
C=\sqrt{(\alpha-\gamma)(\beta-\gamma)}\\
D=\sqrt{(\alpha-\delta)(\beta-\delta)}\ ;
\end{gathered}
\]
for of the quantities $A$, $B$, $C$, $D$ one could be determined at will.

From the same equations, the one being divided by the other, we obtain:
\[
\frac{1+\sqrt{k}}{1-\sqrt{k}}=\frac{\sqrt[4]{(\alpha-\gamma)(\beta-\delta)}}{\sqrt[4]{(\alpha-\delta)(\beta-\gamma)}};
\]
whence:
\[
\sqrt{k}=\frac{\sqrt[4]{(\alpha-\gamma)(\beta-\delta)}-\sqrt[4]{(\alpha-\delta)(\beta-\gamma)}}{\sqrt[4]{(\alpha-\gamma)(\beta-\delta)}+\sqrt[4]{(\alpha-\delta)(\beta-\gamma)}}.
\]
Let the formula be further noted:
\[
\sqrt{k}+\frac{1}{\sqrt{k}}=2.\frac{\sqrt{(\alpha-\gamma)(\beta-\delta)}+\sqrt{(\alpha-\delta)(\beta-\gamma)}}{\sqrt{(\alpha-\gamma)(\beta-\delta)}-\sqrt{(\alpha-\delta)(\beta-\gamma)}},
\]

\origpage{9}

whence:
\[
\begin{gathered}
(1-\sqrt{k})\left(1-\frac{1}{\sqrt{k}}\right)=\frac{-4\sqrt{(\alpha-\delta)(\beta-\gamma)}}{\sqrt{(\alpha-\gamma)(\beta-\delta)}-\sqrt{(\alpha-\delta)(\beta-\gamma)}}\\
(1+\sqrt{k})\left(1+\frac{1}{\sqrt{k}}\right)=\frac{+4\sqrt{(\alpha-\gamma)(\beta-\delta)}}{\sqrt{(\alpha-\gamma)(\beta-\delta)}-\sqrt{(\alpha-\delta)(\beta-\gamma)}}.
\end{gathered}
\]
In order that the Constants $A$, $B$ may be defined, I observe that from equations 1), 2), 3), with $x=\dfrac{1}{\sqrt{k}}$ put, by which $U=\delta V$, there is obtained:\ednote{The author's Corrigenda (printed p. 189, for p. 9) direct that $U$ and $V$ be interchanged throughout this page; the print is reproduced as printed.}
\[
\begin{gathered}
\frac{\delta-\alpha}{\delta-\gamma}=\frac{A(1-\sqrt{k})\left(1-\sqrt{\dfrac{1}{k}}\right)}{4\sqrt{(\alpha-\gamma)(\beta-\delta)}}\\
\frac{\delta-\beta}{\delta-\gamma}=\frac{B(1+\sqrt{k})\left(1+\sqrt{\dfrac{1}{k}}\right)}{4\sqrt{(\alpha-\gamma)(\beta-\gamma)}},
\end{gathered}
\]\ednote{The denominator of the first fraction is printed $4\sqrt{(\alpha-\gamma)(\beta-\delta)}$, while that of the second is $4\sqrt{(\alpha-\gamma)(\beta-\gamma)}$; the derivation would suggest $(\beta-\gamma)$ in the first as well.}
whence:
\[
\begin{gathered}
A=\frac{-\sqrt{(\alpha-\gamma)(\alpha-\delta)}}{\gamma-\delta}\left\{\sqrt{(\alpha-\gamma)(\beta-\delta)}-\sqrt{(\alpha-\delta)(\beta-\gamma)}\right\}\\
B=\frac{\sqrt{(\beta-\gamma)(\beta-\delta)}}{\gamma-\delta}\left\{\sqrt{(\alpha-\gamma)(\beta-\delta)}-\sqrt{(\alpha-\delta)(\beta-\gamma)}\right\}.
\end{gathered}
\]

\subsection*{7.}

From the general principles established above by us it follows that in our example the expression $V\dfrac{dU}{dx}-U\dfrac{dV}{dx}$ will be equal to the product $(1+\sqrt{k}x)(1-\sqrt{k}x)$ multiplied into a constant quantity, which is confirmed by the calculation made as follows.

There results, as is established by carrying out the expansion:
\[
(\gamma-\delta)\left(V\frac{dU}{dx}-U\frac{dV}{dx}\right)=(V-\gamma U)\frac{d(V-\delta U)}{dx}-(V-\delta U)\frac{d(V-\gamma U)}{dx}
\]
But we have obtained:
\[
\begin{gathered}
V-\gamma U=C(1+\sqrt{k}.x)^{2}\\
V-\delta U=D(1-\sqrt{k}.x)^{2},
\end{gathered}
\]
whence:
\[
\begin{gathered}
\frac{d(V-\gamma U)}{dx}=2C(1+\sqrt{k}.x)\sqrt{k}\\
\frac{d(V-\delta U)}{dx}=-2D(1-\sqrt{k}.x)\sqrt{k}.
\end{gathered}
\]

\origpage{10}

Whence results:
\[
(\gamma-\delta)\left(V\frac{dU}{dx}-U\frac{dV}{dx}\right)=+4\sqrt{k}.CD(1+\sqrt{k}.x)(1-\sqrt{k}.x).
\]\ednote{The author's Corrigenda (printed p. 190) direct that $-\sqrt{k}$ be replaced by $+\sqrt{k}$ in lines 2, 4 and 6 of this page; the scan shows $+4\sqrt{k}$ in all three displays, as given here.}
All these things being duly collected, we obtain:
\[
\frac{dy}{\sqrt{-(y-\alpha)(y-\beta)(y-\gamma)(y-\delta)}}=\frac{+4\sqrt{k}}{\gamma-\delta}.\sqrt{\frac{CD}{-AB}}.\frac{dx}{\sqrt{(1-x^{2})(1-k^{2}x^{2})}};
\]
whence:
\[
\begin{aligned}
M&=\frac{\gamma-\delta}{+4\sqrt{k}}\sqrt{\frac{-AB}{CD}}=\frac{\sqrt{(\alpha-\gamma)(\beta-\delta)}-\sqrt{(\alpha-\delta)(\beta-\gamma)}}{4\sqrt{k}}\\
&=\left\{\frac{\sqrt[4]{(\alpha-\gamma)(\beta-\delta)}+\sqrt[4]{(\alpha-\delta)(\beta-\gamma)}}{2}\right\}^{2},
\end{aligned}
\]
whence:
\[
\begin{aligned}
\frac{dy}{\sqrt{-(y-\alpha)(y-\beta)(y-\gamma)(y-\delta)}}&=\frac{dx}{M\sqrt{(1-x^{2})(1-k^{2}x^{2})}}=\\
&\frac{dx}{\sqrt{1-xx}\sqrt{\left(\left(\dfrac{\sqrt[4]{(\alpha-\gamma)(\beta-\delta)}+\sqrt[4]{(\alpha-\delta)(\beta-\gamma)}}{2}\right)^{4}-\left(\dfrac{\sqrt[4]{(\alpha-\gamma)(\beta-\delta)}-\sqrt[4]{(\alpha-\delta)(\beta-\gamma)}}{2}\right)^{4}xx\right)}}.
\end{aligned}
\]
With $(\alpha-\gamma).(\beta-\delta)=G$, $(\alpha-\delta)(\beta-\gamma)=G'$ put, there results:
\[
\frac{dx}{M.\sqrt{(1-x^{2})(1-k^{2}x^{2})}}=\frac{dx}{\sqrt{1-x^{2}}\sqrt{\left(\dfrac{\sqrt[4]{G}+\sqrt[4]{G'}}{2}\right)^{4}-\left(\dfrac{\sqrt[4]{G}-\sqrt[4]{G'}}{2}\right)^{4}x^{2}}}.
\]
Let $G=mm$, $G'=nn$, let further:
\[
\begin{gathered}
m'=\frac{1}{2}(m+n),\ n'=\sqrt{mn}\\
m''=\frac{1}{2}(m'+n'),\ n''=\sqrt{m'n'},
\end{gathered}
\]
there will be, with $x=\operatorname{Sin}\phi$ put:
\[
\frac{dx}{M\sqrt{(1-x^{2})(1-k^{2}x^{2})}}=\frac{d\phi}{\sqrt{m''m''\operatorname{Cos}\phi^{2}+n''n''\operatorname{Sin}\phi^{2}}}.
\]
Moreover the value of $x$ is most easily computed by the aid of the formula:
\[
\frac{1-\sqrt{k}.x}{1+\sqrt{k}.x}=\sqrt[4]{\frac{(\alpha-\gamma)(\beta-\gamma)}{(\alpha-\delta)(\beta-\delta)}}.\sqrt{\frac{y-\delta}{y-\gamma}};
\]
where:
\[
\sqrt{k}=\frac{\sqrt[4]{G}-\sqrt[4]{G'}}{\sqrt[4]{G}+\sqrt[4]{G'}}=\sqrt[4]{\frac{m''m''-n''n''}{m''m''}}.
\]

\origpage{11}

\subsection*{8.}

The quantities $\alpha$, $\beta$, $\gamma$, $\delta$ in the proposed formulae may be permuted among themselves at will. What is placed at our discretion becomes fixed and defined as soon as the condition is added that, if it can be done, the transformation succeed by a real substitution. Let us examine this more accurately.

Let us suppose the quantities $\alpha$, $\beta$, $\gamma$, $\delta$ all to be real; let further $\alpha>\beta>\gamma>\delta$, so that $\alpha-\beta$, $\alpha-\gamma$, $\alpha-\delta$ are positive quantities. Now a distinction will have to be made according to the limits between which the value of the argument $y$ is contained:
\[
\text{1)}\ \delta\ \text{and}\ \gamma,\quad\text{2)}\ \gamma\ \text{and}\ \beta,\quad\text{3)}\ \beta\ \text{and}\ \alpha,\quad\text{4)}\ \alpha\ \text{and}\ \delta.
\]
In the last case suppose the passage from $\alpha$ to $\delta$ to be made through infinity. We see the expression $\dfrac{1}{\sqrt{(y-\alpha)(y-\beta)(y-\gamma)(y-\delta)}}$ become real only in the second and fourth case, the expression $\dfrac{1}{\sqrt{-(y-\alpha)(y-\beta)(y-\gamma)(y-\delta)}}$ only in the first and third case. The real substitutions which correspond to those four cases will be indicated by Table I. Then let us embrace in Table II the formulae which serve for transforming the expression $\dfrac{dy}{\sqrt{\pm(y-\alpha)(y-\beta)(y-\gamma)}}$ into a simpler one by a real substitution, for the limits between which the value of the argument $y$ is contained:
\[
\text{1)}\ -\infty\ \text{and}\ \gamma,\quad\text{2)}\ \gamma\ \text{and}\ \beta,\quad\text{3)}\ \beta\ \text{and}\ \alpha,\quad\text{4)}\ \alpha\ \text{and}\ +\infty.
\]
These formulae can easily be derived from Table I by dividing by $\delta$ and then putting $\delta=-\infty$.

\origpage{12}

\subsection*{TABLE I.}

A.
\[
\frac{dy}{\sqrt{(y-\alpha)(y-\beta)(y-\gamma)(y-\delta)}}=\frac{dx}{\sqrt{1-x^{2}}\sqrt{L^{4}-N^{4}x^{2}}}
\]
\[
\begin{gathered}
L=\frac{\sqrt[4]{(\alpha-\gamma)(\beta-\delta)}+\sqrt[4]{(\alpha-\beta)(\gamma-\delta)}}{2}\\
N=\frac{\sqrt[4]{(\alpha-\gamma)(\beta-\delta)}-\sqrt[4]{(\alpha-\beta)(\gamma-\delta)}}{2}
\end{gathered}
\]
\[
\begin{gathered}
\text{I.}\quad\text{Limits:}\ \alpha.\,.\pm\infty.\,.\delta:\quad\frac{L-Nx}{L+Nx}=\sqrt[4]{\frac{(\alpha-\beta)(\beta-\delta)}{(\alpha-\gamma)(\gamma-\delta)}}.\sqrt{\frac{y-\gamma}{y-\beta}}\\
\text{II.}\quad\text{Limits:}\ \gamma.\,.\,.\,.\beta:\quad\frac{L-Nx}{L+Nx}=\sqrt[4]{\frac{(\beta-\delta)(\gamma-\delta)}{(\alpha-\beta)(\alpha-\gamma)}}\ \sqrt{\frac{\alpha-y}{y-\delta}}.
\end{gathered}
\]

B.
\[
\frac{dy}{\sqrt{-(y-\alpha)(y-\beta)(y-\gamma)(y-\delta)}}=\frac{dx}{\sqrt{1-x^{2}}\sqrt{L^{4}-N^{4}x^{2}}}
\]
\[
\begin{gathered}
L=\frac{\sqrt[4]{(\alpha-\gamma)(\beta-\delta)}+\sqrt[4]{(\alpha-\delta)(\beta-\gamma)}}{2}\\
N=\frac{\sqrt[4]{(\alpha-\gamma)(\beta-\delta)}-\sqrt[4]{(\alpha-\delta)(\beta-\gamma)}}{2}
\end{gathered}
\]
\[
\begin{gathered}
\text{I.}\quad\text{Limits:}\ \beta.\,.\,.\,.\alpha:\quad\frac{L-Nx}{L+Nx}=\sqrt[4]{\frac{(\alpha-\gamma)(\beta-\gamma)}{(\alpha-\delta)(\beta-\delta)}}.\sqrt{\frac{y-\delta}{y-\gamma}}\\
\text{II.}\quad\text{Limits:}\ \delta.\,.\,.\,.\gamma:\quad\frac{L-Nx}{L+Nx}=\sqrt[4]{\frac{(\alpha-\gamma)(\alpha-\delta)}{(\beta-\gamma)(\beta-\delta)}}\ \sqrt{\frac{\beta-y}{\alpha-y}}.
\end{gathered}
\]

\origpage{13}

\subsection*{TABLE II.}

A.
\[
\frac{dy}{\sqrt{(y-\alpha)(y-\beta)(y-\gamma)}}=\frac{dx}{\sqrt{1-x^{2}}\sqrt{L^{4}-N^{4}x^{2}}}
\]
\[
\begin{gathered}
L=\frac{\sqrt[4]{\alpha-\gamma}+\sqrt[4]{\alpha-\beta}}{2}\\
N=\frac{\sqrt[4]{\alpha-\gamma}-\sqrt[4]{\alpha-\beta}}{2}
\end{gathered}
\]
\[
\begin{gathered}
\text{I.}\quad\text{Limits}\ \alpha.\,.\,.\,.+\infty:\quad\frac{L-Nx}{L+Nx}=\sqrt[4]{\frac{\alpha-\beta}{\alpha-\gamma}}\ \sqrt{\frac{y-\gamma}{y-\beta}}\\
\text{II.}\quad\text{Limits}\ \gamma.\,.\,.\,.\beta:\quad\frac{L-Nx}{L+Nx}=\frac{\sqrt{\alpha-y}}{\sqrt[4]{(\alpha-\beta)(\alpha-\gamma)}}.
\end{gathered}
\]

B.
\[
\frac{dy}{\sqrt{-(y-\alpha)(y-\beta)(y-\gamma)}}=\frac{dx}{\sqrt{1-x^{2}}\sqrt{L^{4}-N^{4}x^{2}}}
\]
\[
\begin{gathered}
L=\frac{\sqrt[4]{\alpha-\gamma}+\sqrt[4]{\beta-\gamma}}{2}\\
N=\frac{\sqrt[4]{\alpha-\gamma}-\sqrt[4]{\beta-\gamma}}{2}
\end{gathered}
\]
\[
\begin{gathered}
\text{I.}\quad\text{Limits}\ \beta.\,.\,.\,.\alpha:\quad\frac{L-Nx}{L+Nx}=\frac{\sqrt[4]{(\alpha-\gamma)(\beta-\gamma)}}{\sqrt{y-\gamma}}\\
\text{II.}\quad\text{Limits}\ -\infty\ldots\gamma:\quad\frac{L-Nx}{L+Nx}=\sqrt[4]{\frac{\alpha-\gamma}{\beta-\gamma}}\ \sqrt{\frac{\beta-y}{\alpha-y}}.
\end{gathered}
\]

\origpage{14}

\subsection*{9.}

In these formulae, for the assigned limits, $x$ passes from $-1$ up to $+1$ and at the same time $y$ from one limit to the other. But with the limits which correspond to formulae I and II interchanged with one another, we see an imaginary value of the form $\pm iR$ produced for the expression $\frac{L-Nx}{L+Nx}$, with $i=\sqrt{-1}$ put and $R$ denoting some real quantity; and for $x$ itself the form is procured
\[
\frac{Le^{i\Phi}}{N}=\frac{e^{i\Phi}}{\sqrt{k}};\ \text{whence}\ \frac{L-Nx}{L+Nx}=\frac{1-e^{i\Phi}}{1+e^{i\Phi}}=\frac{e^{-\frac{i\Phi}{2}}-e^{\frac{i\Phi}{2}}}{e^{-\frac{i\Phi}{2}}+e^{\frac{i\Phi}{2}}}=-i\operatorname{tang}\frac{\Phi}{2}.
\]

The form to which we have been carried on this occasion, $x=\frac{e^{i\Phi}}{\sqrt{k}}$, let us substitute in the expression $\frac{dx}{\sqrt{(1-x^{2})(1-k^{2}x^{2})}}$. There results:
\[
\begin{gathered}
\frac{dx}{\sqrt{(1-x^{2})(1-k^{2}x^{2})}}=\frac{i e^{i\Phi}d\Phi}{\sqrt{k}.\sqrt{\left(1-\frac{e^{2i\Phi}}{k}\right)\left(1-ke^{2i\Phi}\right)}}=\frac{d\Phi}{\sqrt{\left(1-ke^{2i\Phi}\right)\left(1-ke^{-2i\Phi}\right)}}\\
=\frac{d\Phi}{\sqrt{1-2k\cos2\Phi+kk}}=\frac{d\Phi}{\sqrt{(1-k)^{2}\cos\Phi^{2}+(1+k)^{2}\sin\Phi^{2}}}.
\end{gathered}
\]
This substitution indeed seems to us sufficiently memorable. From it there also flows the following more general formula, by putting $x=\sin\psi$:
\[
\frac{k^{n}\sin\psi^{2n}d\psi}{\sqrt{1-k^{2}\sin\psi^{2}}}=\frac{(\cos2n\Phi+i\sin2n\Phi)\,d\Phi}{\sqrt{1-2k\cos2\Phi+kk}},
\]
whence for the limits $0$ and $\pi$ is obtained, the imaginary part vanishing:
\[
\int_{0}^{\pi}\frac{k^{n}\sin\psi^{2n}d\psi}{\sqrt{1-k^{2}\sin\psi^{2}}}=\int_{0}^{\pi}\frac{\cos2n\Phi\,d\Phi}{\sqrt{1-2k\cos2\Phi+kk}}=\int_{0}^{\pi}\frac{\cos n\Phi\,d\Phi}{\sqrt{1-2k\cos\Phi+kk}},
\]
which is a succinct demonstration of the memorable formula put forth by Cl.\ Legendre. From Tables I and II two others may be derived as follows, the limits between which the value of $y$ is contained being interchanged, and with $x=\frac{Le^{i\Phi}}{N}$ put. For the assigned limits the angle $\Phi$ increases from $0$ up to $\pi$, while $y$ passes from one limit to the other.

\origpage{15}

\subsection*{TABLE III.}

A.
\[
\frac{dy}{\sqrt{(y-\alpha)(y-\beta)(y-\gamma)(y-\delta)}}=\frac{d\Phi}{\sqrt{mm\cos\Phi^{2}+nn\sin\Phi^{2}}}
\]
\[
\begin{gathered}
m=\sqrt[4]{(\alpha-\gamma)(\beta-\delta)(\alpha-\beta)(\gamma-\delta)}\\
n=\frac{\sqrt{(\alpha-\gamma)(\beta-\delta)}+\sqrt{(\alpha-\beta)(\gamma-\delta)}}{2}
\end{gathered}
\]
\[
\begin{gathered}
\text{I.}\quad\text{Limits}\ \gamma\ldots\beta:\quad\operatorname{tg}\frac{\Phi}{2}=\sqrt[4]{\frac{(\alpha-\beta)(\beta-\delta)}{(\alpha-\gamma)(\gamma-\delta)}}\ \sqrt{\frac{y-\gamma}{\beta-y}}\\
\text{II.}\quad\text{Limits}\ \alpha\ldots\delta:\quad\operatorname{tg}\frac{\Phi}{2}=\sqrt[4]{\frac{(\alpha-\beta)(\alpha-\gamma)}{(\beta-\delta)(\gamma-\delta)}}\ \sqrt{\frac{y-\alpha}{y-\delta}}.
\end{gathered}
\]

B.
\[
\frac{dy}{\sqrt{-(y-\alpha)(y-\beta)(y-\gamma)(y-\delta)}}=\frac{d\Phi}{\sqrt{mm\cos\Phi^{2}+nn\sin\Phi^{2}}}
\]
\[
\begin{gathered}
m=\sqrt[4]{(\alpha-\gamma)(\beta-\delta)(\alpha-\delta)(\beta-\gamma)}\\
n=\frac{\sqrt{(\alpha-\gamma)(\beta-\delta)}+\sqrt{(\alpha-\delta)(\beta-\gamma)}}{2}
\end{gathered}
\]
\[
\begin{gathered}
\text{I.}\quad\text{Limits}\ \delta\ldots\gamma:\quad\operatorname{tg}\frac{\Phi}{2}=\sqrt[4]{\frac{(\alpha-\gamma)(\beta-\gamma)}{(\alpha-\delta)(\beta-\delta)}}\ \sqrt{\frac{y-\delta}{\gamma-y}}\\
\text{II.}\quad\text{Limits}\ \beta\ldots\alpha:\quad\operatorname{tg}\frac{\Phi}{2}=\sqrt[4]{\frac{(\alpha-\gamma)(\alpha-\delta)}{(\beta-\gamma)(\beta-\delta)}}\ \sqrt{\frac{y-\beta}{\alpha-y}}.
\end{gathered}
\]

\origpage{16}

\subsection*{TABLE IV.}

A.
\[
\frac{dy}{\sqrt{(y-\alpha)(y-\beta)(y-\gamma)}}=\frac{d\Phi}{\sqrt{mm\cos\Phi^{2}+nn\sin\Phi^{2}}}
\]
\[
\begin{gathered}
m=\sqrt[4]{(\alpha-\gamma)(\alpha-\beta)}\\
n=\frac{\sqrt{\alpha-\gamma}+\sqrt{\alpha-\beta}}{2}
\end{gathered}
\]
\[
\begin{gathered}
\text{I.}\quad\text{Limits}\ \gamma\ldots\beta:\quad\operatorname{tg}\frac{\Phi}{2}=\sqrt[4]{\frac{\alpha-\beta}{\alpha-\gamma}}\ \sqrt{\frac{y-\gamma}{\beta-y}}\\
\text{II.}\quad\text{Limits}\ \alpha\ldots+\infty:\quad\operatorname{tg}\frac{\Phi}{2}=\frac{\sqrt{y-\alpha}}{\sqrt[4]{(\alpha-\beta)(\alpha-\gamma)}}.
\end{gathered}
\]

B.
\[
\frac{dy}{\sqrt{-(y-\alpha)(y-\beta)(y-\gamma)}}=\frac{d\Phi}{\sqrt{mm\cos\Phi^{2}+nn\sin\Phi^{2}}}
\]
\[
\begin{gathered}
m=\sqrt[4]{(\alpha-\gamma)(\beta-\gamma)}\\
n=\frac{\sqrt{\alpha-\gamma}+\sqrt{\beta-\gamma}}{2}
\end{gathered}
\]
\[
\begin{gathered}
\text{I.}\quad\text{Limits}\ -\infty\ldots\gamma:\quad\operatorname{tg}\frac{\Phi}{2}=\frac{\sqrt{(\alpha-\gamma)(\beta-\gamma)}}{\sqrt{\gamma-y}}\\
\text{II.}\quad\text{Limits}\ \beta\ldots\alpha:\quad\operatorname{tg}\frac{\Phi}{2}=\sqrt[4]{\frac{\alpha-\gamma}{\beta-\gamma}}\ \sqrt{\frac{y-\beta}{\alpha-y}}.
\end{gathered}
\]

\origpage{17}

We have treated this question at some length, so that a worked example may be at hand. There remain the cases in which of the quantities $\alpha$, $\beta$, $\gamma$, $\delta$ either two or four are imaginary. The former case itself admits a real solution, which however suffers from an imaginary appearance. The latter case admits no such real solution at all. Hence, in order that everything may be brought back to real quantities, new transformations will be needed, whereby the elegance of the formulae perishes. We therefore forgo this question.

To the proposed substitution there corresponds another, its inverse, of the form
\[
x=\frac{a+a'y+a''y}{b+b'y+b''y},
\]\ednote{Printed $a''y$ and $b''y$ without the exponent 2 (the Corrigenda read $y^{2}$); a small hand-written 2 in ink is not transcribed.}
which also furnishes most elegant formulae. But since we perhaps already seem to dwell too long on this question, let us relegate its investigation to another occasion. We return to the general questions.

\section*{ON THE TRANSFORMATION OF THE EXPRESSION $\dfrac{dy}{\sqrt{1-y^{2}}\sqrt{1-\lambda^{2}y^{2}}}$ INTO ANOTHER SIMILAR TO IT $\dfrac{dx}{M\sqrt{1-x^{2}}\sqrt{1-k^{2}x^{2}}}$.}

\subsection*{10.}

We have seen that the given expression:
\[
\frac{dy}{\sqrt{A'+B'y+C'y^{2}+D'y^{3}+E'y^{4}}}
\]
can be transformed, by the employed substitution of the kind:
\[
y=\frac{a+a'x+a''x^{2}+.\,.\,.\,.+a^{(p)}x^{p}}{b+b'x+b''x^{2}+.\,.\,.\,.+b^{(p)}x^{p}}=\frac{U}{V},
\]
whatever the number $p$ may be, into another similar to it:
\[
\frac{dx}{\sqrt{A+Bx+Cx^{2}+Dx^{3}+Ex^{4}}}.
\]

Such a substitution depends both on the given Coefficients $A'$, $B'$, $C'$, $D'$, $E'$, and above all on the number $p$, since this denotes the exponent of the highest power that is found in the rational functions $U$, $V$. Wherefore in what follows we shall say that such

\origpage{18}
a substitution or transformation \emph{is of the $p^{\mathrm{ti}}$ order or belongs to the $p^{\mathrm{tum}}$ order, or more simply pertains to the number $p$.}

Now, being about to examine the nature of these substitutions more accurately, let us set aside that more complex form:
\[
\frac{dy}{\sqrt{A+By+Cy^{2}+Dy^{3}+Ey^{4}}},
\]
and let us inquire into this simpler one $\frac{dy}{\sqrt{1-y^{2}}.\sqrt{1-\lambda^{2}y^{2}}}$, to which we both see, and it is known, that the former can be brought back, as it is to be transformed into another similar to it $\frac{dx}{M\sqrt{1-x^{2}}\sqrt{1-k^{2}x^{2}}}$.

When the nature of the proposed question has been duly weighed, the problem is found to be satisfied, provided that of the functions $U$, $V$ one is set to be odd and the other even, which the examples hitherto explored by Analysts already indicate. In this matter we shall have above all to distinguish between the case in which the order of the odd function is less than the order of the even one and that in which it is greater than the order of the even one; or between the case in which the transformation pertains to an even number and that in which it pertains to an odd number.

Let us then prove \emph{first} that the transformation succeeds when a substitution of even order is employed, or of the form:
\[
y=\frac{x\left(a+a'x^{2}+a''x^{4}+\ldots+a^{(m-1)}x^{2m-2}\right)}{1+b'x^{2}+b''x^{4}+\ldots+b^{(m)}x^{2m}}=\frac{U}{V}.
\]
Here the functions $V+U$, $V-U$, $V+\lambda U$, $V-\lambda U$ will themselves be of even order, whence let us put:
\[
\begin{gathered}
\text{1)}\quad V+U=(1+x)(1+kx)AA\\
\text{2)}\quad V-U=(1-x)(1-kx)BB\\
\text{3)}\quad V+\lambda U=CC\\
\text{4)}\quad V-\lambda U=DD;
\end{gathered}
\]
where $A$, $B$, $C$, $D$ denote rational integral functions of the element $x$. As soon as these equations are satisfied, there will be obtained, as we have proved:
\[
\frac{dy}{\sqrt{1-y^{2}}\sqrt{1-\lambda^{2}y}}=\frac{dx}{M\sqrt{1-x^{2}}\sqrt{1-k^{2}x^{2}}}.
\]\ednote{Printed $\lambda^{2}y$ instead of $\lambda^{2}y^{2}$ under the second root on the left; an apparent misprint.}
Since, when $x$ is changed into $-x$, $U$ passes into $-U$ while $V$ is not changed, from the equations 1), 3) the remaining ones 2), 4) flow of themselves. In order that the equations 1), 3) may be satisfied, $V+\lambda U$ must have $m$ times,

\origpage{19}
and $V+U$ $(m-1)$ times, two equal linear factors; moreover to $V+U$ itself the factor $1+x$ must also be assigned. All of which demand of themselves Equations of Condition in number $m+m-1+1=2m$, which is itself the number of the Indeterminates $a$, $a'$, $\ldots a^{(m-1)}$; $b'$, $b''$, $\ldots b^{(m)}$. Whence the proposed problem is determinate.

In the \emph{second} place let us prove that the transformation also succeeds when a substitution of this kind is employed:
\[
y=\frac{x\left(a+a'x^{2}+a''x^{4}+\ldots+a^{(m)}x^{2m}\right)}{1+b'x^{2}+b''x^{4}+\ldots+b^{(m)}x^{2m}}=\frac{U}{V};
\]
which pertains to an odd number. Here $V+U$, $V-U$, $V+\lambda U$, $V-\lambda U$ are themselves of odd order, whence let us put:
\[
\begin{gathered}
\text{1)}\quad V+U=(1+x)AA\\
\text{2)}\quad V-U=(1-x)BB\\
\text{3)}\quad V+\lambda U=(1+kx)CC\\
\text{4)}\quad V-\lambda U=(1-kx)DD.
\end{gathered}
\]
Here too only the equations 1), 3) will have to be satisfied, since from them, by changing $x$ into $-x$, the two remaining ones flow of themselves. In order that they may be satisfied, it is necessary that both $V+U$ and $V+\lambda U$ each have $m$ times two equal linear factors, to which end $2m$ Equations of Condition will have to be satisfied, to which one is added, that moreover $V+U$ acquire $(1+x)$ as a factor. Hence we see that the number of Equations of Condition is $2m+1$, which is itself the number of the Indeterminates $a$, $a'$, $a''$, $.\,.\,a^{(m)}$; $b'$, $b''$, $\ldots b^{(m)}$; whence in this case too the problem is determinate.

\subsection*{11.}

Let $U'$, $V'$ denote rational integral functions of the element $y$ of such a kind that, on putting $z=\frac{U'}{V'}$, there results:
\[
\frac{dz}{\sqrt{1-z^{2}}\sqrt{1-\mu^{2}z^{2}}}=\frac{dy}{M'\sqrt{1-y^{2}}\sqrt{1-\lambda^{2}y^{2}}}.
\]
Let the substitution $z=\frac{U'}{V'}$ which has been employed be of the order $p'^{\mathrm{ti}}$; and by another substitution

\origpage{20}
$y=\frac{U}{V}$, (where $U$, $V$ denote, as above, rational integral functions of the element $x$,) which may be of the order $p^{\mathrm{ti}}$, there will result, as above:
\[
\frac{dy}{\sqrt{1-y^{2}}\sqrt{1-\lambda^{2}y^{2}}}=\frac{dx}{M\sqrt{1-x^{2}}\sqrt{1-k^{2}x^{2}}}.
\]
Now when the value $y=\frac{U}{V}$ is substituted in the expression $z=\frac{U'}{V'}$, let there arise $z=\frac{U''}{V''}$: this will be one substitution $z=\frac{U''}{V''}$, by the employment of which there results:
\[
\frac{dz}{\sqrt{1-z^{2}}\sqrt{1-\mu^{2}z^{2}}}=\frac{dx}{MM'\sqrt{1-x^{2}}\sqrt{1-k^{2}x^{2}}}
\]
of the order $(pp')^{\mathrm{ti}}$. Thus we see that from several transformations, which pertain respectively to the numbers $p$, $p'$, $p''$, $\ldots$, employed successively, one can be composed which pertains to the number $pp'p''\ldots$. And conversely, which however we shall not prove at present, a transformation which pertains to some composite number $pp'p''.\,.$ can always be composed from others employed successively, which pertain respectively to the numbers $p$, $p'$, $p''.\,.$. Wherefore only those transformations need be investigated which pertain to a \emph{prime} number, since from them all the remaining ones can be composed. Now therefore in what follows let us set aside the first case, which regards an even order of the transformation, since it can always be composed from a transformation of odd order and the transformation which pertains to the number 2, repeated again and again where needed. But the \emph{second} case, or the transformations of odd order, let us now examine more closely.

\subsection*{12.}

We see that in that case two functions, the one $V$ even of the $2m^{\mathrm{ti}}$ order, the other $U$ odd of the $(2m+1)^{\mathrm{ti}}$ order, are to be so determined that there be:
\[
\begin{gathered}
V+U=(1+x)AA\\
V+\lambda U=(1+kx)CC.
\end{gathered}
\]
\emph{Now I say, if the functions $U$, $V$ are so determined that, on putting $\frac{1}{kx}$ in place of $x$, $y=\frac{U}{V}$ passes into $\frac{1}{\lambda y}=\frac{V}{\lambda U}$: then those equations follow, the one from the other, of themselves.}

\origpage{21}

Let us put $V=\phi(x^{2})$, $U=xF(x^{2})$; we see that the expression $y=\frac{xF(x^{2})}{\phi(x^{2})}$, on putting $\frac{1}{kx}$ in place of $x$, passes into
\[
\frac{F\left(\frac{1}{k^{2}x^{2}}\right)}{kx\,\phi\left(\frac{1}{k^{2}x^{2}}\right)}=\frac{x^{2m}F\left(\frac{1}{k^{2}x^{2}}\right)}{kx.x^{2m}\,\phi\left(\frac{1}{k^{2}x^{2}}\right)},
\]
where $x^{2m}F\left(\frac{1}{k^{2}x^{2}}\right)$, $x^{2m}\phi\left(\frac{1}{k^{2}x^{2}}\right)$ are integral functions. In order that this may become equal to the expression $\frac{1}{\lambda y}=\frac{V}{\lambda U}=\frac{\phi x^{2}}{\lambda xF(x^{2})}$\ednote{Printed $\phi x^{2}$ without parentheses in the numerator; an apparent misprint for $\phi(x^{2})$.}, the following equations must hold:
\[
\begin{gathered}
\phi(x^{2})=px^{2m}F\left(\frac{1}{k^{2}x^{2}}\right)\\
\lambda F(x^{2})=pkx^{2m}\phi\left(\frac{1}{k^{2}x^{2}}\right),
\end{gathered}
\]
where $p$ denotes a Constant quantity. When in these equations we again put $\frac{1}{kx}$ in place of $x$, we obtain:
\[
\begin{gathered}
\phi\left(\frac{1}{k^{2}x^{2}}\right)=\frac{p}{k^{2m}x^{2m}}F(x^{2})\\
\lambda F\left(\frac{1}{k^{2}x^{2}}\right)=\frac{pk}{k^{2m}x^{2m}}\phi(x^{2}).
\end{gathered}
\]
When these are compared with the earlier equations, we obtain $\frac{p}{k^{2m}}=\frac{\lambda}{pk}$, whence
\[
p=\sqrt{\lambda k^{2m-1}}.
\]
Hence there results:
\[
\begin{gathered}
\phi(x^{2})=x^{2m}\sqrt{\lambda k^{2m-1}}\,F\left(\frac{1}{k^{2}x^{2}}\right)\\
F(x^{2})=x^{2m}\sqrt{\frac{k^{2m+1}}{\lambda}}\,\phi\left(\frac{1}{k^{2}x^{2}}\right),
\end{gathered}
\]
of which equations the one follows from the other.

Now whenever the expression:
\[
\frac{V+U}{1+x}=\frac{\phi(x^{2})+xF(x^{2})}{1+x}
\]
is the square of a rational integral function of the element $x$, the same will also hold of another, which

\origpage{22}
is derived from that one by putting $\frac{1}{kx}$ in place of $x$ and multiplying by $x^{2m}\sqrt{\lambda k^{2m-1}}$. When this is done we obtain, \emph{provided $\frac{V+U}{1+x}$ be a square, the function:}
\[
\begin{gathered}
x^{2m}\sqrt{\lambda k^{2m-1}}\,\frac{\phi\left(\frac{1}{k^{2}x^{2}}\right)+\frac{1}{kx}F\left(\frac{1}{k^{2}x^{2}}\right)}{1+\frac{1}{kx}}\\
=\frac{\sqrt{\lambda k^{2m-1}}\,x^{2m}F\left(\frac{1}{k^{2}x^{2}}\right)+\sqrt{\lambda k^{2m+1}}\,x^{2m+1}\phi\left(\frac{1}{k^{2}x^{2}}\right)}{1+kx}\\
=\frac{\phi(x^{2})+\lambda xF(x^{2})}{1+kx}=\frac{V+\lambda U}{1+kx},
\end{gathered}
\]
\emph{and this too will be a square.} Q.\ E.\ D.

Thus the problem has been brought back to this, that the expression:
\[
\frac{\phi(x^{2})+\sqrt{\frac{k^{2m+1}}{\lambda}}\,x^{2m+1}\phi\left(\frac{1}{k^{2}x^{2}}\right)}{1+x}=\frac{V+U}{1+x}.
\]
be rendered a square, where $\phi(x^{2})$ denotes an expression of this kind:
\[
\phi(x^{2})=V=b+b'x^{2}+b''x^{4}+\ldots+b^{(m)}x^{2m}.
\]
Now, on putting $U=xF(x^{2})=x\left(a+a'x^{2}+a''x^{4}+\ldots+a^{(m)}x^{2m}\right)$, since there is
\[
U=xF(x^{2})=\sqrt{\frac{x^{2m+1}}{\lambda}}\,x^{2m+1}\phi\left(\frac{1}{k^{2}x^{2}}\right):
\]\ednote{Under the root, the numerator is printed $x^{2m+1}$ where $k^{2m+1}$ is expected, as in the preceding formula for $F(x^{2})$; a faint hand-written $k$ in the copy is not transcribed.}
it comes about that:
\[
\mathcal{H}\left\{
\begin{array}{l}
a=\sqrt{\dfrac{k}{\lambda}}.\dfrac{b^{(m)}}{k^{m}};\qquad a'=\sqrt{\dfrac{k}{\lambda}}.\dfrac{b^{(m-1)}}{k^{m-4}},\qquad a''=\sqrt{\dfrac{k}{\lambda}}.\dfrac{b^{(m-2)}}{k^{m-4}},\ \ldots\\[2.5ex]
a^{(m)}=\sqrt{\dfrac{k}{\lambda}}.\,k^{m};\qquad a^{(m-1)}=\sqrt{\dfrac{k}{\lambda}}.\,b'k^{m-2},\qquad a^{(m-2)}=\sqrt{\dfrac{k}{\lambda}}.\,b''k^{m-4},\ \ldots
\end{array}
\right.
\]\ednote{In the value of $a'$ the exponent of $k$ in the denominator is printed $m-4$ where $m-2$ is expected; a hand correction to 2 in the copy is not transcribed.} \ednote{In the value of $a^{(m)}$ the factor $b$ is not printed, which gives $k^{m}$ where $bk^{m}$ is expected, as in the neighbouring values $b'k^{m-2}$ and $b''k^{m-4}$; a gap is left in the print before $k^{m}$.}
Now we descend to the examples.

\origpage{23}

\subsection*{THE TRANSFORMATION OF THE THIRD ORDER IS PROPOSED.}

\subsection*{13.}

Let $m=1$, which is the simplest case, $V=1+b'x^{2}$, $U=x(a+a'x^{2})$. On putting $A=(1+\alpha x)$, we find:
\[
\begin{gathered}
AA=(1+\alpha x)^{2}=1+2\alpha x+\alpha^{2}x^{2},\ \text{whence:}\\
V+U=(1+x)AA=1+(1+2\alpha)x+\alpha(2+\alpha)x^{2}+\alpha^{2}x^{3}.
\end{gathered}
\]\ednote{In the first line the closing parenthesis after $1+x$ is printed as a small comma-like mark, so that the text reads $(1+x,\ AA$; read here as $(1+x)AA$.}
Hence there results:
\[
b'=\alpha(2+\alpha),\ a=1+2\alpha,\ a'=\alpha^{2}.
\]
The equations $\mathcal{H}$ §.\ 12 pass into the following:
\[
a=\sqrt{\frac{k}{\lambda}}.\frac{b'}{k};\ a'=\sqrt{\frac{k^{3}}{\lambda}};
\]
whence we obtain:
\[
1+2\alpha=\frac{\alpha(2+\alpha)}{\sqrt{k\lambda}};\ \alpha^{2}=\frac{\sqrt{k^{3}}}{\sqrt{\lambda}},\ \text{whence}\ \alpha=\sqrt[4]{\frac{k^{3}}{\lambda}}.
\]
Let $\sqrt[4]{k}=u$, $\sqrt[4]{\lambda}=v$; there will be $\alpha=\frac{u^{3}}{v}$, $1+2\alpha=\frac{v+2u^{3}}{v}$, $\alpha(2+\alpha)=\frac{u^{3}(2v+u^{3})}{v^{2}}$. Hence the equation:
\[
1+2\alpha=\alpha(2+\alpha)\sqrt{\frac{1}{k\lambda}}
\]
passes into the following:
\[
\frac{v+2u^{3}}{v}=\frac{u(2v+v^{3})}{v^{4}},
\]\ednote{Printed $v^{3}$ in the numerator on the right where $u^{3}$ is expected; a hand correction to $u^{3}$ in the copy is not transcribed.}
or:
\[
\text{1)}\quad u^{4}-v^{4}+2uv(1-u^{2}v^{2})=0.
\]
Moreover there results:
\[
\begin{gathered}
a=(1+2\alpha)=\frac{v+2u^{3}}{v}\\
a'=\alpha\alpha=\frac{u^{6}}{v^{2}}\\
b'=\alpha(2+\alpha)=u^{3}\left(\frac{2v+u^{3}}{v^{2}}\right)=vu^{2}(v+2u^{3}).
\end{gathered}
\]
Hence we obtain:
\[
\text{2)}\quad y=\frac{(v+2u^{3})vx+u^{6}x^{3}}{vv+v^{3}u^{2}(v+2u^{3})x^{2}}.
\]

\origpage{24}

Moreover we obtain, since $1+y=\frac{(1+x)AA}{V}$:
\[
\begin{gathered}
\text{3)}\quad 1+y=\frac{(1+x)(v+u^{3}x)^{2}}{vv+v^{3}u^{2}(v+2u^{3})x^{2}}\\[1.5ex]
\text{4)}\quad 1-y=\frac{(1-x)(v-u^{3}x)^{2}}{vv+v^{3}u^{2}(v+2u^{3})x^{2}}\\[1.5ex]
\text{5)}\quad \sqrt{\frac{1-y}{1+y}}=\sqrt{\frac{1-x}{1+x}}.\frac{v-u^{3}x}{v+u^{3}x}\\[1.5ex]
\text{6)}\quad \sqrt{1-yy}=\frac{\sqrt{1-xx}(v^{2}-u^{6}x^{2})}{vv+v^{3}u^{3}(v+2u^{3})x^{2}}.
\end{gathered}
\]\ednote{In formula 6) the denominator is printed with $u^{3}$ where $u^{2}$ is expected, as in 3) and 4); an apparent misprint.}
Further, on putting in place of $x$ the quantity $\frac{1}{kx}=\frac{1}{u^{4}x}$, since $y$ passes into $\frac{1}{\lambda y}=\frac{1}{v^{4}y}$, we find the following system of formulae:
\[
\begin{gathered}
\text{7)}\quad 1+v^{4}y=\frac{(1+u^{4}x)(1+uvx)^{2}}{1+vu^{2}(v+2u^{3})x^{2}}\\[1.5ex]
\text{8)}\quad 1-v^{4}y=\frac{(1-u^{4}x)(1-uvx)^{2}}{1+vu^{2}(v+2u^{3})x^{2}}\\[1.5ex]
\text{9)}\quad \sqrt{\frac{1-v^{4}y}{1+v^{4}y}}=\sqrt{\frac{1-u^{4}x}{1+u^{4}x}}.\frac{1-uvx}{1+uvx}\\[1.5ex]
\text{10)}\quad \sqrt{1-v^{8}y^{2}}=\frac{\sqrt{1-u^{8}x^{2}}(1-u^{2}v^{2}x^{2})}{1+vu^{2}(v+2u^{3})x^{2}}.
\end{gathered}
\]

\subsection*{14.}

On putting $V+U=(1+x)AA$, $V+\lambda U=(1+k)CC$\ednote{Printed $(1+k)CC$ without the $x$; an apparent misprint for $(1+kx)CC$.}, $V-U=(1-x)BB$, $V-\lambda U=(1-kx)DD$, we have seen that there results:
\[
ABCD=M\left\{V\frac{dU}{dx}-U\frac{dV}{dx}\right\},
\]
where $M$ denotes a Constant quantity; which one may find out by a comparison of the Coefficient of one and the same power, set up in each of the two expressions $ABCD$, $V\frac{dU}{dx}-U\frac{dV}{dx}$. Now on putting $V=b+b'x^{2}+$ etc., $U=ax+a'x^{3}+$ etc., in the single expressions $A$, $B$, $C$, $D$ the Constant becomes $\sqrt{b}$, whence in the product compounded from them it is $bb$; but in the expression $V\frac{dU}{dx}-U\frac{dV}{dx}$ we see the Constant become $ab$; whence:
\[
M=\frac{b}{a}.
\]

\origpage{25}

Hence in our example it results, since $b=1$,
\[
a=\frac{v+2u^{3}}{v}=\frac{u(2v+u^{3})}{v^{4}}:
\]
\[
M=\frac{v}{v+2u^{3}}=\frac{v^{4}}{u(2v+u^{3})},
\]
whence:
\[
\frac{dy}{\sqrt{(1-y^{2})(1-v^{8}y^{2})}}=\frac{(v+2u^{3})\,dx}{v\sqrt{(1-x^{2})(1-u^{8}x^{2})}}
\]
The Moduli $k$, $\lambda$, which we have seen depend on each other through an equation of the fourth degree §.\ 13.\ 1), are easily expressed rationally by the same quantity $\alpha$. For from the formulae brought forward above:
\[
\alpha=\frac{u^{3}}{v};\ 1+2\alpha=\frac{\alpha(2+\alpha)}{\sqrt{k\lambda}}=\frac{\alpha(2+\alpha)}{u^{2}v^{2}}
\]
it follows:
\[
\alpha=\frac{u^{3}}{v};\ u^{2}v^{2}=\frac{\alpha(2+\alpha)}{1+2\alpha},
\]
whence:
\[
\begin{gathered}
u^{8}=\frac{\alpha^{3}(2+\alpha)}{1+2\alpha}=k^{2}\\
v^{8}=\alpha\left(\frac{2+\alpha}{1+2\alpha}\right)^{3}=\lambda^{2}.
\end{gathered}
\]
There results moreover: $M=\dfrac{1}{1+2\alpha}$, whence, on putting $y=\sin T'$, $x=\sin T$, the equation:
\[
\frac{dy}{\sqrt{1-y^{2}}\sqrt{1-\lambda^{2}y^{2}}}=\frac{dx}{M\sqrt{1-x^{2}}\sqrt{1-k^{2}x^{2}}},
\]
passes into the following:
\[
\frac{dT'}{\sqrt{(1+2\alpha)^{2}-\alpha(2+\alpha)^{3}\sin T'^{2}}}=\frac{dT}{\sqrt{1+2\alpha-\alpha^{3}(2+\alpha)\sin T^{2}}},
\]\ednote{The print has the exponent 2 on the first factor under the first root, where 3 is expected; the author's Corrigenda corrects it to 3.}
or into this:
\[
\frac{dT'}{\sqrt{(1+2\alpha)^{3}\operatorname{Cos}T'^{2}+(1-\alpha)^{3}(1+\alpha)\sin T'^{2}}}=\frac{dT}{\sqrt{(1+2\alpha)\operatorname{Cos}T^{2}+(1+\alpha)^{3}(1-\alpha)\sin T^{2}}},
\]
which is arrived at by the substitution made:
\[
\operatorname{Sin}T'=\frac{(1+2\alpha)\sin T+\alpha\alpha\sin T^{3}}{1+\alpha(2+\alpha)\sin T^{2}}.
\]

\origpage{26}

\subsection*{THE TRANSFORMATION OF THE FIFTH ORDER IS PROPOSED.}

\subsection*{15.}

Now let us pass to the example which is next in simplicity, in which $m=2$,
\[
V=1+b'x^{2}+b''x^{4},\quad U=x(a+a'x^{2}+a''x^{4}),\quad A=1+\alpha x+\beta x^{2}.
\]
We find:
\[
AA=1+2\alpha x+(2\beta+\alpha\alpha)x^{2}+2\alpha\beta x^{3}+\beta\beta x^{4}
\]
whence:
\[
AA(1+x)=1+x(1+2\alpha)+x^{2}(2\alpha+2\beta+\alpha\alpha)+x^{3}(2\beta+\alpha\alpha+2\alpha\beta)+x^{4}(2\alpha\beta+\beta\beta)+\beta\beta x^{5}.
\]
Hence we obtain:
\[
\begin{gathered}
b'=2\alpha+2\beta+\alpha\alpha,\quad b''=\beta(2\alpha+\beta)\\
a=1+2\alpha,\quad a'=2\beta+\alpha\alpha+2\alpha\beta,\quad a''=\beta\beta.
\end{gathered}
\]
The equations $\mathcal{H}$ §.\ 12 become:
\[
a=\sqrt{\frac{k}{\lambda}}.\frac{b''}{k^{2}},\quad a'=\sqrt{\frac{k}{\lambda}}.b',\quad a''=\sqrt{\frac{k^{5}}{\lambda}}.
\]
From these it follows:
\[
\frac{a'a'}{aa''}=\frac{b'b'}{b''},
\]
or, since we have $b'=(2\alpha+\beta)+(\beta+\alpha\alpha)$, $a'=\beta(1+2\alpha)+(\beta+\alpha\alpha)$:
\[
\frac{\left\{(2\alpha+\beta)+(\beta+\alpha\alpha)\right\}^{2}}{2\alpha+\beta}=\frac{\left\{\beta(1+2\alpha)+(\beta+\alpha\alpha)\right\}^{2}}{\beta(1+2\alpha)},
\]
whence:
\[
2\alpha+\beta+\frac{(\beta+\alpha\alpha)^{2}}{2\alpha+\beta}=\beta(1+2\alpha)+\frac{(\beta+\alpha\alpha)^{2}}{\beta(1+2\alpha)}.
\]
Hence it easily follows:
\[
\beta(1+2\alpha)(2\alpha+\beta)=(\beta+\alpha\alpha)^{2},
\]
which, expanded and divided by $\alpha$, passes into:
\[
\alpha^{3}=2\beta(1+\alpha+\beta).
\]
This equation can also be represented in these two ways:
\[
\begin{gathered}
(\alpha\alpha+\beta)(\alpha-2\beta)=\beta(2-\alpha)(1+2\alpha)\\
(\alpha\alpha+\beta)(2-\alpha)=(\alpha-2\beta)(2\alpha+\beta),
\end{gathered}
\]

\origpage{27}
whence it follows:
\[
\left(\frac{2-\alpha}{\alpha-2\beta}\right)^{2}=\frac{2\alpha+\beta}{\beta(1+2\alpha)}.
\]
With these preparations made, the rest is easily accomplished. For we find, on putting $k=u^{4}$, $\lambda=v^{4}$:
\[
\frac{2\alpha+\beta}{\beta(1+2\alpha)}=\frac{b''}{aa''}=\frac{b'b'}{a'a'}=\frac{\lambda}{k}=\frac{v^{4}}{u^{4}},
\]
whence also:
\[
\frac{2-\alpha}{\alpha-2\beta}=\frac{v^{2}}{u^{2}}.
\]
There is moreover $\beta=\sqrt{a''}=\sqrt[4]{\dfrac{k^{5}}{\lambda}}=\dfrac{u^{5}}{v}$, whence the equations:
\[
\frac{v^{4}}{u^{4}}=\left(\frac{2-\alpha}{\alpha-2\beta}\right)^{2}=\frac{2\alpha+\beta}{\beta(1+2\alpha)};\ \frac{2-\alpha}{\alpha-2\beta}=\frac{v^{2}}{u^{2}}
\]
pass into the following:
\[
\begin{gathered}
2\alpha v+u^{5}=uv^{4}(1+2\alpha)\\
u^{2}(2-\alpha)=v(v\alpha-2u^{5}),
\end{gathered}
\]
or:
\[
\begin{gathered}
2\alpha v(1-uv^{3})=u(v^{4}-u^{4})\\
\alpha(vv+uu)=2u^{2}(1+u^{3}v),
\end{gathered}
\]
whence:
\[
(u^{2}+v^{2})(u^{4}-v^{4})+4uv(1+u^{3}v)(1-uv^{3})=0.
\]
When the expansion is made there results:
\[
\text{1)}\quad u^{6}-v^{6}+5u^{2}v^{2}(u^{2}-v^{2})+4uv(1-u^{4}v^{4})=0.
\]

The rest are found thus. From the equations:
\[
\begin{gathered}
2\alpha v(1-uv^{3})=u(v^{4}-u^{4})\\
\alpha(uu+vv)=2u^{2}(1+u^{3}v),
\end{gathered}
\]
it follows:
\[
\alpha=\frac{u(v^{4}-u^{4})}{2v(1-uv^{3})}=\frac{2u^{2}(1+u^{3}v)}{u^{2}+v^{2}}.
\]
Hence there results:
\[
\begin{gathered}
a=1+2\alpha=\frac{1}{v}\left(\frac{v-u^{5}}{1-uv^{3}}\right)\\
\beta+2\alpha=\frac{u^{5}}{v}+2\alpha=uv^{2}\left(\frac{v-u^{5}}{1-uv^{3}}\right)
\end{gathered}
\]

\origpage{28}
\[
\begin{gathered}
\alpha-2\beta=\alpha-\frac{2u^{5}}{v}=\frac{2u^{2}}{v}\left(\frac{v-u^{5}}{u^{2}+v^{2}}\right)\\
2-\alpha=2v\left(\frac{v-u^{5}}{u^{2}+v^{2}}\right)\\
\alpha\alpha+\beta=\frac{(\alpha-2\beta)(2\alpha+\beta)}{2-\alpha}=u^{3}\left(\frac{v-u^{5}}{1-uv^{3}}\right).
\end{gathered}
\]
Hence at length is deduced:
\[
\begin{gathered}
b'=\beta+2\alpha+\alpha\alpha+\beta=\frac{u(u^{2}+v^{2})(v-u^{5})}{1-uv^{3}}\\
b''=\frac{u^{5}}{v}(2\alpha+\beta)=u^{6}v\left(\frac{v-u^{5}}{1-uv^{3}}\right)\\
a=\frac{1}{v}\left(\frac{v-u^{5}}{1-uv^{3}}\right)\\
a'=\frac{u^{2}}{v^{2}}.\,b'=u^{3}\left(\frac{u^{2}+v^{2}}{v^{2}}\right)\left(\frac{v-u^{5}}{1-uv^{3}}\right)\\
a''=\frac{u^{10}}{v^{2}}.
\end{gathered}
\]
Now since $M=\dfrac{1}{a}=v\left(\dfrac{1-uv^{3}}{v-u^{5}}\right)$, the transformation of the fifth order will be contained in the following theorem:

\subsection*{THEOREM.}

On putting:
\[
\begin{gathered}
\text{1)}\quad u^{6}-v^{6}+5u^{2}v^{2}(u^{2}-v^{2})+4uv(1-u^{4}v^{4})=0\\
\text{2)}\quad y=\frac{v(v-u^{5})x+u^{3}(u^{2}+v^{2})(v-u^{5})x^{3}+u^{10}(1-uv^{3})x^{5}}{v^{2}(1-uv^{3})+uv^{2}(u^{2}+v^{2})(v-u^{5})x^{2}+u^{6}v^{3}(v-u^{5})x^{4}}
\end{gathered}
\]
there results:
\[
\frac{v(1-uv^{3})\,dy}{\sqrt{1-y^{2}}\sqrt{1-v^{8}y^{2}}}=\frac{(v-u^{5})\,dx}{\sqrt{1-x^{2}}\sqrt{1-u^{8}x^{2}}}.
\]

\origpage{29}

\subsection*{HOW, BY THE TRANSFORMATION APPLIED TWICE, ONE ARRIVES AT MULTIPLICATION.}

\subsection*{16.}

To one inspecting the equations between $u$ and $v$ found in the two proposed examples:
\[
\begin{gathered}
u^{4}+v^{4}+2uv(1-u^{2}v^{2})=0\\
u^{6}-v^{6}+5u^{2}v^{2}(u^{2}-v^{2})+4uv(1-u^{4}v^{4})=0.
\end{gathered}
\]\ednote{In the first equation the sign before $v^{4}$ is printed as $+$ instead of $-$ (as in the next equation and the following examples); an apparent misprint, struck through by a hand correction in the scan.}
it cannot escape notice that they remain unchanged when $v$ is put in place of $u$, and $-v$ moreover in place of $u$. Hence from the theorem found in the first example, namely on putting:
\[
\begin{gathered}
u^{4}-v^{4}+2uv(1-u^{2}v^{2})=0\\
y=\frac{v(v+2u^{3})x+u^{6}x^{3}}{v^{2}+v^{3}u^{2}(v+2u^{3})x^{2}},
\end{gathered}
\]
there results:
\[
\frac{dy}{\sqrt{1-y^{2}}\sqrt{1-v^{8}y^{2}}}=\frac{v+2u^{3}}{v}.\,\frac{dx}{\sqrt{1-x^{2}}\sqrt{1-u^{8}x^{2}}}
\]
another is immediately derived from this, that on putting:
\[
z=\frac{u(u-2v^{3})y+v^{6}y^{3}}{u^{2}+u^{3}v^{2}(u-2v^{3})y^{2}},
\]
there results:
\[
\frac{dz}{\sqrt{1-z^{2}}\sqrt{1-u^{8}z^{2}}}=\frac{u-2v^{3}}{u}.\,\frac{dy}{\sqrt{1-y^{2}}\sqrt{1-v^{8}y^{2}}}.
\]
Now indeed there is:
\[
\left(\frac{v+2u^{3}}{v}\right)\left(\frac{u-2v^{3}}{u}\right)=\frac{2(u^{4}-v^{4})+uv(1-u^{2}v^{2})}{uv}=-3,
\]\ednote{The factor $(1-u^{2}v^{2})$ is printed; the product of the two brackets gives $uv(1-4u^{2}v^{2})$, so $1-4u^{2}v^{2}$ is expected. A hand-written 4 is added above in the scan.}
whence it follows:
\[
\frac{dz}{\sqrt{1-z^{2}}\sqrt{1-u^{8}z^{2}}}=\frac{-3\,dx}{\sqrt{1-x^{2}}\sqrt{1-u^{8}x^{2}}}.
\]
In order that $+3$ may result in place of $-3$, either $z$ must be changed into $-z$, or $x$ into $-x$.

In a similar way from the theorem proposed in the second example another is deduced, namely on putting:
\[
z=\frac{u(u+v^{5})y+v^{3}(u^{2}+v^{2})(u+v^{5})y^{3}+v^{10}(1+u^{3}v)y^{5}}{u^{2}(1+u^{3}v)+u^{2}v(u^{2}+v^{2})(u+v^{5})y^{2}+u^{3}v^{6}(u+v^{5})y^{4}}
\]\ednote{Printed with $+$ before $v^{3}(u^{2}+v^{2})$ in the numerator and before $u^{2}v(u^{2}+v^{2})$ in the denominator, where $-$ is expected (substituting $v$ for $u$ and $-u$ for $v$ in the preceding theorem); an apparent misprint, struck through by hand corrections in the scan.}

\origpage{30}
there results:
\[
\frac{dz}{\sqrt{1-z^{2}}\sqrt{1-u^{8}z^{2}}}=\frac{u+v^{5}}{u(1+u^{3}v)}.\,\frac{dy}{\sqrt{1-y^{2}}\sqrt{1-v^{8}y^{2}}}.
\]
Now since it follows from the equation:
\[
\begin{gathered}
u^{6}-v^{6}+5u^{2}v^{2}(u^{2}-v^{2})+4uv(1-u^{4}v^{4})=0,\\
\frac{(u+v^{5})(v-u^{5})}{uv(1+u^{3}v)(1-uv^{3})}=\frac{uv(1-u^{4}v^{4})-(u^{6}-v^{6})}{uv(1+u^{3}v)(1-uv^{3})}=5,
\end{gathered}
\]
we see that there results:
\[
\frac{dz}{\sqrt{1-z^{2}}\sqrt{1-u^{8}z^{2}}}=\frac{5\,dx}{\sqrt{1-x^{2}}\sqrt{1-u^{8}x^{2}}}.
\]
Thus, by the transformation applied twice, one arrives at Multiplication.

These two examples, viz.\ the transformations of the third and fifth order, I presented earlier in letters which I gave to Cl.\ Schumacher in the month of June of the year 1827. See Nova Astron.\ l.\ l. And in the same place I also proclaimed the generality of the method by which they were found. The latter had already been found two years earlier by Cl.\ Legendre.

\section*{ON THE NEW NOTATION OF ELLIPTIC FUNCTIONS.}

\subsection*{17.}

The algebraic questions having been set aside, let us inquire more accurately into the analytic nature of our functions. But first it is necessary that we indicate the manner of notation of which use will be made in what follows.

On putting $\displaystyle\int_{0}^{\phi}\frac{d\phi}{\sqrt{1-k^{2}\sin\phi^{2}}}=u$, Geometers have been accustomed to call the angle $\phi$ the \emph{amplitude} of the function $u$. This angle therefore we shall denote in what follows by: ampl.\ $u$ or, more briefly, by:
\[
\phi=\operatorname{am}.u
\]
Thus, where $\displaystyle\int_{0}^{x}\frac{dx}{\sqrt{1-x^{2}}\sqrt{1-k^{2}x^{2}}}=u$, there will be:
\[
x=\sin.\operatorname{am}.u.
\]

\origpage{31}
Moreover, on putting:
\[
\int_{0}^{1}\frac{dx}{\sqrt{1-x^{2}}\sqrt{1-k^{2}x^{2}}}=\int_{0}^{\frac{\pi}{2}}\frac{dx}{\sqrt{1-k^{2}\sin\phi^{2}}}=K,
\]\ednote{In the second integral the numerator is printed as $dx$ where $d\phi$ is expected (the author's Corrigenda give $d\phi$); a hand-written $\phi$ is added over the $x$ in the scan.}
we shall call $K-u$ the Complement of the function $u$; the amplitude of the Complement we shall denote by \emph{coam}, so that there is:
\[
\operatorname{am}(K-u)=\operatorname{coam}.u.
\]

The expression $\sqrt{1-k^{2}\sin^{2}\operatorname{am}u}=\dfrac{d.\operatorname{am}u}{du}$, following the lead of Cl.\ Legendre, we shall denote by
\[
\Delta\operatorname{am}u=\sqrt{1-k^{2}\sin^{2}\operatorname{am}u}.
\]
The Complement, as it is called by Cl.\ Legendre, of the Modulus $k$ I shall designate by $k'$, so that there is:
\[
kk+k'k'=1.
\]

Further, from our notation there will be:
\[
K'=\int_{0}^{\frac{\pi}{2}}\frac{d\phi}{\sqrt{1-k'k'\sin\phi^{2}}}.
\]
The Modulus, which must be understood where needed, will be added either enclosed in brackets, or appended in the margin. Where the Modulus is not added, in what follows you are to understand everywhere the same Modulus $k$.

The expressions themselves $\sin\operatorname{am}u$, $\sin\operatorname{coam}u$, $\cos\operatorname{am}.u$, $\cos\operatorname{coam}u$, $\Delta\operatorname{am}u$, $\Delta\operatorname{coam}u$, etc.\ and in general the \emph{trigonometric functions of the amplitude}, it is convenient in what follows to distinguish by the name of \emph{Elliptic Functions}; so that we assign to that name a certain other notion than has hitherto been done by Analysts. The quantity $u$ itself we shall call the \emph{Argument of the Elliptic Function}, so that on putting $x=\sin\operatorname{am}u$, $u=\operatorname{Arg}.\sin\operatorname{am}x$. From the proposed notation there will be:
\[
\begin{gathered}
\sin\operatorname{coam}u=\frac{\cos\operatorname{am}u}{\Delta\operatorname{am}u}\\
\cos\operatorname{coam}u=\frac{k'\sin\operatorname{am}u}{\Delta\operatorname{am}u}\\
\Delta\operatorname{coam}u=\frac{k'}{\Delta\operatorname{am}u}\\
\operatorname{tg}\operatorname{coam}u=\frac{1}{k'\operatorname{tg}\operatorname{am}u}\\
\operatorname{cotg}\operatorname{coam}u=\frac{k'}{\operatorname{cotg}\operatorname{am}u}.
\end{gathered}
\]

\origpage{32}

\subsection*{FUNDAMENTAL FORMULAE IN THE ANALYSIS OF ELLIPTIC FUNCTIONS.}

\subsection*{18.}

Let us put $\operatorname{am}.u=a$, $\operatorname{am}.v=b$, $\operatorname{am}(u+v)=\sigma$, $\operatorname{am}.(u-v)=\vartheta$; the fundamental formulae for the addition and subtraction of Elliptic Functions are known:
\[
\begin{gathered}
\sin\sigma=\frac{\sin a\cos b\,\Delta b+\sin b\cos a\,\Delta a}{1-k^{2}\sin a^{2}\sin b^{2}}\\[1ex]
\cos.\sigma=\frac{\cos a\cos b-\sin a\sin b\,\Delta a\,\Delta b}{1-k^{2}\sin a^{2}\sin b^{2}}\\[1ex]
\Delta\sigma=\frac{\Delta a\,\Delta b-k^{2}\sin a\sin b\cos a\cos b}{1-k^{2}\sin a^{2}\sin b^{2}}\\[1ex]
\sin\vartheta=\frac{\sin a\cos b\,\Delta b-\sin b\cos a\,\Delta a}{1-k^{2}\sin a^{2}\sin b^{2}}\\[1ex]
\cos\vartheta=\frac{\cos a\cos b+\sin a\sin b\,\Delta a\,\Delta b}{1-k^{2}\sin a^{2}\sin b^{2}}\\[1ex]
\Delta\vartheta=\frac{\Delta a\,\Delta b+k^{2}\sin a\sin b\cos a\cos b}{1-k^{2}\sin a^{2}\sin b^{2}}
\end{gathered}
\]

That everything of which there will be use hereafter may be at hand, let us note down in addition the following formulae, which are easily demonstrated, and whose number is easily increased:
\[
\begin{gathered}
\text{1)}\quad\sin\sigma+\sin\vartheta=\frac{2.\sin a\operatorname{Cos}b\,\Delta b}{1-k^{2}\sin a^{2}\sin b^{2}}\\[1ex]
\text{2)}\quad\cos\sigma+\cos\vartheta=\frac{2\cos a.\cos b}{1-k^{2}\sin a^{2}\sin b^{2}}\\[1ex]
\text{3)}\quad\Delta\sigma+\Delta\vartheta=\frac{2\Delta a.\Delta b}{1-k^{2}\sin a^{2}\sin b^{2}}\\[1ex]
\text{4)}\quad\sin\sigma-\sin\vartheta=\frac{2\sin b\cos a\,\Delta a}{1-k^{2}\sin a^{2}\sin b^{2}}\\[1ex]
\text{5)}\quad\cos.\vartheta-\cos\sigma=\frac{2\sin a\sin b\,\Delta a\,\Delta b}{1-k^{2}\sin a^{2}.\sin b^{2}}\\[1ex]
\text{6)}\quad\Delta\vartheta-\Delta\sigma=\frac{2k^{2}\sin a.\sin b\cos a.\cos b}{1-k^{2}\sin a^{2}\sin b^{2}}\\[1ex]
\text{7)}\quad\sin\sigma.\sin\vartheta=\frac{\sin a^{2}-\sin b^{2}}{1-k^{2}\sin a^{2}\sin b^{2}}\\[1ex]
\text{8)}\quad1+k^{2}\sin\sigma.\sin\vartheta=\frac{\Delta b^{2}+k^{2}\sin a^{2}.\cos b^{2}}{1-k^{2}.\sin a^{2}\sin b^{2}}\\[1ex]
\text{9)}\quad1+\sin\sigma.\sin\vartheta=\frac{\cos b^{2}+\sin a^{2}\Delta b^{2}}{1-k^{2}\sin a^{2}\sin b^{2}}
\end{gathered}
\]\ednote{In formula 1) the factor is printed as capital Cos b, where plain cos b is expected; an apparent misprint.}

\origpage{33}
\[
\begin{gathered}
\text{10)}\quad1+\cos\sigma.\cos\vartheta=\frac{\cos a^{2}+\cos b^{2}}{1-k^{2}\sin a^{2}\sin b^{2}}\\[1ex]
\text{11)}\quad1+\Delta\sigma.\Delta\vartheta=\frac{\Delta a^{2}+\Delta b^{2}}{1-k^{2}\sin a^{2}\sin b^{2}}\\[1ex]
\text{12)}\quad1-k^{2}\sin\sigma\sin\vartheta=\frac{\Delta a^{2}+k^{2}\sin b^{2}\cos a^{2}}{1-k^{2}\sin a^{2}\sin b^{2}}\\[1ex]
\text{13)}\quad1-\sin\sigma\sin\vartheta=\frac{\cos a^{2}+\sin b^{2}\Delta a^{2}}{1-k^{2}\sin a^{2}\sin b^{2}}\\[1ex]
\text{14)}\quad1-\cos\sigma\cos.\vartheta=\frac{\sin a^{2}\Delta b^{2}+\sin b^{2}\Delta a^{2}}{1-k^{2}\sin a^{2}\sin b^{2}}\\[1ex]
\text{15)}\quad1-\Delta\sigma\,\Delta\vartheta=\frac{k^{2}(\sin a^{2}\cos b^{2}+\sin b^{2}\cos a^{2})}{1-k^{2}\sin a^{2}\sin b^{2}}\\[1ex]
\text{16)}\quad(1\pm\sin\sigma)(1\pm\sin\vartheta)=\frac{(\cos b\pm\sin a\,\Delta b)^{2}}{1-k^{2}\sin a^{2}\sin b^{2}}\\[1ex]
\text{17)}\quad(1\pm\sin\sigma)(1\mp\sin\vartheta)=\frac{(\cos a\pm\sin b\,\Delta a)^{2}}{1-k^{2}\sin a^{2}\sin b^{2}}\\[1ex]
\text{18)}\quad(1\pm k\sin\sigma)(1\pm k\sin\vartheta)=\frac{(\Delta b\pm k\sin a\cos b)^{2}}{1-k^{2}\sin a^{2}\sin b^{2}}\\[1ex]
\text{19)}\quad(1\pm k\sin\sigma)(1\mp k\sin\vartheta)=\frac{(\Delta a\pm k\sin b\cos a)^{2}}{1-k^{2}\sin a^{2}\sin b^{2}}\\[1ex]
\text{20)}\quad(1\pm\cos\sigma)(1\pm\cos\vartheta)=\frac{(\cos a\pm\cos b)^{2}}{1-k^{2}\sin a^{2}\sin b^{2}}\\[1ex]
\text{21)}\quad(1\pm\cos\sigma)(1\mp\cos\vartheta)=\frac{(\sin a\,\Delta b\mp\sin b\,\Delta a)^{2}}{1-k^{2}\sin a^{2}\sin b^{2}}\\[1ex]
\text{22)}\quad(1\pm\Delta\sigma)(1\pm\Delta\vartheta)=\frac{(\Delta a+\Delta b)^{2}}{1-k^{2}\sin a^{2}\sin b^{2}}\\[1ex]
\text{23)}\quad(1\pm\Delta\sigma)(1\mp\Delta\vartheta)=\frac{k^{2}\sin^{2}(a\mp b)}{1-k^{2}\sin a^{2}\sin b^{2}}\\[1ex]
\text{24)}\quad\sin\sigma\cos\vartheta=\frac{\sin a\cos a\,\Delta b+\sin b\cos b\,\Delta a}{1-k^{2}\sin a^{2}\sin b^{2}}\\[1ex]
\text{25)}\quad\sin\vartheta\cos\sigma=\frac{\sin a\cos a\,\Delta b-\sin b\cos b\,\Delta a}{1-k^{2}\sin a^{2}\sin b^{2}}\\[1ex]
\text{26)}\quad\sin\sigma\,\Delta\vartheta=\frac{\cos b\sin a\,\Delta a+\cos a\sin b\,\Delta b}{1-k^{2}\sin a^{2}\sin b^{2}}\\[1ex]
\text{27)}\quad\sin\vartheta\,\Delta\sigma=\frac{\cos b\sin a\,\Delta a-\cos a\sin b\,\Delta b}{1-k^{2}\sin a^{2}\sin b^{2}}\\[1ex]
\text{28)}\quad\cos\sigma\,\Delta\vartheta=\frac{\cos a\cos b\,\Delta a\,\Delta b-k'k'\sin a\sin b}{1-k^{2}\sin a^{2}\sin b^{2}}\\[1ex]
\text{29)}\quad\cos\vartheta\,\Delta\sigma=\frac{\cos a\cos b\,\Delta a\,\Delta b+k'k'\sin a\sin b}{1-k^{2}\sin a^{2}\sin b^{2}}
\end{gathered}
\]\ednote{In formula 22) the numerator is printed $(\Delta a+\Delta b)^{2}$ with a plain plus sign; since the left side carries the double sign, $(\Delta a\pm\Delta b)^{2}$ is expected. An apparent misprint.}

\origpage{34}
\[
\begin{gathered}
\text{30)}\quad\sin(\sigma+\vartheta)=\frac{2\sin a\cos a\,\Delta b}{1-k^{2}\sin a^{2}\sin b^{2}}\\[1ex]
\text{31)}\quad\sin(\sigma-\vartheta)=\frac{2\sin b.\cos b\,\Delta a}{1-k^{2}\sin a^{2}\sin b^{2}}\\[1ex]
\text{32)}\quad\cos(\sigma+\vartheta)=\frac{\cos a^{2}-\sin a^{2}\Delta b^{2}}{1-k^{2}\sin a^{2}\sin b^{2}}\\[1ex]
\text{33)}\quad\cos(\sigma-\vartheta)=\frac{\cos b^{2}-\sin b^{2}\Delta a^{2}}{1-k^{2}\sin a^{2}\sin b^{2}}
\end{gathered}
\]

\subsection*{ON THE IMAGINARY VALUES OF ELLIPTIC FUNCTIONS. THE PRINCIPLE OF THE DOUBLE PERIOD.}

\subsection*{19.}

Let us put $\sin\phi=i\operatorname{tg}\psi$, where $i$ is put in place of $\sqrt{-1}$ according to the usage of most Geometers; then $\cos\phi=\sec\psi=\dfrac{1}{\cos\psi}$, whence $d\phi=\dfrac{i\,d\psi}{\cos\psi}$. Hence there results:
\[
\frac{d\phi}{\sqrt{1-k^{2}\sin\phi^{2}}}=\frac{i\,d\psi}{\sqrt{\cos\psi^{2}+k^{2}\sin\psi^{2}}}=\frac{i\,d\psi}{\sqrt{1-k'k'\sin\psi^{2}}}.
\]
We see that this passes over, in our notation, into the equation:
\[
\text{1)}\quad\sin\operatorname{am}iu=i\operatorname{tang}\operatorname{am}(u,k').
\]
Hence follows:
\[
\begin{gathered}
\text{2)}\quad\cos\operatorname{am}(iu,k)=\sec\operatorname{am}(u,k')\\
\text{3)}\quad\operatorname{tang}\operatorname{am}(iu,k)=i\sin\operatorname{am}(u,k')\\
\text{4)}\quad\Delta\operatorname{am}(iu,k)=\frac{\Delta\operatorname{am}(u,k')}{\operatorname{Cos}\operatorname{am}(u,k')}=\frac{1}{\sin\operatorname{coam}(u,k')}\\
\text{5)}\quad\sin\operatorname{coam}(iu,k)=\frac{1}{\Delta\operatorname{am}(u,k')}\\
\text{6)}\quad\cos\operatorname{coam}(iu,k)=i\frac{k'}{k}\cos\operatorname{coam}(u,k')\\
\text{7)}\quad\operatorname{tg}\operatorname{coam}(iu,k)=\frac{-i}{k'\sin\operatorname{am}(u,k')}\\
\text{8)}\quad\Delta\operatorname{coam}(iu,k)=k'\sin\operatorname{coam}(u,k').
\end{gathered}
\]\ednote{In formula 4) the denominator is printed as capital Cos am, where plain cos am is expected; an apparent misprint.}
Another system of formulae, which flows from this, is the following:
\[
\begin{gathered}
\text{9)}\quad\sin\operatorname{am}2iK'=0\\
\text{10)}\quad\sin\operatorname{am}iK'=\infty,\ \text{or if you prefer }\pm i\infty.
\end{gathered}
\]

\origpage{35}
\[
\begin{gathered}
\text{11)}\quad\sin\operatorname{am}(u+2iK')=+\sin\operatorname{am}u\\[1ex]
\text{12)}\quad\cos\operatorname{am}(u+2iK')=-\cos\operatorname{am}u\\[1ex]
\text{13)}\quad\Delta\operatorname{am}(u+2iK')=-\Delta\operatorname{am}u\\[1ex]
\text{14)}\quad\sin\operatorname{am}(u+iK')=\frac{1}{k\sin\operatorname{am}u}\\[1ex]
\text{15)}\quad\cos\operatorname{am}(u+iK')=\frac{-i\Delta\operatorname{am}u}{k\sin\operatorname{am}u}=\frac{-ik'}{k\cos\operatorname{coam}u}\\[1ex]
\text{16)}\quad\operatorname{tg}\operatorname{am}(u+iK')=\frac{+i}{\Delta\operatorname{am}u}\\[1ex]
\text{17)}\quad\Delta\operatorname{am}(u+iK')=-i\operatorname{cotg}\operatorname{am}u\\[1ex]
\text{18)}\quad\sin\operatorname{coam}(u+iK')=\frac{\Delta\operatorname{am}u}{k\cos\operatorname{am}u}=\frac{1}{k\sin\operatorname{coam}u}\\[1ex]
\text{19)}\quad\cos\operatorname{coam}(u+iK')=\frac{+k'i}{k\cos\operatorname{am}u}\\[1ex]
\text{20)}\quad\operatorname{tg}\operatorname{coam}(u+iK')=\frac{-i}{k}\Delta\operatorname{am}u\\[1ex]
\text{21)}\quad\Delta\operatorname{coam}(u+iK')=+ik\operatorname{tg}\operatorname{am}u.
\end{gathered}
\]\ednote{In formulae 20) and 21) the print has $k$ without the prime (the author's Corrigenda give $k'$); the faint primes visible in the scan are hand corrections in ink and are not transcribed.}
From the preceding formulae, which themselves also ought to be considered as fundamental in the Analysis of elliptic functions, it is evident:

\emph{a.} that elliptic functions of an imaginary argument $iv$, with Modulus $k$, can be transformed into others of a real argument $v$, with Modulus $k'=\sqrt{1-kk}$; whence in general the elliptic functions of an imaginary argument $u+iv$, with Modulus $k$, may be composed from elliptic functions of the argument $u$, with Modulus $k$, and others of the argument $v$, with Modulus $k'$;

\emph{b.} that elliptic functions enjoy a double period, one real, the other imaginary, provided the Modulus $k$ is real. Both become imaginary when the Modulus itself is imaginary. This we shall call the \emph{Principle of the Double Period}. Since it embraces the whole Analytic Periodicity that can be conceived, it is evident from it that elliptic functions ought not to be numbered among other transcendents, which enjoy certain elegances, perhaps more numerous or greater than these, but that there is in them a certain appearance of being perfect\ednote{Printed profecti, where perfecti is expected; an apparent misprint, corrected by hand in the scan.} and absolute.

\origpage{36}

\subsection*{ANALYTICAL THEORY OF THE TRANSFORMATION OF ELLIPTIC FUNCTIONS.}

\subsection*{20.}

We have seen in what precedes, whenever rational integral functions $A$, $B$, $C$, $D$, $U$, $V$ of the element $x$ are so determined that
\[
\begin{gathered}
V+U=(1+x)AA\\
V-U=(1-x)BB\\
V+\lambda U=(1+kx)CC\\
V-\lambda U=(1-kx)DD,
\end{gathered}
\]
then, on putting $y=\dfrac{U}{V}$, it will be that
\[
\frac{dy}{\sqrt{1-y^{2}}\sqrt{1-\lambda^{2}y^{2}}}=\frac{dx}{M\sqrt{1-x^{2}}\sqrt{1-k^{2}x^{2}}},
\]
where $M$ denotes a Constant quantity. Let us now set forth the general analytical expressions of those functions.

Let $n$ be any odd number, let $m$, $m'$ be any integer numbers, positive or negative, which however have no common factor that itself also divides the number $n$: let us put
\[
\omega=\frac{mK+m'iK'}{n},
\]
then there results:
\[
\begin{gathered}
U=\frac{x}{M}\left(1-\frac{xx}{\sin^{2}\operatorname{am}4\omega}\right)\left(1-\frac{xx}{\sin^{2}\operatorname{am}8\omega}\right)\ldots\left(1-\frac{xx}{\sin^{2}\operatorname{am}2(n-1)\omega}\right)\\[1.5ex]
V=\left(1-k^{2}\sin^{2}\operatorname{am}4\omega.xx\right)\left(1-k^{2}\sin^{2}\operatorname{am}8\omega.xx\right).\,.\left(1-k^{2}\sin^{2}\operatorname{am}2(n-1)\omega.xx\right)\\[1.5ex]
A=\left(1+\frac{x}{\sin\operatorname{coam}4\omega}\right)\left(1+\frac{x}{\sin\operatorname{coam}8\omega}\right).\,.\,.\,.\left(1+\frac{x}{\sin\operatorname{coam}2(n-1)\omega}\right)\\[1.5ex]
B=\left(1-\frac{x}{\sin\operatorname{coam}4\omega}\right)\left(1-\frac{x}{\sin\operatorname{coam}8\omega}\right)\ldots\left(1-\frac{x}{\sin\operatorname{coam}2(n-1)\omega}\right)\\[1.5ex]
C=\left(1+k\sin\operatorname{coam}4\omega.x\right)\left(1+k\sin\operatorname{coam}8\omega.x\right)\ldots\left(1+k\sin\operatorname{coam}2(n-1)\omega.x\right)\\[1.5ex]
D=\left(1-k\sin\operatorname{coam}4\omega.x\right)\left(1-k\sin\operatorname{coam}8\omega.x\right).\,.\,.\,.\left(1-k\sin\operatorname{coam}2(n-1)\omega.x\right)
\end{gathered}
\]

\origpage{37}
\[
\begin{gathered}
\lambda=k^{n}\left\{\sin\operatorname{coam}4\omega.\sin\operatorname{coam}8\omega.\,.\,.\,.\sin\operatorname{coam}2(n-1)\omega\right\}^{4}\\[1.5ex]
M=(-1)^{\frac{n-1}{2}}\frac{\left\{\sin\operatorname{coam}4\omega\sin\operatorname{coam}8\omega.\,.\,.\,.\sin\operatorname{coam}2(n-1)\omega\right\}^{2}}{\left\{\sin\operatorname{am}4\omega\sin\operatorname{am}8\omega.\,.\,.\,.\sin\operatorname{am}2(n-1)\omega\right\}^{2}}.
\end{gathered}
\]
With these put, where $x=\sin\operatorname{am}u$, there results $y=\frac{U}{V}=\sin\operatorname{am}\left(\frac{u}{M},\lambda\right)$.

Before we approach the demonstration of the formulae itself, we shall indicate a certain transformation of them. To this end we note down the following formulae, which run at once from the formulae of §.\ 18.:
\[
\begin{gathered}
\text{1)}\quad\sin\operatorname{am}(u+\alpha)\sin\operatorname{am}(u-\alpha)=\frac{\sin^{2}\operatorname{am}u-\sin^{2}\operatorname{am}\alpha}{1-k^{2}\sin^{2}\operatorname{am}u\sin^{2}\operatorname{am}\alpha}\\[1.5ex]
\text{2)}\quad\frac{\left(1+\sin\operatorname{am}(u+\alpha)\right)\left(1+\sin\operatorname{am}(u-\alpha)\right)}{\cos^{2}\operatorname{am}\alpha}=\frac{\left(1+\frac{\sin\operatorname{am}u}{\sin\operatorname{coam}\alpha}\right)^{2}}{1-k^{2}\sin^{2}\operatorname{am}u\sin^{2}\operatorname{am}\alpha}\\[1.5ex]
\text{3)}\quad\frac{\left(1-\sin\operatorname{am}(u+\alpha)\right)\left(1-\sin\operatorname{am}(u-\alpha)\right)}{\cos^{2}\operatorname{am}\alpha}=\frac{\left(1-\frac{\sin\operatorname{am}u}{\sin\operatorname{coam}\alpha}\right)^{2}}{1-k^{2}\sin^{2}\operatorname{am}u\sin^{2}\operatorname{am}\alpha}\\[1.5ex]
\text{4)}\quad\frac{\left(1+k\sin\operatorname{am}(u+\alpha)\right)\left(1+k\sin\operatorname{am}(u-\alpha)\right)}{\Delta^{2}\operatorname{am}\alpha}=\frac{\left(1+k\sin\operatorname{am}u\sin\operatorname{coam}\alpha\right)^{2}}{1-k^{2}\sin^{2}\operatorname{am}u\sin^{2}\operatorname{am}\alpha}\\[1.5ex]
\text{5)}\quad\frac{\left(1-k\sin\operatorname{am}(u+\alpha)\right)\left(1-k\sin\operatorname{am}(u-\alpha)\right)}{\Delta^{2}\operatorname{am}\alpha}=\frac{\left(1-k\sin\operatorname{am}u\sin\operatorname{coam}\alpha\right)^{2}}{1-k^{2}\sin^{2}\operatorname{am}u\sin^{2}\operatorname{am}\alpha}
\end{gathered}
\]
From these formulae there also follows:
\[
\begin{gathered}
\text{6)}\quad\frac{\cos\operatorname{am}(u+\alpha)\cos\operatorname{am}(u-\alpha)}{\cos^{2}\operatorname{am}\alpha}=\frac{1-\frac{\sin^{2}\operatorname{am}u}{\sin^{2}\operatorname{coam}\alpha}}{1-k^{2}\sin^{2}\operatorname{am}u\sin^{2}\operatorname{am}\alpha}\\[1.5ex]
\text{7)}\quad\frac{\Delta\operatorname{am}(u+\alpha)\Delta\operatorname{am}(u-\alpha)}{\Delta^{2}\operatorname{am}\alpha}=\frac{1-k^{2}\sin^{2}\operatorname{am}u\sin^{2}\operatorname{coam}\alpha}{1-k^{2}\sin^{2}\operatorname{am}u\sin^{2}\operatorname{am}\alpha}
\end{gathered}
\]
On putting $x=\sin\operatorname{am}u$, we obtain from formula 1):
\[
\frac{1-\frac{xx}{\sin^{2}\operatorname{am}\alpha}}{1-k^{2}\sin^{2}\operatorname{am}\alpha\,xx}=\frac{-\sin\operatorname{am}(u+\alpha)\sin\operatorname{am}(u-\alpha)}{\sin^{2}\operatorname{am}\alpha}
\]
from formulae 2), 3):
\[
\frac{\left(1\pm\frac{x}{\sin\operatorname{coam}\alpha}\right)^{2}}{1-k^{2}x^{2}\sin^{2}\operatorname{am}\alpha}=\frac{\left(1\pm\sin\operatorname{am}(u+\alpha)\right)\left(1\pm\sin\operatorname{am}(u-\alpha)\right)}{\cos^{2}\operatorname{am}\alpha};
\]

\origpage{38}
from formulae 4), 5):
\[
\frac{\left(1\pm kx\sin\operatorname{coam}\alpha\right)^{2}}{1-k^{2}x^{2}\sin^{2}\operatorname{am}\alpha}=\frac{\left(1\pm k\sin\operatorname{am}(u+\alpha)\right)\left(1\pm k\sin\operatorname{am}(u-\alpha)\right)}{\Delta^{2}\operatorname{am}\alpha}
\]
Hence, where in place of $\alpha$ there is put successively $4\omega$, $8\omega$, \ldots $2(n-1)\omega$, and in place of $-\alpha$ however $4n\omega-\alpha$, we obtain:
\[
\begin{gathered}
\text{8)}\quad\frac{U}{V}=\frac{\frac{x}{M}\left(1-\frac{xx}{\sin^{2}\operatorname{am}4\omega}\right)\left(1-\frac{xx}{\sin^{2}\operatorname{am}8\omega}\right)\ldots\left(1-\frac{xx}{\sin^{2}\operatorname{am}2(n-1)\omega}\right)}{\left(1-k^{2}x^{2}\sin^{2}\operatorname{am}4\omega\right)\left(1-k^{2}x^{2}\sin^{2}\operatorname{am}8\omega\right).\,.\,.\,.\left(1-k^{2}x^{2}\sin^{2}\operatorname{am}2(n-1)\omega\right)}\\[1.5ex]
=\frac{\sin\operatorname{am}u.\sin\operatorname{am}(u+4\omega)\sin\operatorname{am}(u+8\omega).\,.\,.\,.\sin\operatorname{am}(u+4(n-1)\omega)}{\left\{\sin\operatorname{coam}4\omega\sin\operatorname{coam}8\omega.\,.\,.\,.\sin\operatorname{coam}2(n-1)\omega\right\}^{2}}
\end{gathered}
\]
\[
\begin{gathered}
\text{9)}\quad\frac{(1+x)AA}{V}=\frac{(1+x)\left\{\left(1+\frac{x}{\sin\operatorname{coam}4\omega}\right)\left(1+\frac{x}{\sin\operatorname{coam}8\omega}\right).\,.\,.\,.\,.\,.\left(1+\frac{x}{\sin\operatorname{coam}2(n-1)\omega}\right)\right\}^{2}}{\left(1-k^{2}x^{2}\sin^{2}\operatorname{am}4\omega\right)\left(1-k^{2}x^{2}\sin^{2}\operatorname{am}8\omega\right).\,.\,.\,.\left(1-k^{2}x^{2}\sin^{2}\operatorname{am}2(n-1)\omega\right)}\\[1.5ex]
=\frac{\left(1+\sin\operatorname{am}u\right)\left(1+\sin\operatorname{am}(u+4\omega)\right)\left(1+\sin\operatorname{am}(u+8\omega)\right).\,.\,.\,.\left(1+\sin\operatorname{am}(u+4(n-1)\omega)\right)}{\left\{\cos\operatorname{am}4\omega.\cos\operatorname{am}8\omega.\,.\,.\,.\cos\operatorname{am}2(n-1)\omega\right\}^{2}}
\end{gathered}
\]
\[
\begin{gathered}
\text{10)}\quad\frac{(1-x)BB}{V}=\frac{(1-x)\left\{\left(1-\frac{x}{\sin\operatorname{coam}4\omega}\right)\left(1-\frac{x}{\sin\operatorname{coam}8\omega}\right).\,.\,.\,.\left(1-\frac{x}{\sin\operatorname{coam}2(n-1)\omega}\right)\right\}^{2}}{\left(1-k^{2}x^{2}\sin^{2}\operatorname{am}4\omega\right)\left(1-k^{2}x^{2}\sin^{2}\operatorname{am}8\omega\right).\,.\,.\,.\left(1-k^{2}x^{2}\sin^{2}\operatorname{am}2(n-1)\omega\right)}\\[1.5ex]
=\frac{\left(1-\sin\operatorname{am}u\right)\left(1-\sin\operatorname{am}(u+4\omega)\right)\left(1-\sin\operatorname{am}(u+8\omega)\right).\,.\,.\,.\left(1-\sin\operatorname{am}(u+4(n-1)\omega)\right)}{\left\{\cos\operatorname{am}4\omega.\cos\operatorname{am}8\omega.\,.\,.\,.\cos\operatorname{am}2(n-1)\omega\right\}^{2}}
\end{gathered}
\]
\[
\begin{gathered}
\text{11)}\quad\frac{(1+kx)CC}{V}=\frac{(1+kx)\left\{\left(1+kx\sin\operatorname{coam}4\omega\right)\left(1+kx\sin\operatorname{coam}8\omega\right).\,.\,.\,.\left(1+kx\sin\operatorname{coam}2(n-1)\omega\right)\right\}^{2}}{\left(1-k^{2}x^{2}\sin^{2}\operatorname{am}4\omega\right)\left(1-k^{2}x^{2}\sin^{2}\operatorname{am}8\omega\right).\,.\,.\,.\left(1-k^{2}x^{2}\sin^{2}\operatorname{am}2(n-1)\omega\right)}\\[1.5ex]
=\frac{\left(1+k\sin\operatorname{am}u\right)\left(1+k\sin\operatorname{am}(u+4\omega)\right)\left(1+k\sin\operatorname{am}(u+8\omega)\right)\ldots\left(1+k\sin\operatorname{am}(u+4(n-1)\omega)\right)}{\left\{\Delta\operatorname{am}4\omega\Delta\operatorname{am}8\omega\ldots\Delta\operatorname{am}2(n-1)\omega\right\}^{2}}
\end{gathered}
\]
\[
\begin{gathered}
\text{12)}\quad\frac{(1-kx)DD}{V}=\frac{(1-kx)\left\{\left(1-kx\sin\operatorname{coam}4\omega\right)\left(1-kx\sin\operatorname{coam}8\omega\right).\,.\,.\,.\left(1-kx\sin\operatorname{coam}2(n-1)\omega\right)\right\}^{2}}{\left(1-k^{2}x^{2}\sin^{2}\operatorname{am}4\omega\right)\left(1-k^{2}x^{2}\sin^{2}\operatorname{am}8\omega\right).\,.\,.\,.\left(1-k^{2}x^{2}\sin^{2}\operatorname{am}2(n-1)\omega\right)}\\[1.5ex]
=\frac{\left(1-k\sin\operatorname{am}u\right)\left(1-k\sin\operatorname{am}(u+4\omega)\right)\left(1-k\sin\operatorname{am}(u+8\omega)\right).\,.\,.\,.\left(1-k\sin\operatorname{am}(u+4(n-1)\omega)\right)}{\left\{\Delta\operatorname{am}4\omega\Delta\operatorname{am}8\omega\ldots\Delta\operatorname{am}2(n-1)\omega\right\}^{2}}
\end{gathered}
\]

\origpage{39}
Hence there also follow the formulae:
\[
\begin{gathered}
\text{13)}\quad\frac{\sqrt{1-xx}\,AB}{V}=\sqrt{1-xx}\,\frac{\left(1-\frac{xx}{\sin^{2}\operatorname{coam}4\omega}\right)\left(1-\frac{xx}{\sin^{2}\operatorname{coam}8\omega}\right).\,.\,.\,.\left(1-\frac{xx}{\sin^{2}\operatorname{coam}2(n-1)\omega}\right)}{\left(1-k^{2}x^{2}\sin^{2}\operatorname{am}4\omega\right)\left(1-k^{2}x^{2}\sin^{2}\operatorname{am}8\omega\right).\,.\,.\,.\left(1-k^{2}x^{2}\sin^{2}\operatorname{am}2(n-1)\omega\right)}\\[1.5ex]
=\frac{\cos\operatorname{am}u.\cos\operatorname{am}(u+4\omega)\cos\operatorname{am}(u+8\omega).\,.\,.\,.\cos\operatorname{am}(u+4(n-1)\omega)}{\left\{\cos\operatorname{am}4\omega.\cos\operatorname{am}8\omega.\,.\,.\,.\cos\operatorname{am}2(n-1)\omega\right\}^{2}}
\end{gathered}
\]
\[
\begin{gathered}
\text{14)}\quad\frac{\sqrt{1-k^{2}x^{2}}.CD}{V}=\sqrt{1-k^{2}x^{2}}.\frac{\left(1-k^{2}x^{2}\sin^{2}\operatorname{coam}4\omega\right)\left(1-k^{2}x^{2}\sin^{2}\operatorname{coam}8\omega\right).\,.\,.\,.\left(1-k^{2}x^{2}\sin^{2}\operatorname{coam}2(n-1)\omega\right)}{\left(1-k^{2}x^{2}\sin^{2}\operatorname{am}4\omega\right)\left(1-k^{2}x^{2}\sin^{2}\operatorname{am}8\omega\right).\,.\,.\,.\left(1-k^{2}x^{2}\sin^{2}\operatorname{am}2(n-1)\omega\right)}\\[1.5ex]
=\frac{\Delta\operatorname{am}u\Delta\operatorname{am}(u+4\omega)\Delta\operatorname{am}(u+8\omega).\,.\,.\,.\Delta\operatorname{am}(u+4(n-1)\omega)}{\left\{\Delta\operatorname{am}4\omega\Delta\operatorname{am}8\omega.\,.\,.\,.\Delta\operatorname{am}2(n-1)\omega\right\}^{2}}.
\end{gathered}
\]

\subsection*{DEMONSTRATION OF THE ANALYTICAL FORMULAE FOR THE TRANSFORMATION.}

\subsection*{21.}

Let us now demonstrate, on putting:
\[
\begin{gathered}
1-y=(1-x).\frac{\left\{\left(1-\frac{x}{\sin\operatorname{coam}4\omega}\right)\left(1-\frac{x}{\sin\operatorname{coam}8\omega}\right).\,.\,.\,.\left(1-\frac{x}{\sin\operatorname{coam}2(n-1)\omega}\right)\right\}^{2}}{\left(1-k^{2}x^{2}\sin^{2}\operatorname{am}4\omega\right)\left(1-k^{2}x^{2}\sin^{2}\operatorname{am}8\omega\right)\ldots\left(1-k^{2}x^{2}\sin^{2}\operatorname{am}2(n-1)\omega\right)}\\[1.5ex]
=\frac{\left(1-\sin\operatorname{am}u\right)\left(1-\sin\operatorname{am}(u+4\omega)\right)\left(1-\sin\operatorname{am}(u+8\omega)\right)\ldots\left(1-\sin\operatorname{am}(u+4(n-1)\omega)\right)}{\left\{\cos\operatorname{am}4\omega.\cos\operatorname{am}8\omega.\cos\operatorname{am}12\omega.\,.\,.\,.\cos\operatorname{am}2(n-1)\omega\right\}^{2}},
\end{gathered}
\]
that the remaining formulae can be derived, and this one:
\[
\frac{dy}{\sqrt{1-y^{2}}\sqrt{1-\lambda^{2}y^{2}}}=\frac{dx}{M\sqrt{1-x^{2}}\sqrt{1-k^{2}x^{2}}},
\]
provided that:
\[
\begin{gathered}
\lambda=k^{n}\left\{\sin\operatorname{coam}4\omega.\sin\operatorname{coam}8\omega.\,.\,.\,.\sin\operatorname{coam}2(n-1)\omega\right\}^{4}\\[1.5ex]
M=\frac{\left\{\sin\operatorname{coam}4\omega.\sin\operatorname{coam}8\omega.\,.\,.\,.\sin\operatorname{coam}2(n-1)\omega\right\}^{2}}{\left\{\sin\operatorname{am}4\omega.\sin\operatorname{am}8\omega.\,.\,.\,.\sin\operatorname{am}2(n-1)\omega\right\}^{2}}.
\end{gathered}
\]\ednote{The factor $(-1)^{\frac{n-1}{2}}$ before the fraction for $M$ is not printed here (the same formula for $M$ in the preceding article carries it); the copy supplies it by a faint hand-written addition in ink, which is not transcribed.}
From the proposed formula it appears that $y$ is not changed at all whenever $u$ passes over into $u+4\omega$. For then every factor passes over into the following one, but the last into the first. Whence in general $y$ is not changed,

\origpage{40}
provided that in place of $u$ there is put $u+4p\omega$, where $p$ denotes a positive or negative integer number. But where $u=0$, there results:
\[
1-y=\frac{\left(1-\sin\operatorname{am}4\omega\right)\left(1-\sin\operatorname{am}8\omega\right)\ldots\left(1-\sin\operatorname{am}4(n-1)\omega\right)}{\left\{\cos\operatorname{am}4\omega.\cos\operatorname{am}8\omega\ldots\cos\operatorname{am}2(n-1)\omega\right\}^{2}}=1,
\]
or $y=0$. For it is easily evident that it will be:
\[
\begin{gathered}
-\sin\operatorname{am}4(n-1)\omega=+\sin\operatorname{am}4\omega\\
-\sin\operatorname{am}4(n-2)\omega=+\sin\operatorname{am}8\omega,\\
\ldots
\end{gathered}
\]
whence:
\[
\begin{gathered}
\left(1-\sin\operatorname{am}4\omega\right)\left(1-\sin\operatorname{am}4(n-1)\omega\right)=\cos^{2}\operatorname{am}4\omega\\
\left(1-\sin\operatorname{am}8\omega\right)\left(1-\sin\operatorname{am}4(n-2)\omega\right)=\cos^{2}\operatorname{am}8\omega\\
\ldots\\
\left(1-\sin\operatorname{am}2(n-1)\omega\right)\left(1-\sin\operatorname{am}2(n+1)\omega\right)=\cos^{2}\operatorname{am}2(n-1)\omega.
\end{gathered}
\]
Now since $y=0$ whenever $u=0$, and $y$ is not changed where in place of $u$ there is put $u+4p\omega$, in general $y$ vanishes whenever $u$ takes on the values:
\[
0,\ 4\omega,\ 8\omega,\ .\,.\,.\,.\,.,\ 4(n-2)\omega,\ 4(n-1)\omega,
\]
to which correspond the values of the quantity $x=\sin\operatorname{am}\omega$\ednote{The print has $\omega$ here, where $u$ is evidently meant ($x=\sin\operatorname{am}u$); kept as printed.}:
\[
0,\ \sin\operatorname{am}4\omega,\ \sin\operatorname{am}8\omega,\ \ldots\sin\operatorname{am}4(n-2)\omega,\ \sin\operatorname{am}4(n-1)\omega,
\]
which it is permitted to exhibit also thus:
\[
0,\ \pm\sin\operatorname{am}4\omega,\ \pm\sin\operatorname{am}8\omega,\ \ldots,\ \pm\sin\operatorname{am}2(n-1)\omega,
\]
or even in this manner:
\[
0,\ \pm\sin\operatorname{am}2\omega,\ \pm\sin\operatorname{am}4\omega,\ \ldots\pm\sin\operatorname{am}(n-1)\omega.
\]
These values of the element $x$, which it can assume when $y$ vanishes, will all be different from one another, and their number will be $=n$. Now from the supposed equation between $x$ and $y$, from which we set out, it is evident, on putting:
\[
\begin{gathered}
V=\left(1-k^{2}x^{2}\sin^{2}\operatorname{am}4\omega\right)\left(1-k^{2}x^{2}\sin^{2}\operatorname{am}8\omega\right).\,.\,.\,.\left(1-k^{2}x^{2}\sin^{2}\operatorname{am}2(n-1)\omega\right)\\[1.5ex]
=\left(1-k^{2}x^{2}\sin^{2}\operatorname{am}2\omega\right)\left(1-k^{2}x^{2}\sin^{2}\operatorname{am}4\omega\right).\,.\,.\,.\left(1-k^{2}x^{2}\sin^{2}\operatorname{am}(n-1)\omega\right),
\end{gathered}
\]
$y=\frac{U}{V}$, that $U$ becomes a rational integral function of the element $x$ of the $n^{\mathrm{ti}}$ order. Since this vanishes together with $y$ for the following values of the quantity $x$, $n$ in number and different from one another:
\[
0,\ \pm\sin\operatorname{am}2\omega,\ \pm\sin\operatorname{am}4\omega,\ \ldots\pm\sin\operatorname{am}(n-1)\omega,
\]

\origpage{41}
it necessarily assumes the form:
\[
\begin{gathered}
U=\frac{x}{M}\left(1-\frac{xx}{\sin^{2}\operatorname{am}2\omega}\right)\left(1-\frac{xx}{\sin^{2}\operatorname{am}4\omega}\right).\,.\,.\,.\left(1-\frac{xx}{\sin^{2}\operatorname{am}(n-1)\omega}\right)\\[1.5ex]
=\frac{x}{M}\left(1-\frac{xx}{\sin^{2}\operatorname{am}4\omega}\right)\left(1-\frac{xx}{\sin^{2}\operatorname{am}8\omega}\right).\,.\,.\,.\left(1-\frac{xx}{\sin^{2}\operatorname{am}2(n-1)\omega}\right),
\end{gathered}
\]
where $M$ denotes a Constant. Since on putting $x=1$ there results $1-y=0$, $y=1$, we obtain from the equation $y=\frac{U}{V}$:
\[
\begin{gathered}
1=\frac{\left(1-\frac{1}{\sin^{2}\operatorname{am}2\omega}\right)\left(1-\frac{1}{\sin^{2}\operatorname{am}4\omega}\right)\ldots\left(1-\frac{1}{\sin^{2}\operatorname{am}(n-1)\omega}\right)}{M\left(1-k^{2}\sin^{2}\operatorname{am}2\omega\right)\left(1-k^{2}\sin^{2}\operatorname{am}4\omega\right).\,.\,.\,.\left(1-k^{2}\sin^{2}\operatorname{am}(n-1)\omega\right)}\\[1.5ex]
=\frac{(-1)^{\frac{n-1}{2}}\left\{\sin\operatorname{coam}2\omega.\sin\operatorname{coam}4\omega.\,.\,.\,.\sin\operatorname{coam}(n-1)\omega\right\}^{2}}{M\left\{\sin\operatorname{am}2\omega.\sin\operatorname{am}4\omega.\,.\,.\,.\sin\operatorname{am}(n-1)\omega\right\}^{2}}
\end{gathered}
\]
whence:
\[
M=\frac{(-1)^{\frac{n-1}{2}}\left\{\sin\operatorname{coam}2\omega.\sin\operatorname{coam}4\omega.\,.\,.\,.\sin\operatorname{coam}(n-1)\omega\right\}^{2}}{\left\{\sin\operatorname{am}2\omega.\sin\operatorname{am}4\omega.\,.\,.\,.\sin\operatorname{am}(n-1)\omega\right\}^{2}}.
\]

Between the functions $U$, $V$ there intervenes a memorable correlation, that I mean which was mentioned above, by whose benefit it comes about that, on putting $\frac{1}{kx}$ in place of $x$, $y$ at the same time passes over into $\frac{1}{\lambda y}$, where $\lambda$ denotes a Constant.

For on putting $\frac{1}{kx}$ in place of $x$ there passes over:
\[
U=\frac{x}{M}\left(1-\frac{xx}{\sin^{2}\operatorname{am}2\omega}\right)\left(1-\frac{xx}{\sin^{2}\operatorname{am}4\omega}\right).\,.\,.\,.\left(1-\frac{xx}{\sin^{2}\operatorname{am}(n-1)\omega}\right)
\]
into this expression:
\[
(-1)^{\frac{n-1}{2}}\frac{V}{Mx^{n}}.\frac{1}{k^{n}\left(\sin\operatorname{am}2\omega.\sin\operatorname{am}4\omega.\,.\,.\,.\sin\operatorname{am}(n-1)\omega\right)^{2}}.
\]
But on the other hand, the same substitution being made,
\[
V=\left(1-k^{2}x^{2}\sin^{2}\operatorname{am}2\omega\right)\left(1-k^{2}x^{2}\sin^{2}\operatorname{am}4\omega\right).\,.\,.\,.\left(1-k^{2}x^{2}\sin^{2}\operatorname{am}(n-1)\omega\right)
\]
passes over into this expression:
\[
(-1)^{\frac{n-1}{2}}\frac{U}{x^{n}}.M\left\{\sin\operatorname{am}2\omega.\sin\operatorname{am}4\omega.\,.\,.\,.\sin\operatorname{am}(n-1)\omega\right\}^{2}.
\]

\origpage{42}
Whence, with $\frac{1}{kx}$ put in place of $x$, $y=\frac{U}{V}$ passes over into:
\[
\frac{V}{U}.\frac{1}{MM.k^{n}\left\{\sin\operatorname{am}2\omega.\sin\operatorname{am}4\omega\ldots\sin\operatorname{am}(n-1)\omega\right\}^{4}},
\]
or $y$ into $\frac{1}{\lambda y}$, provided that there is put:
\[
\begin{gathered}
\lambda=MMk^{n}\left\{\sin\operatorname{am}2\omega.\sin\operatorname{am}4\omega.\,.\,.\,.\sin\operatorname{am}(n-1)\omega\right\}^{4}\\[1.5ex]
=k^{n}\left\{\sin\operatorname{coam}2\omega.\sin\operatorname{coam}4\omega.\,.\,.\,.\sin\operatorname{coam}(n-1)\omega\right\}^{4}
\end{gathered}
\]
This is what was to be demonstrated.

From the proposed equation:
\[
1-y=(1-x)\frac{\left\{\left(1-\frac{x}{\sin\operatorname{coam}4\omega}\right)\left(1-\frac{x}{\sin\operatorname{coam}8\omega}\right).\,.\,.\,.\left(1-\frac{x}{\sin\operatorname{coam}2(n-1)\omega}\right)\right\}^{2}}{\left(1-k^{2}x^{2}\sin^{2}\operatorname{am}4\omega\right)\left(1-k^{2}x^{2}\sin^{2}\operatorname{am}8\omega\right).\,.\,.\,.\left(1-k^{2}x^{2}\sin^{2}\operatorname{am}2(n-1)\omega\right)},
\]
on putting $\frac{1}{kx}$ in place of $x$, $\frac{1}{\lambda y}$ in place of $y$, which is allowed by what precedes, we derive:
\[
\frac{1}{\lambda y}-1=\frac{1-kx}{\lambda U}\left\{\left(1-kx\sin\operatorname{coam}4\omega\right)\left(1-kx\sin\operatorname{coam}8\omega\right).\,.\,.\,.\left(1-kx\sin\operatorname{coam}2(n-1)\omega\right)\right\}^{2},
\]
which multiplied by $\lambda y=\frac{\lambda U}{V}$, yields:
\[
1-\lambda y=(1-kx)\frac{\left\{\left(1-kx\sin\operatorname{coam}4\omega\right)\left(1-kx\sin\operatorname{coam}8\omega\right)\ldots\left(1-kx\sin\operatorname{coam}2(n-1)\omega\right)\right\}^{2}}{V}
\]
Moreover it is evident that $y=\frac{U}{V}$ passes over into $-y$ where $x$ is changed into $-x$; this being done, we therefore obtain at once also $1+y$, $1+\lambda y$ from $1-y$, $1-\lambda y$.

Now therefore we have found such rational integral functions $U$, $V$ of the element $x$ that it is:
\[
\begin{gathered}
V+U=V(1+y)=(1+x)AA\\
V-U=V(1-y)=(1-x)BB\\
V+\lambda U=V(1+\lambda y)=(1+kx)CC\\
V-\lambda U=V(1-\lambda y)=(1-kx)DD,
\end{gathered}
\]
where $A$, $B$, $C$, $D$ themselves denote rational integral functions of the element $x$. Hence, however, according to the Principles of Transformation established at the beginning, there follows at once:
\[
\frac{dy}{\sqrt{1-y^{2}}\sqrt{1-\lambda^{2}y^{2}}}=\frac{dx}{M\sqrt{1-x^{2}}\sqrt{1-k^{2}x^{2}}}.
\]

\origpage{43}
The Multiplier $M$, as we shall call it, we obtain from the observation made in §.\ 15. Whence all the general analytical formulae that concern the theory of the transformation of elliptic functions have now been demonstrated.

\subsection*{22.}

The proposed demonstration is derived from the one which we gave in the New Astronomical Notices published by Cl.\ Schumacher, No.\ 127, where $\omega$ is put in place of $\frac{K}{n}$, all else remaining unchanged. The general analytical theorem on Transformation itself, in a somewhat different form, I had already communicated earlier to analysts in the same journal, No.\ 123. Cl.\ Legendre, the foremost judge in this doctrine, kindly and splendidly wished to review the demonstration in the same place, No.\ 130. The man, venerable on many counts, observes there that the equation:
\[
V\frac{dU}{dx}-U\frac{dV}{dx}=\frac{ABCD}{M}=\frac{T}{M},
\]
by whose aid the demonstration is completed, and which followed for us from purely algebraic principles of transformation, can be proved analytically even without them. Since, in the opinion of the distinguished man himself, this casts excellent light on our theorem, let us, following his lead, demonstrate it briefly in the following manner.

The proposed equation:
\[
V\frac{dU}{dx}-U\frac{dV}{dx}=\frac{ABCD}{M}=\frac{T}{M}
\]
can also be exhibited thus:
\[
\frac{dU}{U\,dx}-\frac{dV}{V\,dx}=\frac{d\log U}{dx}-\frac{d\log V}{dx}=\frac{ABCD}{MUV}=\frac{T}{MUV}.
\]
But we have found:
\[
\begin{gathered}
U=\frac{x}{M}\left(1-\frac{xx}{\sin^{2}\operatorname{am}2\omega}\right)\left(1-\frac{xx}{\sin^{2}\operatorname{am}4\omega}\right).\,.\,.\,.\left(1-\frac{xx}{\sin^{2}\operatorname{am}(n-1)\omega}\right)\\[1.5ex]
V=\left(1-k^{2}x^{2}\sin^{2}\operatorname{am}2\omega\right)\left(1-k^{2}x^{2}\sin^{2}\operatorname{am}4\omega\right)\ldots\left(1-k^{2}x^{2}\sin^{2}\operatorname{am}(n-1)\omega\right),
\end{gathered}
\]
whence:
\[
\frac{d\log U}{dx}-\frac{d\log V}{dx}=\frac{1}{x}+\Sigma\left\{\frac{-2x}{\sin^{2}\operatorname{am}2q\omega-xx}+\frac{2k^{2}x\sin^{2}\operatorname{am}2q\omega}{1-k^{2}x^{2}\sin^{2}\operatorname{am}2q\omega}\right\},
\]

\origpage{44}
the number $q$ designated in the sum being assigned the values $1,2,3,\ldots,\frac{n-1}{2}$. Further we have found:
\[
\begin{gathered}
AB=\left(1-\frac{xx}{\sin^{2}\operatorname{coam}2\omega}\right)\left(1-\frac{xx}{\sin^{2}\operatorname{coam}4\omega}\right)\ldots\left(1-\frac{xx}{\sin^{2}\operatorname{coam}(n-1)\omega}\right)\\[1.5ex]
CD=\left(1-k^{2}x^{2}\sin^{2}\operatorname{coam}2\omega\right)\left(1-k^{2}x^{2}\sin^{2}\operatorname{coam}4\omega\right).\,.\,.\,.\left(1-k^{2}x^{2}\sin^{2}\operatorname{coam}(n-1)\omega\right),
\end{gathered}
\]
whence:
\[
\frac{T}{MUV}=\frac{ABCD}{MUV}=\frac{x\,\Pi\left(1-\frac{xx}{\sin^{2}\operatorname{coam}2p\omega}\right)\left(1-k^{2}x^{2}\sin^{2}\operatorname{coam}2p\omega\right)}{x^{2}\,\Pi\left(1-\frac{xx}{\sin^{2}\operatorname{am}2p\omega}\right)\left(1-k^{2}x^{2}\sin^{2}\operatorname{am}2p\omega\right)},
\]
provided that in the products, denoted for brevity by the prefixed sign $\Pi$, the element $p$ is assigned the values $1,2,3,.\,.\,.\,.,\frac{n-1}{2}$. This expression can be resolved into simple fractions, so that it takes the form:
\[
\frac{1}{x}+\Sigma\left\{\frac{A^{(q)}x}{\sin^{2}\operatorname{am}2q\omega-xx}+\frac{B^{(q)}x}{1-k^{2}x^{2}\sin^{2}\operatorname{am}2q\omega}\right\},
\]
and having done this, in order that we may have established what was proposed, it must be demonstrated that:
\[
A^{(q)}=-2,\quad B^{(q)}=2k^{2}\sin^{2}\operatorname{am}2q\omega.
\]

In what follows we shall denote by the prefixed sign $\Pi^{(q)}$ a product so formed that the element $p$ is assigned the values $1,2,3,\ldots,\frac{n-1}{2}$, the value $p=q$ however being omitted. Hence, from the rules of the theory of simple fractions, abundantly known, it follows that:
\[
A^{(q)}=\left(1-k^{2}\sin^{2}\operatorname{am}2q\omega.\sin^{2}\operatorname{coam}2q\omega\right)\frac{\Pi\left(\frac{1-\frac{\sin^{2}\operatorname{am}2q\omega}{\sin^{2}\operatorname{coam}2p\omega}}{1-k^{2}\sin^{2}\operatorname{am}2q\omega.\sin^{2}\operatorname{am}2p\omega}\right)}{\Pi^{(q)}\left(\frac{1-\frac{\sin^{2}\operatorname{am}2q\omega}{\sin^{2}\operatorname{am}2p\omega}}{1-k^{2}\sin^{2}\operatorname{am}2q\omega.\sin^{2}\operatorname{coam}2p\omega}\right)}.
\]

Now from the formulae exhibited by us above there results:
\[
\begin{gathered}
\frac{1-\frac{\sin^{2}\operatorname{am}2q\omega}{\sin^{2}\operatorname{coam}2p\omega}}{1-k^{2}\sin^{2}\operatorname{am}2q\omega\sin^{2}\operatorname{am}2p\omega}=\frac{\cos\operatorname{am}(2q+2p)\omega.\cos\operatorname{am}(2q-2p)\omega}{\cos^{2}\operatorname{am}2p\omega}\\[1.5ex]
\frac{1-\frac{\sin^{2}\operatorname{am}2q\omega}{\sin^{2}\operatorname{am}2p\omega}}{1-k^{2}\sin^{2}\operatorname{am}2q\omega.\sin^{2}\operatorname{coam}2p\omega}=\frac{\cos\operatorname{coam}(2p+2q)\omega.\cos\operatorname{coam}(2p-2q)\omega}{\cos^{2}\operatorname{coam}2p\omega}.
\end{gathered}
\]

\origpage{45}
But it is easily evident, the factors that are found the same in the denominator and numerator having been removed, that there results:
\[
\begin{gathered}
\Pi\frac{\cos\operatorname{am}(2q+2p)\omega\cos\operatorname{am}(2q-2p)\omega}{\cos^{2}\operatorname{am}2p\omega}=\frac{\pm1}{\cos\operatorname{am}2q\omega}\\[1.5ex]
\Pi^{(q)}\frac{\cos\operatorname{coam}(2q+2p)\omega\cos\operatorname{coam}(2p-2q)\omega}{\cos^{2}\operatorname{coam}2p\omega}=\frac{\mp1}{\cos\operatorname{coam}2q\omega}.\frac{\cos^{2}\operatorname{coam}2q\omega}{\cos\operatorname{coam}4q\omega}=\frac{\mp\cos\operatorname{coam}2q\omega}{\cos\operatorname{coam}4q\omega},
\end{gathered}
\]
whence:
\[
A^{(q)}=\frac{-\left(1-k^{2}\sin^{2}\operatorname{am}2q\omega\sin^{2}\operatorname{coam}2q\omega\right)\cos\operatorname{coam}4q\omega}{\cos\operatorname{am}2q\omega\cos\operatorname{coam}2q\omega}.
\]
But by the known formula for duplication there results:
\[
\begin{gathered}
\cos\operatorname{coam}4q\omega=\frac{2k'\sin\operatorname{am}2q\omega\cos\operatorname{am}2q\omega\Delta\operatorname{am}2q\omega}{1-2k^{2}\sin^{2}\operatorname{am}2q\omega+k^{2}\sin^{4}\operatorname{am}2q\omega}\\[1.5ex]
=\frac{2k'\sin\operatorname{am}2q\omega\cos\operatorname{am}2q\omega\Delta\operatorname{am}2q\omega}{\Delta^{2}\operatorname{am}2q\omega-k^{2}\sin^{2}\operatorname{am}2q\omega\cos^{2}\operatorname{am}2q\omega}\\[1.5ex]
=\frac{2\cos\operatorname{am}2q\omega\cos\operatorname{coam}2q\omega}{1-k^{2}\sin^{2}\operatorname{am}2q\omega\sin^{2}\operatorname{coam}2q\omega},
\end{gathered}
\]
whence at length, which was to be demonstrated, $A^{(q)}=-2$. In an entirely similar manner the other equation: $B^{(q)}=2k^{2}\sin^{2}\operatorname{am}2q\omega$ can be proved; this however, $A^{(q)}=-2$ having now been found, is done more easily as follows.

It is easily evident that, when $\frac{1}{kx}$ is put in place of $x$, the expression is not changed:
\[
\Pi\frac{\left(1-\frac{x^{2}}{\sin^{2}\operatorname{coam}2p\omega}\right)\left(1-k^{2}x^{2}\sin^{2}\operatorname{coam}2p\omega\right)}{\left(1-k^{2}x^{2}\sin^{2}\operatorname{am}2p\omega\right)\left(1-\frac{x^{2}}{\sin^{2}\operatorname{am}2p\omega}\right)},
\]
which we saw can be set equal to the expression:
\[
1+\Sigma\frac{-2x^{2}}{\sin^{2}\operatorname{am}2q\omega-x^{2}}+\Sigma\frac{B^{(q)}x^{2}}{1-k^{2}\sin^{2}\operatorname{am}2q\omega\,x^{2}}.
\]
But this expression, $\frac{1}{kx}$ being put in place of $x$, passes into this:
\[
\begin{gathered}
1+\Sigma\frac{2}{1-k^{2}x^{2}\sin^{2}\operatorname{am}2q\omega}+\Sigma\frac{-B^{(q)}}{k^{2}\left(\sin^{2}\operatorname{am}2q\omega-x^{2}\right)}=\\[1.5ex]
1+\Sigma\left(2-\frac{B^{(q)}}{k^{2}\sin^{2}\operatorname{am}2q\omega}\right)+\Sigma\frac{2k^{2}x^{2}\sin^{2}\operatorname{am}2q\omega}{1-k^{2}x^{2}\sin^{2}\operatorname{am}2q\omega}+\Sigma\frac{-B^{(q)}}{k^{2}\sin^{2}\operatorname{am}2q\omega}.\,\frac{x^{2}}{\sin^{2}\operatorname{am}2q\omega-x^{2}},
\end{gathered}
\]
whence, in order that it may remain unchanged, as it must, it is necessary that there result:
\[
B^{(q)}=2k^{2}\sin^{2}\operatorname{am}2q\omega.
\]
Q.\ E.\ D.

\origpage{46}
\subsection*{23.}

From formula 14) of §.\ 20 it follows:
\[
\sqrt{1-\lambda^{2}y^{2}}=\sqrt{1-k^{2}x^{2}}\,\frac{CD}{V}=\sqrt{1-k^{2}x^{2}}\,\frac{\left(1-k^{2}x^{2}\sin^{2}\operatorname{coam}2\omega\right)\left(1-k^{2}x^{2}\sin^{2}\operatorname{coam}4\omega\right).\,.\,.\,.\left(1-k^{2}x^{2}\sin^{2}\operatorname{coam}(n-1)\omega\right)}{\left(1-k^{2}x^{2}\sin^{2}\operatorname{am}2\omega\right)\left(1-k^{2}x^{2}\sin^{2}\operatorname{am}4\omega\right).\,.\,.\,.\left(1-k^{2}x^{2}\sin^{2}\operatorname{am}(n-1)\omega\right)}.
\]
Putting $x=1$, whence also $y=1$, and $\sqrt{1-\lambda\lambda}=\lambda'$, there results:
\[
\lambda'=k'\left\{\frac{\Delta\operatorname{coam}2\omega\Delta\operatorname{coam}4\omega\ldots\Delta\operatorname{coam}(n-1)\omega}{\Delta\operatorname{am}2\omega\Delta\operatorname{am}4\omega\ldots\Delta\operatorname{am}(n-1)\omega}\right\}^{2}
\]
Now however we have:
\[
\Delta\operatorname{coam}u=\frac{k'}{\Delta\operatorname{am}u},
\]
whence:
\[
\text{1)}\quad\lambda'=\frac{k'^{n}}{\left\{\Delta\operatorname{am}2\omega.\Delta\operatorname{am}4\omega.\,.\,.\,.\Delta\operatorname{am}(n-1)\omega\right\}^{4}}.
\]
Further, the formulae being brought into use:
\[
\begin{gathered}
\text{2)}\quad\lambda=k^{n}\left\{\sin\operatorname{coam}2\omega.\sin\operatorname{coam}4\omega.\,.\,.\,.\sin\operatorname{coam}(n-1)\omega\right\}^{4}\\[1.5ex]
\text{3)}\quad M=(-1)^{\frac{n-1}{2}}\frac{\left\{\sin\operatorname{coam}2\omega.\sin\operatorname{coam}4\omega.\,.\,.\,.\sin\operatorname{coam}(n-1)\omega\right\}^{2}}{\left\{\sin\operatorname{am}2\omega.\sin\operatorname{am}4\omega.\,.\,.\,.\sin\operatorname{am}(n-1\,\omega^{2}\right\}}
\end{gathered}
\]\ednote{In the denominator of formula 3) the closing parenthesis after $n-1$ is missing and the exponent 2, which belongs after the closing brace, is printed on $\omega$ instead; an apparent misprint for $(n-1)\omega$ with the exponent 2 on the brace.}
we obtain:
\[
\begin{gathered}
\text{4)}\quad\frac{(-1)^{\frac{n-1}{2}}}{M}\sqrt{\frac{\lambda}{k^{n}}}=\left\{\sin\operatorname{am}2\omega.\sin\operatorname{am}4\omega.\,.\,.\,.\sin\operatorname{am}(n-1)\omega\right\}^{2}\\[1.5ex]
\text{5)}\quad\sqrt{\frac{\lambda k'^{n}}{\lambda'k^{n}}}=\left\{\cos\operatorname{am}2\omega\cos\operatorname{am}4\omega.\,.\,.\,.\cos\operatorname{am}(n-1)\omega\right\}^{2}\\[1.5ex]
\text{6)}\quad\sqrt{\frac{k'^{n}}{\lambda'}}=\left\{\Delta\operatorname{am}2\omega\Delta\operatorname{am}4\omega\ldots\Delta\operatorname{am}(n-1)\omega\right\}^{2}\\[1.5ex]
\text{7)}\quad\frac{(-1)^{\frac{n-1}{2}}}{M}\sqrt{\frac{\lambda'}{k'^{n}}}=\left\{\operatorname{tg}\operatorname{am}2\omega.\operatorname{tg}\operatorname{am}4\omega\ldots\operatorname{tg}\operatorname{am}(n-1)\omega\right\}^{2}\\[1.5ex]
\text{8)}\quad\sqrt{\frac{\lambda}{k^{n}}}=\left\{\sin\operatorname{coam}2\omega.\sin\operatorname{coam}4\omega\ldots\sin\operatorname{coam}(n-1)\omega\right\}^{2}\\[1.5ex]
\text{9)}\quad\frac{(-1)^{\frac{n-1}{2}}}{M}\sqrt{\frac{\lambda\lambda'k'^{n}}{k'k'k^{n}}}=\left\{\cos\operatorname{coam}2\omega\cos\operatorname{coam}4\omega.\,.\,.\,.\cos\operatorname{coam}(n-1)\omega\right\}^{2}
\end{gathered}
\]

\origpage{47}
\[
\begin{gathered}
\text{10)}\quad\sqrt{\lambda'k'^{n-2}}=\left\{\Delta\operatorname{coam}2\omega\Delta\operatorname{coam}4\omega\ldots\Delta\operatorname{coam}(n-1)\omega\right\}^{2}\\[1.5ex]
\text{11)}\quad(-1)^{\frac{n-1}{2}}M\sqrt{\frac{1}{\lambda'k'^{n-2}}}=\left\{\operatorname{tg}\operatorname{coam}2\omega\operatorname{tg}\operatorname{coam}4\omega\ldots\operatorname{tg}\operatorname{coam}(n-1)\omega\right\}^{2}.
\end{gathered}
\]
By the aid of these formulae, formulae 1), 4), 5) pass into the following:
\[
\begin{gathered}
\text{12)}\quad\sin\operatorname{am}\left(\frac{u}{M},\lambda\right)=\sqrt{\frac{k^{n}}{\lambda}}\sin\operatorname{am}u\sin\operatorname{am}(u+4\omega)\sin\operatorname{am}(u+8\omega).\,.\,.\,.\sin\operatorname{am}(u+4(n-1)\omega)\\[1.5ex]
\text{13)}\quad\cos\operatorname{am}\left(\frac{u}{M},\lambda\right)=\sqrt{\frac{\lambda'k^{n}}{\lambda k'^{n}}}\cos\operatorname{am}u\cos\operatorname{am}(u+4\omega)\cos\operatorname{am}(u+8\omega).\,.\,.\,.\cos\operatorname{am}(u+4(n-1)\omega)\\[1.5ex]
\text{14)}\quad\Delta\operatorname{am}\left(\frac{u}{M},\lambda\right)=\sqrt{\frac{\lambda'}{k'^{n}}}\Delta\operatorname{am}u\Delta\operatorname{am}(u+4\omega)\Delta\operatorname{am}(u+8\omega).\,.\,.\,.\Delta\operatorname{am}(u+4(n-1)\omega),
\end{gathered}
\]
whence also:
\[
\text{15)}\quad\operatorname{tg}\operatorname{am}\left(\frac{u}{M},\lambda\right)=\sqrt{\frac{k'^{n}}{\lambda'}}\operatorname{tg}\operatorname{am}u\operatorname{tg}\operatorname{am}(u+4\omega)\operatorname{tg}\operatorname{am}(u+8\omega).\,.\,.\,.\operatorname{tg}\operatorname{am}(u+4(n-1)\omega).
\]

Another system of formulae is found as follows. From equation 4) it follows:
\[
\frac{\lambda}{M^{2}k^{n}}=\left\{\sin\operatorname{am}2\omega\sin\operatorname{am}4\omega.\,.\,.\,.\sin\operatorname{am}(n-1)\omega\right\}^{2},
\]\ednote{The exponent is printed 2, where 4 is expected (squaring formula 4) gives the fourth power); a hand correction to 4 in ink is not transcribed.}
whence:
\[
y=\sin\operatorname{am}\left(\frac{u}{M},\lambda\right)=\frac{x}{M}\,\Pi\frac{1-\frac{xx}{\sin^{2}\operatorname{am}2p\omega}}{1-k^{2}x^{2}\sin^{2}\operatorname{am}2p\omega}=\frac{kM}{\lambda}\,x\,\Pi\frac{xx-\sin^{2}\operatorname{am}2p\omega}{xx-\frac{1}{k^{2}\sin^{2}\operatorname{am}2p\omega}},
\]
or:
\[
0=x\,\Pi\left(x^{2}-\sin^{2}\operatorname{am}2p\omega\right)-\frac{\lambda}{kM}\sin\operatorname{am}\left(\frac{u}{M},\lambda\right)\Pi\left(x^{2}-\frac{1}{k^{2}\sin^{2}\operatorname{am}2p\omega}\right).
\]
The roots of this equation of the $n^{\mathrm{ti}}$ order are:
\[
x=\sin\operatorname{am}u,\ \sin\operatorname{am}(u+4\omega),\ \sin\operatorname{am}(u+8\omega),\ .\,.\,.\,.,\ \sin\operatorname{am}(u+4(n-1)\omega),
\]
whence we obtain the identical equation:
\[
\begin{gathered}
x\,\Pi\left(x^{2}-\sin^{2}\operatorname{am}2p\omega\right)-\frac{\lambda}{kM}\sin\operatorname{am}\left(\frac{u}{M},\lambda\right)\Pi\left(x^{2}-\frac{1}{k^{2}\sin^{2}\operatorname{am}2p\omega}\right)=\\[1.5ex]
\left(x-\sin\operatorname{am}u\right)\left(x-\sin\operatorname{am}(u+4\omega)\right)\left(x-\sin\operatorname{am}(u+8\omega)\right)\ldots\left(x-\sin\operatorname{am}(u+4(n-1)\omega)\right).
\end{gathered}
\]

\origpage{48}
Hence comes forth the sum of the roots
\[
\text{16)}\quad\Sigma\sin\operatorname{am}(u+4q\omega)=\frac{\lambda}{kM}\sin\operatorname{am}\left(\frac{u}{M},\lambda\right).
\]
In the same way is found:
\[
\begin{gathered}
\text{17)}\quad\Sigma\cos\operatorname{am}(u+4q\omega)=\frac{(-1)^{\frac{n-1}{2}}\lambda}{kM}\cos\operatorname{am}\left(\frac{u}{M},\lambda\right)\\[1.5ex]
\text{18)}\quad\Sigma\Delta\operatorname{am}(u+4q\omega)=\frac{(-1)^{\frac{n-1}{2}}}{M}\Delta\operatorname{am}\left(\frac{u}{M},\lambda\right)\\[1.5ex]
\text{19)}\quad\Sigma\operatorname{tg}\operatorname{am}(u+4q\omega)=\frac{\lambda'}{k'M}\operatorname{tg}\operatorname{am}\left(\frac{u}{M},\lambda\right),
\end{gathered}
\]
in which formulae the number $q$ is assigned the values $0,1,2,3,\ldots n-1$. These formulae it is convenient to represent also in this manner:
\[
\begin{gathered}
\frac{\lambda}{kM}\sin\operatorname{am}\left(\frac{u}{M},\lambda\right)=\sin\operatorname{am}u+\Sigma\left\{\sin\operatorname{am}(u+4q\omega)+\sin\operatorname{am}(u-4q\omega)\right\}\\[1.5ex]
\frac{(-1)^{\frac{n-1}{2}}\lambda}{kM}\cos\operatorname{am}\left(\frac{u}{M},\lambda\right)=\cos\operatorname{am}u+\Sigma\left\{\cos\operatorname{am}(u+4q\omega)+\cos\operatorname{am}(u-4q\omega)\right\}\\[1.5ex]
\frac{(-1)^{\frac{n-1}{2}}}{M}\Delta\operatorname{am}\left(\frac{u}{M},\lambda\right)=\Delta\operatorname{am}u+\Sigma\left\{\Delta\operatorname{am}(u+4q\omega)+\Delta\operatorname{am}(u-4q\omega)\right\}\\[1.5ex]
\frac{\lambda'}{k'M}\operatorname{tg}\operatorname{am}\left(\frac{u}{M},\lambda\right)=\operatorname{tg}\operatorname{am}u+\Sigma\left\{\operatorname{tg}\operatorname{am}(u+4q\omega)+\operatorname{tg}\operatorname{am}(u-4q\omega)\right\},
\end{gathered}
\]
where the number $q$ is assigned the values $1,2,3,\ldots\frac{n-1}{2}$. Now let the following formulae be noted:
\[
\begin{gathered}
\sin\operatorname{am}(u+4q\omega)+\sin\operatorname{am}(u-4q\omega)=\frac{2\cos\operatorname{am}4q\omega\,\Delta\operatorname{am}4q\omega\sin\operatorname{am}u}{1-k^{2}\sin^{2}\operatorname{am}4q\omega\sin^{2}\operatorname{am}u}\\[1.5ex]
\cos\operatorname{am}(u+4q\omega)+\cos\operatorname{am}(u-4q\omega)=\frac{2\cos\operatorname{am}4q\omega\cos\operatorname{am}u}{1-k^{2}\sin^{2}\operatorname{am}4q\omega\sin^{2}\operatorname{am}u}\\[1.5ex]
\Delta\operatorname{am}(u+4q\omega)+\Delta\operatorname{am}(u-4q\omega)=\frac{2\Delta\operatorname{am}4q\omega\Delta\operatorname{am}u}{1-k^{2}\sin^{2}\operatorname{am}4q\omega\sin^{2}\operatorname{am}u}\\[1.5ex]
\operatorname{tg}\operatorname{am}(u+4q\omega)+\operatorname{tg}\operatorname{am}(u-4q\omega)=\frac{2\Delta\operatorname{am}4q\omega\sin\operatorname{am}u\cos\operatorname{am}u}{\cos^{2}\operatorname{am}4q\omega-\Delta^{2}\operatorname{am}4q\omega\sin^{2}\operatorname{am}u}\,{}^{*)},
\end{gathered}
\]

\textbf{*)} Cf.\ §.\ 18 formulae 1), 2), 3); the last formula flows from formulae 10), 30), where you reflect that $\operatorname{tg}\sigma+\operatorname{tg}d=\frac{\sin(\sigma+d)}{\cos\sigma\cos d}$.

\origpage{49}
by whose aid formulae 16) --- 19) pass into these:
\[
\begin{gathered}
\text{20)}\quad\frac{\lambda}{kM}\sin\operatorname{am}\left(\frac{u}{M},\lambda\right)=\sin\operatorname{am}u+\Sigma\frac{2\cos\operatorname{am}4q\omega\,\Delta\operatorname{am}4q\omega\sin\operatorname{am}u}{1-k^{2}\sin^{2}\operatorname{am}4q\omega\sin^{2}\operatorname{am}u}\\[1.5ex]
\text{21)}\quad\frac{(-1)^{\frac{n-1}{2}}\lambda}{kM}\cos\operatorname{am}\left(\frac{u}{M},\lambda\right)=\cos\operatorname{am}u+\Sigma\frac{2\cos\operatorname{am}4q\omega\cos\operatorname{am}u}{1-k^{2}\sin^{2}\operatorname{am}4q\omega\sin^{2}\operatorname{am}u}\\[1.5ex]
\text{22)}\quad\frac{(-1)^{\frac{n-1}{2}}}{M}\Delta\operatorname{am}\left(\frac{u}{M},\lambda\right)=\Delta\operatorname{am}u+\Sigma\frac{2\Delta\operatorname{am}4q\omega\,\Delta\operatorname{am}u}{1-k^{2}\sin^{2}\operatorname{am}4q\omega\sin^{2}\operatorname{am}u}\\[1.5ex]
\text{23)}\quad\frac{\lambda'}{k'M}\operatorname{tg}\operatorname{am}\left(\frac{u}{M},\lambda\right)=\operatorname{tg}\operatorname{am}u+\Sigma\frac{2\Delta\operatorname{am}4q\omega\sin\operatorname{am}u\cos\operatorname{am}u}{1-k^{2}\sin^{2}\operatorname{am}4q\omega\sin^{2}\operatorname{am}u},
\end{gathered}
\]
which are also obtained when the formulae proposed above are resolved into simple fractions by known methods.

\section*{ON THE VARIOUS TRANSFORMATIONS OF THE SAME ORDER.}

\subsection*{TWO REAL TRANSFORMATIONS, OF THE GREATER MODULUS INTO THE LESSER AND OF THE LESSER INTO THE GREATER.}

\subsection*{24.}

We have seen that the element $\omega$ can be assigned any value of the scheme $\dfrac{mK+m'iK'}{n}$, $m$, $m'$ denoting positive or negative integers, which however, whenever $n$ is a composite number, have no common factor with $n$. But it is easily evident that, when $q$ is prime to $n$, the values $\omega=\dfrac{qmK+qm'iK'}{n}$ will not exhibit different substitutions. Hence, when $n$ itself is a prime number, the values of the element $\omega$ which furnish different transformations will be all of these:
\[
\frac{K}{n},\;\frac{iK'}{n},\;\frac{K+iK'}{n},\;\frac{K+2iK'}{n},\;\frac{K+3iK'}{n},\;.\,.\,.\,.\;\frac{K+(n-1)iK'}{n},
\]
or also:
\[
\frac{K}{n},\;\frac{iK'}{n},\;\frac{K+iK'}{n},\;\frac{2K+iK'}{n},\;\frac{3K+iK'}{n},\;.\,.\,.\,.\;\frac{(n-1)K+iK'}{n},
\]
or, if it pleases:
\[
\frac{K}{n},\;\frac{iK'}{n},\;\frac{K\pm iK'}{n},\;\frac{K\pm2iK'}{n},\;\frac{K\pm3iK'}{n},\;.\,.\,.\,.\;\frac{K\pm\frac{n-1}{2}iK}{n},
\]\ednote{Printed $iK$ instead of $iK'$ in the last term; the prime is missing, an apparent misprint or dropped ink.}

\origpage{50}
or also:

\[
\frac{K}{n},\;\frac{iK'}{n},\;\frac{K\pm iK'}{n},\;\frac{2K\pm iK'}{n},\;\frac{3K\pm iK'}{n},\;.\,.\,.\,.\;\frac{\frac{n-1}{2}K\pm iK'}{n},
\]
of which the number is $n+1$. And in fact we saw that in the transformations of the third and fifth order, proposed above as examples, the equations between $u=\sqrt[4]{k}$ and $v=\sqrt[4]{\lambda}$, which we shall call \emph{Modular Equations}, ascended respectively to the fourth and sixth degree. But whenever $n$ is a composite number, this number is greatly increased; for there are added the cases in which either $m$, or $m'$, or even both have a factor in common with $n$, provided only that it is not the same factor that $m$ and $m'$ both have in common with $n$. In general, however, the theorem holds:

„\emph{the number of substitutions of the $n^{\mathrm{ti}}$ order different from one another, by whose aid it is permitted to transform elliptic functions, equals the sum of the factors of $n$, which number, however, whenever $n$ is divisible by a square, embraces also substitutions mixed of transformation and multiplication; and so, whenever $n$ itself is a square, multiplication itself.}"\ednote{Printed as such: numerum substitutionem (where substitutionum is expected) and substisutiones (where substitutiones is expected).}

That sum of factors, therefore, will designate the degree to which, for a given number $n$, the Modular Equation will ascend, where it is to be noted that, whenever $n$ is a square number, one of the roots will present $k=\lambda$, and in general, whenever $n=m^{2}v$, $m^{2}$ denoting the least square by which it is permitted to divide the number $n$, among the roots will also be all the roots of the Modular Equation which belongs to $v$ itself.

Among the values of the element $\omega$ proposed above, which, in the case where $n$ is prime (to which, since the remaining ones return into it, it is convenient to attend either solely or before the rest), supply the entire multitude of transformations, only two, generally speaking ${}^{*)}$, are found which furnish real transformations; these I call $\omega=\dfrac{K}{n}$, $\omega=\dfrac{iK'}{n}$. The former we shall call in what follows the \emph{first} transformation, the latter the \emph{second}; and the moduli which correspond to them we shall designate respectively by $\lambda$, $\lambda_{,}$, and their Complements by $\lambda'$, $\lambda_{,}'$. The arguments of amplitude $\dfrac{\pi}{2}$ which correspond to these moduli, (which Cl.\ Legendre calls complete functions,) we shall designate by $\Lambda$, $\Lambda_{,}$, $\Lambda'$, $\Lambda_{,}'$. Our general formulae for these cases turn out to be the following.

\textbf{*)} For in infinitely many cases, for special Moduli, it happens that a pair of imaginary roots of the Modular Equations becomes equal to itself and so becomes real.

\origpage{51}
\subsection*{I.}

\subsection*{FORMULAE FOR THE FIRST REAL TRANSFORMATION OF THE MODULUS $k$ INTO THE MODULUS $\lambda$.}

\[
\begin{gathered}
\lambda=k^{n}\left\{\sin\operatorname{coam}\frac{2K}{n}\sin\operatorname{coam}\frac{4K}{n}.\,.\,.\,.\,.\sin\operatorname{coam}\frac{(n-1)K}{n}\right\}^{*}\\[2ex]
\lambda'=\frac{k'^{n}}{\left\{\Delta\operatorname{am}\frac{2K}{n}\Delta\operatorname{am}\frac{4K}{n}.\,.\,.\,.\,.\Delta\operatorname{am}\frac{(n-1)K}{n}\right\}^{4}}\\[2ex]
M=\left\{\frac{\sin\operatorname{coam}\frac{2K}{n}\sin\operatorname{coam}\frac{4K}{n}.\,.\,.\,.\sin\operatorname{coam}\frac{(n-1)K}{n}}{\sin\operatorname{am}\frac{2K}{n}\sin\operatorname{am}\frac{4K}{n}.\,.\,.\,.\sin\operatorname{am}\frac{(n-1)K}{n}}\right\}^{2}\\[2.5ex]
\sin\operatorname{am}\left(\frac{u}{M},\lambda\right)=\frac{\frac{\sin\operatorname{am}u}{M}\left(1-\frac{\sin^{2}\operatorname{am}u}{\sin^{2}\operatorname{am}\frac{2K}{n}}\right)\left(1-\frac{\sin^{2}\operatorname{am}u}{\sin^{2}\operatorname{am}\frac{4K}{n}}\right).\,.\,.\,.\left(1-\frac{\sin^{2}\operatorname{am}u}{\sin^{2}\operatorname{am}\frac{(n-1)K}{n}}\right)}{\left(1-k^{2}\sin^{2}\operatorname{am}\frac{2K}{n}\sin^{2}\operatorname{am}u\right)\left(1-k^{2}\sin^{2}\operatorname{am}\frac{4K}{n}\sin^{2}\operatorname{am}u\right).\,.\,.\,.\left(1-k^{2}\sin^{2}\operatorname{am}\frac{(n-1)K}{n}\sin^{2}\operatorname{am}u\right)}\\[2.5ex]
=(-1)^{\frac{n-1}{2}}\sqrt{\frac{k^{n}}{\lambda}}\sin\operatorname{am}u\sin\operatorname{am}\left(u+\frac{4K}{n}\right)\sin\operatorname{am}\left(u+\frac{8K}{n}\right).\,.\,.\,.\sin\operatorname{am}\left(u+\frac{4(n-1)K}{n}\right)\\[2.5ex]
\cos\operatorname{am}\left(\frac{u}{M},\lambda\right)=\sqrt{\frac{\lambda'k^{n}}{\lambda k'^{n}}}.\frac{\cos\operatorname{am}u\left(1-\frac{\sin^{2}\operatorname{am}u}{\sin^{2}\operatorname{coam}\frac{2K}{n}}\right)\left(1-\frac{\sin^{2}\operatorname{am}u}{\sin^{2}\operatorname{coam}\frac{4K}{n}}\right).\,.\,.\,.\left(1-\frac{\sin^{2}\operatorname{am}u}{\sin^{2}\operatorname{coam}\frac{(n-1)K}{n}}\right)}{\left(1-k^{2}\sin^{2}\operatorname{am}\frac{2K}{n}\sin^{2}\operatorname{am}u\right)\left(1-k^{2}\sin^{2}\operatorname{am}\frac{4K}{n}\sin^{2}\operatorname{am}u\right).\,.\,.\,.\left(1-k^{2}\sin^{2}\operatorname{am}\frac{(n-1)K}{n}\sin^{2}\operatorname{am}u\right)}
\end{gathered}
\]\ednote{The exponent after the closing brace of the first formula (for $\lambda$) is printed as an asterisk-like mark where the exponent 4 is expected (as in the formula for $\lambda'$); no footnote exists. The author's Corrigenda for this page say that in the expression for $\cos\operatorname{am}\left(\frac{u}{M},\lambda\right)$ the factor $\sqrt{\frac{\lambda'k^{n}}{\lambda k'^{n}}}$ is to be deleted and a second expression with that factor and the product of the cosines of $u$, $u+\frac{4K}{n}$, up to $u+\frac{4(n-1)K}{n}$ is to be added; the print has the factor and lacks the second expression.}

\[
\Delta\operatorname{am}\left(\frac{u}{M},\lambda\right)=\sqrt{\frac{\lambda'}{k'^{n}}}.\frac{\Delta\operatorname{am}u\left(1-k^{2}\sin^{2}\operatorname{coam}\frac{2K}{n}\sin^{2}\operatorname{am}u\right)\left(1-k^{2}\sin^{2}\operatorname{coam}\frac{2K}{n}\sin^{2}\operatorname{am}u\right).\,.\,.\,.\left(1-k^{2}\sin^{2}\operatorname{coam}\frac{(n-1)K}{n}\sin^{2}\operatorname{am}u\right)}{\left(1-k^{2}\sin^{2}\operatorname{am}\frac{2K}{n}\sin^{2}\operatorname{am}u\right)\left(1-k^{2}\sin^{2}\operatorname{am}\frac{4K}{n}\sin^{2}\operatorname{am}u\right).\,.\,.\,.\left(1-k^{2}\sin^{2}\operatorname{am}\frac{(n-1)K}{n}\sin^{2}\operatorname{am}u\right)}
\]\ednote{In the numerator the second factor is printed with $\frac{2K}{n}$, repeating the first factor, where $\frac{4K}{n}$ is expected; an apparent misprint.}

\[
\begin{gathered}
\sqrt{\frac{1\mp\sin\operatorname{am}\left(\frac{u}{M},\lambda\right)}{1\pm\sin\operatorname{am}\left(\frac{u}{M},\lambda\right)}}=\sqrt{\frac{1-\sin\operatorname{am}u}{1+\sin\operatorname{am}u}}.\frac{\left(1-\frac{\sin\operatorname{am}u}{\sin\operatorname{coam}\frac{4K}{n}}\right)\left(1-\frac{\sin\operatorname{am}u}{\sin\operatorname{coam}\frac{8K}{n}}\right).\,.\,.\,.\left(1-\frac{\sin\operatorname{am}u}{\sin\operatorname{coam}\frac{2(n-1)K}{n}}\right)}{\left(1+\frac{\sin\operatorname{am}u}{\sin\operatorname{coam}\frac{4K}{n}}\right)\left(1+\frac{\sin\operatorname{am}u}{\sin\operatorname{coam}\frac{8K}{n}}\right).\,.\,.\,.\left(1+\frac{\sin\operatorname{am}u}{\sin\operatorname{coam}\frac{2(n-1)K}{n}}\right)}\\[2.5ex]
\sqrt{\frac{1\mp\lambda\sin\operatorname{am}\left(\frac{u}{M},\lambda\right)}{1\pm\lambda\sin\operatorname{am}\left(\frac{u}{M},\lambda\right)}}=\\[2.5ex]
\sqrt{\frac{1-k\sin\operatorname{am}u}{1+k\sin\operatorname{am}u}}.\frac{\left(1-k\sin\operatorname{coam}\frac{4K}{n}\sin\operatorname{am}u\right)\left(1-k\sin\operatorname{coam}\frac{8K}{n}\sin\operatorname{am}u\right)\ldots\left(1-k\sin\operatorname{coam}\frac{2(n-1)K}{n}\sin\operatorname{am}u\right)}{\left(1+k\sin\operatorname{coam}\frac{4K}{n}\sin\operatorname{am}u\right)\left(1+k\sin\operatorname{coam}\frac{8K}{n}\sin\operatorname{am}u\right)\ldots\left(1+k\sin\operatorname{coam}\frac{2(n-1)K}{n}\sin\operatorname{am}u\right)}
\end{gathered}
\]

\origpage{52}
\[
\begin{gathered}
\frac{\lambda}{kM}\sin\operatorname{am}\left(\frac{u}{M},\lambda\right)=\sin\operatorname{am}u+2\Sigma\frac{(-1)^{q}\cos\operatorname{am}\frac{2qK}{n}\Delta\operatorname{am}\frac{2qK}{n}\sin\operatorname{am}u}{1-k^{2}\sin^{2}\operatorname{am}\frac{2qK}{n}\sin^{2}\operatorname{am}u}\\[2.5ex]
\frac{\lambda}{kM}\cos\operatorname{am}\left(\frac{u}{M},\lambda\right)=\cos\operatorname{am}u+2\Sigma\frac{(-1)^{q}\cos\operatorname{am}\frac{2qK}{n}\cos\operatorname{am}u}{1-k^{2}\sin^{2}\operatorname{am}\frac{2qK}{n}\sin^{2}\operatorname{am}u}\\[2.5ex]
\frac{1}{M}\Delta\operatorname{am}\left(\frac{u}{M},\lambda\right)=\Delta\operatorname{am}u+2\Sigma\frac{\Delta\operatorname{am}\frac{2qK}{n}\Delta\operatorname{am}u}{1-k^{2}\sin^{2}\operatorname{am}\frac{2qK}{n}\sin^{2}\operatorname{am}u}\\[2.5ex]
\frac{\lambda'}{k'M}\operatorname{tg}\operatorname{am}\left(\frac{u}{M},\lambda\right)=\operatorname{tg}\operatorname{am}u+2\Sigma\frac{\Delta\operatorname{am}\frac{2qK}{n}\sin\operatorname{am}u\cos\operatorname{am}u}{\cos^{2}\operatorname{am}\frac{2qK}{n}-\Delta^{2}\operatorname{am}\frac{2qK}{n}\sin^{2}\operatorname{am}u}.
\end{gathered}
\]

\subsection*{II.}

\subsection*{A. FORMULAE FOR THE SECOND REAL TRANSFORMATION, OF THE MODULUS $k$ INTO THE MODULUS $\lambda_{,}$, IN IMAGINARY FORM.}

\[
\begin{gathered}
\lambda_{,}=k^{n}\left\{\sin\operatorname{coam}\frac{2iK'}{n}\sin\operatorname{coam}\frac{4iK'}{n}.\,.\,.\,.\sin\operatorname{coam}\frac{(n-1)iK'}{n}\right\}^{4}\\[2ex]
\lambda_{,}'=\frac{k'^{n}}{\left\{\Delta\operatorname{am}\frac{2iK'}{n}\Delta\operatorname{am}\frac{4iK'}{n}.\,.\,.\,.\Delta\operatorname{am}\frac{(n-1)iK'}{n}\right\}^{4}}\\[2.5ex]
M_{,}=(-1)^{\frac{n-1}{2}}\left\{\frac{\sin\operatorname{coam}\frac{2iK'}{n}\sin\operatorname{coam}\frac{4iK'}{n}\ldots\sin\operatorname{coam}\frac{(n-1)iK'}{n}}{\sin\operatorname{am}\frac{2iK'}{n}\sin\operatorname{am}\frac{4iK'}{n}\ldots\sin\operatorname{am}\frac{(n-1)iK'}{n}}\right\}^{2}\\[2.5ex]
\sin\operatorname{am}\left(\frac{u}{M_{,}},\lambda_{,}\right)=\frac{\sin\operatorname{am}u\left(1-\frac{\sin^{2}\operatorname{am}u}{\sin^{2}\operatorname{am}\frac{2iK'}{n}}\right)\left(1-\frac{\sin^{2}\operatorname{am}u}{\sin^{2}\operatorname{am}\frac{4iK'}{n}}\right).\,.\,.\,.\left(1-\frac{\sin^{2}\operatorname{am}u}{\sin^{2}\operatorname{am}\frac{(n-1)iK'}{n}}\right)}{\left(1-\frac{\sin^{2}\operatorname{am}u}{\sin^{2}\operatorname{am}\frac{iK'}{n}}\right)\left(1-\frac{\sin^{2}\operatorname{am}u}{\sin^{2}\operatorname{am}\frac{3iK'}{n}}\right).\,.\,.\,.\left(1-\frac{\sin^{2}\operatorname{am}u}{\sin^{2}\operatorname{am}\frac{(n-2)iK'}{n}}\right)}\\[2.5ex]
=\sqrt{\frac{k^{n}}{\lambda_{,}}}.\sin\operatorname{am}u\sin\operatorname{am}(u+4iK')\sin\operatorname{am}(u+8iK').\,.\,.\,.\sin\operatorname{am}(u+4(n-1)iK')
\end{gathered}
\]\ednote{The arguments $u+4iK'$, $u+8iK'$, up to $u+4(n-1)iK'$ are printed without the denominator $n$, where $u+\frac{4iK'}{n}$ and so on are expected; left as printed.}

\origpage{53}
\[
\begin{gathered}
\cos\operatorname{am}\left(\frac{u}{M_{,}},\lambda_{,}\right)=\frac{\cos\operatorname{am}u\left(1-\frac{\sin^{2}\operatorname{am}u}{\sin^{2}\operatorname{coam}\frac{2iK'}{n}}\right)\left(1-\frac{\sin^{2}\operatorname{am}u}{\sin^{2}\operatorname{coam}\frac{4iK'}{n}}\right).\,.\,.\,.\left(1-\frac{\sin^{2}\operatorname{am}u}{\sin^{2}\operatorname{coam}\frac{(n-1)iK'}{n}}\right)}{\left(1-\frac{\sin^{2}\operatorname{am}u}{\sin^{2}\operatorname{am}\frac{iK'}{n}}\right)\left(1-\frac{\sin^{2}\operatorname{am}u}{\sin^{2}\operatorname{am}\frac{3iK'}{n}}\right).\,.\,.\,.\left(1-\frac{\sin^{2}\operatorname{am}u}{\sin^{2}\operatorname{am}\frac{(n-2)iK'}{n}}\right)}\\[2.5ex]
=\sqrt{\frac{\lambda_{,}'k^{n}}{\lambda_{,}k'^{n}}}.\cos\operatorname{am}u\cos\operatorname{am}\left(u+\frac{4iK'}{n}\right)\cos\operatorname{am}\left(u+\frac{8iK'}{n}\right).\,.\,.\,.\cos\operatorname{am}\left(u+\frac{4(n-1)iK'}{n}\right)\\[2.5ex]
\Delta\operatorname{am}\left(\frac{u}{M_{,}},\lambda_{,}\right)=\frac{\Delta\operatorname{am}u\left(1-\frac{\sin^{2}\operatorname{am}u}{\sin^{2}\operatorname{coam}\frac{iK'}{n}}\right)\left(1-\frac{\sin^{2}\operatorname{am}u}{\sin^{2}\operatorname{coam}\frac{3iK'}{n}}\right).\,.\,.\,.\left(1-\frac{\sin^{2}\operatorname{am}u}{\sin^{2}\operatorname{coam}\frac{(n-2)iK'}{n}}\right)}{\left(1-\frac{\sin^{2}\operatorname{am}u}{\sin^{2}\operatorname{am}\frac{iK'}{n}}\right)\left(1-\frac{\sin^{2}\operatorname{am}u}{\sin^{2}\operatorname{am}\frac{3iK'}{n}}\right).\,.\,.\,.\left(1-\frac{\sin^{2}\operatorname{am}u}{\sin^{2}\operatorname{am}\frac{(n-2)iK'}{n}}\right)}\\[2.5ex]
=\sqrt{\frac{\lambda'}{k'^{n}}}\Delta\operatorname{am}u\,\Delta\operatorname{am}(u+4iK')\Delta\operatorname{am}(u+8iK').\,.\,.\,.\Delta\operatorname{am}(u+4(n-1)iK')
\end{gathered}
\]\ednote{As in the sine formula on the previous page, the arguments $u+4iK'$, $u+8iK'$, and so on are printed without the denominator $n$; left as printed.}

\[
\begin{gathered}
\sqrt{\frac{1-\sin\operatorname{am}\left(\frac{u}{M_{,}},\lambda_{,}\right)}{1+\sin\operatorname{am}\left(\frac{u}{M_{,}},\lambda_{,}\right)}}=\\[2.5ex]
\sqrt{\frac{1-\sin\operatorname{am}u}{1+\sin\operatorname{am}u}}.\frac{\left(1-\frac{\sin\operatorname{am}u}{\sin\operatorname{coam}\frac{2iK'}{n}}\right)\left(1-\frac{\sin\operatorname{am}u}{\sin\operatorname{coam}\frac{4iK'}{n}}\right).\,.\,.\,.\left(1-\frac{\sin\operatorname{am}u}{\sin\operatorname{coam}\frac{(n-1)iK'}{n}}\right)}{\left(1+\frac{\sin\operatorname{am}u}{\sin\operatorname{coam}\frac{2iK'}{n}}\right)\left(1+\frac{\sin\operatorname{am}u}{\sin\operatorname{coam}\frac{4iK'}{n}}\right).\,.\,.\,.\left(1+\frac{\sin\operatorname{am}u}{\sin\operatorname{coam}\frac{(n-1)iK'}{n}}\right)}\\[2.5ex]
\sqrt{\frac{1-\lambda_{,}\sin\operatorname{am}\left(\frac{u}{M_{,}},\lambda_{,}\right)}{1+\lambda_{,}\sin\operatorname{am}\left(\frac{u}{M_{,}},\lambda_{,}\right)}}=\\[2.5ex]
\sqrt{\frac{1-k\sin\operatorname{am}u}{1+k\sin\operatorname{am}u}}.\frac{\left(1-\frac{\sin\operatorname{am}u}{\sin\operatorname{coam}\frac{iK'}{n}}\right)\left(1-\frac{\sin\operatorname{am}u}{\sin\operatorname{coam}\frac{3iK'}{n}}\right)\ldots\left(1-\frac{\sin\operatorname{am}u}{\sin\operatorname{coam}\frac{(n-2)iK'}{n}}\right)}{\left(1+\frac{\sin\operatorname{am}u}{\sin\operatorname{coam}\frac{iK'}{n}}\right)\left(1+\frac{\sin\operatorname{am}u}{\sin\operatorname{coam}\frac{3iK'}{n}}\right)\ldots\left(1+\frac{\sin\operatorname{am}u}{\sin\operatorname{coam}\frac{(n-2)iK'}{n}}\right)}\\[2.5ex]
\frac{\lambda_{,}}{kM_{,}}\sin\operatorname{am}\left(\frac{u}{M_{,}},\lambda_{,}\right)=\sin\operatorname{am}u-\frac{2}{k}\Sigma\frac{\cos\operatorname{am}\frac{(2q-1)iK'}{n}\Delta\operatorname{am}\frac{(2q-1)iK'}{n}\sin\operatorname{am}u}{\sin^{2}\operatorname{am}\frac{(2q-1)iK'}{n}-\sin^{2}\operatorname{am}u}
\end{gathered}
\]

\origpage{54}
\[
\begin{gathered}
\frac{(-1)^{\frac{n-1}{2}}\lambda_{,}}{kM_{,}}\cos\operatorname{am}\left(\frac{u}{M_{,}},\lambda_{,}\right)=\cos\operatorname{am}u+\frac{2(-1)^{\frac{n-1}{2}}}{ik}\Sigma\frac{(-1)^{q}\sin\operatorname{am}\frac{(2q-1)iK'}{n}\Delta\operatorname{am}\frac{(2q-1)iK'}{n}\cos\operatorname{am}u}{\sin^{2}\operatorname{am}\frac{(2q-1)iK'}{n}-\sin^{2}\operatorname{am}u}\\[2.5ex]
\frac{(-1)^{\frac{n-1}{2}}}{M_{,}}\Delta\operatorname{am}\left(\frac{u}{M_{,}},\lambda_{,}\right)=\Delta\operatorname{am}u+\frac{2(-1)^{\frac{n-1}{2}}}{i}\Sigma\frac{(-1)^{q}\sin\operatorname{am}\frac{(2q-1)iK'}{n}\cos\operatorname{am}\frac{(2q-1)iK'}{n}\Delta\operatorname{am}u}{\sin^{2}\operatorname{am}\frac{(2q-1)iK'}{n}-\sin^{2}\operatorname{am}u}\\[2.5ex]
-\frac{\lambda_{,}'}{k'M_{,}}\operatorname{tg}\operatorname{am}\left(\frac{u}{M_{,}},\lambda_{,}\right)=\operatorname{tg}\operatorname{am}u+2\Sigma\frac{(-1)^{q}\Delta\operatorname{am}\frac{2qiK'}{n}\sin\operatorname{am}u\cos\operatorname{am}u}{\cos^{2}\operatorname{am}\frac{2qiK'}{n}-\Delta^{2}\operatorname{am}\frac{2qiK'}{n}\sin^{2}\operatorname{am}u}.
\end{gathered}
\]

\subsection*{B. FORMULAE FOR THE SECOND REAL TRANSFORMATION IN REAL FORM.}

\[
\begin{gathered}
\lambda_{,}=\frac{k^{n}}{\left\{\Delta\operatorname{am}\left(\frac{2K'}{n},k'\right)\Delta\operatorname{am}\left(\frac{4K'}{n},k'\right).\,.\,.\,.\Delta\operatorname{am}\left(\frac{(n-1)K'}{n},k'\right)\right\}^{4}}\\[2.5ex]
\lambda_{,}'=k'^{n}\left\{\sin\operatorname{coam}\left(\frac{2K'}{n},k'\right)\sin\operatorname{coam}\left(\frac{4K'}{n},k'\right).\,.\,.\,.\sin\operatorname{coam}\left(\frac{(n-1)K'}{n},k'\right)\right\}^{4}\\[2.5ex]
M_{,}=\left\{\frac{\sin\operatorname{coam}\left(\frac{2K'}{n},k'\right)\sin\operatorname{coam}\left(\frac{4K'}{n},k'\right).\,.\,.\,.\sin\operatorname{coam}\left(\frac{(n-1)K'}{n},k'\right)}{\sin\operatorname{am}\left(\frac{2K'}{n},k'\right)\sin\operatorname{am}\left(\frac{4K'}{n},k'\right).\,.\,.\,.\sin\operatorname{am}\left(\frac{(n-1)K'}{n},k'\right)}\right\}^{2}
\end{gathered}
\]

\[
\sin\operatorname{am}\left(\frac{u}{M_{,}},\lambda_{,}\right)=\frac{\frac{\sin\operatorname{am}u}{M_{,}}\left(1+\frac{\sin^{2}\operatorname{am}u}{\operatorname{tg}^{2}\operatorname{am}\left(\frac{2K'}{n},k'\right)}\right)\left(1+\frac{\sin^{2}\operatorname{am}u}{\operatorname{tg}^{2}\operatorname{am}\left(\frac{4K'}{n},k'\right)}\right)\ldots\left(1+\frac{\sin^{2}\operatorname{am}u}{\operatorname{tg}^{2}\operatorname{am}\left(\frac{(n-1)K'}{n},k'\right)}\right)}{\left(1+\frac{\sin^{2}\operatorname{am}u}{\operatorname{tg}^{2}\operatorname{am}\left(\frac{K'}{n},k'\right)}\right)\left(1+\frac{\sin^{2}\operatorname{am}u}{\operatorname{tg}^{2}\operatorname{am}\left(\frac{3K'}{n},k'\right)}\right)\ldots\left(1+\frac{\sin^{2}\operatorname{am}u}{\operatorname{tg}^{2}\operatorname{am}\left(\frac{(n-2)K'}{n},k\right)}\right)}
\]\ednote{In the last denominator factor the modulus is printed $k$ instead of $k'$; an apparent misprint.}

\[
\begin{gathered}
\cos\operatorname{am}\left(\frac{u}{M_{,}},\lambda_{,}\right)=\frac{\cos\operatorname{am}u\left(1+\sin^{2}\operatorname{am}u\,\Delta^{2}\operatorname{am}\left(\frac{2K'}{n},k'\right)\right)\left(1+\sin^{2}\operatorname{am}u\,\Delta^{2}\operatorname{am}\left(\frac{4K'}{n},k'\right)\right)\ldots\left(1+\sin^{2}\operatorname{am}u\,\Delta^{2}\operatorname{am}\left(\frac{(n-1)K'}{n},k'\right)\right)}{\left(1+\frac{\sin^{2}\operatorname{am}u}{\operatorname{tg}^{2}\operatorname{am}\left(\frac{K'}{n},k'\right)}\right)\left(1+\frac{\sin^{2}\operatorname{am}u}{\operatorname{tg}^{2}\operatorname{am}\left(\frac{3K'}{n},k'\right)}\right)\ldots\left(1+\frac{\sin^{2}\operatorname{am}u}{\operatorname{tg}^{2}\operatorname{am}\left(\frac{(n-2)K'}{n},k'\right)}\right)}\\[2.5ex]
\Delta\operatorname{am}\left(\frac{u}{M_{,}},\lambda_{,}\right)=\frac{\Delta\operatorname{am}u\left(1+\sin^{2}\operatorname{am}u\,\Delta^{2}\operatorname{am}\left(\frac{K'}{n},k'\right)\right)\left(1+\sin^{2}\operatorname{am}u\,\Delta^{2}\operatorname{am}\left(\frac{3K'}{n},k'\right)\right)\ldots\left(1+\sin^{2}\operatorname{am}u\,\Delta^{2}\operatorname{am}\left(\frac{(n-2)K'}{n},k'\right)\right)}{\left(1+\frac{\sin^{2}\operatorname{am}u}{\operatorname{tg}^{2}\operatorname{am}\left(\frac{K'}{n},k'\right)}\right)\left(1+\frac{\sin^{2}\operatorname{am}u}{\operatorname{tg}^{2}\operatorname{am}\left(\frac{3K'}{n},k'\right)}\right)\ldots\left(1+\frac{\sin^{2}\operatorname{am}u}{\operatorname{tg}^{2}\operatorname{am}\left(\frac{(n-2)K'}{n},k'\right)}\right)}
\end{gathered}
\]

\origpage{55}
\[
\begin{gathered}
\sqrt{\frac{1-\sin\operatorname{am}\left(\frac{u}{M_{,}},\lambda_{,}\right)}{1+\sin\operatorname{am}\left(\frac{u}{M_{,}},\lambda_{,}\right)}}=\\[2.5ex]
\sqrt{\frac{1-\sin\operatorname{am}u}{1+\sin\operatorname{am}u}}.\frac{\left(1-\sin\operatorname{am}u\,\Delta\operatorname{am}\left(\frac{2K'}{n},k'\right)\right)\left(1-\sin\operatorname{am}u\,\Delta\operatorname{am}\left(\frac{4K'}{n},k'\right)\right)\ldots\left(1-\sin\operatorname{am}u\,\Delta\operatorname{am}\left(\frac{(n-1)K'}{n},k'\right)\right)}{\left(1+\sin\operatorname{am}u\,\Delta\operatorname{am}\left(\frac{2K'}{n},k'\right)\right)\left(1+\sin\operatorname{am}u\,\Delta\operatorname{am}\left(\frac{4K'}{n},k'\right)\right)\ldots\left(1+\sin\operatorname{am}u\,\Delta\operatorname{am}\left(\frac{(n-1)K'}{n},k'\right)\right)}\\[2.5ex]
\sqrt{\frac{1-\lambda_{,}\sin\operatorname{am}\left(\frac{u}{M_{,}},\lambda_{,}\right)}{1+\lambda_{,}\sin\operatorname{am}\left(\frac{u}{M_{,}},\lambda_{,}\right)}}=\\[2.5ex]
\sqrt{\frac{1-k\sin\operatorname{am}u}{1+k\sin\operatorname{am}u}}.\frac{\left(1-\Delta\operatorname{am}\left(\frac{K'}{n},k'\right)\sin\operatorname{am}u\right)\left(1-\Delta\operatorname{am}\left(\frac{3K'}{n},k'\right)\sin\operatorname{am}u\right)\ldots\left(1-\Delta\operatorname{am}\left(\frac{(n-2)K'}{n},k'\right)\sin\operatorname{am}u\right)}{\left(1+\Delta\operatorname{am}\left(\frac{K'}{n},k'\right)\sin\operatorname{am}u\right)\left(1+\Delta\operatorname{am}\left(\frac{3K'}{n},k'\right)\sin\operatorname{am}u\right)\ldots\left(1+\Delta\operatorname{am}\left(\frac{(n-2)K'}{n},k'\right)\sin\operatorname{am}u\right)}\\[2.5ex]
\frac{\lambda_{,}}{kM_{,}}\sin\operatorname{am}\left(\frac{u}{M_{,}},\lambda_{,}\right)=\sin\operatorname{am}u+\frac{2}{k}\Sigma\frac{\Delta\operatorname{am}\left(\frac{2(q-1)K'}{n},k'\right)\sin\operatorname{am}u}{\sin^{2}\operatorname{am}\left(\frac{2(q-1)K'}{n},k'\right)+\cos^{2}\operatorname{am}\left(\frac{2(q-1)K'}{n},k'\right)\sin^{2}\operatorname{am}u}\\[2.5ex]
\frac{(-1)^{\frac{n-1}{2}}\lambda_{,}}{kM_{,}}\cos\operatorname{am}\left(\frac{u}{M_{,}},\lambda_{,}\right)=\cos\operatorname{am}u-\frac{2(-1)^{\frac{n-1}{2}}}{k}\Sigma\frac{(-1)^{q}\sin\operatorname{am}\left(\frac{2(q-1)K'}{n},k'\right)\Delta\operatorname{am}\left(\frac{2(q-1)K'}{n},k'\right)\cos\operatorname{am}u}{\sin^{2}\operatorname{am}\left(\frac{2(q-1)K'}{n},k'\right)+\cos^{2}\operatorname{am}\left(\frac{2(q-1)K'}{n},k'\right)\sin^{2}\operatorname{am}u}\\[2.5ex]
\frac{(-1)^{\frac{n-1}{2}}}{M_{,}}\Delta\operatorname{am}\left(\frac{u}{M_{,}},\lambda_{,}\right)=\Delta\operatorname{am}u-2(-1)^{\frac{n-1}{2}}\Sigma\frac{(-1)^{q}\sin\operatorname{am}\left(\frac{2(q-1)K'}{n},k'\right)\Delta\operatorname{am}u}{\sin^{2}\operatorname{am}\left(\frac{2(q-1)K'}{n},k'\right)+\cos^{2}\operatorname{am}\left(\frac{2(q-1)K'}{n},k'\right)\sin^{2}\operatorname{am}u}\\[2.5ex]
\frac{\lambda_{,}'}{k'M_{,}}\operatorname{tg}\operatorname{am}\left(\frac{u}{M_{,}},\lambda'\right)=\operatorname{tg}\operatorname{am}u+2\Sigma\frac{(-1)^{q}\cos\operatorname{am}\left(\frac{2qK'}{n},k'\right)\Delta\operatorname{am}\left(\frac{2qK'}{n},k'\right)\sin\operatorname{am}u\cos\operatorname{am}u}{1-\Delta^{2}\operatorname{am}\left(\frac{2qK'}{n},k'\right)\sin^{2}\operatorname{am}u}.
\end{gathered}
\]

\origpage{56}
In the formulae for the first transformation, $(-1)^{\frac{n-1}{2}}M$ has been put in place of $M$. I have chosen to present the formulae for the second transformation in two ways, both under the imaginary form and under the real form, in which moreover, in place of $k\sin\operatorname{am}\dfrac{2miK'}{n}$, $k\sin\operatorname{coam}\dfrac{2miK'}{n}$, etc., there is everywhere written $\dfrac{-1}{\sin\operatorname{am}\dfrac{(n-2m)iK'}{n}}$, $\dfrac{1}{\sin\operatorname{coam}\dfrac{(n-2m)iK'}{n}}$ etc. This, like the reduction to the real form, is easily accomplished by means of the formulae of §$^{\mathrm{i}}$ 19. Where the ambiguous sign $\pm$ is put, the upper $+$ is to be chosen when $\dfrac{n-1}{2}$ is an even number, the lower $-$ when $\dfrac{n-1}{2}$ is an odd number; for the sign $\mp$ the opposite holds. In the sums designated by the prefixed $\Sigma$, the number $q$ is to be given the values $1$, $2$, $3$, $.\,.\,.\,.$ $\dfrac{n-1}{2}$.

From the formulae set forth for the first transformation it is evident that, whenever $u$ becomes successively:

\[
0,\;\frac{K}{n},\;\frac{2K}{n},\;\frac{3K}{n},\;\frac{4K}{n},\;.\,.\,.\,.,
\]
$\operatorname{am}\left(\dfrac{u}{M},\lambda\right)$ will be:

\[
0,\;\frac{\pi}{2},\;\pi,\;\frac{3\pi}{2},\;2\pi,\;.\,.\,.\,.
\]
whence we obtain:

\[
\frac{K}{nM}=\Lambda.
\]
On the other hand we see that in the second transformation, whenever $u$ becomes: $0$, $K$, $2K$, $3K$, $\ldots$ or $\operatorname{am}u$: $0$, $\dfrac{\pi}{2}$, $\pi$, $\dfrac{3\pi}{2}$, $\ldots$, $\operatorname{am}\left(\dfrac{u}{M_{,}},\lambda_{,}\right)$ also becomes $=0$, $\dfrac{\pi}{2}$, $\pi$, $\dfrac{3\pi}{2}$, $\ldots$, whence in this case:

\[
\frac{K}{M_{,}}=\Lambda_{,}.
\]

Moreover, from the formulae exhibited for the Moduli $\lambda$, $\lambda'$, $\lambda_{,}$, $\lambda_{,}'$ it is clear that, as $n$ increases, the Moduli $\lambda$, $\lambda_{,}'$ converge rapidly to zero, and thus at the same time the Moduli $\lambda'$, $\lambda_{,}$ approach very closely to unity. It is therefore fitting to call the first transformation of the Modulus that \emph{of the greater into the less}, the second that \emph{of the less into the greater.}

\origpage{57}
\section*{ON THE COMPLEMENTARY TRANSFORMATIONS}

\subsection*{§. HOW FROM THE TRANSFORMATION OF A MODULUS INTO A MODULUS ANOTHER IS DERIVED, OF THE COMPLEMENT INTO THE COMPLEMENT.}

\subsection*{25.}

In the formula found above:

\[
\operatorname{tg}\operatorname{am}\left(\frac{u}{M},\lambda\right)=\sqrt{\frac{k'^{n}}{\lambda'}}\operatorname{tg}\operatorname{am}u\operatorname{tg}\operatorname{am}(u+4\omega)\operatorname{tg}\operatorname{am}(u+8\omega).\,.\,.\,.\operatorname{tg}\operatorname{am}(u+4(n-1)\omega)
\]
let us put $u=iu'$, $\omega=i\omega'$, so that $\omega=mK+m'iK'$, $\omega'=m'K'-miK$. Now indeed (§.\ 19) we have

\[
\begin{gathered}
\operatorname{tg}\operatorname{am}(iu',k)=i\sin\operatorname{am}(u',k')\\
\operatorname{tg}\operatorname{am}(iu',\lambda)=i\sin\operatorname{am}(u',\lambda'),
\end{gathered}
\]
whence we see the cited formula pass into the following:

\[
\sin\operatorname{am}\left(\frac{u'}{M},\lambda\right)=(-1)^{\frac{n-1}{2}}\sqrt{\frac{k'^{n}}{\lambda'}}\sin\operatorname{am}u'\sin\operatorname{am}(u'+4\omega')\sin\operatorname{am}(u'+8\omega').\,.\,.\,.\sin\operatorname{am}(u'+4(n-1)\omega').\left\{\operatorname{Mod}k'\right\}.
\]
Further, we have found the formulae:

\[
\begin{gathered}
\lambda'=\frac{k'^{n}}{\left\{\Delta\operatorname{am}2\omega\,\Delta\operatorname{am}4\omega.\,.\,.\,.\Delta\operatorname{am}(n-1)\omega\right\}^{4}}\\[2ex]
M=(-1)^{\frac{n-1}{2}}\frac{\left\{\sin\operatorname{coam}2\omega\sin\operatorname{coam}4\omega.\,.\,.\,.\sin\operatorname{coam}(n-1)\omega\right\}^{2}}{\left\{\sin\operatorname{am}2\omega\sin\operatorname{am}4\omega.\,.\,.\,.\sin\operatorname{am}(n-1)\omega\right\}^{2}},
\end{gathered}
\]
which by means of the formulae:

\[
\begin{gathered}
\Delta\operatorname{am}(iu,k)=\frac{1}{\sin\operatorname{coam}(u,k')}\\[1.5ex]
\sin\operatorname{coam}(iu,k)=\frac{1}{\Delta\operatorname{am}(u,k')},
\end{gathered}
\]
whence there also follows:

\[
\frac{\sin\operatorname{coam}(iu,k)}{\sin\operatorname{am}(iu,k)}=\frac{-i}{\operatorname{tg}\operatorname{am}(u,k')\Delta\operatorname{am}(u,k')}=-i\,\frac{\sin\operatorname{coam}(u,k')}{\sin\operatorname{am}(u,k')},
\]
pass into the following:

\[
\begin{gathered}
\lambda'=k'^{n}\left\{\sin\operatorname{coam}2\omega'\sin\operatorname{coam}4\omega'.\,.\,.\,.\sin\operatorname{coam}(n-1)\omega'\right\}^{4}\quad\left\{\operatorname{Mod}k'\right\}\\[2ex]
M=\frac{\left\{\sin\operatorname{coam}2\omega'\sin\operatorname{coam}4\omega'.\,.\,.\,.\sin\operatorname{coam}(n-1)\omega'\right\}^{2}}{\left\{\sin\operatorname{am}2\omega'\sin\operatorname{am}4\omega'.\,.\,.\,.\sin\operatorname{am}(n-1)\omega'\right\}^{2}}\quad\left\{\operatorname{Mod}k'\right\}.
\end{gathered}
\]

\origpage{58}
Comparing these formulae with those which serve for the transformation of the Modulus $k$ into the Modulus $\lambda$:

\[
\begin{gathered}
\sin\operatorname{am}\left(\frac{u}{M},\lambda\right)=\sqrt{\frac{k^{n}}{\lambda}}\sin\operatorname{am}u\sin\operatorname{am}(u+4\omega)\sin\operatorname{am}(u+8\omega).\,.\,.\,.\sin\operatorname{am}(u+4(n-1)\omega)\\[1.5ex]
\lambda=k^{n}\left\{\sin\operatorname{coam}2\omega\sin\operatorname{coam}4\omega.\,.\,.\,.\sin\operatorname{coam}(n-1)\omega\right\}^{4}\\[1.5ex]
M=(-1)^{\frac{n-1}{2}}\left\{\frac{\sin\operatorname{coam}2\omega\sin\operatorname{coam}4\omega.\,.\,.\,.\sin\operatorname{coam}(n-1)\omega}{\sin\operatorname{am}2\omega\sin\operatorname{am}4\omega.\,.\,.\,.\sin\operatorname{am}(n-1)\omega}\right\}^{2},
\end{gathered}
\]
there emerges a Theorem which must be reckoned of the greatest moment in the Theory of Transformation:

„\emph{Whatever formulae can be set forth concerning the Transformation of the Modulus $k$ into the Modulus $\lambda$, the same hold, on changing $k$ into $k'$, $\lambda$ into $\lambda'$, $\omega$ into $\omega'=\dfrac{\omega}{i}$, $M$ into $(-1)^{\frac{n-1}{2}}M$.}"
The transformation of the Complement into the Complement, derived in the said manner from the proposed transformation, we shall call the \emph{Complementary Transformation.}

It is readily evident that the complementary transformations of the real transformations of the Modulus $k$ are the real transformations of the Modulus $k'$, in such a way, however, that the complementary of the first of the Modulus $k$ is the second of the Modulus $k'$, and the complementary of the second of the Modulus $k$ is the first of the Modulus $k'$. For where in the theorem just set forth one puts $\omega=\dfrac{\pm K}{n}$, $\omega=\dfrac{\pm iK'}{n}$, which correspond to the first and second transformations of the Modulus $k$, there results $\omega'=\dfrac{\omega}{i}=\dfrac{\pm iK}{n}$, $\omega'=\dfrac{\omega}{i}=\dfrac{\pm K'}{n}$, which correspond to the transformations of the Modulus $k'$, the second and first respectively. Also, since as the Modulus increases the Complement decreases and vice versa, the transformation of Modulus into Modulus, where it is that of the greater into the less, must be the transformation of Complement into Complement, or the complementary transformation, of the less into the greater, and vice versa. We see therefore that, on changing $k$ into $k'$, $\lambda$ passes into $\lambda_{,}'$, and $\lambda_{,}$ into $\lambda'$. Also the Multiplier $M$, common to the first transformation and its complementary ${}^{*)}$, will pass

\textbf{*)} This holds generally only if the sign be neglected; for we have seen that what was $M$ in the one transformation is $(-1)^{\frac{n-1}{2}}M$ in the complementary; but in our cases, because in the first transformation $(-1)^{\frac{n-1}{2}}M$ has been put in place of $M$ (see above), the ambiguity of sign is removed, so that for the complementary real transformations the Multiplier $M$ is wholly the same.

\origpage{59}
into $M_{,}$, which belongs to the second transformation and its complementary, and vice versa $M_{,}$ into $M$. Hence from the formulae found above:

\[
\Lambda=\frac{K}{nM},\quad\Lambda_{,}=\frac{K}{M_{,}},
\]
these follow:

\[
\Lambda_{,}'=\frac{K'}{nM_{,}},\quad\Lambda'=\frac{K'}{M},
\]
whence arise the formulae of the highest moment in this theory:

\[
\frac{\Lambda'}{\Lambda}=n\frac{K'}{K};\quad\frac{\Lambda_{,}'}{\Lambda_{,}}=\frac{1}{n}.\,\frac{K'}{K}.
\]
These formulae constitute the genuine character of the proposed transformation, whence it is evident that we have rightly referred the individual transformations to the individual numbers $n$. I shall remark that, whenever $n$ is a composite number $=n'n''$, from the individual real roots of the Modular Equations, or from the individual real Moduli into which a given Modulus $k$ may be transformed by a substitution of the $n^{\mathrm{ti}}$ order, equations of this kind arise:

\[
\frac{\Lambda'}{\Lambda}=\frac{n'}{n''}\frac{K'}{K},
\]
which correspond to the individual resolutions of the number $n$ into two factors. Among the number of these, therefore, whenever $n$ is a square number, there will also be this one:

\[
\frac{\Lambda'}{\Lambda}=\frac{K'}{K},\;\text{whence }\lambda=k,
\]
which teaches that, in the case where $n$ is a square, among the number of substitutions there is one which supplies a multiplication.

\origpage{60}
\section*{ON THE TRANSFORMATIONS SUPPLEMENTARY TO MULTIPLICATION.}

\subsection*{26.}

Let us recall the formulae:

\[
\frac{\Lambda'}{\Lambda}=n\frac{K'}{K},\quad\frac{\Lambda_{,}'}{\Lambda_{,}}=\frac{1}{n}\frac{K'}{K},
\]
which, written in this manner:

\[
\begin{gathered}
\frac{\Lambda'}{\Lambda}=n\frac{K'}{K}\\[1.5ex]
\frac{K'}{K}=n.\,\frac{\Lambda_{,}'}{\Lambda_{,}},
\end{gathered}
\]
make it clear that \emph{the Modulus $\lambda$ depends on the Modulus $k$ in the same way as the Modulus $k$ on the Modulus $\lambda_{,}$, or that the Modulus $k$ depends on the Modulus $\lambda$ in the same way as the Modulus $\lambda_{,}$ on the Modulus $k$.} Therefore by the first transformation, or that of the greater into the less, by which $k$ is transformed into $\lambda$, $\lambda_{,}$ will be transformed into $k$; by the second transformation, or that of the less into the greater, by which $k$ is transformed into $\lambda_{,}$, $\lambda$ will be transformed into $k$. Therefore \emph{after the first transformation the second is applied, or after the second the first, the Modulus $k$ returns into itself, or the first and second transformations applied successively, in whichever order one pleases, yield Multiplication.}

Let us call $M'$ the Multiplier which depends on $\lambda$ in the same way as $M_{,}$ on $k$; $M_{,}'$ the Multiplier which depends on $\lambda_{,}$ in the same way as $M$ on $k$; so that the equations are obtained:

\[
\begin{gathered}
\frac{dy}{\sqrt{1-y^{2}}}.\,\frac{1}{\sqrt{1-\lambda^{2}y^{2}}}=\frac{dx}{M\sqrt{1-x^{2}}\sqrt{1-k^{2}x^{2}}}\\[2ex]
\frac{dz}{\sqrt{1-z^{2}}}.\,\frac{1}{\sqrt{1-k^{2}z^{2}}}=\frac{dy}{M'\sqrt{1-y^{2}}\sqrt{1-\lambda^{2}z^{2}}},
\end{gathered}
\]\ednote{In the last radical of the second equation the print has $z^{2}$ where $y^{2}$ is expected; an apparent misprint.}
of which the one corresponds to the transformation of the Modulus $k$ into the Modulus $\lambda$ by the first transformation, the other to the transformation of the Modulus $\lambda$ into the Modulus $k$ by the second transformation. From these equations there results:

\[
\frac{dz}{\sqrt{(1-z^{2})(1-k^{2}z^{2})}}=\frac{dx}{MM'\sqrt{(1-x^{2})(1-k^{2}x^{2})}};\;\text{whence }z=\sin\operatorname{am}\left(\frac{u}{MM'}\right).
\]

\origpage{61}
But from the equation $\Lambda_{,}=\dfrac{K}{M_{,}}$, on changing $k$ into $\lambda$, by which change $K$ passes into $\Lambda$, $\lambda_{,}$ into $k$, $\Lambda_{,}$ into $K$, $M_{,}$ into $M'$, there is obtained $K=\dfrac{\Lambda}{M'}$, and this equation, compared with the equation $\Lambda=\dfrac{K}{nM}$, gives $\dfrac{1}{MM'}=n$, whence:

\[
\frac{dz}{\sqrt{(1-z^{2})(1-k^{2}z^{2})}}=\frac{n\,dx}{\sqrt{(1-x^{2})(1-k^{2}x^{2})}}.
\]

In the same way from the equation $\Lambda=\dfrac{K}{nM}$, on changing $k$ into $\lambda_{,}$, by which change $K$ passes into $\Lambda_{,}$, $\lambda$ into $k$, $\Lambda$ into $K$, $M_{,}$ into $M_{,}'$, there results $K=\dfrac{\Lambda_{,}}{nM_{,}'}$, and this equation, compared with the equation $\Lambda_{,}=\dfrac{K}{M_{,}}$, gives $\dfrac{1}{M_{,}M_{,}'}=n$; whence we see that in those two cases, after two transformations applied successively, the Argument is multiplied by the number $n$.

Where after the transformation of the Modulus $k$ into the Modulus $\lambda$ the Modulus $\lambda$ is transformed again into the Modulus $k$, so that Multiplication results, we shall call this transformation the one \emph{supplementary to multiplication} of that former, or simply \emph{supplementary.}

Let us set down, both by way of example and for use in what follows, the formulae for the transformation \emph{supplementary to the first}, that is, of the Modulus $\lambda$ into the Modulus $k$, which will be the second of $\lambda$ itself, yet only under the one imaginary form, since the reduction to the real is ready to hand. We obtain these formulae at once when in those which were set forth above for the second transformation of the Modulus $k$ (see Table II. A. §.\ 24) we put $\lambda$ in place of $k$, $k$ in place of $\lambda$, $\dfrac{u}{M}$ in place of $u$, $M'=\dfrac{1}{nM}$ in place of $M_{,}$, whence $\dfrac{u}{MM'}=nu$ in place of $\dfrac{u}{M}$. In these formulae, but in these only, the Modulus $\lambda$ will be understood, unless the Modulus $k$ is expressly added; moreover, for brevity's sake I have put $y=\sin\operatorname{am}\left(\dfrac{u}{M},\lambda\right)$; the number $q$ is to be given, as above, the values: $1,\,2,\,3,\,\ldots,\,\dfrac{n-1}{2}$.\ ---

\origpage{62}
\section*{FORMULAE FOR THE TRANSFORMATION OF THE MODULUS $\lambda$ INTO THE MODULUS $k$, OR THE ONE SUPPLEMENTARY TO THE FIRST.}

\subsection*{27.}

\[
k=\lambda^{n}\left\{\sin\operatorname{coam}\frac{2i\Lambda'}{n}\sin\operatorname{coam}\frac{4i\Lambda'}{n}.\,.\,.\,.\sin\operatorname{coam}\frac{(n-1)i\Lambda'}{n}\right\}{}^{*}
\]

\[
k'=\frac{\lambda'^{n}}{\left\{\Delta\operatorname{am}\dfrac{2i\Lambda'}{n}\Delta\operatorname{am}\dfrac{4i\Lambda'}{n}.\,.\,.\,.\Delta\operatorname{am}\dfrac{(n-1)i\Lambda'}{n}\right\}{}^{*}}
\]

\[
\frac{1}{nM}=\left\{\frac{\sin\operatorname{coam}\dfrac{2i\Lambda'}{n}\sin\operatorname{coam}\dfrac{4i\Lambda'}{n}.\,.\,.\,.\sin\operatorname{coam}\dfrac{(n-1)i\Lambda'}{n}}{\sin\operatorname{am}\dfrac{2i\Lambda'}{n}\sin\operatorname{am}\dfrac{4i\Lambda'}{n}\ldots\sin\operatorname{am}\dfrac{(n-1)i\Lambda'}{n}}\right\}^{2}
\]

\[
\sin\operatorname{am}(nu,k)=\frac{nMy\left(1-\dfrac{yy}{\sin^{2}\operatorname{am}\dfrac{2i\Lambda'}{n}}\right)\left(1-\dfrac{yy}{\sin^{2}\operatorname{am}\dfrac{4i\Lambda'}{n}}\right)\ldots\left(1-\dfrac{yy}{\sin^{2}\operatorname{am}\dfrac{(n-1)i\Lambda'}{n}}\right)}{\left(1-\dfrac{yy}{\sin^{2}\operatorname{am}\dfrac{i\Lambda'}{n}}\right)\left(1-\dfrac{yy}{\sin^{2}\operatorname{am}\dfrac{3i\Lambda'}{n}}\right)\ldots\left(1-\dfrac{yy}{\sin^{2}\operatorname{am}\dfrac{(n-2)i\Lambda'}{n}}\right)}
\]

\[
=\sqrt{\frac{\lambda^{n}}{k}}\sin\operatorname{am}\frac{u}{M}\sin\operatorname{am}\left(\frac{u}{M}+\frac{4i\Lambda'}{n}\right)\sin\operatorname{am}\left(\frac{u}{M}+\frac{8i\Lambda'}{n}\right)\ldots\sin\operatorname{am}\left(\frac{u}{M}+\frac{4(n-1)i\Lambda'}{n}\right)
\]

\[
\cos\operatorname{am}(nu,k)=\frac{\sqrt{1-yy}\left(1-\dfrac{yy}{\sin^{2}\operatorname{coam}\dfrac{2i\Lambda'}{n}}\right)\left(1-\dfrac{yy}{\sin^{2}\operatorname{coam}\dfrac{4i\Lambda'}{n}}\right)\ldots\left(1-\dfrac{yy}{\sin^{2}\operatorname{coam}\dfrac{(n-1)i\Lambda'}{n}}\right)}{\left(1-\dfrac{yy}{\sin^{2}\operatorname{am}\dfrac{i\Lambda'}{n}}\right)\left(1-\dfrac{yy}{\sin^{2}\operatorname{am}\dfrac{3i\Lambda'}{n}}\right)\ldots\left(1-\dfrac{yy}{\sin^{2}\operatorname{am}\dfrac{(n-2)i\Lambda'}{n}}\right)}
\]

\[
=\sqrt{\frac{k'\lambda^{n}}{k\lambda'^{n}}}\cos\operatorname{am}\frac{u}{M}\cos\operatorname{am}\left(\frac{u}{M}+\frac{4i\Lambda'}{n}\right)\cos\operatorname{am}\left(\frac{u}{M}+\frac{8i\Lambda'}{n}\right)\ldots\cos\operatorname{am}\left(\frac{u}{M}+\frac{4(n-1)i\Lambda'}{n}\right)
\]

\[
\Delta\operatorname{am}(nu,k)=\frac{\sqrt{1-\lambda^{2}yy}\left(1-\dfrac{yy}{\sin^{2}\operatorname{coam}\dfrac{i\Lambda'}{n}}\right)\left(1-\dfrac{yy}{\sin^{2}\operatorname{coam}\dfrac{3i\Lambda'}{n}}\right)\ldots\left(1-\dfrac{yy}{\sin^{2}\operatorname{coam}\dfrac{(n-2)i\Lambda'}{n}}\right)}{\left(1-\dfrac{yy}{\sin^{2}\operatorname{am}\dfrac{i\Lambda'}{n}}\right)\left(1-\dfrac{yy}{\sin^{2}\operatorname{am}\dfrac{3i\Lambda'}{n}}\right)\ldots\left(1-\dfrac{yy}{\sin^{2}\operatorname{am}\dfrac{(n-2)i\Lambda'}{n}}\right)}
\]

\[
=\sqrt{\frac{k'}{\lambda'^{n}}}\Delta\operatorname{am}\frac{u}{M}\Delta\operatorname{am}\left(\frac{u}{M}+\frac{4i\Lambda'}{n}\right)\Delta\operatorname{am}\left(\frac{u}{M}+\frac{8i\Lambda'}{n}\right)\ldots\Delta\operatorname{am}\left(\frac{u}{M}+\frac{4(n-1)i\Lambda'}{n}\right)
\]

\origpage{63}
\[
\sqrt{\frac{1-\sin\operatorname{am}(nu,k)}{1+\sin\operatorname{am}(nu,k)}}=\sqrt{\frac{1-y}{1+y}}.\frac{\left(1-\dfrac{y}{\sin\operatorname{coam}\dfrac{2i\Lambda'}{n}}\right)\left(1-\dfrac{y}{\sin\operatorname{coam}\dfrac{4i\Lambda'}{n}}\right)\ldots\left(1-\dfrac{y}{\sin\operatorname{coam}\dfrac{(n-1)i\Lambda'}{n}}\right)}{\left(1+\dfrac{y}{\sin\operatorname{coam}\dfrac{2i\Lambda'}{n}}\right)\left(1+\dfrac{y}{\sin\operatorname{coam}\dfrac{4i\Lambda'}{n}}\right)\ldots\left(1+\dfrac{y}{\sin\operatorname{coam}\dfrac{(n-1)i\Lambda'}{n}}\right)}
\]

\[
\sqrt{\frac{1-k\sin\operatorname{am}(nu,k)}{1+k\sin\operatorname{am}(nu,k)}}=\sqrt{\frac{1-\lambda y}{1+\lambda y}}.\frac{\left(1-\dfrac{y}{\sin\operatorname{coam}\dfrac{i\Lambda'}{n}}\right)\left(1-\dfrac{y}{\sin\operatorname{coam}\dfrac{3i\Lambda'}{n}}\right)\ldots\left(1-\dfrac{y}{\sin\operatorname{coam}\dfrac{(n-2)i\Lambda'}{n}}\right)}{\left(1+\dfrac{y}{\sin\operatorname{coam}\dfrac{i\Lambda'}{n}}\right)\left(1+\dfrac{y}{\sin\operatorname{coam}\dfrac{3i\Lambda'}{n}}\right)\ldots\left(1+\dfrac{y}{\sin\operatorname{coam}\dfrac{(n-2)i\Lambda'}{n}}\right)}
\]

\[
\sin\operatorname{am}(nu,k)=\frac{\lambda y}{knM}-\frac{2y}{knM}\Sigma\frac{\cos\operatorname{am}\dfrac{(2q-1)i\Lambda'}{n}\Delta\operatorname{am}\dfrac{(2q-1)i\Lambda'}{n}}{\sin^{2}\operatorname{am}\dfrac{(2q-1)i\Lambda'}{n}-yy}
\]

\[
\cos\operatorname{am}(nu,k)=\frac{(-1)^{\frac{n-1}{2}}\lambda\sqrt{1-yy}}{knM}+\frac{2\sqrt{1-yy}}{iknM}\Sigma\frac{(-1)^{q}\sin\operatorname{am}\dfrac{(2q-1)i\Lambda'}{n}\Delta\operatorname{am}\dfrac{(2q-1)i\Lambda'}{n}}{\sin^{2}\operatorname{am}\dfrac{(2q-1)i\Lambda'}{n}-yy}
\]

\[
\Delta\operatorname{am}(nu,k)=\frac{(-1)^{\frac{n-1}{2}}}{nM}\sqrt{1-\lambda^{2}yy}+\frac{2\sqrt{1-\lambda^{2}yy}}{inM}\Sigma\frac{(-1)^{q}\sin\operatorname{am}\dfrac{(2q-1)i\Lambda'}{n}\cos\operatorname{am}\dfrac{(2q-1)i\Lambda'}{n}}{\sin^{2}\operatorname{am}\dfrac{(2q-1)i\Lambda'}{n}-yy}
\]

\[
\operatorname{tg}\operatorname{am}(nu,k)=\frac{\lambda'}{k'nM}.\frac{y}{\sqrt{1-yy}}+\frac{2y\sqrt{1-yy}}{k'nM}\Sigma\frac{(-1)^{q}\Delta\operatorname{am}\dfrac{2qi\Lambda'}{n}}{\cos^{2}\operatorname{am}\dfrac{2qi\Lambda'}{n}-\Delta^{2}\operatorname{am}\dfrac{2qi\Lambda'}{n}\sin^{2}\operatorname{am}u}.
\]\ednote{The last term of the denominator is printed with $\sin^{2}\operatorname{am}u$, where $yy$ is expected (the author's Corrigenda give $yy$); a hand correction in ink is not transcribed.}

The general analytic Theorem concerning that transformation supplementary to the first I communicated to Cl.\ Legendre as early as the beginning of August 1827, and he also wished to make mention of it in the Note cited above (Nova Astr.\ a.\ 1827.\ no.\ 130). A similar system of formulae could have been established for the other transformation, supplementary to the second, or the transformation of the Modulus $\lambda_{,}$ into the Modulus $k$. That all these things may become clearer, I have chosen to add a table setting before the eye the fundamental formulae for the first and second transformations, and for their complementary and supplementary ones.

\origpage{64}
Moreover, of the number of imaginary transformations, each one has its supplementary one pertaining to Multiplication. Let us suppose, as is permitted, that the numbers $m$, $m'$ of §.\ 20 have no common factor; let further $m\mu'-\mu m'=1$, with $\mu$, $\mu'$ denoting positive or negative integers. Now if in our general formulae proposed in §.\ 20 sqq.\ on transformation one puts $\omega=\dfrac{\mu K+\mu' iK'}{nM}$, and $k$ and $\lambda$ are interchanged, one obtains the formulae which pertain to the supplementary of the transformation. With $m=1$, $m'=0$ put, there results $\mu=0$, $\mu'=1$, whence $\dfrac{\mu K+\mu' iK'}{nM}=\dfrac{iK'}{nM}=\dfrac{i\Lambda'}{n}$, which affords the supplementary of the first, as we have seen.

\section*{GENERAL ANALYTIC FORMULAE FOR THE MULTIPLICATION OF ELLIPTIC FUNCTIONS.}

\subsection*{28.}

From the two Supplementary Transformations the formulae for Multiplication itself can be composed, or the formulae by which the elliptic functions of the Argument $nu$ are expressed by elliptic functions of the Argument $u$. To demonstrate this by an example, let us compose the Multiplication from the first transformation and its supplementary. To this end let the formula be recalled:

\[
\sin\operatorname{am}\left(\frac{u}{M},\lambda\right)=(-1)^{\frac{n-1}{2}}\sqrt{\frac{k^{n}}{\lambda}}\sin\operatorname{am}u\sin\operatorname{am}\left(u+\frac{4K}{n}\right)\sin\operatorname{am}\left(u+\frac{8K}{n}\right)\ldots\sin\operatorname{am}\left(u+\frac{4(n-1)K}{n}\right),
\]

which can also be represented in this manner:

\[
(-1)^{\frac{n-1}{2}}\sin\operatorname{am}\left(\frac{u}{M},\lambda\right)=\sqrt{\frac{k^{n}}{\lambda}}\,\Pi\sin\operatorname{am}\left(u+\frac{2mK}{n}\right),
\]

where $m$ denotes the numbers $0$, $\pm1$, $\pm2$, $\ldots$, $\pm\dfrac{n-1}{2}$. In this formula let us put $u+\dfrac{2m'iK'}{n}$ in place of $u$, whence $\dfrac{u}{M}$ passes into $\dfrac{u}{M}+\dfrac{2m'iK'}{n}=\dfrac{u}{M}+\dfrac{2m'i\Lambda'}{n}$: there results
\ednote{Printed $\frac{2m'iK'}{n}$ in the first term, where $\frac{2m'iK'}{nM}$ is expected so that it equals $\frac{2m'i\Lambda'}{n}$; an apparent misprint (the factor $M$ is missing).}

\[
(-1)^{\frac{n-1}{2}}\sin\operatorname{am}\left(\frac{u}{M}+\frac{2m'i\Lambda'}{n},\lambda\right)=\sqrt{\frac{k^{n}}{\lambda}}\,\Pi\sin\operatorname{am}\left(u+\frac{2mK+2m'iK'}{n}\right).
\]

\origpage{65}
Now when the values $0$, $\pm1$, $\pm2$, $\ldots$, $\pm\dfrac{n-1}{n}$\ednote{Printed $\pm\frac{n-1}{n}$ instead of $\pm\frac{n-1}{2}$; an apparent misprint.} are assigned to $m'$ itself as well, so that to both $m$, $m'$ those values are appropriate, on forming the product we obtain:

\[
(-1)^{\frac{n-1}{2}}\,\Pi\sin\operatorname{am}\left(\frac{u}{M}+\frac{2m'i\Lambda'}{n},\lambda\right)=\sqrt{\frac{k^{nn}}{\lambda^{n}}}\,\Pi\sin\operatorname{am}\left(u+\frac{2mK+2m'iK'}{n}\right);
\]
where in the one product the number $m'$, in the other both $m$, $m'$, are to be assigned the values $0$, $\pm1$, $\pm2$, $.\,.$, $\pm\dfrac{n-1}{2}$.

But we saw in the preceding §$^{\mathrm{o}}$ that:

\[
\sin\operatorname{am}(nu,k)=\sqrt{\frac{\lambda^{n}}{k}}\sin\operatorname{am}\left(\frac{u}{M}\right)\sin\operatorname{am}\left(\frac{u}{M}+\frac{4i\Lambda'}{n}\right)\sin\operatorname{am}\left(\frac{u}{M}+\frac{8i\Lambda'}{n}\right)\ldots\sin\operatorname{am}\left(\frac{u}{M}+\frac{4(n-1)i\Lambda'}{n}\right)\left\{\operatorname{Mod}\lambda\right\},
\]
a formula which can also be represented thus:

\[
\sin\operatorname{am}(nu,k)=\sqrt{\frac{\lambda^{n}}{k}}\,\Pi\sin\operatorname{am}\left(\frac{u}{M}+\frac{2m'i\Lambda'}{n},\lambda\right),
\]
whence now:

\[
\text{1)}\quad\sin\operatorname{am}nu=(-1)^{\frac{n-1}{2}}\sqrt{k^{nn-1}}\,\Pi\sin\operatorname{am}\left(u+\frac{2mK+2m'iK'}{n}\right).
\]
In the same way one finds:

\[
\begin{gathered}
\text{2)}\quad\cos\operatorname{am}nu=\sqrt{\left(\frac{k}{k'}\right)^{nn-1}}\,\Pi\cos\operatorname{am}\left(u+\frac{2mK+2m'iK'}{n}\right)\\[1.5ex]
\text{3)}\quad\Delta\operatorname{am}nu=\sqrt{\left(\frac{1}{k'}\right)^{nn-1}}\,\Pi\Delta\operatorname{am}\left(u+\frac{2mK+2m'iK'}{n}\right).
\end{gathered}
\]
These formulae are also easily reduced to this form:

\[
\begin{gathered}
\text{4)}\quad\sin\operatorname{am}nu=u\sin\operatorname{am}u\,\Pi\frac{\left(1-\dfrac{\sin^{2}\operatorname{am}u}{\sin^{2}\operatorname{am}\dfrac{2mK+2m'iK'}{n}}\right)}{\left(1-k^{2}\sin^{2}\operatorname{am}\dfrac{2mK+2m'iK'}{n}\sin^{2}\operatorname{am}u\right)}\\[2ex]
\text{5)}\quad\cos\operatorname{am}nu=\cos\operatorname{am}u\,\Pi\frac{1-\dfrac{\sin^{2}\operatorname{am}u}{\sin^{2}\operatorname{coam}\dfrac{2mK+2m'iK'}{n}}}{1-k^{2}\sin^{2}\operatorname{am}\dfrac{2mK+2m'iK'}{n}\sin^{2}\operatorname{am}u}
\end{gathered}
\]\ednote{In formula 4) the print has $u$ where $n$ is expected before the first $\sin\operatorname{am}u$; an apparent misprint.}

\origpage{66}
\[
\begin{gathered}
\text{6)}\quad\Delta\operatorname{am}nu=\Delta\operatorname{am}u\,\Pi\frac{1-k^{2}\sin^{2}\operatorname{coam}\dfrac{2mK+2m'iK'}{n}\sin^{2}\operatorname{am}u}{1-k^{2}\sin^{2}\operatorname{am}\dfrac{2mK+2m'iK'}{n}\sin^{2}\operatorname{am}u}.
\end{gathered}
\]
To these it pleases me to add the following:

\[
\begin{gathered}
\text{7)}\quad\Pi\sin^{2}\operatorname{am}\frac{2mK+2m'iK'}{n}=\frac{(-1)^{\frac{n-1}{2}n}}{k^{\frac{nn-1}{2}}}\\[1.5ex]
\text{8)}\quad\Pi\cos^{2}\operatorname{am}\frac{2mK+2m'iK'}{n}=\left(\frac{k'}{k}\right)^{\frac{nn-1}{2}}\\[1.5ex]
\text{9)}\quad\Pi\Delta^{2}\operatorname{am}\frac{2mK+2m'iK'}{n}=k'^{\frac{nn-1}{2}}.
\end{gathered}
\]
In the last six formulae only the positive values $0$, $1$, $2$, $3$, $\ldots$, $\dfrac{n-1}{2}$ are appropriate to the number $m$, yet so that whenever $m=0$, only the positive values $1$, $2$, $3$, $\ldots$, $\dfrac{n-1}{2}$ are assigned to $m'$ itself. Both these and other formulae for Multiplication had already been proposed earlier by Cl.\ \emph{Abel}, with the necessary changes made, whence it was permitted to us to be briefer.

\section*{ON THE PROPERTIES OF MODULAR EQUATIONS.}

\subsection*{29.}

Since in the same manner $\lambda$ depends on $k$ and $k$ on $\lambda_{,}$, and likewise $\lambda_{,}'$ on $k'$, $k'$ on $\lambda_{,}'$; it is evident, when according to the same law one constructs scales of Moduli which can be transformed into one another, the one containing the Modulus $k$, the other its Complement $k'$, that in them the terms will follow one another in the same order:

\[
\begin{gathered}
.\,.\,.\,.,\ \lambda,\ k,\ \lambda_{,},\ \ldots\\
\ldots,\ \lambda_{,}',\ k',\ \lambda',\ .\,.\,.\,.
\end{gathered}
\]
This was already observed earlier by Cl.\ Legendre in the transformations of the second and third order and confirmed by calculation. Since similar things hold for all transformed and imaginary Moduli, it is evident that, with $\lambda$ denoting any transformed Modulus, the algebraic equations between $k$ and $\lambda$, or between $u=\sqrt[4]{k}$ and $v=\sqrt[4]{\lambda}$, which we have called \emph{Modular Equations}, remain unchanged,

\origpage{67}
1) when $k$ and $\lambda$ are interchanged,

2) when $k'$ is put in place of $k$, $\lambda'$ in place of $\lambda$.

The former we have already observed above in the Modular equations which pertain to the transformations of the third and fifth order:

\[
\begin{gathered}
\text{1)}\quad u^{4}-v^{4}+2uv(1-u^{2}v^{2})=0\\
\text{2)}\quad u^{6}-v^{6}+5u^{2}v^{2}(u^{2}-v^{2})+4uv(1-u^{4}v^{4})=0
\end{gathered}
\]
and by means of that observation we have exhibited the algebraic expressions for the supplementary transformations. That the latter also may be verified by these examples, let us transform those equations into others between $kk=u^{8}$ and $\lambda\lambda=v^{8}$, which is not done without a lengthy calculation. This having been performed, the equations are obtained:

\[
\begin{gathered}
\text{1)}\quad(k^{2}-\lambda^{2})^{4}=128k^{2}\lambda^{2}(1-k^{2})(1-\lambda^{2})(2-k^{2}-\lambda^{2}+2k^{2}\lambda^{2})\\
\text{2)}\quad(k^{2}-\lambda^{2})^{6}=512k^{2}\lambda^{2}(1-k^{2})(1-\lambda^{2})\left\{L-L'k^{2}+L''k^{4}-L'''k^{6}\right\},
\end{gathered}
\]
provided that in the second one puts:

\[
\begin{aligned}
L&=128-192\lambda^{2}+\ 78\lambda^{4}-\ 7\lambda^{6}\\
L'&=192+252\lambda^{2}-423\lambda^{4}-\ 78\lambda^{6}\\
L''&=\ 78+423\lambda^{2}-252\lambda^{4}-192\lambda^{6}\\
L'''&=\ \ 7-\ 78\lambda^{2}-192\lambda^{4}-128\lambda^{6}.
\end{aligned}
\]\ednote{In the last line the print has $-192\lambda^{4}$; the author's Corrigenda give $+192\lambda^{4}$ (a hand correction in ink appears in this copy).}
These equations pass into a much more convenient form on introducing the quantities $q=1-2k^{2}$, $l=1-2\lambda^{2}$. This done, the proposed equations come out as:

\[
\begin{gathered}
\text{1)}\quad(q-l)^{4}=64(1-qq)(1-ll)\left\{3+ql\right\}\\[1ex]
\text{2)}\quad(q-l)^{6}=256(1-qq)(1-ll)\left\{16ql(9-ql)^{2}+9(45-ql)(q-l)^{2}\right\}\\
=256(1-qq)(1-ll)\left\{405(qq+ll)+486ql-9ql(qq+ll)-270qqll+16q^{3}l^{3}\right\}.
\end{gathered}
\]
These equations remain unchanged when $k'$ is put in place of $k$, $\lambda'$ in place of $\lambda$, whence $q$ passes into $-q$, $l$ into $-l$; which was to be demonstrated.

\emph{Corollary.} Since we have seen that the Modular Equations proposed between $q=1-2k^{2}$ and $l=1-2\lambda^{2}$ assume a sufficiently convenient form, it may be of interest also to expand the functions $K$, $K'$ themselves according to the quantity $q$. This is done not inelegantly by the series:

\[
\begin{aligned}
K&=J\left(1+\frac{q^{2}}{2.4}+\frac{5.5.q^{4}}{2.4.6.8}+\frac{5.5.9.9.q^{6}}{2.4.6.8.10.12}+\ldots\right)\\
&\quad-\frac{\pi}{2J}\left(\frac{q}{2}+\frac{3.3.q^{3}}{2.4.6}+\frac{3.3.7.7.q^{5}}{2.4.6.8.10}+\frac{3.3.7.7.11.11.q^{7}}{2.4.6.8.10.12.14}+.\,.\,.\,.\right)
\end{aligned}
\]

\origpage{68}
\[
\begin{aligned}
K'&=J\left(1+\frac{q^{2}}{2.4}+\frac{5.5.q^{4}}{2.4.6.8}+\frac{5.5.9.9.q^{6}}{2.4.6.8.10.12}+.\,.\,.\,.\right)\\
&\quad+\frac{\pi}{2J}\left(\frac{q}{2}+\frac{3.3.q^{3}}{2.4.6}+\frac{3.3.7.7.q^{5}}{2.4.6.8.10}+\frac{3.3.7.7.11.11.q^{7}}{2.4.6.8.10.12.14}+.\,.\,.\,.\right)
\end{aligned}
\]
where for the sake of brevity it is put

\[
\int_{0}^{\frac{\pi}{2}}\frac{d\phi}{\sqrt{1-\dfrac{1}{2}\sin\phi^{2}}}=J.
\]

\subsection*{30.}

With easier labour, for the transformation of the third order the equation:

\[
u^{4}-v^{4}+2uv(1-u^{2}v^{2})=0
\]
can be so transformed that that correlation between the Moduli and Complements shines forth. For we obtain from it:

\[
\begin{gathered}
(1-u^{4})(1+v^{4})=1-u^{4}v^{4}+2uv(1-u^{2}v^{2})=(1-u^{2}v^{2})(1+uv)^{2}\\
(1+u^{4})(1-v^{4})=1-u^{4}v^{4}-2uv(1-u^{2}v^{2})=(1-u^{2}v^{2})(1-uv)^{2},
\end{gathered}
\]
and on multiplying these equations into one another there results:

\[
(1-u^{8})(1-v^{8})=(1-u^{2}v^{2})^{4}.
\]
Now let:

\[
\begin{gathered}
1-u^{8}=k'k'=u'^{8}\\
1-v^{8}=\lambda'\lambda'=v'^{8},
\end{gathered}
\]
roots having been extracted, there results:

\[
u'^{2}v'^{2}=1-u^{2}v^{2},
\]
or

\[
u^{2}v^{2}+u'^{2}v'^{2}=\sqrt{k\lambda}+\sqrt{k'\lambda'}=1,
\]
a most elegant formula which Cl.\ Legendre has already exhibited. Nor is it proved inelegantly by our analytic formulae. For from them in the case $n=3$ there flows:

\[
\lambda=k^{3}\sin^{4}\operatorname{coam}4\omega;\ \lambda'=\frac{k'^{3}}{\Delta^{4}\operatorname{am}4\omega},
\]

\origpage{69}
whence:

\[
\begin{gathered}
\sqrt{k\lambda}=k^{2}\sin^{2}\operatorname{coam}4\omega=\frac{k^{2}\cos^{2}\operatorname{am}4\omega}{\Delta^{2}\operatorname{am}4\omega}\\[1.5ex]
\sqrt{k'\lambda'}=\frac{k'}{\Delta^{2}\operatorname{am}4\omega},
\end{gathered}
\]
whence, since

\[
k'k'+kk\cos^{2}\operatorname{am}4\omega=1-kk\sin^{2}\operatorname{am}4\omega=\Delta^{2}\operatorname{am}4\omega,
\]
we obtain, what was to be demonstrated:

\[
\sqrt{k\lambda}+\sqrt{k'\lambda'}=1.
\]

That I may derive a simpler equation between $u$, $v$, $u'$, $v'$ for the second example, I proceed thus. The proposed equation:

\[
u^{6}-v^{6}+5u^{2}v^{2}(u^{2}-v^{2})+4uv(1-u^{4}v^{4})=0
\]
I exhibit as follows:

\[
(u^{2}-v^{2})(u^{4}+6u^{2}v^{2}+v^{4})+4uv(1-u^{4}v^{4})=0,
\]
which, it is easily evident, can assume the two following forms:

\[
\begin{gathered}
(u^{2}-v^{2})(u+v)^{4}=-4uv(1-u^{4})(1+v^{4})\\
(u^{2}-v^{2})(u-v)^{4}=-4uv(1+u^{4})(1-v^{4}),
\end{gathered}
\]
and on multiplying these equations into one another there results:

\[
(u^{2}-v^{2})^{6}=16u^{2}v^{2}(1-u^{8})(1-v^{8})=16u^{2}v^{2}u'^{8}v'^{8}.
\]
Since at the same time, as was proved above, $u^{8}$ passes into $u'^{8}$, $v^{8}$ into $v'^{8}$, we obtain also:

\[
(v'^{2}-u'^{2})^{6}=16u'^{2}v'^{2}(1-u'^{8})(1-v'^{8})=16u'^{2}v'^{2}u^{8}v^{8}.
\]
Hence, on division and extraction of roots, is derived:

\[
\frac{u^{2}-v^{2}}{v'^{2}-u'^{2}}=\frac{u'v'}{uv},\ \text{or }uv(u^{2}-v^{2})=u'v'(v'^{2}-u'^{2}),
\]
or

\[
\sqrt[4]{k\lambda}\left(\sqrt{k}-\sqrt{\lambda}\right)=\sqrt[4]{k'\lambda'}\left(\sqrt{\lambda'}-\sqrt{k'}\right).
\]

\origpage{70}
\subsection*{31.}

Another remarkable property of the Modular equations

\[
\begin{gathered}
u^{4}-v^{4}+2uv(1-u^{2}v^{2})=0\\
u^{6}-v^{6}+5u^{2}v^{2}(u^{2}-v^{2})+4uv(1-u^{4}v^{4})=0
\end{gathered}
\]
is found even at first sight, viz.\ that they remain unchanged if in place of $u$, $v$ one puts $\dfrac{1}{u}$, $\dfrac{1}{v}$. That this may be demonstrated generally for Modular equations, let the following be noted, which can be of use for other questions also.

When $y=kx$ is put, one obtains:

\[
\frac{dy}{\sqrt{(1-y^{2})\left(1-\dfrac{y^{2}}{k^{2}}\right)}}=\frac{k\,dx}{\sqrt{(1-x^{2})(1-k^{2}x^{2})}},
\]
whence, since simultaneously $x=0$, $y=0$:

\[
\int_{0}^{y}\frac{dy}{\sqrt{(1-y^{2})\left(1-\dfrac{y^{2}}{k^{2}}\right)}}=k\int_{0}^{x}\frac{dx}{\sqrt{(1-x^{2})(1-k^{2}x^{2})}}.
\]
Hence on putting

\[
\int_{0}^{x}\frac{dx}{\sqrt{(1-x^{2})(1-k^{2}x^{2})}}=u,\ \text{there results:}
\]

\[
\int_{0}^{y}\frac{dy}{\sqrt{(1-y^{2})\left(1-\dfrac{y^{2}}{k^{2}}\right)}}=ku,
\]
whence $x=\sin\operatorname{am}(u,k)$, $y=\sin\operatorname{am}\left(ku,\dfrac{1}{k}\right)$. Hence arises the equation:

\[
\begin{gathered}
\sin\operatorname{am}\left(ku,\frac{1}{k}\right)=k\sin\operatorname{am}(u,k),\ \text{whence also}\\[1.5ex]
\cos\operatorname{am}\left(ku,\frac{1}{k}\right)=\Delta\operatorname{am}(u,k)\\[1.5ex]
\Delta\operatorname{am}\left(ku,\frac{1}{k}\right)=\cos\operatorname{am}(u,k)\\[1.5ex]
\operatorname{tg}\operatorname{am}\left(ku,\frac{1}{k}\right)=\frac{k}{k'}\cos\operatorname{coam}(u,k)
\end{gathered}
\]

\origpage{71}
\[
\begin{gathered}
\sin\operatorname{coam}\left(ku,\frac{1}{k}\right)=\frac{1}{\sin\operatorname{coam}(u,k)}\\[1.5ex]
\cos\operatorname{coam}\left(ku,\frac{1}{k}\right)=ik'\operatorname{tg}\operatorname{am}(u,k)\\[1.5ex]
\Delta\operatorname{coam}\left(ku,\frac{1}{k}\right)=\frac{ik'}{k\cos\operatorname{am}(u,k)}\\[1.5ex]
\operatorname{tg}\operatorname{coam}\left(ku,\frac{1}{k}\right)=\frac{-i}{\cos\operatorname{coam}(u,k)}
\end{gathered}
\]
Further, by putting $iu$ in place of $u$, since the Complement of the Modulus $\dfrac{1}{k}$ becomes $\dfrac{ik'}{k}$, we obtain with the aid of the formulae of §$^{\mathrm{i}}$ 19:

\[
\begin{gathered}
\sin\operatorname{am}\left(ku,\frac{ik'}{k}\right)=\cos\operatorname{coam}(u,k')\\[1.5ex]
\cos\operatorname{am}\left(ku,\frac{ik'}{k}\right)=\sin\operatorname{coam}(u,k')\\[1.5ex]
\Delta\operatorname{am}\left(ku,\frac{ik'}{k}\right)=\frac{1}{\Delta\operatorname{am}(u,k)}\\[1.5ex]
\operatorname{tg}\operatorname{am}\left(ku,\frac{ik'}{k}\right)=\operatorname{cotg}\operatorname{coam}(u,k')\\[1.5ex]
\sin\operatorname{coam}\left(ku,\frac{ik'}{k}\right)=\cos\operatorname{am}(u,k')\\[1.5ex]
\cos\operatorname{coam}\left(ku,\frac{ik'}{k}\right)=\sin\operatorname{am}(u,k')\\[1.5ex]
\Delta\operatorname{coam}\left(ku,\frac{ik'}{k}\right)=\frac{\Delta\operatorname{am}(u,k')}{k}\\[1.5ex]
\operatorname{tg}\operatorname{coam}\left(ku,\frac{ik'}{k}\right)=\operatorname{cotg}\operatorname{am}(u,k').
\end{gathered}
\]
Now let us investigate what $K$, $K'$ or arg.\ am$\left(\dfrac{\pi}{2},k\right)$, arg.\ am$\left(\dfrac{\pi}{2},k'\right)$ become, if in place of $k$ one puts $\dfrac{1}{k}$; or let us investigate the value of the expressions arg.\ am$\left(\dfrac{\pi}{2},\dfrac{1}{k}\right)$, arg.\ am$\left(\dfrac{\pi}{2},\dfrac{ik'}{k}\right)$, which expressions, in the notation employed by Cl.\ Legendre, would be $F^{\mathrm{I}}\left(\dfrac{1}{k}\right)$, $F^{\mathrm{I}}\left(\dfrac{ik'}{k}\right)$. First, however, there results:

\[
\text{arg.\ am}\left(\frac{\pi}{2},\frac{1}{k}\right)=\int_{0}^{1}\frac{dy}{\sqrt{(1-y^{2})\left(1-\dfrac{y^{2}}{k^{2}}\right)}}=\int_{0}^{k}\frac{dy}{\sqrt{(1-y^{2})\left(1-\dfrac{y^{2}}{k^{2}}\right)}}+\int_{k}^{1}\frac{dy}{\sqrt{(1-y^{2})\left(1-\dfrac{y^{2}}{k^{2}}\right)}}.
\]

\origpage{72}
With $y=kx$ put, there results

\[
\int_{0}^{k}\frac{dy}{\sqrt{(1-y^{2})\left(1-\dfrac{y^{2}}{k^{2}}\right)}}=k\int_{0}^{1}\frac{dx}{\sqrt{(1-x^{2})(1-k^{2}x^{2})}}=kK.
\]
That the other integral $\displaystyle\int_{k}^{1}\frac{dy}{\sqrt{(1-y^{2})\left(1-\frac{y^{2}}{k^{2}}\right)}}$ may be derived, let us put $y=\sqrt{1-k'k'x^{2}}$, whence

\[
\frac{dy}{\sqrt{(1-y^{2})\left(\dfrac{y^{2}}{k^{2}}-1\right)}}=\frac{-k\,dx}{\sqrt{(1-x^{2})(1-k'k'x^{2})}}.
\]
Now since $x$ increases from $0$ up to $1$ while $y$ decreases from $1$ down to $k$, we obtain:

\[
\int_{k}^{1}\frac{dy}{\sqrt{(1-y^{2})\left(1-\dfrac{y^{2}}{k^{2}}\right)}}=-i\int_{k}^{1}\frac{dy}{\sqrt{(1-y^{2})\left(\dfrac{y^{2}}{k^{2}}-1\right)}}=i\int_{0}^{1}\frac{k\,dx}{\sqrt{(1-x^{2})(1-k'k'x^{2})}}=ikK'.
\]
Hence results $\text{arg.\ am}\left(\dfrac{\pi}{2},\dfrac{1}{k}\right)=k\left\{\text{arg.\ am}\left(\dfrac{\pi}{2},k\right)+i\,\text{arg.\ am}\left(\dfrac{\pi}{2},k'\right)\right\}=k\left\{K+iK'\right\}$, or when $k$ is changed into $\dfrac{1}{k}$, $K$ passes into $k\left\{K+iK'\right\}$.

Secondly, with $y=\cos\phi$ put, there results:

\[
\int_{0}^{1}\frac{dy}{\sqrt{(1-y^{2})\left(1+\dfrac{k'k'}{kk}y^{2}\right)}}=k\int_{0}^{\frac{\pi}{2}}\frac{d\phi}{\sqrt{1-k'k'\sin\phi^{2}}}=kK',\ \text{whence:}
\]

\[
\text{arg.\ am}\left(\frac{\pi}{2},\frac{ik'}{k}\right)=k\,\text{arg.\ am}\left(\frac{\pi}{2},k'\right)=kK',
\]
or when $k$ is changed into $\dfrac{1}{k}$, $K'$ passes into $kK'$.

In general, therefore, when $k$ is changed into $\dfrac{1}{k}$, $mK+im'K'$ passes into $k\left\{mK+(m+m')iK'\right\}$, whence $\sin\operatorname{coam}\left\{\dfrac{p(mK+m'iK')}{n},k\right\}$ into $\sin\operatorname{coam}\left\{\dfrac{kp\left(mK+(m+m')iK'\right)}{n},\dfrac{1}{k}\right\}$, which comes about from the formula $\sin\operatorname{coam}\left(ku,\dfrac{1}{k}\right)=\dfrac{1}{\sin\operatorname{coam}(u,k)}$:

\[
\sin\operatorname{coam}\left\{\frac{kp\left(mK+(m+m')iK'\right)}{n},\frac{1}{k}\right\}=\frac{1}{\sin\operatorname{coam}\left\{\dfrac{p\left(mK+(m+m')iK'\right)}{n},k\right\}}.
\]

\origpage{73}

Now therefore, with $\dfrac{mK+m'iK'}{n}=\omega$, $\dfrac{mK+(m+m')iK'}{n}=\omega_{,}$ put, the expression

\[
\lambda=k^{n}\left\{\sin\operatorname{coam}2\omega\sin\operatorname{coam}4\omega\sin\operatorname{coam}6\omega.\,.\,.\,.\sin\operatorname{coam}(n-1)\omega\right\}^{4},
\]
when $k$ is changed into $\dfrac{1}{k}$, passes into this:

\[
\frac{1}{k^{n}\left\{\sin\operatorname{coam}2\omega_{,}\,\sin\operatorname{coam}4\omega_{,}\,\sin\operatorname{coam}6\omega_{,}\,\ldots\,\sin\operatorname{coam}(n-1)\omega_{,}\right\}^{4}}=\frac{1}{\mu},
\]
where $\mu$ is itself a root of the Modular equation, or is among the number of Moduli into which the proposed Modulus $k$ can be transformed by a transformation of the $n^{\mathrm{ti}}$ order. For among the values which $\omega$ can assume, in order that a transformed Modulus result, that $\omega_{,}$ will also be. Whence now the reason is evident why in general Modular Equations must remain unchanged when $k$ is changed into $\dfrac{1}{k}$, $\lambda$ into $\dfrac{1}{\lambda}$.

I shall note in addition that, when according to the same law of transformation $k$ is simultaneously transformed into $k^{(m)}$, $\lambda$ into $\lambda^{(m)}$, whenever $k^{(m)}$ is put in place of $k$, $\lambda$ also passes into $\lambda^{(m)}$; whence the Modular equations, when $k$ is simultaneously changed into $k^{(m)}$ and $\lambda$ into $\lambda^{(m)}$, must remain unchanged. Thus e.g.\ the equation $\sqrt{k\lambda}+\sqrt{k'\lambda'}=1$, which is for the transformation of the third order, must remain unchanged when in place of $k$, $\lambda$ there is put resp.\ $\dfrac{1-k'}{1+k'}$, $\dfrac{1-\lambda'}{1+\lambda'}$, whence in place of $k'$, $\lambda'$ there will be put $\dfrac{2\sqrt{k'}}{1+k'}$, $\dfrac{2\sqrt{\lambda'}}{1+\lambda'}$, which it is known comes about by the transformation of the second order.
For the equation $\sqrt{k\lambda}+\sqrt{k'\lambda'}=1$ passes into this:

\[
\begin{gathered}
\sqrt{\frac{(1-k')(1-\lambda')}{(1+k')(1+\lambda')}}+\frac{2\sqrt[4]{k'\lambda'}}{\sqrt{(1+k')(1+\lambda')}}=1,\quad\text{or}\\[1.5ex]
2\sqrt[4]{k'\lambda'}=\sqrt{(1+k')(1+\lambda')}-\sqrt{(1-k')(1-\lambda')}.
\end{gathered}
\]
On multiplying this into itself there results:

\[
4\sqrt{k'\lambda'}=2(1+k'\lambda')-2k\lambda,\quad\text{or }k\lambda=1+k'\lambda'-2\sqrt{k'\lambda'},
\]
which, roots having been extracted, returns to the proposed:

\[
\sqrt{k\lambda}=1-\sqrt{k'\lambda'}\quad\text{or }\sqrt{k\lambda}+\sqrt{k'\lambda'}=1.
\]
This example was already proposed by Cl.\ Legendre. In general, however, concerning the composition of transformations it can be proved that, when two or more transformations are applied successively, one arrives at the same result, in whatever order they are applied.

\origpage{74}

\subsection*{32.}

But among the properties of the Modular Equations that seems to me the most memorable and singular which I think I perceive, \emph{that they all satisfy the same Differential Equation of the Third Order.} The investigation of this, however, will have to be taken up from somewhat further back.

It is sufficiently well known ${}^{*)}$, with $aK+bK'=Q$ put, that:

\[
k(1-k^{2})\frac{d^{2}Q}{dk^{2}}+(1-3k^{2})\frac{dQ}{dk}=kQ,
\]
with $a$, $b$ denoting any Constants whatever. Thus also with $a'K+b'K'=Q'$ put, with $a'$, $b'$ denoting any other Constants whatever, there will be
\[
k(1-k^{2})\frac{d^{2}Q'}{dk^{2}}+(1-3k^{2})\frac{dQ'}{dk}=kQ'.
\]
These equations being combined, one obtains:

\[
k(1-k^{2})\left\{Q\frac{d^{2}Q'}{dk^{2}}-Q'\frac{d^{2}Q}{dk^{2}}\right\}+(1-3k^{2})\left\{Q\frac{dQ'}{dk}-Q'\frac{dQ}{dk}\right\}=0,
\]
whence, integration having been performed:

\[
k(1-k^{2})\left\{Q\frac{dQ'}{dk}-Q'\frac{dQ}{dk}\right\}=(ab'-a'b)k(1-k^{2})\left\{K\frac{dK'}{dk}-K'\frac{dK}{dk}\right\}=(ab'-a'b)\,C.
\]
The Constant $C$ was found by Cl.\ Legendre from a special case to be $=-\dfrac{\pi}{2}$, whence now

\[
Q\frac{dQ'}{dk}-Q'\frac{dQ}{dk}=\frac{-\frac{1}{2}\pi(ab'-a'b)}{k(1-k^{2})},\quad\text{or}
\]

\[
d\frac{Q'}{Q}=\frac{-\frac{1}{2}\pi(ab'-a'b)\,dk}{k(1-k^{2})QQ}.
\]
Similarly, with $\lambda$ denoting any other Modulus whatever, with $\alpha\Lambda+\beta\Lambda'=L$, $\alpha'\Lambda+\beta'\Lambda'=L'$ put, there will be,

\[
d\frac{L'}{L}=\frac{-\frac{1}{2}\pi(\alpha\beta'-\alpha'\beta)\,d\lambda}{\lambda(1-\lambda^{2})LL}.
\]
Let $\lambda$ be the Modulus into which $k$ is transformed by the first transformation of the $n^{\mathrm{ti}}$ order; let further $Q=K$, $Q'=K'$, $L=\Lambda$, $L'=\Lambda'$; there will be:

\[
\frac{L'}{L}=\frac{\Lambda'}{\Lambda}=\frac{nK'}{K}=\frac{nQ'}{Q},
\]

\textbf{*)} Cf.\ Legendre Traité des F.\ E.\ Tom.\ I.\ Cap.\ XIII.

\origpage{75}

whence

\[
\frac{n\,dk}{k(1-k^{2})KK}=\frac{d\lambda}{\lambda(1-\lambda^{2})\Lambda\Lambda}.
\]
But we have found for that transformation $\Lambda=\dfrac{K}{uM}$\ednote{The print has $u$ in the denominator where $n$ is expected; an apparent misprint.}, whence now:

\[
MM=\frac{1}{n}.\,\frac{\lambda(1-\lambda^{2})\,dk}{k(1-k^{2})\,d\lambda}.
\]
In the second transformation we saw that $\dfrac{\Lambda_{,}'}{\Lambda_{,}}=\dfrac{1}{n}\,\dfrac{K'}{K}$, $\Lambda_{,}=\dfrac{K}{M_{,}}$, whence:

\[
\frac{dk}{k(1-k^{2})KK}=\frac{n\,d\lambda_{,}}{\lambda_{,}(1-\lambda_{,}^{2})\Lambda_{,}\Lambda_{,}},
\]
whence here too:

\[
M_{,}M_{,}=\frac{1}{n}.\,\frac{\lambda_{,}(1-\lambda_{,}^{2})\,dk}{k(1-k^{2})\,d\lambda_{,}}.
\]
In general, however, whatever the Modulus $\lambda$ be, whether real or imaginary, into which the proposed Modulus $k$ can be transformed by a transformation of the $n^{\mathrm{ti}}$ order, the equation will hold:

\[
MM=\frac{1}{n}.\,\frac{k(1-k^{2})\,d\lambda}{\lambda(1-\lambda^{2})\,dk}.
\]
That this may be proved, I shall note that in general equations of the form are obtained:

\[
\begin{gathered}
\alpha\Lambda+i\beta\Lambda'=\frac{aK+ibK'}{nM}\\[1.5ex]
\alpha'\Lambda'+i\beta'\Lambda=\frac{a'K'+ib'K}{nM},
\end{gathered}
\]
with $a$, $a'$, $\alpha$, $\alpha'$ denoting odd numbers, $b$, $b'$, $\beta$, $\beta'$ even numbers, both kinds positive or negative, of such a sort that $aa'+bb'=1$, $\alpha\alpha'+\beta\beta'=1$ ${}^{*)}$. Hence, with the following put:

\[
\begin{gathered}
aK+ibK'=Q,\quad a'K'+ib'K=Q'\\[1ex]
\alpha\Lambda+i\beta\Lambda'=L,\quad\alpha'\Lambda'+i\beta'\Lambda=L',
\end{gathered}
\]
we obtain, because $aa'+bb'=1$, $\alpha\alpha'+\beta\beta'=1$:

\[
d\frac{Q'}{Q}=\frac{-n\pi\,dk}{2k(1-k^{2})QQ},\quad d\frac{L'}{L}=\frac{-\pi\,d\lambda}{2\lambda(1-\lambda^{2})LL},
\]

\textbf{*)} A more accurate determination of the numbers $a$, $a'$, $b$, $b'$ etc.\ etc.\ for the individual transformations of the same order seems to labour under grave difficulties. Indeed this determination, unless we are greatly mistaken, depends chiefly on the limits between which the Modulus $k$ lies, so that for different limits it comes out wholly different. How intricate this renders the question, he who has tried it will recognise. Above all, however, it seems that the nature of imaginary Moduli must be inquired into more accurately; this question as yet lies wholly untouched.

\origpage{76}
whence, since $\dfrac{Q'}{Q}=\dfrac{L'}{L}$, $L=\dfrac{Q}{nM}$, in general there results:

\[
MM=\frac{1}{n}.\,\frac{\lambda(1-\lambda^{2})\,dk}{k(1-k^{2})\,d\lambda}.
\]
I shall note in addition that the equation found can also be exhibited thus:

\[
MM=\frac{1}{n}.\,\frac{\lambda^{2}(1-\lambda^{2})\,d(k^{2})}{k^{2}(1-k^{2})\,d(\lambda^{2})}=\frac{1}{n}.\,\frac{\lambda'^{2}(1-\lambda'^{2})\,d(k'^{2})}{k^{2}(1-k'^{2})\,d(\lambda'^{2})},
\]\ednote{The print sets $k^{2}$ in the last denominator, where $k'^{2}$ is expected; the author's Corrigenda list this misprint.}
whence we see that the expression $MM$ is not changed when in place of $k$, $\lambda$ the Complements $k'$, $\lambda'$ are put, or, what we demonstrated above, that for complementary transformations, no account being taken of sign, the multiplier $M$ is the same. Further, by changing $k$ into $\lambda$, $\lambda$ into $k$, whereby the transformation passes into the supplementary, $MM$ is changed into

\[
\frac{1}{n}.\,\frac{k(1-k^{2})\,d\lambda}{\lambda(1-\lambda^{2})\,dk}=\frac{1}{nnMM},\quad\text{or }M\text{ into }\frac{1}{nM},
\]
which too was itself proved above.

\subsection*{33.}

With $Q=aK+bK'$, $L=\alpha\Lambda+\beta\Lambda'$ put, the Constants $a$, $b$, $\alpha$, $\beta$ can always be so determined that $L=\dfrac{Q}{M}$, or $Q=ML$. Further, one has the equations:

\[
\begin{gathered}
\text{1)}\quad(k-k^{3})\frac{d^{2}Q}{dk^{2}}+(1-3k^{2})\frac{dQ}{dk}-kQ=0\\[1.5ex]
\text{2)}\quad(\lambda-\lambda^{3})\frac{d^{2}L}{d\lambda^{2}}+(1-3\lambda^{2})\frac{dL}{d\lambda}-\lambda L=0,
\end{gathered}
\]
which can also be represented in this manner:

\[
\begin{gathered}
\text{3)}\quad d\,\frac{\dfrac{(k-k^{3})\,dQ}{dk}}{dk}-kQ=0\\[2ex]
\text{4)}\quad d\,\frac{\dfrac{(\lambda-\lambda^{3})\,dL}{d\lambda}}{d\lambda}-\lambda L=0.
\end{gathered}
\]
Let us substitute in the equation:

\[
(k-k^{3})\frac{d^{2}Q}{dk^{2}}+(1-3k^{2})\frac{dQ}{dk}-kQ=0
\]

\origpage{77}
$Q=ML$; there results:

\[
L\left\{(k-k^{3})\frac{d^{2}M}{dk^{2}}+(1-3k^{2})\frac{dM}{dk}-kM\right\}+\frac{dL}{dk}\left\{2(k-k^{3})\frac{dM}{dk}+(1-3k^{2})M\right\}+(k-k^{3})M\frac{d^{2}L}{dk^{2}}=0,
\]
which, multiplied by $M$, gives us:

\[
\text{5)}\quad LM\left\{(k-k^{3})\frac{d^{2}M}{dk^{2}}+(1-3k^{2})\frac{dM}{dk}-kM\right\}+d\,\frac{\dfrac{(k-k^{3})M^{2}dL}{dk}}{dk}=0.
\]
But from the antecedent §$^{\mathrm{o}}$ there results:

\[
M^{2}=\frac{(\lambda-\lambda^{3})\,dk}{n(k-k^{3})\,d\lambda},\quad\text{whence }\frac{(k-k^{3})M^{2}\,dL}{dk}=\frac{(\lambda-\lambda^{3})\,dL}{n\,d\lambda}.
\]
Further, from equation 4) there results:

\[
d\left\{\frac{(\lambda-\lambda^{3})\,dL}{d\lambda}\right\}=\lambda L\,d\lambda,\quad\text{whence}
\]

\[
d\,\frac{\dfrac{(k-k^{3})M^{2}\,dL}{dk}}{dk}=d\,\frac{\dfrac{(\lambda-\lambda^{3})\,dL}{d\lambda}}{n\,dk}=\frac{\lambda L\,d\lambda}{n\,dk}.
\]
Hence equation 5), divided by $L$, passes into this:

\[
\text{6)}\quad M\left\{(k-k^{3})\frac{d^{2}M}{dk^{2}}+(1-3k^{2})\frac{dM}{dk}-kM\right\}+\frac{\lambda\,d\lambda}{n\,dk}=0.
\]
When in this equation the value of $M$ from the equation $M^{2}=\dfrac{(\lambda-\lambda^{3})\,dk}{n(k-k^{3})\,d\lambda}$ is substituted, a differential equation between the Moduli $k$, $\lambda$ themselves is obtained, which, it is easily evident, ascends to the third order. By a somewhat troublesome calculation there is found:

\[
\text{7)}\quad\frac{3d^{2}\lambda^{2}}{dk^{4}}-\frac{2d\lambda}{dk}.\frac{d^{3}\lambda}{dk^{3}}+\frac{d\lambda^{2}}{dk^{2}}\left\{\left(\frac{1+k^{2}}{k-k^{3}}\right)^{2}-\left(\frac{1+\lambda^{2}}{\lambda-\lambda^{3}}\right)^{2}.\frac{d\lambda^{2}}{dk^{2}}\right\}=0.
\]
In this equation $dk$ has been considered as a constant differential. Where it pleases to transform it into another, in which no differential is put constant, there will have to be put:

\[
\begin{gathered}
\frac{d^{2}\lambda}{dk^{2}}=\frac{d^{2}\lambda}{dk^{2}}-\frac{d\lambda\,d^{2}k}{dk^{3}}\\[1.5ex]
\frac{d^{3}\lambda}{dk^{3}}=\frac{d^{3}\lambda}{dk^{3}}-\frac{3d^{2}\lambda\,d^{2}k}{dk^{4}}-\frac{d\lambda\,d^{3}k}{dk^{4}}+\frac{3d\lambda\,d^{2}k^{2}}{dk^{5}}
\end{gathered}
\]
whence:

\[
\frac{3d^{2}\lambda^{2}}{dk^{4}}-\frac{2d\lambda}{dk}.\frac{d^{3}\lambda}{dk^{3}}=\frac{3d^{2}\lambda^{2}}{dk^{4}}-\frac{3d\lambda^{2}d^{2}k^{2}}{dk^{6}}+\frac{2d\lambda^{2}d^{3}k}{dk^{5}}-\frac{2d\lambda\,d^{3}\lambda}{dk^{4}}.
\]

\origpage{78}
Hence equation 7), multiplied by $dk^{6}$, passes into the following, in which no differential is put constant, or in which any one whatever that one pleases may be considered as such:

\[
\text{8)}\quad3\left\{dk^{2}d^{2}\lambda^{2}-d\lambda^{2}d^{2}k^{2}\right\}-2dk\,d\lambda\left\{dk\,d^{3}\lambda-d\lambda\,d^{3}k\right\}+dk^{2}d\lambda^{2}\left\{\left(\frac{1+k^{2}}{k-k^{3}}\right)^{2}dk^{2}-\left(\frac{1+\lambda^{2}}{\lambda-\lambda^{3}}\right)^{2}d\lambda^{2}\right\}=0.
\]
It is evident that this equation, the elements $k$ and $\lambda$ being interchanged, remains unchanged, which we proved above of the Modular Equations.

It is worth the labour to investigate that differential equation of the third order by yet another method. To this end let us introduce into the equation from which we set out:

\[
(k-k^{3})\frac{d^{2}Q}{dk^{2}}+(1-3k^{2})\frac{dQ}{dk}-kQ=0
\]
the quantity $(k-k^{3})QQ=s$. There results

\[
\begin{gathered}
\frac{ds}{dk}=(1-3k^{2})QQ+2(k-k^{3})Q\frac{dQ}{dk}\\[1.5ex]
\frac{d^{2}s}{dk^{2}}=-6kQQ+4(1-3k^{2})Q\frac{dQ}{dk}+2(k-k^{3})\left(\frac{dQ}{dk}\right)^{2}+2(k-k^{3})Q\frac{d^{2}Q}{dk^{2}}.
\end{gathered}
\]
In this equation, when one puts:

\[
(k-k^{3})\frac{d^{2}Q}{dk^{2}}=kQ-(1-3k^{2})\frac{dQ}{dk},\quad\text{there results}
\]

\[
\begin{gathered}
\frac{d^{2}s}{dk^{2}}=-4kQQ+2(1-3k^{2})Q\frac{dQ}{dk}+2(k-k^{3})\left(\frac{dQ}{dk}\right)^{2}\\[1.5ex]
=2\frac{dQ}{dk}\left\{(1-3k^{2})Q+(k-k^{3})\frac{dQ}{dk}\right\}-4kQQ.
\end{gathered}
\]
This equation, multiplied into $2s=2(k-k^{3})QQ$, gives:

\[
\frac{2s\,d^{2}s}{dk^{2}}=2(k-k^{3})Q\frac{dQ}{dk}\left\{2(1-3k^{2})QQ+2(k-k^{3})Q\frac{dQ}{dk}\right\}-8k^{2}(1-k^{2})Q^{4},
\]
or, since

\[
\begin{gathered}
2(k-k^{3})Q\frac{dQ}{dk}=\frac{ds}{dk}-(1-3k^{2})QQ\\[1.5ex]
2(1-3k^{2})QQ+2(k-k^{3})Q\frac{dQ}{dk}=\frac{ds}{dk}+(1-3k^{2})QQ,
\end{gathered}
\]
we obtain:

\[
\frac{2s\,d^{2}s}{dk^{2}}=\left(\frac{ds}{dk}\right)^{2}-(1-3k^{2})^{2}Q^{4}-8k^{2}(1-k^{2})Q^{4}=\left(\frac{ds}{dk}\right)^{2}-(1+k^{2})^{2}Q^{4},\quad\text{or}
\]

\[
\text{9)}\quad\frac{2s\,d^{2}s}{dk^{2}}-\left(\frac{ds}{dk}\right)^{2}+\left(\frac{1+k^{2}}{k-k^{3}}\right)^{2}ss=0.
\]

\origpage{79}
Now indeed, having put $a'K+b'K'=Q'$, $\dfrac{Q'}{Q}=t$, we have seen that $\dfrac{dt}{dk}=\dfrac{m}{(k-k^{3})QQ}=\dfrac{m}{s}$, $m$ denoting a Constant; whence $s=\dfrac{m\,dk}{dt}$. Let us transform equation 9) into another, in which $dt$ is put constant. There will be $\dfrac{ds}{dk}=\dfrac{m\,d^{2}k}{dt\,dk}$, $\dfrac{d^{2}s}{dk^{2}}=\dfrac{m\,d^{3}k}{dt\,dk^{2}}-\dfrac{m\,d^{2}k^{2}}{dt\,dk^{3}}$; and when these are substituted, there results from equation 9):

\[
\frac{2d^{3}k}{dt^{2}dk}-\frac{3d^{2}k^{2}}{dt^{2}dk^{2}}+\left(\frac{1+k^{2}}{k-k^{3}}\right)^{2}\frac{dk^{2}}{dt^{2}}=0,\quad\text{or}
\]

\[
\text{10)}\quad2d^{3}k\,dk-3d^{2}k^{2}+\left(\frac{1+k^{2}}{k-k^{3}}\right)^{2}dk^{4}=0;
\]
where differentiation is to be made with respect to $t$, which has come out of the equation.

Putting $\dfrac{\alpha'\Lambda+\beta'\Lambda'}{\alpha\Lambda+\beta\Lambda'}=\omega$, the Constants $\alpha$, $\beta$, $\alpha'$, $\beta'$, whenever $\lambda$ is a transformed Modulus, can be so determined that $t=\omega$; and in a similar way we obtain:

\[
\text{11)}\quad2d^{3}\lambda\,d\lambda-3d^{2}\lambda^{2}+\left(\frac{1+\lambda^{2}}{\lambda-\lambda^{3}}\right)d\lambda^{4}=0,
\]\ednote{The print omits the exponent 2 after the parenthesis; compare 10) and the author's Corrigenda.}
in which equation too differentiation will have to be made with respect to $\omega=t$. Let equation 10) be multiplied by $d\lambda^{2}$, equation 11) by $dk^{2}$: on subtraction one obtains:

\[
\text{12)}\quad2dk\,d\lambda\left\{d\lambda\,d^{3}k-dk\,d^{3}\lambda\right\}-3\left\{d\lambda^{2}d^{2}k^{2}-dk^{2}d^{2}\lambda^{2}\right\}+dk^{2}d\lambda^{2}\left\{\left(\frac{1+k^{2}}{k-k^{3}}\right)^{2}dk^{2}-\left(\frac{1+\lambda^{2}}{\lambda-\lambda^{3}}\right)^{2}d\lambda^{2}\right\}=0.
\]
But this equation agrees with equation 8), in which we know that any differential one pleases can be considered as constant, and therefore, although it has been found on the supposition that $dt$ is a constant differential, it will also hold whichever other is considered as such.

Behold, therefore, a differential equation of the third order, which has innumerable algebraic solutions, yet particular ones, namely the Equations we have called Modular. But the complete Integral depends on elliptic functions; for it is $t=\omega$, or $\dfrac{a'K+b'K'}{aK+bK'}=\dfrac{\alpha'\Lambda+\beta'\Lambda'}{\alpha\Lambda+\beta\Lambda'}$, which equation can also be represented thus:

\[
mK\Lambda+m'K'\Lambda'+m''K\Lambda'+m'''K'\Lambda=0,
\]
$m$, $m'$, $m''$, $m'''$ denoting Arbitrary Constants. This integration we deem to be of the deepest investigation.

\origpage{80}

We could inquire whether the Modular Equations for the transformations of the third and fifth orders in fact satisfy, as they must, our differential equation of the third order. But since this seems to demand overly long calculations, let it suffice to demonstrate the same for the transformation of the second order, where $\lambda=\dfrac{1-k'}{1+k'}$.

Let $dk'$ be considered as constant; there results:

\[
\begin{gathered}
\lambda=\frac{1-k'}{1+k'}=-1+\frac{2}{1+k'}\qquad kk+k'k'=1\\[1.5ex]
\frac{d\lambda}{dk'}=\frac{-2}{(1+k')^{2}}\qquad\frac{dk}{dk'}=\frac{-k'}{k}\\[1.5ex]
\frac{d^{2}\lambda}{dk'^{2}}=\frac{4}{(1+k')^{3}}\qquad\frac{d^{2}k}{dk'^{2}}=\frac{-1}{k}-\frac{k'k'}{k^{3}}=\frac{-1}{k^{3}}\\[1.5ex]
\frac{d^{3}\lambda}{dk'^{3}}=\frac{-12}{(1+k')^{4}}\qquad\frac{d^{3}k}{dk'^{3}}=\frac{3k'}{k^{5}}.
\end{gathered}
\]\ednote{The last fraction is printed without the minus sign, as $\frac{3k'}{k^{5}}$; the author's Corrigenda give $-\frac{3k'}{k^{5}}$.}
Hence there is:

\[
\begin{gathered}
\frac{dk^{2}d^{2}\lambda^{2}-d\lambda^{2}d^{2}k^{2}}{dk'^{6}}=\frac{16k'k'}{k^{2}(1+k')^{6}}-\frac{4}{k^{6}(1+k')^{4}}\\[1.5ex]
=\frac{4\left\{k^{4}k'^{2}-(1+k')^{2}\right\}}{k^{6}(1+k')^{6}}=\frac{4\left\{k'^{2}(1-k')^{2}-1\right\}}{k^{6}(1+k')^{4}}.
\end{gathered}
\]\ednote{Printed with the factor 4 outside the braces and not before $k^{4}k'^{2}$ and $k'^{2}(1-k')^{2}$ inside them; the author's Corrigenda give $4k^{4}k'^{2}$ and $4k'^{2}(1-k')^{2}$.}

Further one obtains:

\[
\begin{gathered}
\frac{dk\,d^{3}\lambda-d\lambda\,d^{3}k}{dk'^{4}}=\frac{12k'}{k(1+k')^{4}}+\frac{6k'}{k^{5}(1+k')^{2}}=\frac{6k'\left\{2(1-k')^{2}+1\right\}}{k^{5}(1+k')^{2}}\\[1.5ex]
\frac{dk\,d\lambda\left\{dk\,d^{3}\lambda-d\lambda\,d^{3}k\right\}}{dk'^{6}}=\frac{12k'^{2}\left\{2(1-k')^{2}-1\right\}}{k^{6}(1+k')^{4}},
\end{gathered}
\]
whence

\[
\frac{3\left\{dk^{2}d^{2}\lambda^{2}-d\lambda^{2}d^{2}k^{2}\right\}-2dk\,d\lambda\left\{dk\,d^{3}\lambda-d\lambda\,d^{3}k\right\}}{dk'^{6}}=\frac{12(2k'^{2}-1)}{k^{6}(1+k')4}.
\]\ednote{Printed $(1+k')4$ with the exponent set on the baseline instead of raised; an apparent misprint for $(1+k')^{4}$.}
Further there is

\[
\begin{gathered}
\left(\frac{1+k^{2}}{k-k^{3}}\right)^{2}\frac{dk^{2}}{dk'^{2}}=\frac{(1+k^{2})^{2}}{k^{4}k'^{2}}\\[1.5ex]
\left(\frac{1+\lambda^{2}}{\lambda-\lambda^{3}}\right)^{2}\frac{d\lambda^{2}}{dk'^{2}}=\frac{4}{(1+k')^{4}}\left\{\frac{1+k'}{1-k'}\right\}^{2}\left\{\frac{1+k'^{2}}{2k'}\right\}^{2}=\frac{(1+k'^{2})^{2}}{k'^{2}k^{4}},
\end{gathered}
\]
whence:

\[
\begin{gathered}
\left\{\frac{1+k^{2}}{k-k^{3}}\right\}^{2}\frac{dk^{2}}{dk'^{2}}-\left\{\frac{1+\lambda^{2}}{\lambda-\lambda^{3}}\right\}^{2}\frac{d\lambda^{2}}{d\lambda'^{2}}=\frac{3(1-2k'^{2})}{k^{4}k'^{2}}\\[1.5ex]
\frac{dk^{2}d\lambda^{2}}{dk'^{4}}\left\{\left(\frac{1+k^{2}}{k-k^{3}}\right)^{2}\frac{dk^{2}}{dk'^{2}}-\left(\frac{1+\lambda^{2}}{\lambda-\lambda^{3}}\right)\frac{d\lambda^{2}}{d\lambda'^{2}}\right\}=\frac{12(1-2k'^{2})}{k^{6}(1+k')^{4}}.
\end{gathered}
\]\ednote{Printed $d\lambda'^{2}$ for $dk'^{2}$ in both denominators, and in the second row the exponent 2 after the bracket of the $\lambda$ term is missing; apparent misprints.}

\origpage{81}
Hence at last there results, as it must:

\[
\begin{gathered}
\frac{3\left\{dk^{2}d^{2}\lambda^{2}-d\lambda^{2}d^{2}k^{2}\right\}-2dk\,d\lambda\left\{dk\,d^{3}\lambda-d\lambda\,d^{3}k\right\}}{dk'^{6}}\\[1.5ex]
+\frac{dk^{2}d\lambda^{2}}{dk'^{4}}\left\{\left(\frac{1+k^{2}}{k-k^{3}}\right)^{2}\frac{dk^{2}}{dk'^{2}}-\left(\frac{1+\lambda^{2}}{\lambda-\lambda^{3}}\right)^{2}\frac{d\lambda^{2}}{dk'^{2}}\right\}=\frac{12(2k'^{2}-1)}{k^{6}(1+k')^{4}}+\frac{12(1-2k'^{2})}{k^{6}(1+k')^{4}}=0.
\end{gathered}
\]
If expeditious methods were at hand for bringing out all the algebraic solutions that a differential equation has, then from the differential equation proposed by us alone we could elicit all the Modular Equations that concern the single orders of transformations. Yet apart from Cl.\ \emph{Condorcet,} I know of no one who has touched this arduous subject, which is worthy of the attention of Analysts.

\subsection*{34.}

The equation found above:

\[
MM=\frac{1}{n}.\,\frac{\lambda(1-\lambda\lambda)}{k(1-kk)}.\,\frac{dk}{d\lambda},
\]
by whose aid, from a Modular Equation that has been found, the quantity $M$ too can at once be determined, seems worthy of our dwelling on it a little longer. It is not evident at first sight how the values of the quantity $M$ found in the transformations of the third and fifth orders agree with that equation. Let us therefore examine this more accurately:

\emph{a)} In the transformation of the \emph{third} order, putting $u=\sqrt[4]{k}$, $v=\sqrt[4]{\lambda}$, we found:

\[
\text{1)}\quad u^{4}-v^{4}+2uv(1-u^{2}v^{2})=0,
\]
which equation we also exhibited thus in §.\ 16:

\[
\text{2)}\quad\left(\frac{v+2u^{3}}{v}\right)\left(\frac{u-2v^{3}}{u}\right)=-3.
\]
Further we saw that there results:

\[
\text{3)}\quad M=\frac{v}{v+2u^{3}}=\frac{2v^{3}-u}{3u}.
\]
Equation 1) being differentiated, we obtain:

\[
\frac{du}{dv}=\frac{2v^{3}-u+3u^{3}v^{2}}{2u^{3}+v-3u^{2}v^{3}},
\]

\origpage{82}
or, putting in place of $3$ the expression $\left(\dfrac{v+2u^{3}}{v}\right)\left(\dfrac{2v^{3}-u}{u}\right)$:

\[
\text{4)}\quad\frac{du}{dv}=\frac{2v^{3}-u}{2u^{3}+v}.\,\frac{1+u^{2}v^{2}+2u^{5}v}{1+u^{2}v^{2}-2uv^{5}}.
\]
From equation 1) it follows:

\[
\begin{gathered}
1-u^{8}=(1+u^{4})\left\{1-v^{4}+2uv(1-u^{2}v^{2})\right\}\\[1ex]
=1-u^{4}v^{4}+u^{4}-v^{4}+2uv(1+u^{4})(1-u^{2}v^{2}\\[1ex]
=1-u^{4}v^{4}+2u^{5}v(1-u^{2}v^{2})=(1-u^{2}v^{2})(1+u^{2}v^{2}+2u^{5}v).
\end{gathered}
\]\ednote{The closing bracket after $u^{2}v^{2}$ is missing in the second row of the print.}
In the same way one finds:

\[
1-v^{8}=(1-u^{2}v^{2})(1+u^{2}v^{2}-2uv^{5}),
\]
whence

\[
\frac{1-v^{8}}{1-u^{8}}=\frac{1+u^{2}v^{2}-2uv^{5}}{1+u^{2}v^{2}+2u^{5}v},\quad\text{or from equation 4):}
\]

\[
\frac{1-v^{8}}{1-u^{8}}.\frac{du}{dv}=\frac{2v^{3}-u}{2u^{3}+v}.
\]
This equation being multiplied by

\[
\frac{v}{3u}=\frac{vv}{(2u^{3}+v)(2v^{3}-u)},
\]
there results:

\[
\frac{1}{3}.\,\frac{v(1-v^{8})}{u(1-u^{8})}.\,\frac{du}{dv}=\frac{1}{3}.\,\frac{\lambda(1-\lambda\lambda)}{k(1-kk)}.\,\frac{dk}{d\lambda}=\left(\frac{v}{v+2u^{3}}\right)^{2}=MM,
\]
Q.\ E.\ D.

\emph{b)} In the transformation of the \emph{fifth} order, putting $u=\sqrt[4]{k}$, $v=\sqrt[4]{\lambda}$, we found:

\[
\text{1)}\quad u^{6}-v^{6}+5u^{2}v^{2}(u^{2}-v^{2})+4uv(1-u^{4}v^{4})=0,
\]
which equation we also exhibited in these ways in §§.\ 16.\ 30:

\[
\begin{gathered}
\text{2)}\quad\frac{u+v^{5}}{u(1+u^{3}v)}.\,\frac{v-u^{5}}{v(1-uv^{3})}=5\\[1.5ex]
\text{3)}\quad(u^{2}-v^{2})^{6}=16u^{2}v^{2}(1-u^{8})(1-v^{8}).
\end{gathered}
\]
Further we found:

\[
\text{4)}\quad M=\frac{v(1-uv^{3})}{v-u^{5}}=\frac{u+v^{5}}{5u(1+u^{3}v)}.
\]

\origpage{83}
Equation 3) being differentiated, we obtain:

\[
\begin{gathered}
6uv(1-u^{8})(1-v^{8})(u\,du-v\,dv)=\\[1ex]
u(u^{2}-v^{2})(1-u^{8})(1-5v^{8})\,dv+v(u^{2}-v^{2})(1-v^{8})(1-5u^{8})\,du,
\end{gathered}
\]
or:

\[
\text{5)}\quad v(1-v^{8})\left\{5u^{2}-u^{10}+v^{2}-5u^{8}v^{2}\right\}du=u(1-u^{8})\left\{5v^{2}-v^{10}+u^{2}-5u^{2}v^{8}\right\}dv.
\]
Equation 1) being multiplied by $u^{4}$, $v^{4}$, there is elicited:

\[
\begin{gathered}
5u^{2}-u^{10}+v^{2}-5u^{8}v^{2}=(1-u^{4}v^{4})(v^{2}+5u^{2}+4uv^{5})\\[1ex]
5v^{2}-v^{10}+u^{2}-5u^{2}v^{8}=(1-u^{4}v^{4})(u^{2}+5v^{2}-4uv^{5}),
\end{gathered}
\]
whence equation 5) passes into this:

\[
\text{6)}\quad\frac{v(1-v^{8})}{u(1-u^{8})}.\,\frac{du}{dv}=\frac{u^{2}+5v^{2}-4uv^{5}}{v^{2}+5u^{2}+4u^{5}v}.
\]
Let $u+v^{5}=A$, $u+u^{4}v=B$, $v-u^{5}=C$, $v-uv^{4}=D$, so that:

\[
\begin{gathered}
\frac{AC}{BD}=5,\quad\text{or }AC=5BD\\[1.5ex]
\frac{D}{C}=\frac{A}{5B}=M:\\[1.5ex]
u^{2}+5v^{2}-4uv^{5}=uA+5vD\\[1ex]
v^{2}+5u^{2}+4u^{5}v=vC+5uB,
\end{gathered}
\]
there will be:

\[
\begin{gathered}
\text{7)}\quad\frac{v(1-v^{8})}{u(1-u^{8})}\,\frac{du}{dv}=\frac{uA+5vD}{vC+5uB}=\frac{uAB+vAC}{vCD+uAC}.\,\frac{D}{B}\\[1.5ex]
=\frac{uB+vC}{vD+uA}.\,\frac{AD}{BC}=\frac{AD}{BC}=5MM.
\end{gathered}
\]
For there is:

\[
uB+vC=vD+uA=uu+vv.
\]
Whence also:

\[
MM=\frac{1}{5}.\,\frac{v(1-v^{8})}{u(1-u^{8})}.\,\frac{du}{dv}=\frac{1}{5}.\,\frac{\lambda(1-\lambda\lambda)}{k(1-kk)}.\,\frac{dk}{d\lambda}.
\]
Q.\ E.\ D.

\origpage{84}

\section*{THEORY OF THE EXPANSION OF ELLIPTIC FUNCTIONS.}

\subsection*{ON THE EXPANSION OF ELLIPTIC FUNCTIONS INTO INFINITE PRODUCTS.}

\subsection*{35.}

A real Modulus $k$, less than unity, being proposed, we see that the Modulus

\[
\lambda=k^{n}\left\{\sin\operatorname{coam}\frac{2K}{n}.\sin\operatorname{coam}\frac{4K}{n}.\,.\sin\operatorname{coam}\frac{(n-1)K}{n}\right\}^{4},
\]
into which that Modulus is changed by the first transformation of the $n^{\mathrm{ti}}$ order, converges most rapidly to nothing as the number $n$ increases; and hence for the limit $n=\infty$ there results $\lambda=0$. Then there will be $\Lambda=\dfrac{\pi}{2}$, $\operatorname{am}(u,\lambda)=u$, whence from the formulae $\Lambda=\dfrac{K}{nM}$, $\Lambda'=\dfrac{K'}{M}$ we shall obtain:

\[
nM=\frac{2K}{\pi},\quad\frac{\Lambda'}{n}=\frac{K'}{nM}=\frac{\pi K'}{2K}.
\]
Let us now put, in the formulae for the supplementary transformation of the first §.\ 26, $\dfrac{u}{n}$ in place of $u$, and $n=\infty$: $\operatorname{am}\left(\dfrac{u}{M},\lambda\right)$ passes into $\operatorname{am}\left(\dfrac{u}{nM},\lambda\right)=\dfrac{\pi u}{2K}$, $y=\sin\operatorname{am}\left(\dfrac{u}{M},\lambda\right)$ into $\sin\dfrac{\pi u}{2K}$; further $\operatorname{am}(nu)$ into $\operatorname{am}(u)$. Hence from those formulae we arrive at the following:

\[
\begin{gathered}
\sin\operatorname{am}u=\frac{2Ky}{\pi}.\frac{\left(1-\dfrac{yy}{\sin^{2}.\dfrac{i\pi K'}{K}}\right)\left(1-\dfrac{yy}{\sin^{2}.\dfrac{2i\pi K'}{K}}\right)\left(1-\dfrac{yy}{\sin^{2}.\dfrac{3i\pi K'}{K}}\right)\ldots}{\left(1-\dfrac{yy}{\sin^{2}.\dfrac{i\pi K'}{2K}}\right)\left(1-\dfrac{yy}{\sin^{2}.\dfrac{3i\pi K'}{2K}}\right)\left(1-\dfrac{yy}{\sin^{2}.\dfrac{5i\pi K'}{2K}}\right)\ldots}\\[3ex]
\cos\operatorname{am}u=\sqrt{1-yy}.\frac{\left(1-\dfrac{yy}{\cos^{2}.\dfrac{i\pi K'}{K}}\right)\left(1-\dfrac{yy}{\cos^{2}.\dfrac{2i\pi K'}{K}}\right)\left(1-\dfrac{yy}{\cos^{2}.\dfrac{3i\pi K'}{K}}\right)\ldots}{\left(1-\dfrac{yy}{\sin^{2}.\dfrac{i\pi K'}{2K}}\right)\left(1-\dfrac{yy}{\sin^{2}.\dfrac{3i\pi K'}{2K}}\right)\left(1-\dfrac{yy}{\sin^{2}.\dfrac{5i\pi K'}{2K}}\right)\ldots}
\end{gathered}
\]

\origpage{85}
\[
\begin{gathered}
\Delta\operatorname{am}u=\frac{\left(1-\dfrac{yy}{\cos^{2}.\dfrac{i\pi K'}{2K}}\right)\left(1-\dfrac{yy}{\cos^{2}.\dfrac{3i\pi K'}{2K}}\right)\left(1-\dfrac{yy}{\cos^{2}.\dfrac{5i\pi K'}{2K}}\right)\ldots}{\left(1-\dfrac{yy}{\sin^{2}.\dfrac{i\pi K'}{2K}}\right)\left(1-\dfrac{yy}{\sin^{2}.\dfrac{3i\pi K'}{2K}}\right)\left(1-\dfrac{yy}{\sin^{2}.\dfrac{5i\pi K'}{2K}}\right)\ldots}
\\[2ex]
\sqrt{\frac{1-\sin\operatorname{am}u}{1+\sin\operatorname{am}u}}=\sqrt{\frac{1-y}{1+y}}.\frac{\left(1-\dfrac{y}{\cos.\dfrac{i\pi K'}{K}}\right)\left(1-\dfrac{y}{\cos.\dfrac{2i\pi K'}{K}}\right)\left(1-\dfrac{y}{\cos.\dfrac{3i\pi K'}{K}}\right)\ldots}{\left(1+\dfrac{y}{\cos.\dfrac{i\pi K'}{K}}\right)\left(1+\dfrac{y}{\cos.\dfrac{2i\pi K'}{K}}\right)\left(1+\dfrac{y}{\cos.\dfrac{3i\pi K'}{K}}\right)\ldots}
\\[2ex]
\sqrt{\frac{1-k\sin\operatorname{am}u}{1+k\sin\operatorname{am}u}}=\frac{\left(1-\dfrac{y}{\cos.\dfrac{i\pi K'}{2K}}\right)\left(1-\dfrac{y}{\cos.\dfrac{3i\pi K'}{2K}}\right)\left(1-\dfrac{y}{\cos.\dfrac{5i\pi K'}{2K}}\right)\ldots}{\left(1+\dfrac{y}{\cos.\dfrac{i\pi K'}{2K}}\right)\left(1+\dfrac{y}{\cos.\dfrac{3i\pi K'}{2K}}\right)\left(1+\dfrac{y}{\cos.\dfrac{5i\pi K'}{2K}}\right)\ldots}
\\[2ex]
\sin\operatorname{am}u=-\frac{\pi y}{kK}\left(\frac{\cos.\dfrac{i\pi K'}{2K}}{\sin^{2}.\dfrac{i\pi K'}{2K}-yy}+\frac{\cos.\dfrac{3i\pi K'}{2K}}{\sin^{2}.\dfrac{3i\pi K'}{2K}-yy}+\frac{\cos.\dfrac{5i\pi K'}{2K}}{\sin^{2}.\dfrac{5i\pi K'}{2K}-yy}+.\,.\right)
\\[2ex]
\cos\operatorname{am}u=\frac{i\sqrt{1-yy}.\pi}{kK}\left(\frac{\sin.\dfrac{i\pi K'}{2K}}{\sin^{2}.\dfrac{i\pi K'}{2K}-yy}-\frac{\sin.\dfrac{3i\pi K'}{2K}}{\sin^{2}.\dfrac{3i\pi K'}{2K}-yy}+\frac{\sin.\dfrac{5i\pi K'}{2K}}{\sin^{2}.\dfrac{5i\pi K'}{2K}-yy}-.\,.\right).
\end{gathered}
\]

In what follows let us put $e^{-\frac{\pi K'}{K}}=q$, $\dfrac{\pi u}{2K}=x$, or $u=\dfrac{2Kx}{\pi}$, whence $y=\sin\dfrac{\pi u}{2K}=\sin x$; there is:
\[
\begin{gathered}
\sin.\dfrac{mi\pi K'}{K}=\dfrac{q^{m}-q^{-m}}{2i}=\dfrac{i(1-q^{2m})}{2q^{m}}
\\[2ex]
\cos.\dfrac{mi\pi K'}{K}=\dfrac{q^{m}+q^{-m}}{2}=\dfrac{1+q^{2m}}{2q^{m}},
\end{gathered}
\]
whence:
\[
1-\frac{yy}{\sin^{2}.\dfrac{mi\pi K'}{K}}=1+\frac{4q^{2m}\sin x^{2}}{(1-q^{2m})^{2}}=\frac{1-2q^{2m}\cos 2x+q^{4m}}{(1-q^{2m})^{2}}
\]

\origpage{86}
\[
\begin{gathered}
1-\frac{yy}{\cos^{2}.\dfrac{mi\pi K'}{K}}=1-\frac{4q^{2m}\sin x^{2}}{(1+q^{2m})^{2}}=\frac{1+2q^{2m}\cos 2x+q^{4m}}{(1+q^{2m})^{2}}
\\[2ex]
1\pm\frac{y}{\cos.\dfrac{mi\pi K'}{K}}=1\pm\frac{2q^{m}\sin x}{1+q^{2m}}=\frac{1\pm 2q^{m}\sin x+q^{2m}}{1+q^{2m}}
\\[2ex]
\frac{-\cos.\dfrac{mi\pi K'}{K}}{\sin^{2}.\dfrac{mi\pi K'}{K}-yy}=\frac{2q^{m}(1+q^{2m})}{1-2q^{2m}\cos 2x+q^{4m}}
\\[2ex]
\frac{i.\sin.\dfrac{mi\pi K'}{K}}{\sin^{2}.\dfrac{mi\pi K'}{K}-yy}=\frac{2q^{m}(1-q^{2m})}{1-2q^{2m}\cos 2x+q^{4m}}.
\end{gathered}
\]
These things having been prepared, and having put for the sake of brevity:
\[
\begin{gathered}
A=\left\{\frac{(1-q)(1-q^{3})(1-q^{5})\,.\,.}{(1-q^{2})(1-q^{4})(1-q^{6})\,.\,.}\right\}^{2}
\\[2ex]
B=\left\{\frac{(1-q)(1-q^{3})(1-q^{5})\,.\,.}{(1+q^{2})(1+q^{4})(1+q^{6})\,.\,.}\right\}^{2}
\\[2ex]
C=\left\{\frac{(1-q)(1-q^{3})(1-q^{5})\,.\,.}{(1+q)(1+q^{3})(1+q^{5})\,.\,.}\right\}^{2},
\end{gathered}
\]
there come forth the fundamental expansions of the Elliptic Functions into Infinite Products:
\[
\begin{gathered}
\text{1)}\quad\sin\operatorname{am}\frac{2Kx}{\pi}=\frac{2AK}{\pi}\sin x.\frac{(1-2q^{2}\cos 2x+q^{4})(1-2q^{4}\cos 2x+q^{8})(1-2q^{6}\cos 2x+q^{12})\,.\,.}{(1-2q\cos 2x+q^{2})(1-2q^{3}\cos 2x+q^{6})(1-2q^{5}\cos 2x+q^{10})\,.\,.}
\\[2ex]
\text{2)}\quad\cos\operatorname{am}\frac{2Kx}{\pi}=B\cos x.\frac{(1+2q^{2}\cos 2x+q^{4})(1+2q^{4}\cos 2x+q^{8})(1+2q^{6}\cos 2x+q^{12})\,.\,.}{(1-2q\cos 2x+q^{2})(1-2q^{3}\cos 2x+q^{6})(1-2q^{5}\cos 2x+q^{10})\,.\,.}
\\[2ex]
\text{3)}\quad\Delta\operatorname{am}\frac{2Kx}{\pi}=C.\frac{(1+2q\cos 2x+q^{2})(1+2q^{3}\cos 2x+q^{6})(1+2q^{5}\cos 2x+q^{10})\,.\,.}{(1-2q\cos 2x+q^{2})(1-2q^{3}\cos 2x+q^{6})(1-2q^{5}\cos 2x+q^{10})\,.\,.}
\\[2ex]
\text{4)}\quad\sqrt{\frac{1-\sin\operatorname{am}\dfrac{2Kx}{\pi}}{1+\sin\operatorname{am}\dfrac{2Kx}{\pi}}}=\sqrt{\frac{1-\sin x}{1+\sin x}}.\frac{(1-2q\sin x+q^{2})(1-2q^{2}\sin x+q^{4})(1-2q^{3}\sin x+q^{6})\,.\,.}{(1+2q\sin x+q^{2})(1+2q^{2}\sin x+q^{4})(1+2q^{3}\sin x+q^{6})\,.\,.}
\\[2ex]
\text{5)}\quad\sqrt{\frac{1-k\sin\operatorname{am}\dfrac{2Kx}{\pi}}{1+k\sin\operatorname{am}\dfrac{2Kx}{\pi}}}=\frac{(1-2\sqrt{q}\sin x+q)(1-2\sqrt{q^{3}}\sin x+q^{3})(1-2\sqrt{q^{5}}\sin x+q^{5})\,.\,.}{(1+2\sqrt{q}\sin x+q)(1+2\sqrt{q^{3}}\sin x+q^{3})(1+2\sqrt{q^{5}}\sin x+q^{5})\,.\,.}
\end{gathered}
\]

\origpage{87}
And also another system of formulae, which supplies the resolution of the proposed ones into simple fractions:
\[
\begin{gathered}
\text{6)}\quad\sin\operatorname{am}\frac{2Kx}{\pi}=\frac{2\pi}{kK}\sin x\left(\frac{\sqrt{q}(1+q)}{1-2q\cos 2x+q^{2}}+\frac{\sqrt{q^{3}}(1+q^{3})}{1-2q^{3}\cos 2x+q^{6}}+\frac{\sqrt{q^{5}}(1+q^{5})}{1-2q^{5}\cos 2x+q^{10}}+.\,.\right)
\\[2ex]
\text{7)}\quad\cos\operatorname{am}\frac{2Kx}{\pi}=\frac{2\pi}{kK}\cos x\left(\frac{\sqrt{q}(1-q)}{1-2q\cos 2x+q^{2}}-\frac{\sqrt{q^{3}}(1-q^{3})}{1-2q^{3}\cos 2x+q^{6}}+\frac{\sqrt{q^{5}}(1-q^{5})}{1-2q^{5}\cos 2x+q^{10}}-.\,.\right).
\end{gathered}
\]

To these we add, flowing from the same source:
\[
\begin{gathered}
\text{8)}\quad 1-\Delta\operatorname{am}\frac{2Kx}{\pi}=\frac{4\pi\sin x^{2}}{K}\left(\frac{q\left(\dfrac{1+q}{1-q}\right)}{1-2q\cos 2x+q^{2}}-\frac{q^{3}\left(\dfrac{1+q^{3}}{1-q^{3}}\right)}{1-2q^{3}\cos 2x+q^{6}}+\frac{q^{5}\left(\dfrac{1+q^{5}}{1-q^{5}}\right)}{1-2q^{5}\cos 2x+q^{10}}-.\,.\right)
\\[2ex]
\text{9)}\quad\operatorname{am}\frac{2Kx}{\pi}=\pm x+2\operatorname{Arc}\operatorname{tg}.\frac{(1+q)\operatorname{tg}x}{1-q}-2\operatorname{Arc}\operatorname{tg}.\frac{(1+q^{3})\operatorname{tg}x}{1-q^{3}}+2\operatorname{Arc}\operatorname{tg}.\frac{(1+q^{5})\operatorname{tg}x}{1-q^{5}}-.\,.
\end{gathered}
\]
In the last formula the upper sign is to be chosen whenever you stop the computation at a negative term, the lower whenever at a positive term.

\subsection*{36.}

Let us contemplate formulae 1)---3), in which first of all the values of the quantities that we have designated by $A$, $B$, $C$ are to be elicited. It is indeed easily found, putting $x=\dfrac{\pi}{2}$, from formulae 3), 1):
\[
k'=C\left\{\frac{(1-q)(1-q^{3})(1-q^{5})\ldots}{(1+q)(1+q^{3})(1+q^{5})\ldots}\right\}^{2}=CC,
\]
whence $C=\sqrt{k'}$;
\[
1=\frac{2AK}{\pi}\left\{\frac{(1+q^{2})(1+q^{4})(1+q^{6})\ldots}{(1+q)(1+q^{3})(1+q^{5})\ldots}\right\}^{2}=\frac{2AK}{\pi}.\frac{C}{B}=\frac{2\sqrt{k'}AK}{\pi B},
\]
whence $B=\dfrac{2\sqrt{k'}AK}{\pi}$. But in order that the value of $A$ itself be elicited, we must take refuge in other devices.

Let us put $e^{ix}=U$: where $x$ is changed into $x+\dfrac{i\pi K'}{2K}$, $U$ passes into $\sqrt{q}U$, and $\sin\operatorname{am}\dfrac{2Kx}{\pi}$ into
\[
\sin\operatorname{am}\left(\frac{2Kx}{\pi}+iK'\right)=\frac{1}{k\sin\operatorname{am}\dfrac{2Kx}{\pi}}\,.
\]

\origpage{88}
From formula 1), however, we obtain:
\[
\sin\operatorname{am}\frac{2Kx}{\pi}=\frac{AK}{\pi}\left(\frac{U-U^{-1}}{i}\right)\frac{\left\{(1-q^{2}U^{2})(1-q^{4}U^{2})\,.\,.\right\}\left\{(1-q^{2}U^{-2})(1-q^{4}U^{-2})\,.\,.\right\}}{\left\{(1-qU^{2})(1-q^{3}U^{2})\,.\,.\right\}\left\{(1-qU^{-2})(1-q^{3}U^{-2})\,.\,.\right\}},
\]
whence, changing $x$ into $x+\dfrac{i\pi K'}{2K}$:
\[
\frac{1}{k\sin\operatorname{am}\dfrac{2Kx}{\pi}}=\frac{AK}{\pi}\left(\frac{\sqrt{q}U-\sqrt{q^{-1}}U^{-1}}{i}\right)\frac{\left\{(1-q^{3}U^{2})(1-q^{5}U^{2})\,.\,.\right\}\left\{(1-qU^{-2})(1-q^{3}U^{-2})\,.\,.\right\}}{\left\{(1-q^{2}U^{2})(1-q^{4}U^{2})\,.\,.\right\}\left\{(1-U^{-2})(1-q^{2}U^{-2})\,.\,.\right\}},
\]
and these equations being multiplied together, since there is:
\[
\frac{\sqrt{q}U-\sqrt{q^{-1}}U^{-1}}{1-U^{2}}=\frac{-1}{\sqrt{q}}.\frac{1-qU^{2}}{U-U^{-1}},
\]
there results:
\[
\frac{1}{k}=\frac{1}{\sqrt{q}}\left(\frac{AK}{\pi}\right)^{2},\quad\text{or}\quad A=\frac{\pi\sqrt[4]{q}}{\sqrt{k.K}};\quad\text{whence}\quad\frac{2KA}{\pi}=\frac{2\sqrt[4]{q}}{\sqrt{k}}.
\]
Hence there will be $B=\dfrac{2\sqrt{k'}AK}{\pi}=2\sqrt[4]{q}\sqrt{\dfrac{k'}{k}}$. Now therefore there results:
\[
\begin{gathered}
\sin\operatorname{am}\frac{2Kx}{\pi}=\frac{1}{\sqrt{k}}.\,\frac{2\sqrt[4]{q}\sin x(1-2q^{2}\cos 2x+q^{4})(1-2q^{4}\cos 2x+q^{8})(1-2q^{6}\cos 2x+q^{12})\ldots}{(1-2q\cos 2x+q^{2})(1-2q^{3}\cos 2x+q^{6})(1-2q^{5}\cos 2x+q^{10})\ldots}
\\[2ex]
\cos\operatorname{am}\frac{2Kx}{\pi}=\sqrt{\frac{k'}{k}}.\frac{2\sqrt[4]{q}\cos x(1+2q^{2}\cos 2x+q^{4})(1+2q^{4}\cos 2x+q^{8})(1+2q^{6}\cos 2x+q^{12})\ldots}{(1-2q\cos 2x+q^{2})(1-2q^{3}\cos 2x+q^{6})(1-q^{5}\cos 2x+q^{10})\ldots}
\\[2ex]
\Delta\operatorname{am}\frac{2Kx}{\pi}=\sqrt{k'}.\frac{(1+2q\cos 2x+q^{2})(1+2q^{3}\cos 2x+q^{6})(1+2q^{5}\cos 2x+q^{10})\ldots}{(1-2q\cos 2x+q^{2})(1-2q^{3}\cos 2x+q^{6})(1-2q^{5}\cos 2x+q^{10})\ldots}.
\end{gathered}
\]\ednote{In the second of these three formulae the print has the last denominator factor as $(1-q^{5}\cos 2x+q^{10})$, the factor 2 before $q^{5}$ being missing (author's Corrigenda, p. 88 line 11: read $-2q^{5}$).}

The equations being multiplied together:
\[
\begin{aligned}
B&=2\sqrt[4]{q}\sqrt{\frac{k'}{k}}&&=\left\{\frac{(1-q)(1-q^{3})(1-q^{5})\,.\,.}{(1+q^{2})(1+q^{4})(1+q^{6})\,.\,.}\right\}^{2}\\[2ex]
C&=\sqrt{k'}&&=\left\{\frac{(1-q)(1-q^{3})(1-q^{5})\ldots}{(1+q)(1+q^{3})(1+q^{5})\ldots}\right\}^{2},
\end{aligned}
\]
there results:
\[
\frac{2\sqrt[4]{q}.k'}{\sqrt{k}}=\frac{\left\{(1-q)(1-q^{3})(1-q^{5})\,.\,.\right\}^{4}}{\left\{(1+q)(1+q^{2})(1+q^{3})\,.\,.\right\}^{2}}.
\]

\origpage{89}
Now indeed, according to \emph{Euler} in the \emph{Introd.} (\emph{on the Partition of Numbers}), there is:
\[
\begin{aligned}
(1+q)(1+q^{2})(1+q^{3})\ldots&=\frac{(1-q^{2})(1-q^{4})(1-q^{6})\ldots}{(1-q)(1-q^{2})(1-q^{3})\ldots}\\[2ex]
&=\frac{1}{(1-q)(1-q^{3})(1-q^{5})\ldots},
\end{aligned}
\]
whence we obtain:
\[
\text{1)}\quad\left\{(1-q)(1-q^{3})(1-q^{5})(1-q^{7})\,.\,.\right\}^{6}=\frac{2\sqrt[4]{q}\,k'}{\sqrt{k}}.
\]
The formula being called in:
\[
A=\frac{\pi\sqrt[4]{q}}{\sqrt{k.K}}=\left\{\frac{(1-q)(1-q^{3})(1-q^{5})\,.\,.}{(1-q^{2})(1-q^{4})(1-q^{6})\,.\,.}\right\}^{2},
\]
there is:
\[
\begin{gathered}
\text{2)}\quad\left\{(1-q^{2})(1-q^{4})(1-q^{6})(1-q^{8})\,.\,.\right\}^{6}=\frac{2kk'K^{3}}{\pi^{3}\sqrt{q}},\quad\text{whence also:}
\\[2ex]
\text{3)}\quad\left\{(1-q)(1-q^{2})(1-q^{3})(1-q^{4})\,.\,.\right\}^{6}=\frac{4\sqrt{k}\,k'k'K^{3}}{\pi^{3}\sqrt[4]{q}}.
\end{gathered}
\]
To these one may add the formulae which follow easily:
\[
\begin{gathered}
\text{4)}\quad\left\{(1+q)(1+q^{3})(1+q^{5})(1+q^{7})\,.\,.\right\}^{6}=\frac{2\sqrt[4]{q}}{\sqrt{kk'}}
\\[2ex]
\text{5)}\quad\left\{(1+q^{2})(1+q^{4})(1+q^{6})(1+q^{8})\,.\,.\right\}^{6}=\frac{k}{4\sqrt{k'}\sqrt{q}}
\\[2ex]
\text{6)}\quad\left\{(1+q)(1+q^{2})(1+q^{3})(1+q^{4})\,.\,.\right\}^{6}=\frac{\sqrt{k}}{2k'\sqrt[4]{q}}.
\end{gathered}
\]
From these is also gathered:
\[
\begin{gathered}
\text{7)}\quad k=4\sqrt{q}\left\{\frac{(1+q^{2})(1+q^{4})(1+q^{6})\,.\,.}{(1+q)(1+q^{3})(1+q^{5})\,.\,.}\right\}^{4}
\\[2ex]
\text{8)}\quad k'=\left\{\frac{(1-q)(1-q^{3})(1-q^{5})\,.\,.}{(1+q)(1+q^{3})(1+q^{5})\,.\,.}\right\}^{4}
\\[2ex]
\text{9)}\quad\frac{2K}{\pi}=\left\{\frac{(1-q^{2})(1-q^{4})(1-q^{6})\,.\,.}{(1-q)(1-q^{3})(1-q^{5})\,.\,.}\right\}^{2}\left\{\frac{(1+q)(1+q^{3})(1+q^{5})\,.\,.}{(1+q^{2})(1+q^{4})(1+q^{6})\,.\,.}\right\}^{2}
\\[2ex]
\text{10)}\quad\frac{2kK}{\pi}=4\sqrt{q}\left\{\frac{(1-q^{4})(1-q^{8})(1-q^{12})\,.\,.}{(1-q^{2})(1-q^{6})(1-q^{10})\,.\,.}\right\}^{2}
\end{gathered}
\]

\origpage{90}
\[
\begin{gathered}
\text{11)}\quad\frac{2k'K}{\pi}=\left\{\frac{(1-q)(1-q^{2})(1-q^{3})\,.\,.}{(1+q)(1+q^{2})(1+q^{3})\,.\,.}\right\}^{2}
\\[2ex]
\text{12)}\quad\frac{2\sqrt{k.K}}{\pi}=2\sqrt[4]{q}\left\{\frac{(1-q^{2})(1-q^{4})(1-q^{6})\,.\,.}{(1-q)(1-q^{3})(1-q^{5})\,.\,.}\right\}^{2}
\\[2ex]
\text{13)}\quad\frac{2\sqrt{k'}K}{\pi}=\left\{\frac{(1-q^{2})(1-q^{4})(1-q^{6})\ldots}{(1+q^{2})(1+q^{4})(1+q^{6})\ldots}\right\}^{2}.
\end{gathered}
\]

From formulae 7), 8) follows a fairly abstruse identical equation:
\[
\begin{gathered}
\text{14)}\quad\left\{(1-q)(1-q^{3})(1-q^{5})\,.\,.\right\}^{8}+16q\left\{(1+q^{2})(1+q^{4})(1+q^{6})\,.\,.\right\}^{8}
\\[1ex]
=
\\[1ex]
\left\{(1+q)(1+q^{3})(1+q^{5})\,.\,.\right\}^{8}.
\end{gathered}
\]

\subsection*{37.}

We have seen above, where the properties of the Modular equations were dealt with, that when $k$ is changed into $\dfrac{1}{k}$, $K$ passes into $k(K+iK')$, $K'$ into $kK'$; further that there results:
\[
\begin{gathered}
\sin\operatorname{am}\left(ku,\frac{ik'}{k}\right)=\cos\operatorname{coam}(u,k')
\\[2ex]
\cos\operatorname{am}\left(ku,\frac{ik'}{k}\right)=\sin\operatorname{coam}(u,k')
\\[2ex]
\Delta\operatorname{am}\left(ku,\frac{ik'}{k}\right)=\frac{1}{\Delta\operatorname{am}(u,k')}.
\end{gathered}
\]
$k$ and $k'$ being interchanged, it follows hence, where $k'$ passes into $\dfrac{1}{k'}$ or $k$ into $\dfrac{ik}{k'}$, that at the same time $K$ passes into $k'K$, $K'$ into $k'(K'+iK)$; further that there results:
\[
\begin{gathered}
\sin\operatorname{am}\left(k'u,\frac{ik}{k'}\right)=\cos\operatorname{coam}u
\\[2ex]
\cos\operatorname{am}\left(k'u,\frac{ik}{k'}\right)=\sin\operatorname{coam}u
\\[2ex]
\Delta\operatorname{am}\left(k'u,\frac{ik}{k'}\right)=\frac{1}{\Delta\operatorname{am}u},
\end{gathered}
\]
whence also:
\[
\operatorname{am}\left(k'u,\frac{ik}{k'}\right)=\frac{\pi}{2}-\operatorname{coam}u.
\]

\origpage{91}
But $K$ being changed into $k'K$, $K'$ into $k'(K'+iK)$, $q=e^{-\frac{\pi K'}{K}}$ passes into $-q$, whence conversely flows

\subsection*{THEOREM I.}

When $q$ is changed into $-q$ there pass:
\[
\begin{gathered}
k\text{ into }\frac{ik}{k'},\quad k'\text{ into }\frac{1}{k'}
\\[2ex]
K\text{ into }k'K,\quad K'\text{ into }k'(K'+iK)
\\[2ex]
\sin\operatorname{am}\frac{2Kx}{\pi}\text{ into }\cos\operatorname{coam}\frac{2Kx}{\pi}
\\[2ex]
\cos\operatorname{am}\frac{2Kx}{\pi}\text{ into }\sin\operatorname{coam}\frac{2Kx}{\pi}
\\[2ex]
\Delta\operatorname{am}\frac{2Kx}{\pi}\text{ into }\frac{1}{\Delta\operatorname{am}\dfrac{2Kx}{\pi}}
\\[2ex]
\operatorname{am}\frac{2Kx}{\pi}\text{ into }\frac{\pi}{2}-\operatorname{coam}\frac{2Kx}{\pi};
\end{gathered}
\]
when at the same time $q$ is changed into $-q$, $x$ into $\dfrac{\pi}{2}-x$, there pass:
\[
\begin{gathered}
\operatorname{am}\frac{2Kx}{\pi}\text{ into }\frac{\pi}{2}-\operatorname{am}\frac{2Kx}{\pi}
\\[2ex]
\sin\operatorname{am}\frac{2Kx}{\pi}\text{ into }\cos\operatorname{am}\frac{2Kx}{\pi}
\\[2ex]
\cos\operatorname{am}\frac{2Kx}{\pi}\text{ into }\sin\operatorname{am}\frac{2Kx}{\pi}
\\[2ex]
\Delta\operatorname{am}\frac{2Kx}{\pi}\text{ into }\frac{1}{k'}\Delta\operatorname{am}\frac{2Kx}{\pi}.
\end{gathered}
\]

Let us inquire further what changes the Elliptic Functions undergo when $q$ is changed either into $q^{2}$ or into $\sqrt{q}$.

We have seen above that the Modulus $\lambda$, derived from the Modulus $k$ by the real first transformation of the $n^{\mathrm{ti}}$ order, enjoys this remarkable property, that there is:
\[
\frac{\Lambda'}{\Lambda}=n.\frac{K'}{K};
\]
whence, $k$ being changed into $\lambda$, $q=e^{-\frac{K'\pi}{K}}$ passes into $q^{n}$. The same, which we have proved generally for the transformations

\origpage{92}
of odd order, was already long ago proved by Cl.\ Legendre for the transformation of the second order, namely, having put $\lambda=\dfrac{1-k'}{1+k'}$, there results:
\[
\Lambda=\left(\frac{1+k'}{2}\right)K,\quad\Lambda'=(1+k')K',\quad\frac{\Lambda'}{\Lambda}=2.\frac{K'}{K},
\]
whence we see that, $k$ being changed into $\dfrac{1-k'}{1+k'}$, $q$ passes into $q^{2}$. Hence conversely we obtain

\subsection*{THEOREM II.}

„When $q$ is changed into $q^{2}$, $k$ passes into $\dfrac{1-k'}{1+k'}$, $K$ into $\left(\dfrac{1+k'}{2}\right)K$,"

whence also:
\[
\begin{gathered}
k'\text{ into }\frac{2\sqrt{k'}}{1+k'}\qquad\qquad 1+k\text{ into }\frac{2}{1+k'}
\\[2ex]
k'K\text{ into }\sqrt{k'}K\qquad\qquad 1-k\text{ into }\frac{2k'}{1+k'}
\\[2ex]
\sqrt{k}\text{ into }\frac{k}{1+k'}\qquad\qquad 1+k'\text{ into }\frac{(1+\sqrt{k'})^{2}}{1+k'}
\\[2ex]
\sqrt{k}.K\text{ into }\frac{kK}{2}\qquad\qquad 1-k'\text{ into }\frac{(1-\sqrt{k'})^{2}}{1+k'}.
\end{gathered}
\]
From the inversion of this theorem the second is obtained

\subsection*{THEOREM III.}

„When $q$ is changed into $\sqrt{q}$, $k$ passes into $\dfrac{2\sqrt{k}}{1+k}$, $K$ into $(1+k)K$,"

whence also:
\[
\begin{gathered}
k'\text{ into }\frac{1-k}{1+k}\qquad\qquad 1+k\text{ into }\frac{(1+\sqrt{k})^{2}}{1+k}
\\[2ex]
\sqrt{k'}\text{ into }\frac{k'}{1+k}\qquad\qquad 1-k\text{ into }\frac{(1-\sqrt{k})^{2}}{1+k}
\\[2ex]
kK\text{ into }2\sqrt{k}.K\qquad\qquad 1+k'\text{ into }\frac{2}{1+k}
\\[2ex]
\sqrt{k'}K\text{ into }k'K\qquad\qquad 1-k'\text{ into }\frac{2k}{1+k}.
\end{gathered}
\]
These three theorems are confirmed in many ways by the expansions proposed in §§.\ 35. 36, and find in what follows their most frequent application. Indeed by their aid one may either derive other formulae from others, or those found elsewhere are conveniently confirmed.

\origpage{93}
\subsection*{38.}

Let us designate the quantities into which $k$, $k'$, $K$ pass when $q^{m}$ is put in place of $q$ by $k^{(m)}$, $k^{(m)\prime}$, $K^{(m)}$, so that $k^{(m)}$ is the Modulus elicited by the real first transformation of the $n^{\mathrm{ti}}$ order, and $k^{(m)\prime}$ its complement. Let us put in the equation:
\[
\sqrt{k'}=\left\{\frac{(1-q)(1-q^{3})(1-q^{5})(1-q^{7})\ldots}{(1+q)(1+q^{3})(1+q^{5})(1+q^{7})\ldots}\right\}^{2}
\]
in place of $q$ successively $q^{2}$, $q^{4}$, $q^{8}$, $q^{16}$, etc.; there results, after infinite multiplication has been made:
\[
\sqrt{k^{(2)\prime}k^{(4)\prime}k^{(8)\prime}k^{(16)\prime}\ldots}=\left\{\frac{(1-q^{2})(1-q^{4})(1-q^{6})(1-q^{8})\ldots}{(1+q^{2})(1+q^{4})(1+q^{6})(1+q^{8})\ldots}\right\}^{2};
\]
but we have found:
\[
\left\{\frac{(1-q^{2})(1-q^{4})(1-q^{6})(1-q^{8})\ldots}{(1+q^{2})(1+q^{4})(1+q^{6})(1+q^{8})\ldots}\right\}^{2}=\frac{2\sqrt{k'}K}{\pi},
\]
whence:
\[
\text{1)}\quad\frac{2K}{\pi}=\sqrt{\frac{k^{(2)\prime}k^{(4)\prime}k^{(8)\prime}k^{(16)\prime}.\,.\,.\,.}{k'}}.
\]
Since $k^{(2)\prime}=\dfrac{2\sqrt{k'}}{1+k'}$, there is from 1):
\[
\left(\frac{2K}{\pi}\right)^{2}=\frac{1}{k'}.\,\frac{2\sqrt{k'}}{1+k'}.\,\frac{2\sqrt{k^{(2)\prime}}}{1+k^{(2)\prime}}.\,\frac{2\sqrt{k^{(4)\prime}}}{1+k^{(4)\prime}}.\,\frac{2\sqrt{k^{(8)\prime}}}{1+k^{(8)\prime}}\ldots.
\]
whence, division having been made by 1):
\[
\text{2)}\quad\frac{2K}{\pi}=\frac{2}{1+k'}.\,\frac{2}{1+k^{(2)\prime}}.\,\frac{2}{1+k^{(4)\prime}}.\,\frac{2}{1+k^{(8)\prime}}\ldots
\]
This formula is also obtained from the fact that there is:
\[
\begin{gathered}
\frac{2K}{\pi}=\frac{2K^{(2)}}{\pi}.\,\frac{2}{1+k'}
\\[2ex]
\frac{2K^{(2)}}{\pi}=\frac{2K^{(4)}}{\pi}.\,\frac{2}{1+k^{(2)\prime}}
\\[2ex]
\frac{2K^{(4)}}{\pi}=\frac{2K^{(8)}}{\pi}.\,\frac{2}{1+k^{(4)\prime}}
\\[2ex]
\ldots,
\end{gathered}
\]

\origpage{94}
whence, since, as $m$ increases to infinity, the limit of the expression $\dfrac{2K^{(m)}}{\pi}$ is 1, when the infinite product is formed, 2) results. Having put:
\[
\begin{gathered}
m=1,\quad n=k'
\\[2ex]
m'=\frac{m+n}{2},\quad n'=\sqrt{mn}
\\[2ex]
m''=\frac{m'+n'}{2},\quad n''=\sqrt{m'n'}
\\[2ex]
m'''=\frac{m''+n''}{2},\quad n'''=\sqrt{m''n''}
\\[2ex]
\ldots,
\end{gathered}
\]
there results:
\[
\begin{gathered}
k^{(2)\prime}=\frac{2\sqrt{k'}}{1+k'}=\frac{n'}{m'}
\\[2ex]
k^{(4)\prime}=\frac{2\sqrt{k^{(2)\prime}}}{1+k^{(2)\prime}}=\frac{n''}{m''}
\\[2ex]
k^{(6)\prime}=\frac{2\sqrt{k^{(4)\prime}}}{1+k^{(4)\prime}}=\frac{n'''}{m'''}
\\[2ex]
\ldots,
\end{gathered}
\]\ednote{The third line is printed with $k^{(6)\prime}$ on the left, where $k^{(8)\prime}$ is expected (the Corrigenda correct it to $k^{(8)\prime}$; a hand correction in ink is also visible on the scan).}
whence:
\[
\frac{2}{1+k'}=\frac{m}{m'},\quad\frac{2}{1+k^{(2)\prime}}=\frac{m'}{m''},\quad\frac{2}{1+k^{(4)\prime}}=\frac{m''}{m'''},
\]
and therefore:
\[
\frac{2K}{\pi}=\frac{m}{m'}.\,\frac{m'}{m''}.\,\frac{m''}{m'''}.\,\frac{m'''}{m''''}\ldots;
\]
or, if $\mu$ denotes the common limit to which $m^{(p)}$, $n^{(p)}$ converge as $n$ increases to infinity:\ednote{Printed $n$ as the growing index; presumably $p$ (the index of $m^{(p)}$ and $n^{(p)}$) is meant.}
\[
\text{3)}\quad\frac{2K}{\pi}=\frac{1}{\mu}.
\]
These things are abundantly well known.

Let us put again in the formula:
\[
\Delta\operatorname{am}\frac{2Kx}{\pi}=\sqrt{k'}\,\frac{(1+2q\cos 2x+q^{2})(1+2q^{3}\cos 2x+q^{6})(1+2q^{5}\cos 2x+q^{10})\ldots}{(1-2q\cos 2x+q^{2})(1-2q^{3}\cos 2x+q^{6})(1-2q^{5}\cos 2x+q^{10})\ldots}
\]

\origpage{95}
in place of $q$ successively $q^{2}$, $q^{4}$, $q^{8}$, etc.; and further let:
\[
S=\Delta\operatorname{am}\left(\frac{2K^{(2)}x}{\pi},k^{(2)}\right)\Delta\operatorname{am}\left(\frac{2K^{(4)}x}{\pi},k^{(4)}\right)\Delta\operatorname{am}\left(\frac{2K^{(8)}x}{\pi},k^{(8)}\right).\,.\,.\,.
\]
When the infinite product is formed, since:
\[
\frac{2\sqrt{k'}K}{\pi}=\sqrt{k^{(2)\prime}k^{(4)\prime}k^{(8)\prime}k^{(16)\prime}\ldots},
\]
we obtain:
\[
S=\frac{2\sqrt{k'}K}{\pi}\,\frac{(1+2q^{2}\cos 2x+q^{4})(1+2q^{4}\cos 2x+q^{8})(1+2q^{6}\cos 2x+q^{12})\ldots}{(1-2q^{2}\cos 2x+q^{4})(1-2q^{4}\cos 2x+q^{8})(1-2q^{6}\cos 2x+q^{12})\ldots}.
\]
Now from the formulae:
\[
\begin{gathered}
\sin\operatorname{am}\frac{2Kx}{\pi}=
\\[1ex]
\frac{2}{\sqrt{k}}\,\frac{\sqrt[4]{q}\sin x(1-2q^{2}\cos 2x+q^{4})(1-2q^{4}\cos 2x+q^{8})(1-2q^{6}\cos 2x+q^{12})\ldots}{(1-2q\cos 2x+q^{2})(1-2q^{3}\cos 2x+q^{6})(1-2q^{5}\cos 2x+q^{10})\ldots}
\\[2ex]
\cos\operatorname{am}\frac{2Kx}{\pi}=
\\[1ex]
2\sqrt{\frac{k'}{k}}\,\frac{\sqrt[4]{q}\cos x(1+2q^{2}\cos 2x+q^{4})(1+2q^{4}\cos 2x+q^{8})(1+2q^{6}\cos 2x+q^{12})\ldots}{(1-2q\cos 2x+q^{2})(1-2q^{3}\cos 2x+q^{6})(1-2q^{5}\cos 2x+q^{10})\ldots},
\end{gathered}
\]
we obtain:
\[
\begin{gathered}
\operatorname{tang}\operatorname{am}\frac{2Kx}{\pi}=
\\[1ex]
\frac{1}{\sqrt{k'}}\,\frac{\operatorname{tang}.x(1-2q^{2}\cos 2x+q^{4})(1-2q^{4}\cos 2x+q^{8})(1-2q^{6}\cos 2x+q^{12})\ldots}{(1+2q^{2}\cos 2x+q^{4})(1+2q^{4}\cos 2x+q^{8})(1+2q^{6}\cos 2x+q^{12})\ldots},
\end{gathered}
\]
whence the memorable formula results:
\[
\text{4)}\quad\operatorname{tang}x=\frac{S.\operatorname{tg}\operatorname{am}\dfrac{2Kx}{\pi}}{\dfrac{2K}{\pi}}.
\]
To demonstrate the same by known formulae, let us call in the formula for the transformation of the second order, such as Cl.\ \emph{Gauss} has exhibited in the Memoir entitled: „\emph{Determinatio Attractionis}" etc.:
\[
\sin\operatorname{am}\frac{2Kx}{\pi}=\frac{(1+k^{(2)})\sin\operatorname{am}\left(\dfrac{2K^{(2)}x}{\pi},k^{(2)}\right)}{1+k^{(2)}\sin^{2}\operatorname{am}\left(\dfrac{2K^{(2)}x}{\pi},k^{(2)}\right)},
\]

\origpage{96}
which, for the sake of brevity, on putting:
\[
\operatorname{am}\left(\frac{2K^{(m)}x}{\pi},k^{(m)}\right)=\phi^{(m)},\quad\Delta\operatorname{am}\left(\frac{2K^{(m)}x}{\pi},k^{(m)}\right)=\Delta^{(m)},
\]
is thus exhibited:
\[
\sin\phi=\frac{(1+k^{(2)})\sin\phi^{(2)}}{1+k^{(2)}\sin^{2}\phi^{(2)}},
\]
whence also:
\[
\begin{gathered}
\cos\phi=\frac{\cos\phi^{(2)}\Delta^{(2)}}{1+k^{(2)}\sin^{2}\phi^{(2)}}
\\[2ex]
\Delta(\phi)=\frac{1-k^{(2)}\sin^{2}\phi^{(2)}}{1+k^{(2)}\sin^{2}\phi^{(2)}}
\\[2ex]
\operatorname{tg}\phi=\frac{(1+k^{(2)})\operatorname{tg}\phi^{(2)}}{\Delta^{(2)}}.
\end{gathered}
\]
The last formula can also be represented thus:
\[
\frac{\operatorname{tg}\phi}{\dfrac{2K}{\pi}}=\frac{\operatorname{tg}\phi^{(2)}}{\dfrac{2K^{(2)}}{\pi}}.\,\frac{1}{\Delta^{(2)}},
\]
whence, on putting successively $q^{2}$, $q^{4}$, $q^{8}$, $\ldots$ in place of $q$, by which $k$, $K$, $\phi$ pass into $k^{(2)}$, $k^{(4)}$, $k^{(8)}$, $.\,.$; $K^{(2)}$, $K^{(4)}$, $K^{(8)}$, $\ldots$; $\phi^{(2)}$, $\phi^{(4)}$, $\phi^{(8)}$, $\ldots$, we obtain:
\[
\begin{gathered}
\frac{\operatorname{tg}\phi^{(2)}}{\dfrac{2K^{(2)}}{\pi}}=\frac{\operatorname{tg}\phi^{(4)}}{\dfrac{2K^{(4)}}{\pi}}.\,\frac{1}{\Delta^{(4)}}
\\[2ex]
\frac{\operatorname{tg}\phi^{(4)}}{\dfrac{2K^{(4)}}{\pi}}=\frac{\operatorname{tg}\phi^{(8)}}{\dfrac{2K^{(8)}}{\pi}}.\,\frac{1}{\Delta^{(8)}}
\\[2ex]
\frac{\operatorname{tg}\phi^{(8)}}{\dfrac{2K^{(8)}}{\pi}}=\frac{\operatorname{tg}\phi^{(16)}}{\dfrac{2K^{(16)}}{\pi}}.\,\frac{1}{\Delta^{(16)}}
\\[2ex]
\ldots
\end{gathered}
\]
Now the limit of the expression
\[
\frac{\operatorname{tg}\phi^{(p)}}{\dfrac{2K^{(p)}}{\pi}}=\frac{\operatorname{tg}\operatorname{am}\left(\dfrac{2K^{(p)}x}{\pi},k^{(p)}\right)}{\dfrac{2K^{(p)}}{\pi}},
\]

\origpage{97}
as $p$ increases to infinity, is
\[
\operatorname{tang}x;
\]
for then $k^{(p)}=0$, $K^{(p)}=\dfrac{\pi}{2}$, $\operatorname{am}(u,k)=u$; whence, when the infinite product is now formed and, as above, $S=\Delta^{(2)}\Delta^{(4)}\Delta^{(8)}\ldots$ is put, there results:
\[
\frac{\operatorname{tg}\phi}{\dfrac{2K}{\pi}}=\frac{\operatorname{tg}x}{S},
\]
which is the formula to be demonstrated.

From the formula:
\[
\operatorname{tg}x=\frac{S\operatorname{tg}\phi}{\dfrac{2K}{\pi}}
\]
a not inelegant Algorithm can be sought for computing the \emph{indefinite} Elliptic Integrals of the first species; and this by the aid of the formula, easy to prove:
\[
\Delta^{(2)}=\sqrt{\frac{2(\Delta+k')}{(1+k')(1+\Delta)}}.
\]
To this end we propose

\subsection*{THEOREM.}

Having put
\[
\begin{gathered}
\int_{0}^{\Phi}\frac{d\phi}{\sqrt{mm\operatorname{Cos}\phi^{2}+nn\operatorname{Sin}\phi^{2}}}=\Phi\\[1.5ex]
\sqrt{mm\operatorname{Cos}\phi^{2}+nn\operatorname{Sin}\phi^{2}}=\Delta,
\end{gathered}
\]
let the expressions be formed:
\[
\begin{gathered}
\frac{m+n}{2}=m'\qquad\sqrt{mn}=n'\qquad\Delta'=\sqrt{\frac{mm'(\Delta+n)}{m+\Delta}}\\[1.5ex]
\frac{m'+n'}{2}=m''\qquad\sqrt{m'n'}=n''\qquad\Delta''=\sqrt{\frac{m'm''(\Delta'+n')}{m'+\Delta'}}\\[1.5ex]
\frac{m''+n''}{2}=m'''\qquad\sqrt{m''n''}=n'''\qquad\Delta'''=\sqrt{\frac{m''m'''(\Delta''+n'')}{m''+\Delta''}}\\[1.5ex]
.\qquad\qquad.\qquad\qquad.\qquad\qquad.\;,
\end{gathered}
\]

\origpage{98}
and let $\mu$ denote the common limit to which the quantities $m^{(p)}$, $\Delta^{(p)}$, $n^{(p)}$ converge most rapidly as $p$ increases; there will be:
\[
\operatorname{tang}\mu\Phi=\frac{\Delta'\Delta''\Delta'''\ldots}{m\,m'm''\ldots}.\operatorname{tang}\Phi.
\]

By the same methods which we have used in what precedes, the value of the infinite product
\[
\frac{2\sqrt[4]{q}}{\sqrt{k}}.\,\frac{2\sqrt[4]{q^{2}}}{\sqrt{k^{(2)}}}.\,\frac{2\sqrt[4]{q^{4}}}{\sqrt{k^{(4)}}}.\,\frac{2\sqrt[4]{q^{8}}}{\sqrt{k^{(8)}}}\ldots
\]
is also found. To this end we cite the formulae of §. 36, 4), 5):
\[
\begin{gathered}
\left\{(1+q)(1+q^{3})(1+q^{5})(1+q^{7})\ldots\right\}^{6}=\frac{2\sqrt[4]{q}}{\sqrt{kk'}}\\[1.5ex]
\left\{(1+q^{2})(1+q^{4})(1+q^{6})(1+q^{8})\ldots\right\}^{6}=\frac{k}{4\sqrt{k'}\sqrt{q}},
\end{gathered}
\]
of which the latter arises from the former by putting successively $q^{2}$, $q^{4}$, $q^{8}$ etc.\ in place of $q$ and forming the infinite product, whence we obtain:
\[
\frac{k}{4\sqrt{k'}\sqrt{q}}=\frac{2\sqrt[4]{q^{2}}}{\sqrt{k^{(2)}k^{(2)\prime}}}.\,\frac{2\sqrt[4]{q^{4}}}{\sqrt{k^{(4)}k^{(4)\prime}}}.\,\frac{2\sqrt[4]{q^{8}}}{\sqrt{k^{(8)}k^{(8)\prime}}}.\,.
\]
Now we have derived 1):
\[
\frac{2K}{\pi}=\sqrt{\frac{k^{(2)\prime}k^{(4)\prime}k^{(8)\prime}.\,.}{k'}},
\]
whence:
\[
\text{5)}\quad\frac{\sqrt{k}}{2\sqrt[4]{q}}.\,\frac{2K}{\pi}=\frac{2\sqrt[4]{q}}{\sqrt{k}}.\,\frac{2\sqrt[4]{q^{2}}}{\sqrt{k^{(2)}}}.\,\frac{2\sqrt[4]{q^{4}}}{\sqrt{k^{(4)}}}.\,\frac{2\sqrt[4]{q^{8}}}{\sqrt{k^{(8)}}}\ldots
\]
Although these things may seem foreign to our undertaking, since they are not lacking in elegance and contribute greatly to perceiving the nature of the expansions proposed, it is pleasing to have set them down \ednote{So printed (opposuisse); the author's Corrigenda read apposuisse.}.

\origpage{99}
\subsection*{EXPANSION OF ELLIPTIC FUNCTIONS INTO SERIES PROCEEDING ACCORDING TO THE SINES OR COSINES OF MULTIPLES OF THE ARGUMENT.}

\subsection*{39.}

From the formulae delivered above:
\[
\begin{gathered}
\text{1)}\quad\sin\operatorname{am}\frac{2Kx}{\pi}=\frac{2\sqrt[4]{q}}{\sqrt{k}}\sin x.\frac{(1-2q^{2}\cos2x+q^{4})(1-2q^{4}\cos2x+q^{8})(1-2q^{6}\cos2x+q^{12})\ldots}{(1-2q\cos2x+q^{2})(1-2q^{3}\cos2x+q^{6})(1-2q^{5}\cos2x+q^{10})\ldots}\\[1.5ex]
\text{2)}\quad\cos\operatorname{am}\frac{2Kx}{\pi}=\frac{2\sqrt[4]{q}\sqrt{k'}}{\sqrt{k}}\cos x.\frac{(1+2q^{2}\cos2x+q^{4})(1+2q^{4}\cos2x+q^{8})(1+2q^{6}\cos2x+q^{12})\ldots}{(1-2q\cos2x+q^{2})(1-2q^{3}\cos2x+q^{6})(1-2q^{5}\cos2x+q^{10})\ldots}\\[1.5ex]
\text{3)}\quad\Delta\operatorname{am}\frac{2Kx}{\pi}=\sqrt{k'}\,\frac{(1+2q\cos2x+q^{2})(1+2q^{3}\cos2x+q^{6})(1+2q^{5}\cos2x+q^{10})\ldots}{(1-2q\cos2x+q^{2})(1-2q^{3}\cos2x+q^{6})(1-2q^{5}\cos2x+q^{10})\ldots}\\[1.5ex]
\text{4)}\quad\sqrt{\frac{1-\sin\operatorname{am}\dfrac{2Kx}{\pi}}{1+\sin\operatorname{am}\dfrac{2Kx}{\pi}}}=\sqrt{\frac{1-\sin x}{1+\sin x}}.\frac{(1-2q\sin x+q^{2})(1-2q^{2}\sin x+q^{4})(1-2q^{3}\sin x+q^{6})\ldots}{(1+2q\sin x+q^{2})(1+2q^{2}\sin x+q^{4})(1+2q^{3}\sin x+q^{6})\ldots}\\[1.5ex]
\text{5)}\quad\sqrt{\frac{1-k\sin\operatorname{am}\dfrac{2Kx}{\pi}}{1+k\sin\operatorname{am}\dfrac{2Kx}{\pi}}}=\sqrt{k'}\,\frac{(1-2\sqrt{q}\sin x+q)(1-2\sqrt{q^{3}}\sin x+q^{3})(1-2\sqrt{q^{5}}\sin x+q^{5})\ldots}{(1+2\sqrt{q}\sin x+q)(1+2\sqrt{q^{3}}\sin x+q^{3})(1+2\sqrt{q^{5}}\sin x+q^{5})\ldots},
\end{gathered}
\]\ednote{In 4) the exponent 3 of $q$ in the third numerator factor is lost in a blot in the scan; read as $2q^{3}\sin x$ by analogy with the denominator.}
the logarithms of the single factors in either member of the equations being expanded, after obvious reductions, the following ensue:
\[
\begin{gathered}
\text{6)}\quad\log\sin\operatorname{am}\frac{2Kx}{\pi}=\log\left\{\frac{2\sqrt[4]{q}\sin x}{\sqrt{k}}\right\}+\frac{2q\cos2x}{1+q}+\frac{2q^{2}\cos4x}{2(1+q^{2})}+\frac{2q^{3}\cos6x}{3(1+q^{3})}+\ldots\\[1.5ex]
\text{7)}\quad\log\cos\operatorname{am}\frac{2Kx}{\pi}=\log\left\{2\sqrt[4]{q}\sqrt{\frac{k'}{k}}\cos x\right\}+\frac{2q\cos2x}{1-q}+\frac{2q^{2}\cos4x}{2(1+q^{2})}+\frac{2q^{3}\cos6x}{3(1-q^{3})}+\ldots\\[1.5ex]
\text{8)}\quad\log\Delta\operatorname{am}\frac{2Kx}{\pi}=\log\sqrt{k'}+\frac{4q\cos2x}{1-q^{2}}+\frac{4q^{3}\cos6x}{3(1-q^{6})}+\frac{4q^{5}\cos10x}{5(1-q^{10})}+\ldots\\[1.5ex]
\text{9)}\quad\log\sqrt{\frac{1+\sin\operatorname{am}\dfrac{2Kx}{\pi}}{1-\sin\operatorname{am}\dfrac{2Kx}{\pi}}}=\log\sqrt{\frac{1+\sin x}{1-\sin x}}+\frac{4q\sin x}{1-q}-\frac{4q^{3}\sin3x}{3(1-q^{3})}+\frac{4q^{5}\sin5x}{5(1-q^{5})}-\ldots
\end{gathered}
\]

\origpage{100}
\[
\text{10)}\quad\log\sqrt{\frac{1-k\sin\operatorname{am}\dfrac{2Kx}{\pi}}{1+k\sin\operatorname{am}\dfrac{2Kx}{\pi}}}=\frac{4\sqrt{q}\sin x}{1-q}-\frac{4\sqrt{q^{3}}\sin3x}{3(1-q^{3})}+\frac{4\sqrt{q^{5}}\sin5x}{5(1-q^{5})}-\ldots
\]
By differentiating these formulae, where we note the differential formulae, easy to prove,
\[
\begin{gathered}
\frac{d.\log\sin\operatorname{am}\dfrac{2Kx}{\pi}}{dx}=\frac{2k'K}{\pi}.\frac{\cos\operatorname{am}\dfrac{2Kx}{\pi}}{\cos\operatorname{coam}\dfrac{2Kx}{\pi}}\\[2ex]
-\frac{d.\log\cos\operatorname{am}\dfrac{2Kx}{\pi}}{dx}=\frac{2K}{\pi}.\frac{\sin\operatorname{am}\dfrac{2Kx}{\pi}}{\sin\operatorname{coam}\dfrac{2Kx}{\pi}}=\frac{2K}{\pi}.\operatorname{tg}\left(\frac{\operatorname{am}\dfrac{4Kx}{\pi}}{2}\right)\\[2ex]
-\frac{d.\log\Delta\operatorname{am}\dfrac{2Kx}{\pi}}{dx}=\frac{2k^{2}K}{\pi}.\sin\operatorname{am}\frac{2Kx}{\pi}\sin\operatorname{coam}\frac{2Kx}{\pi}\\[2ex]
\frac{d.\log\sqrt{\dfrac{1+\sin\operatorname{am}\dfrac{2Kx}{\pi}}{1-\sin\operatorname{am}\dfrac{2Kx}{\pi}}}}{dx}=\frac{2K}{\pi}.\frac{1}{\sin\operatorname{coam}\dfrac{2Kx}{\pi}}\\[2ex]
\frac{d.\log\sqrt{\dfrac{1+k\sin\operatorname{am}\dfrac{2Kx}{\pi}}{1-k\sin\operatorname{am}\dfrac{2Kx}{\pi}}}}{dx}=\frac{2kK}{\pi}.\sin\operatorname{coam}\frac{2Kx}{\pi},
\end{gathered}
\]
we shall derive the following:
\[
\begin{gathered}
\text{11)}\quad\frac{2k'K}{\pi}.\frac{\cos\operatorname{am}\dfrac{2Kx}{\pi}}{\cos\operatorname{coam}\dfrac{2Kx}{\pi}}=\operatorname{cotg}x-\frac{4q\sin2x}{1+q}-\frac{4q^{2}\sin4x}{1+q^{2}}-\frac{4q^{3}\sin6x}{1+q^{3}}-.\,.\\[2ex]
\text{12)}\quad\frac{2K}{\pi}.\frac{\sin\operatorname{am}\dfrac{2Kx}{\pi}}{\sin\operatorname{coam}\dfrac{2Kx}{\pi}}=\operatorname{tg}x+\frac{4q\sin2x}{1-q}+\frac{4q^{2}\sin4x}{1+q^{2}}+\frac{4q^{3}\sin6x}{1-q^{3}}+.\,.\\[2ex]
\text{13)}\quad\frac{2k^{2}K}{\pi}.\sin\operatorname{am}\frac{2Kx}{\pi}\sin\operatorname{coam}\frac{2Kx}{\pi}=\frac{8q\sin2x}{1-q^{2}}+\frac{8q^{3}\sin6x}{1-q^{6}}+\frac{8q^{5}\sin10x}{1-q^{10}}+.\,.
\end{gathered}
\]

\origpage{101}
\[
\begin{gathered}
\text{14)}\quad\frac{2K}{\pi\sin\operatorname{coam}\dfrac{2Kx}{\pi}}=\frac{1}{\cos x}+\frac{4q\cos x}{1-q}-\frac{4q^{3}\cos3x}{1-q^{3}}+\frac{4q^{5}\cos5x}{1-q^{5}}-.\,.\\[2ex]
\text{15)}\quad\frac{2kK}{\pi}.\sin\operatorname{coam}\frac{2Kx}{\pi}=\frac{4\sqrt{q}\cos x}{1-q}-\frac{4\sqrt{q^{3}}\cos3x}{1-q^{3}}+\frac{4\sqrt{q^{5}}\cos5x}{1-q^{5}}-.\,.
\end{gathered}
\]

When in these formulae $\dfrac{\pi}{2}-x$ is put in place of $x$, there is derived:
\[
\begin{gathered}
\text{16)}\quad\frac{2k'K}{\pi}.\frac{\cos\operatorname{coam}\dfrac{2Kx}{\pi}}{\cos\operatorname{am}\dfrac{2Kx}{\pi}}=\operatorname{tg}x-\frac{4q\sin2x}{1+q}+\frac{4q^{2}\sin4x}{1+q^{2}}-\frac{4q^{3}\sin6x}{1+q^{3}}+.\,.\\[2ex]
\text{17)}\quad\frac{2K}{\pi}.\frac{\sin\operatorname{coam}\dfrac{2Kx}{\pi}}{\sin\operatorname{am}\dfrac{2Kx}{\pi}}=\operatorname{cotg}x+\frac{4q\sin2x}{1-q}-\frac{4q^{2}\sin4x}{1+q^{2}}+\frac{4q^{3}\sin6x}{1-q^{3}}-.\,.\\[2ex]
\text{18)}\quad\frac{2K}{\pi\sin\operatorname{am}\dfrac{2Kx}{\pi}}=\frac{1}{\sin x}+\frac{4q\sin x}{1-q}+\frac{4q^{3}\sin3x}{1-q^{3}}+\frac{4q^{5}\sin5x}{1-q^{5}}+.\,.\\[2ex]
\text{19)}\quad\frac{2kK}{\pi}.\sin\operatorname{am}\frac{2Kx}{\pi}=\frac{4\sqrt{q}\sin x}{1-q}+\frac{4\sqrt{q^{3}}\sin3x}{1-q^{3}}+\frac{4\sqrt{q^{5}}\sin5x}{1-q^{5}}+.\,.
\end{gathered}
\]
Formula 13) remains unchanged on putting $\dfrac{\pi}{2}-x$ in place of $x$.

On changing $q$ into $-q$, by Theorem I. §. 37 the formulae 11), 12) pass into 17), 16); 13) remains unchanged; from the formulae 14), 15), 18), 19) we obtain:
\[
\begin{gathered}
\text{20)}\quad\frac{2k'K}{\pi\cos\operatorname{am}\dfrac{2Kx}{\pi}}=\frac{1}{\cos x}-\frac{4q\cos x}{1+q}+\frac{4q^{3}\cos3x}{1-q^{3}}-\frac{4q^{5}\cos5x}{1+q^{5}}+.\,.\\[2ex]
\text{21)}\quad\frac{2kK}{\pi}.\cos\operatorname{am}\frac{2Kx}{\pi}=\frac{4\sqrt{q}\cos x}{1+q}+\frac{4\sqrt{q^{3}}\cos3x}{1+q^{3}}+\frac{4\sqrt{q^{5}}\cos5x}{1+q^{5}}+.\,.\\[2ex]
\text{22)}\quad\frac{2k'K}{\pi\cos\operatorname{coam}\dfrac{2Kx}{\pi}}=\frac{1}{\sin x}-\frac{4q\sin x}{1+q}-\frac{4q^{3}\sin3x}{1+q^{3}}-\frac{4q^{5}\sin5x}{1+q^{5}}-.\,.\\[2ex]
\text{23)}\quad\frac{2kK}{\pi}.\cos\operatorname{coam}\frac{2Kx}{\pi}=\frac{4\sqrt{q}\sin x}{1+q}-\frac{4\sqrt{q^{3}}\sin3x}{1+q^{3}}+\frac{4\sqrt{q^{5}}\sin5x}{1+q^{5}}-.\,.
\end{gathered}
\]

\origpage{102}

Formulae 19), 21) can also easily be derived, by known expansions, from those which we cited above in §. 35. 6), 7):
\[
\begin{gathered}
\sin\operatorname{am}\frac{2Kx}{\pi}=\frac{2\pi}{kK}\sin x\left(\frac{\sqrt{q}(1+q)}{1-2q\cos2x+q^{2}}+\frac{\sqrt{q^{3}}(1+q^{3})}{1-2q^{3}\cos2x+q^{6}}+\frac{\sqrt{q^{5}}(1+q^{5})}{1-2q^{5}\cos2x+q^{10}}+.\,.\right)\\[2ex]
\cos\operatorname{am}\frac{2Kx}{\pi}=\frac{2\pi}{kK}\cos x\left(\frac{\sqrt{q}(1-q)}{1-2q\cos2x+q^{2}}-\frac{\sqrt{q^{3}}(1-q^{3})}{1-2q^{3}\cos2x+q^{6}}+\frac{\sqrt{q^{5}}(1-q^{5})}{1-2q^{5}\cos2x+q^{10}}-.\,.\right),
\end{gathered}
\]
From formula 9) §. 35:
\[
\begin{gathered}
\operatorname{am}\frac{2Kx}{\pi}=\\[1.5ex]
\pm x+2\operatorname{Arc}\operatorname{tg}\left(\left(\frac{1+q}{1-q}\right)\operatorname{tg}x\right)-2\operatorname{Arc}\operatorname{tg}\left(\left(\frac{1+q^{3}}{1-q^{3}}\right)\operatorname{tg}x\right)+2\operatorname{Arc}\operatorname{tg}\left(\left(\frac{1+q^{5}}{1-q^{5}}\right)\operatorname{tg}x\right)-.\,.
\end{gathered}
\]
there follows further:
\[
\text{24)}\quad\operatorname{am}\frac{2Kx}{\pi}=x+\frac{2q\sin2x}{1+q^{2}}+\frac{2q^{2}\sin4x}{2(1+q^{4})}+\frac{2q^{3}\sin6x}{3(1+q^{6})}+.\,.
\]
For it is permitted to represent the same, according to the nature of the ambiguous sign, thus:
\[
\begin{gathered}
+x+2\operatorname{Arc}\operatorname{tg}.\frac{(1+q)t}{1-q}-2\operatorname{Arc}\operatorname{tg}.\frac{(1+q^{3})t}{1-q^{3}}+2\operatorname{Arc}\operatorname{tg}.\frac{(1+q^{5})t}{1-q^{5}}-.\,.,\\[1.5ex]
-2x\qquad\qquad\qquad\qquad+2x\qquad\qquad\qquad\qquad-2x\qquad\qquad\qquad\qquad+.\,.,
\end{gathered}
\]
provided that for brevity $t=\operatorname{tg}x$. But there results:
\[
\begin{gathered}
\operatorname{Arc}\operatorname{tg}.\frac{(1+q)t}{1-q}-x=\operatorname{Arc}\operatorname{tg}.\left\{\frac{(1+q)t-(1-q)t}{1-q+(1+q)tt}\right\}=\\[1.5ex]
\operatorname{Arc}\operatorname{tg}.\left\{\frac{2qt}{1+tt-q(1-tt)}\right\}=\operatorname{Arc}\operatorname{tg}.\left\{\frac{q\sin2x}{1-q\cos2x}\right\},
\end{gathered}
\]
whence $\operatorname{am}\dfrac{2Kx}{\pi}=$
\[
x+2\operatorname{Arc}\operatorname{tg}.\frac{q\sin2x}{1-q\cos2x}-2\operatorname{Arc}\operatorname{tg}.\frac{q^{3}\sin2x}{1-q^{3}\cos2x}+2\operatorname{Arc}\operatorname{tg}.\frac{q^{5}\sin2x}{1-q^{5}\cos2x}-.\,.,
\]
or, since
\[
\operatorname{Arc}\operatorname{tg}.\frac{q\sin2x}{1-q\cos2x}=q\sin2x+\frac{q^{2}\sin4x}{2}+\frac{q^{3}\sin6x}{3}+.\,.,
\]
there results $\operatorname{am}\dfrac{2Kx}{\pi}=$
\[
x+\frac{2q\sin2x}{1+q^{2}}+\frac{2q^{2}\sin4x}{2(1+q^{4})}+\frac{2q^{3}\sin6x}{3(1+q^{6})}+.\,.
\]

\origpage{103}
which is formula 24). From the differentiation of this there arises:
\[
\text{25)}\quad\frac{2K}{\pi}.\Delta\operatorname{am}\frac{2Kx}{\pi}=1+\frac{4q\cos2x}{1+q^{2}}+\frac{4q^{2}\cos4x}{1+q^{4}}+\frac{4q^{3}\cos6x}{1+q^{6}}+.\,.,
\]
whence also, on putting $-q$ in place of $q$ or $\dfrac{\pi}{2}-x$ in place of $x$:
\[
\text{26)}\quad\frac{2k'K}{\pi\Delta\operatorname{am}\dfrac{2Kx}{\pi}}=1-\frac{4q\cos2x}{1+q^{2}}+\frac{4q^{2}\cos4x}{1+q^{4}}-\frac{4q^{3}\cos6x}{1+q^{6}}+.\,.
\]

\subsection*{40.}

From the formulae proposed, by putting $x=0$ or by other means, the following are easily derived:
\[
\begin{gathered}
\text{1)}\quad\log k=\log4\sqrt{q}-\frac{4q}{1+q}+\frac{4q^{2}}{2(1+q^{2})}-\frac{4q^{3}}{3(1+q^{3})}+\frac{4q^{4}}{4(1+q^{4})}-.\,.\\[2ex]
\text{2)}\quad-\log k'=\frac{8q}{1-q^{2}}+\frac{8q^{3}}{3(1-q^{6})}+\frac{8q^{5}}{5(1-q^{10})}+\frac{8q^{7}}{7(1-q^{14})}+.\,.\\[2ex]
\text{3)}\quad\log\frac{2K}{\pi}=\frac{4q}{1+q}+\frac{4q^{3}}{3(1+q^{3})}+\frac{4q^{5}}{5(1+q^{5})}+\frac{4q^{7}}{7(1+q^{7})}+.\,.\\[2ex]
\text{4)}\quad\frac{2K}{\pi}=1+\frac{4q}{1-q}-\frac{4q^{3}}{1-q^{3}}+\frac{4q^{5}}{1-q^{5}}+.\,.\\[2ex]
=1+\frac{4q}{1+q^{2}}+\frac{4q^{2}}{1+q^{4}}+\frac{4q^{3}}{1+q^{6}}+.\,.\\[2ex]
\text{5)}\quad\frac{2kK}{\pi}=\frac{4\sqrt{q}}{1-q}-\frac{4\sqrt{q^{3}}}{1-q^{3}}+\frac{4\sqrt{q^{5}}}{1-q^{5}}-.\,.\\[2ex]
=\frac{4\sqrt{q}}{1+q}+\frac{4\sqrt{q^{3}}}{1+q^{3}}+\frac{4\sqrt{q^{5}}}{1+q^{5}}+.\,.\\[2ex]
\text{6)}\quad\frac{2k'K}{\pi}=1-\frac{4q}{1+q}+\frac{4q^{3}}{1+q^{3}}-\frac{4q^{5}}{1+q^{5}}+.\,.\\[2ex]
=1-\frac{4q}{1+q^{2}}+\frac{4q^{2}}{1+q^{4}}-\frac{4q^{3}}{1+q^{6}}+.\,.\\[2ex]
\text{7)}\quad\frac{2\sqrt{k'}K}{\pi}=1-\frac{4q^{2}}{1+q^{2}}+\frac{4q^{6}}{1+q^{6}}-\frac{4q^{10}}{1+q^{10}}+.\,.\\[2ex]
=1-\frac{4q^{2}}{1+q^{4}}+\frac{4q^{4}}{1+q^{8}}-\frac{4q^{6}}{1+q^{12}}+.\,.\\[2ex]
\text{8)}\quad\frac{4KK}{\pi\pi}=1+\frac{8q}{1-q}+\frac{16q^{2}}{1+q^{2}}+\frac{24q^{3}}{1-q^{3}}+.\,.\\[2ex]
=1+\frac{8q}{(1-q)^{2}}+\frac{8q^{2}}{(1+q^{2})^{2}}+\frac{8q^{3}}{(1-q^{3})^{2}}+.\,.
\end{gathered}
\]

\origpage{104}
\[
\begin{gathered}
\text{9)}\quad\frac{4kkKK}{\pi\pi}=\frac{16q}{1-q^{2}}+\frac{48q^{3}}{1-q^{6}}+\frac{80q^{5}}{1-q^{10}}+.\,.\\[2ex]
=\frac{16q(1+q^{2})}{(1-q^{2})^{2}}+\frac{16q^{3}(1+q^{6})}{(1-q^{6})^{2}}+\frac{16q^{5}(1+q^{10})}{(1-q^{10})^{2}}+.\,.\\[2ex]
\text{10)}\quad\frac{4k'k'KK}{\pi\pi}=1-\frac{8q}{1+q}+\frac{16q^{2}}{1+q^{2}}-\frac{24q^{3}}{1+q^{3}}+.\,.\\[2ex]
=1-\frac{8q}{(1+q^{2})}+\frac{8q^{2}}{(1+q^{2})^{2}}-\frac{8q^{3}}{(1+q^{3})^{2}}+.\,.\\[2ex]
\text{11)}\quad\frac{4kk'KK}{\pi\pi}=\frac{4\sqrt{q}}{1+q}-\frac{12\sqrt{q^{3}}}{1+q^{3}}+\frac{20\sqrt{q^{5}}}{1+q^{5}}-.\,.\\[2ex]
=\frac{4\sqrt{q}(1+q)}{(1+q)^{2}}+\frac{4\sqrt{q^{3}}(1+q^{3})}{(1+q^{3})^{2}}+\frac{4\sqrt{q^{5}}(1+q^{5})}{(1+q^{5})^{2}}+.\,.\\[2ex]
\text{12)}\quad\frac{4k'KK}{\pi\pi}=1-\frac{8q^{2}}{1+q^{2}}+\frac{16q^{4}}{1+q^{4}}-\frac{24q^{6}}{1+q^{6}}+.\,.\\[2ex]
=1-\frac{8q^{2}}{(1+q^{2})^{2}}+\frac{8q^{4}}{(1+q^{4})^{2}}-\frac{8q^{6}}{(1+q^{6})^{2}}+.\,.\\[2ex]
\text{13)}\quad\frac{4kKK}{\pi\pi}=\frac{4\sqrt{q}}{1-q}+\frac{12\sqrt{q^{3}}}{1-q^{3}}+\frac{20\sqrt{q^{5}}}{1-q^{5}}-.\,.\\[2ex]
=\frac{4\sqrt{q}(1+q)}{(1-q)^{2}}+\frac{4\sqrt{q^{3}}(1+q^{3})}{(1-q^{3})^{2}}+\frac{4\sqrt{q^{5}}(1+q^{5})}{(1-q^{5})^{2}}+.\,.
\end{gathered}
\]\ednote{In the second line of 10) the first denominator is printed $(1+q^{2})$ instead of $(1+q)^{2}$; an apparent misprint.}
Formulae 4) - 13) we have represented in a twofold manner; but the one representation follows easily from the other, when the single denominators are expanded into series. Let us note further that, according to the theorems proposed in §. 37, all of them can be derived from two of their number, the $4^{\mathrm{ta}}$ and the $8^{\mathrm{va}}$. For on putting $\sqrt{q}$ in place of $q$, since $K$ passes into $(1+k)K$, by subtracting from formula 4) there results 5); then on putting $-q$ in place of $q$, $K$ passes into $k'K$, whence from the formulae 4), 8) there result 6), 10); 5) remains unchanged. On putting $q^{2}$ in place of $q$, $k'K$ passes into $\sqrt{k'}K$, whence from 6), 10) there result 7), 12). From 8), 10), because $kk+k'k'=1$, there results 9). On putting $\sqrt{q}$ in place of $q$, $kK$ passes into $2\sqrt{k}K$, whence from 9) there results 13). On putting $-q$ in place of $q$, $kKK$ passes into $ikk'KK$, whence from 13) there results 11). Moreover, for the Modulus itself or the Complement no series of this kind appear to exist.

The formulae proposed being expanded into powers of $q$, we obtain:
\[
\begin{gathered}
\text{14)}\quad\log k=\log4\sqrt{q}-4q+6q^{2}-\frac{16}{3}q^{3}+3q^{4}-\frac{24}{5}q^{5}+8q^{6}-\frac{32}{7}q^{7}+\frac{3}{2}q^{8}-\frac{52}{9}q^{9}+\frac{36}{5}q^{10}\ldots\\[2ex]
\text{15)}\quad-\log k'=8q+\frac{32}{3}q^{3}+\frac{48}{5}q^{5}+\frac{64}{7}q^{7}+\frac{104}{9}q^{9}+\frac{96}{11}q^{11}+\frac{112}{13}q^{13}+\frac{192}{15}q^{15}.\,.\,.\,.
\end{gathered}
\]

\origpage{105}
\[
\begin{gathered}
\text{16)}\quad\log\frac{2K}{\pi}=4q-4q^{2}+\frac{16}{3}q^{3}-4q^{4}+\frac{24}{5}q^{5}-\frac{16}{3}q^{6}+\frac{32}{7}q^{7}-4q^{8}+\frac{52}{9}q^{9}-\frac{24}{10}q^{10}+\ldots\\[2ex]
\text{17)}\quad\frac{2K}{\pi}=1+4q+4q^{2}+4q^{4}+8q^{5}+4q^{8}+4q^{9}+8q^{10}+8q^{13}+4q^{16}+8q^{17}+4q^{18}+.\,.\\[2ex]
\text{18)}\quad\frac{2kK}{\pi}=4\sqrt{q}+8\sqrt{q^{5}}+4\sqrt{q^{9}}+8\sqrt{q^{13}}+8\sqrt{q^{17}}+12\sqrt{q^{25}}+8\sqrt{q^{29}}+8\sqrt{q^{37}}+.\,.\\[2ex]
\text{19)}\quad\frac{2k'K}{\pi}=1-4q+4q^{2}+4q^{4}-8q^{5}+4q^{8}-4q^{9}+8q^{10}-8q^{13}+4q^{16}-8q^{17}+4q^{18}\ldots\\[2ex]
\text{20)}\quad\frac{2\sqrt{k'}K}{\pi}=1-4q^{2}+4q^{4}+4q^{8}-8q^{10}+4q^{16}-4q^{18}+8q^{20}-8q^{26}+4q^{32}\ldots\\[2ex]
\text{21)}\quad\frac{4KK}{\pi\pi}=1+8q+24q^{2}+32q^{3}+24q^{4}+48q^{5}+96q^{6}+64q^{7}+24q^{8}+\ldots\\[2ex]
\text{22)}\quad\frac{4kkKK}{\pi\pi}=16q+64q^{3}+96q^{5}+128q^{7}+208q^{9}+192q^{11}+224q^{13}+384q^{15}+.\,.\\[2ex]
\text{23)}\quad\frac{4k'k'KK}{\pi\pi}=1-8q+24q^{2}-32q^{3}+24q^{4}-48q^{5}+96q^{6}-64q^{7}+24q^{8}\ldots\\[2ex]
\text{24)}\quad\frac{4kk'KK}{\pi\pi}=4\sqrt{q}-16\sqrt{q^{3}}+24\sqrt{q^{5}}-32\sqrt{q^{7}}+52\sqrt{q^{9}}-48\sqrt{q^{11}}+56\sqrt{q^{13}}-\ldots\\[2ex]
\text{25)}\quad\frac{4k'KK}{\pi\pi}=1-8q^{2}+24q^{4}-32q^{6}+24q^{8}-48q^{10}+96q^{12}-64q^{14}+24q^{16}-104q^{18}+\ldots\\[2ex]
\text{26)}\quad\frac{4kKK}{\pi\pi}=4\sqrt{q}+16\sqrt{q^{3}}+24\sqrt{q^{5}}+32\sqrt{q^{7}}+52\sqrt{q^{9}}+48\sqrt{q^{11}}+56\sqrt{q^{13}}+\ldots
\end{gathered}
\]\ednote{In 16) the last coefficient is printed $\frac{24}{10}$; by formula 29) it should be $\frac{24}{5}$, an apparent misprint (the author's Corrigenda list this). A hand correction in ink is not transcribed.}

That the law and principle of these series may be better perceived, we shall denote them by the summatory sign $\Sigma$ prefixed to their general term. Let us suppose that $p$ is an \emph{odd number}, and $\phi(p)$ the \emph{sum of the factors of} $p$. Then there results:
\[
\begin{gathered}
\text{27)}\quad\log k=\log4\sqrt{q}-4\Sigma\frac{\phi(p)}{p}\left\{q^{p}-\frac{3q^{2p}}{2}-\frac{3}{4}q^{4p}-\frac{3}{8}q^{8p}-\frac{3}{16}q^{16p}-.\,.\right\}\\[2ex]
\text{28)}\quad-\log k'=8\Sigma\frac{\phi(p)}{p}q^{p}\\[2ex]
\text{29)}\quad\log\frac{2K}{\pi}=4\Sigma\frac{\phi(p)}{p}\left\{q^{p}-q^{2p}-q^{4p}-q^{8p}-q^{16p}-.\,.\right\}.
\end{gathered}
\]
Further let $n$ be an \emph{odd number, all of whose prime factors have the form} $4a+1$, $\psi(n)$ the \emph{number of factors} of $n$; $l$, $m$ all numbers from $0$ up to $\infty$: we obtain:
\[
\text{30)}\quad\frac{2K}{\pi}=1+4\Sigma\psi(n)q^{2^{l}(4m-1)^{2}n}
\]

\origpage{106}
\[
\begin{gathered}
\text{31)}\quad\frac{2kK}{\pi}=4\Sigma\psi(n)q^{\frac{(4m-1)^{2}n}{2}}\\[2ex]
\text{32)}\quad\frac{2k'K}{\pi}=1-4\Sigma\psi(n)q^{(4m-1)^{2}n}+4\Sigma\psi(n)q^{2^{l+1}(4m-1)^{2}n}\\[2ex]
\text{33)}\quad\frac{2\sqrt{k'}K}{\pi}=1-4\Sigma\psi(n)q^{2(4m-1)^{2}n}+4\Sigma\psi(n)q^{2^{l+2}(4m-1)^{2}n}.
\end{gathered}
\]
$p$ denoting again an odd number, $\phi(p)$ the sum of the factors of $p$: there results
\[
\begin{gathered}
\text{34)}\quad\frac{4KK}{\pi\pi}=1+8\Sigma\phi(p)\left\{q^{p}+3q^{2p}+3q^{4p}+3q^{8p}+3q^{16p}+.\,.\right\}\\[2ex]
\text{35)}\quad\frac{4kkKK}{\pi\pi}=16\Sigma\phi(p)q^{p}\\[2ex]
\text{36)}\quad\frac{4k'k'KK}{\pi\pi}=1+8\Sigma\phi(p)\left\{-q^{p}+3q^{2p}+3q^{4p}+3q^{8p}+3q^{16p}+.\,.\right\}\\[2ex]
\text{37)}\quad\frac{4kk'KK}{\pi\pi}=4\Sigma(-1)^{\frac{p-1}{2}}\phi(p)\sqrt{q^{p}}\\[2ex]
\text{38)}\quad\frac{4k'KK}{\pi\pi}=1+8\Sigma\phi(p)\left\{-q^{2p}+3q^{4p}+3q^{8p}+3q^{16p}+3q^{32p}+.\,.\right\}\\[2ex]
\text{39)}\quad\frac{4kKK}{\pi\pi}=4\Sigma\phi(p)\sqrt{q^{p}}.
\end{gathered}
\]
Let us demonstrate formula 27). We have found 1):
\[
\log k=\log4\sqrt{q}-\frac{4q}{1+q}+\frac{4q^{2}}{2(1+q^{2})}-\frac{4q^{3}}{3(1+q^{3})}+.\,.,
\]
which let us put $=\log4\sqrt{q}+4\Sigma A^{(x)}q^{x}$. Let $x$ be an odd number $p=mm'$; from any term $-\dfrac{q^{m}}{m(1+q^{m})}$ there results $\dfrac{-q^{p}}{m}$, whence it is evident that $A^{(p)}=-\dfrac{\phi(p)}{p}$. Now let $x$ be an even number $=2^{l}p=2^{l}mm'$: from the terms
\[
\frac{-q^{m}}{m(1+q^{m})}+\frac{q^{2m}}{2m(1+q^{2m})}+\frac{q^{4m}}{4m(1+q^{4m})}+\frac{q^{8m}}{8m(1+q^{8m})}+.\,.+\frac{q^{2^{l}m}}{2^{l}m(1+q^{2^{l}m})}
\]
there results
\[
\frac{q^{x}}{m}\left\{1-\frac{1}{2}-\frac{1}{4}-\frac{1}{8}\ldots-\frac{1}{2^{l-1}}+\frac{1}{2^{l}}\right\}=\frac{3q^{x}}{2^{l}m},
\]
whence $A^{(x)}=\dfrac{3\phi(p)}{2^{l}p}$, which furnishes the formula proposed.

\origpage{107}
Let us demonstrate formula 30). We have found 4):
\[
\frac{2K}{\pi}=1+\frac{4q}{1-q}-\frac{4q^{3}}{1-q^{3}}+\frac{4q^{5}}{1-q^{5}}-.\,.=1+4\Sigma A^{(x)}q^{x}.
\]
Let $B^{(x)}$ be the number of factors of $x$ which have the form $4m+1$, $C^{(x)}$ the number of factors which have the form $4m+3$; it is readily evident that $A^{(x)}=B^{(x)}-C^{(x)}$. Let $x=2^{l}nn'$, so that $n$ is an odd number all of whose prime factors have the form $4m+1$, and $n'$ an odd number all of whose prime factors have the form $4m-1$; it is easily proved that, unless $n'$ is a square number, there will always be $B^{(x)}-C^{(x)}=0$, but where $n'$ is a square number, there will be $B^{(x)}-C^{(x)}=B^{(n)}=\psi(n)$, formula 30) flows.\ednote{The author's Corrigenda ask for ``unde formula 30)''; the word unde is added by hand in ink above the line and is not transcribed.}

Lastly let us prove formula 34). We have found 8):
\[
\frac{4KK}{\pi\pi}=1+\frac{8q}{1-q}+\frac{16q^{2}}{1+q^{2}}+\frac{24q^{3}}{1-q^{3}}+\frac{32q^{4}}{1+q^{4}}+.\,.=1+8\Sigma A^{(x)}q^{x}.
\]
$x$ denoting an odd number, it is readily evident that $A^{(x)}=\phi(x)$; but where $x$ is an even number $=2^{l}p$, $p$ denoting an odd number, for every $m$ a factor of $p$, from the terms
\[
8\left\{\frac{mq^{m}}{1-q^{m}}+\frac{2mq^{2m}}{1+q^{2m}}+\frac{4mq^{4m}}{1+q^{4m}}+\frac{8mq^{8m}}{1+q^{8m}}+.\,.+\frac{2^{l}mq^{2^{l}m}}{1+q^{2^{l}m}}\right\}
\]
there results $8mq^{x}\left\{1-2-4-8-.\,.-2^{l-1}+2^{l}\right\}=24mq^{x}$, whence in that case $A^{(x)}=3\phi(p)$, which suggests the formula proposed. The remaining ones are demonstrated similarly or can be deduced from these.

The expressions $\cos\operatorname{am}\dfrac{2Kx}{\pi}$, $\Delta\operatorname{am}\dfrac{2Kx}{\pi}$, $\dfrac{1}{\cos\operatorname{am}\dfrac{2Kx}{\pi}}$ being expanded into powers of $x$, we obtain the Coefficient of $x^{2}$\ednote{The author's Corrigenda ask for ``evolutae'' and ``nanciscuntur'' here, in place of the printed evolutas and nanciscimur; a hand correction in ink is not transcribed.} respectively $-\dfrac{1}{2}\left(\dfrac{2K}{\pi}\right)^{2}$, $-\dfrac{1}{2}\left(\dfrac{2kK}{\pi}\right)^{2}$, $+\dfrac{1}{2}\left(\dfrac{2K}{\pi}\right)^{2}$, whence from the formulae 21), 20), 24) of the preceding §$^{\mathrm{i}}$ we see that the following are derived:
\[
\begin{gathered}
\text{40)}\quad k\left(\frac{2K}{\pi}\right)^{3}=4\left\{\frac{\sqrt{q}}{1+q}+\frac{9\sqrt{q^{3}}}{1+q^{3}}+\frac{25\sqrt{q^{5}}}{1+q^{5}}+\frac{49\sqrt{q^{7}}}{1+q^{7}}+.\,.\right\}=\\[2ex]
4\left\{\frac{\sqrt{q}(1+6q+q^{2})}{(1-q)^{3}}-\frac{\sqrt{q^{3}}(1+6q^{3}+q^{6})}{(1-q^{3})^{3}}+\frac{\sqrt{q^{5}}(1+6q^{5}+q^{10})}{(1-q^{5})^{3}}-.\,.\right\}\\[2ex]
\text{41)}\quad k'\left(\frac{2K}{\pi}\right)^{3}=1+4\left\{\frac{q}{1+q}-\frac{9q^{3}}{1+q^{3}}+\frac{25q^{5}}{1+q^{5}}-\frac{49q^{7}}{1+q^{7}}+.\,.\right\}=\\[2ex]
1+4\left\{\frac{q(1-6q^{2}+q^{4})}{(1+q^{2})^{3}}-\frac{q^{2}(1-6q^{4}+q^{8})}{(1+q^{4})^{3}}+\frac{q^{3}(1-6q^{6}+q^{12})}{(1+q^{6})^{3}}-.\,.\right\}
\end{gathered}
\]

\origpage{108}
\[
\begin{gathered}
\text{42)}\quad kk\left(\frac{2K}{\pi}\right)^{5}=16\left\{\frac{q}{1+q^{2}}+\frac{4q^{2}}{1+q^{4}}+\frac{9q^{3}}{1+q^{6}}+\frac{16q^{4}}{1+q^{8}}+.\,.\right\}=\\[2ex]
16\left\{\frac{q(1+q)}{(1-q)^{3}}-\frac{q^{3}(1+q^{3})}{(1-q^{3})^{3}}+\frac{q^{5}(1+q^{5})}{(1-q^{5})^{3}}+.\,.\right\}.
\end{gathered}
\]\ednote{In 42) the exponent of $\frac{2K}{\pi}$ is printed $5$ instead of $3$; an apparent misprint (the author's Corrigenda list this), with a hand correction in ink that is not transcribed.}
From these, putting $-q$ in place of $q$, we obtain:
\[
\begin{gathered}
\text{43)}\quad kk'k'\left(\frac{2K}{\pi}\right)^{3}=4\left\{\frac{\sqrt{q}}{1-q}-\frac{9\sqrt{q^{3}}}{1-q^{3}}+\frac{25\sqrt{q^{5}}}{1-q^{5}}-\frac{49\sqrt{q^{7}}}{1-q^{7}}+.\,.\right\}\\[2ex]
\text{44)}\quad k'k'\left(\frac{2K}{\pi}\right)^{3}=1-4\left\{\frac{q}{1-q}-\frac{9q^{3}}{1-q^{3}}+\frac{25q^{5}}{1-q^{5}}-\frac{49q^{7}}{1-q^{7}}+.\,.\right\}\\[2ex]
\text{45)}\quad k'kk\left(\frac{2K}{\pi}\right)^{3}=16\left\{\frac{q}{1+q^{2}}-\frac{4q^{2}}{1+q^{4}}+\frac{9q^{3}}{1+q^{6}}-\frac{16q^{4}}{1+q^{8}}+.\,.\right\}.
\end{gathered}
\]
By adding the formulae 42), 44), we obtain $\left(\dfrac{2K}{\pi}\right)^{3}$; by subtracting 40) and 43), 41) and 45), we obtain $\left(\dfrac{2kK}{\pi}\right)^{3}$, $\left(\dfrac{2k'K}{\pi}\right)^{3}$, from which, on putting resp. $\sqrt{q}$, $q^{2}$ in place of $q$, there results $\left(\dfrac{4\sqrt{k}K}{\pi}\right)^{3}$, $\left(\dfrac{2\sqrt{k'}K}{\pi}\right)^{3}$; from $\left(\dfrac{4\sqrt{k}K}{\pi}\right)^{3}$, on putting $-q$ in place of $q$, is obtained $\left(\dfrac{4\sqrt{kk'}K}{\pi}\right)^{3}$.

Finally, putting $k=\sin\vartheta$, let us expand $\vartheta=\operatorname{Arc.}\sin k$ itself. We have seen that on putting $\sqrt{q}$ in place of $q$, $k'$ passes into $\dfrac{1-k}{1+k}$; let us again put $-q$ in place of $q$, then $k$ passes into $\dfrac{ik}{k'}$, or into $i.\operatorname{tang}\vartheta$; so that on putting $i\sqrt{q}$ in place of $q$, the expression $\dfrac{-\log k'}{2i}$ is changed into
\[
-\frac{1}{2i}\log\left(\frac{1-i\operatorname{tang}\vartheta}{1+i\operatorname{tang}\vartheta}\right)=\vartheta.
\]
Hence from formula 2)
\[
-\log k'=\frac{8q}{1-q^{2}}+\frac{8q^{3}}{3(1-q^{6})}+\frac{8q^{5}}{5(1-q^{10})}+\frac{8q^{7}}{7(1-q^{14})}+.\,.
\]
we derive:
\[
\text{46)}\quad\vartheta=\operatorname{Arc}\sin k=\frac{4\sqrt{q}}{1+q}-\frac{4\sqrt{q^{3}}}{3(1+q^{3})}+\frac{4\sqrt{q^{5}}}{5(1+q^{5})}-\frac{4\sqrt{q^{7}}}{7(1+q^{7})}+.\,.,
\]
which is easily transformed into this:
\[
\text{47)}\quad\frac{\vartheta}{4}=\operatorname{Arc}\operatorname{tg}\sqrt{q}-\operatorname{Arc}\operatorname{tg}\sqrt{q^{3}}+\operatorname{Arc}\operatorname{tg}\sqrt{q^{5}}-\operatorname{Arc}\operatorname{tg}\sqrt{q^{7}}+\ldots,
\]
which must be counted among the most elegant formulae.

\origpage{109}
\subsection*{41.}

Let us multiply the equation exhibited above:
\[
\frac{2kK}{\pi}\sin\operatorname{am}\frac{2Kx}{\pi}=\frac{4\sqrt{q}\sin x}{1-q}+\frac{4\sqrt{q^{3}}\sin 3x}{1-q^{3}}+\frac{4\sqrt{q^{5}}\sin 5x}{1-q^{5}}+\ldots
\]
into itself. On substituting everywhere $\cos(m-n)x-\cos(m+n)x$ in place of $2\sin mx\sin nx$, the product assumes the form:
\[
\left(\frac{2kK}{\pi}\right)^{2}\sin^{2}\operatorname{am}\frac{2Kx}{\pi}=A+A'\cos 2x+A''\cos 4x+A'''\cos 6x+.\,.\,.\,.
\]
One finds:
\[
A=\frac{8q}{(1-q)^{2}}+\frac{8q^{3}}{(1-q^{3})^{2}}+\frac{8q^{5}}{(1-q^{5})^{2}}+.\,.
\]
Further there results:
\[
A^{(n)}=16B^{(n)}-8C^{(n)}=8\left\{2B^{(n)}-C^{(n)}\right\},
\]
if one puts:
\[
B^{(n)}=\frac{q^{n+1}}{(1-q)(1-q^{2n+1})}+\frac{q^{n+3}}{(1-q^{3})(1-q^{2n+3})}+\frac{q^{n+5}}{(1-q^{5})(1-q^{2n+5})}+\text{etc.\ to inf.}
\]
\[
C^{(n)}=\frac{q^{n}}{(1-q)(1-q^{2n-1})}+\frac{q^{n}}{(1-q^{3})(1-q^{2n-3})}+\frac{q^{n}}{(1-q^{5})(1-q^{2n-5})}+.\,.+\frac{q^{n}}{(1-q^{2n-1})(1-q)}.
\]
Now since
\[
\frac{q^{n+m}}{(1-q^{m})(1-q^{2n+m})}=\frac{q^{n}}{1-q^{2n}}\left\{\frac{q^{m}}{1-q^{m}}-\frac{q^{2n+m}}{1-q^{2n+m}}\right\},
\]
there results $B^{(n)}=$
\[
\begin{gathered}
\frac{q^{n}}{1-q^{2n}}\left\{\frac{q}{1-q}+\frac{q^{3}}{1-q^{3}}+\frac{q^{5}}{1-q^{5}}+.\,.\right\}\\[1.5ex]
-\frac{q^{n}}{1-q^{2n}}\left\{\frac{q^{2n+1}}{1-q^{2n+1}}+\frac{q^{2n+3}}{1-q^{2n+3}}+\frac{q^{2n+5}}{1-q^{2n+5}}+.\,.\right\},
\end{gathered}
\]
or, the terms which destroy one another being removed:
\[
B^{(n)}=\frac{q^{n}}{1-q^{2n}}\left\{\frac{q}{1-q}+\frac{q^{3}}{1-q^{3}}+.\,.+\frac{q^{2n-1}}{1-q^{2n-1}}\right\}.
\]

\origpage{110}
Further there results:
\[
\frac{q^{n}}{(1-q^{m})(1-q^{2n-m})}=\frac{q^{n}}{1-q^{2n}}\left\{\frac{q^{m}}{1-q^{m}}+\frac{q^{2n-m}}{1-q^{2n-m}}+1\right\},
\]
whence
\[
C^{(n)}=\frac{nq^{n}}{1-q^{2n}}+\frac{2q^{n}}{1-q^{2n}}\left\{\frac{q}{1-q}+\frac{q^{3}}{1-q^{3}}+.\,.+\frac{q^{2n-1}}{1-q^{2n-1}}\right\}.
\]
Hence at length there results:
\[
A^{(n)}=8\left\{2B^{(n)}-C^{(n)}\right\}=\frac{-8nq^{n}}{1-q^{2n}},
\]
whence now:
\[
\text{1)}\quad\left(\frac{2kK}{\pi}\right)^{2}\sin^{2}\operatorname{am}\frac{2Kx}{\pi}=A-8\left\{\frac{q\cos 2x}{1-q^{2}}+\frac{2q^{2}\cos 4x}{1-q^{4}}+\frac{3q^{3}\cos 6x}{1-q^{6}}+.\,.\right\}.
\]
In a similar manner, or from 1), one finds:
\[
\text{2)}\quad\left(\frac{2kK}{\pi}\right)^{2}\cos^{2}\operatorname{am}\frac{2Kx}{\pi}=B+8\left\{\frac{q\cos 2x}{1-q^{2}}+\frac{2q^{2}\cos 4x}{1-q^{4}}+\frac{3q^{3}\cos 6x}{1-q^{6}}+.\,.\right\},
\]
where:
\[
\begin{gathered}
A=8\left\{\frac{q}{(1-q)^{2}}+\frac{q^{3}}{(1-q^{3})^{2}}+\frac{q^{5}}{(1-q^{5})^{2}}+.\,.\right\}\\[1.5ex]
B=8\left\{\frac{q}{(1+q)^{2}}+\frac{q^{3}}{(1+q^{3})^{2}}+\frac{q^{5}}{(1+q^{5})^{2}}+.\,.\right\}.
\end{gathered}
\]
By a well-known theorem of the Integral Calculus, whenever
\[
\phi x=A+A'\cos 2x+A''\cos 4x+A'''\cos 6x+.\,.,
\]
the first or constant term is:
\[
A=\frac{2}{\pi}\int_{0}^{\frac{\pi}{2}}\phi(x).\,dx,
\]
whence in the present case we obtain:
\[
\begin{gathered}
A=\frac{2}{\pi}.\left(\frac{2kK}{\pi}\right)^{2}\int_{0}^{\frac{\pi}{2}}\sin^{2}\operatorname{am}\frac{2Kx}{\pi}.\,dx.\\[1.5ex]
B=\frac{2}{\pi}\left(\frac{2kK}{\pi}\right)^{2}\int_{0}^{\frac{\pi}{2}}\cos^{2}\operatorname{am}\frac{2Kx}{\pi}.\,dx.
\end{gathered}
\]

\origpage{111}
Let us put, with Cl.\ Legendre
\[
E^{\mathrm{I}}=\int_{0}^{\frac{\pi}{2}}d\phi\,\Delta(\phi)=\frac{2K}{\pi}\int_{0}^{\frac{\pi}{2}}dx.\Delta^{2}\operatorname{am}\frac{2Kx}{\pi},
\]
then there will be:
\[
\begin{gathered}
A=\frac{2K}{\pi}.\,\frac{2K}{\pi}-\frac{2K}{\pi}.\,\frac{2E^{\mathrm{I}}}{\pi}\\[1.5ex]
B=\frac{2K}{\pi}.\,\frac{2E^{\mathrm{I}}}{\pi}-\left(\frac{2k'K}{\pi}\right)^{2}.
\end{gathered}
\]
Hence also, since on changing $q$ into $-q$, $A$ passes into $-B$, $K$ into $k'K$, it follows at the same time that $E^{\mathrm{I}}$ passes into $\dfrac{E^{\mathrm{I}}}{k'}$.

Let us remark besides that from formula 1) there follows:
\[
\begin{gathered}
\text{3)}\quad kk\left(\frac{2K}{\pi}\right)^{4}=16\left\{\frac{q}{1-q^{2}}+\frac{2^{3}q^{2}}{1-q^{4}}+\frac{3^{3}q^{3}}{1-q^{6}}+\frac{4^{3}q^{4}}{1-q^{8}}+\ldots\right\}\\[1.5ex]
=16\left\{\frac{q(1+4q+q^{2})}{(1-q)^{4}}+\frac{q^{3}(1+4q^{3}+q^{6})}{(1-q^{3})^{4}}+\frac{q^{5}(1+4q^{5}+q^{10})}{(1-q^{5})^{4}}+.\,.\right\},
\end{gathered}
\]
whence also, on changing $q$ into $-q$:
\[
\begin{gathered}
\text{4)}\quad k^{2}k'^{2}\left(\frac{2K}{\pi}\right)^{4}=16\left\{\frac{q}{1-q^{2}}-\frac{2^{3}q^{2}}{1-q^{4}}+\frac{3^{3}q^{3}}{1-q^{6}}-\frac{4^{3}q^{4}}{1-q^{8}}+\ldots\right\}\\[1.5ex]
=16\left\{\frac{q(1-4q+q^{2})}{(1+q)^{4}}+\frac{q^{3}(1-4q^{3}+q^{6})}{(1+q^{3})^{4}}+\frac{q^{5}(1-4q^{5}+q^{10})}{(1-q^{5})^{4}}+.\,.\right\}.
\end{gathered}
\]\ednote{The third denominator is printed $(1-q^{5})^{4}$ where $(1+q^{5})^{4}$ is expected; an apparent misprint.}
Formula 4) being subtracted from 3), there results:
\[
\begin{gathered}
\text{5)}\quad\left(\frac{2kK}{\pi}\right)^{4}=256\left\{\frac{q^{2}}{1-q^{4}}+\frac{2^{3}q^{4}}{1-q^{8}}+\frac{3^{3}q^{6}}{1-q^{12}}+\frac{4^{3}q^{8}}{1-q^{16}}+.\,.\right\}\\[1.5ex]
=256\left\{\frac{q^{2}(1+4q^{2}+q^{4})}{(1-q^{2})^{4}}+\frac{q^{6}(1+4q^{6}+q^{12})}{(1-q^{6})^{4}}+\frac{q^{10}(1+4q^{10}+q^{20})}{(1-q^{10})^{4}}+.\,.\right\},
\end{gathered}
\]
which\ednote{Printed ``quem''; the author's Corrigenda give ``quam'' here.} you also obtain from 3), on changing $q$ into $q^{2}$.

\origpage{112}
\subsection*{42.}

By a method similar to that by which formula 1) was found, we could inquire into the expansion in series of the expression
\[
\frac{\left(\dfrac{2K}{\pi}\right)^{2}}{\sin^{2}\operatorname{am}\dfrac{2Kx}{\pi}}
\]
if formula 18) §.\ 39 were multiplied into itself. This, however, is accomplished more easily from 1) itself by the following consideration.

For if the formula:
\[
\frac{d.\log\sin\operatorname{am}\dfrac{2Kx}{\pi}}{dx}=\frac{2K}{\pi}.\,\frac{\sqrt{1-(1+kk)\sin^{2}\operatorname{am}\dfrac{2Kx}{\pi}+kk\sin^{4}\operatorname{am}\dfrac{2Kx}{\pi}}}{\sin\operatorname{am}\dfrac{2Kx}{\pi}}
\]
is differentiated again, we obtain, after making the reductions:
\[
\text{1)}\quad\frac{d^{2}.\log\sin\operatorname{am}\dfrac{2Kx}{\pi}}{dx^{2}}=\left(\frac{2K}{\pi}\right)^{2}\left\{kk\sin^{2}\operatorname{am}\frac{2Kx}{\pi}-\frac{1}{\sin^{2}\operatorname{am}\dfrac{2Kx}{\pi}}\right\}.
\]
Now, however, we found in §.\ 39, 6):
\[
\log\sin\operatorname{am}\frac{2Kx}{\pi}=\log\left(\frac{2\sqrt[4]{q}}{\sqrt{k}}\right)+\log\sin x+2\left\{\frac{q\cos 2x}{1+q}+\frac{q^{2}\cos 4x}{2(1+q^{2})}+\frac{q^{3}\cos 6x}{3(1+q^{3})}+.\,.\right\},
\]
whence:
\[
\frac{d^{2}\log\sin\operatorname{am}\dfrac{2Kx}{\pi}}{dx^{2}}=-\frac{1}{\sin^{2}x}-8\left\{\frac{q\cos 2x}{1+q}+\frac{2q^{2}\cos 4x}{1+q^{2}}+\frac{3q^{3}\cos 6x}{1+q^{3}}+.\,.\right\}.
\]
Further we have, §.\ 41, 1):
\[
\begin{gathered}
\left(\frac{2kK}{\pi}\right)^{2}\sin^{2}\operatorname{am}\frac{2Kx}{\pi}=\\[1.5ex]
\frac{2K}{\pi}.\,\frac{2K}{\pi}-\frac{2K}{\pi}.\,\frac{2E^{\mathrm{I}}}{\pi}-8\left\{\frac{q\cos 2x}{1-q^{2}}+\frac{2q^{2}\cos 4x}{1-q^{4}}+\frac{3q^{3}\cos 6x}{1-q^{6}}+.\,.\right\},
\end{gathered}
\]

\origpage{113}
whence, since from formula 1) there is:
\[
\frac{\left(\dfrac{2K}{\pi}\right)^{2}}{\sin^{2}\operatorname{am}\dfrac{2Kx}{\pi}}=\left(\frac{2kK}{\pi}\right)^{2}\sin^{2}\operatorname{am}\frac{2Kx}{\pi}-\frac{d^{2}\log\sin\operatorname{am}\dfrac{2Kx}{\pi}}{dx^{2}},
\]
there results what we seek:
\[
\begin{gathered}
\text{2)}\quad\frac{\left(\dfrac{2K}{\pi}\right)^{2}}{\sin^{2}\operatorname{am}\dfrac{2Kx}{\pi}}=\\[1.5ex]
\frac{2K}{\pi}.\,\frac{2K}{\pi}-\frac{2K}{\pi}.\,\frac{2E^{\mathrm{I}}}{\pi}+\frac{1}{\sin^{2}x}-8\left\{\frac{q^{2}\cos 2x}{1-q^{2}}+\frac{2q^{4}\cos 4x}{1-q^{4}}+\frac{3q^{6}\cos 6x}{1-q^{6}}+.\,.\right\}.
\end{gathered}
\]\ednote{The numerators are printed $q^{2}\cos 2x$, $2q^{4}\cos 4x$, $3q^{6}\cos 6x$ where $q\cos 2x$, $2q^{2}\cos 4x$, $3q^{3}\cos 6x$ are expected (compare formula 1 of §.\ 41); an apparent misprint, repeated in formula 3).}

If $q$ is changed into $-q$ and at the same time $x$ into $\frac{\pi}{2}-x$, whence $K$ passes into $k'K$, $E^{\mathrm{I}}$ into $\dfrac{E^{\mathrm{I}}}{k'}$ §.\ 41, and $\sin\operatorname{am}\dfrac{2Kx}{\pi}$ into $\cos\operatorname{am}\dfrac{2Kx}{\pi}$, there results from 2):
\[
\begin{gathered}
\text{3)}\quad\frac{\left(\dfrac{2k'K}{\pi}\right)^{2}}{\cos^{2}\operatorname{am}\dfrac{2Kx}{\pi}}=\\[1.5ex]
\left(\frac{2k'K}{\pi}\right)^{2}-\frac{2K}{\pi}.\,\frac{2E^{\mathrm{I}}}{\pi}+\frac{1}{\cos^{2}x}+8\left\{\frac{q^{2}\cos 2x}{1-q^{2}}-\frac{2q^{4}\cos 4x}{1-q^{4}}+\frac{3q^{6}\cos 6x}{1-q^{6}}-.\,.\right\}.
\end{gathered}
\]
To these I add the following, which follow easily from §.\ 41.\ 1):
\[
\begin{gathered}
\text{4)}\quad\left(\frac{2K}{\pi}\right)^{2}\Delta^{2}\operatorname{am}\frac{2Kx}{\pi}=\frac{2K}{\pi}.\,\frac{2E^{\mathrm{I}}}{\pi}+8\left\{\frac{q\cos 2x}{1-q^{2}}+\frac{2q^{2}\cos 4x}{1-q^{4}}+\frac{3q^{3}\cos 6x}{1-q^{6}}+.\,.\right\}\\[1.5ex]
\text{5)}\quad\frac{\left(\dfrac{2k'K}{\pi}\right)^{2}}{\Delta^{2}\operatorname{am}\dfrac{2Kx}{\pi}}=\frac{2K}{\pi}.\,\frac{2E^{\mathrm{I}}}{\pi}-8\left\{\frac{q\cos 2x}{1-q^{2}}-\frac{2q^{2}\cos 4x}{1-q^{4}}+\frac{3q^{3}\cos 6x}{1-q^{6}}-.\,.\right\},
\end{gathered}
\]
of which 5) follows from 4), on changing $x$ into $\frac{\pi}{2}-x$ or $q$ into $+q$.

Putting $y=\sin\operatorname{am}\dfrac{2Kx}{\pi}$, $\sqrt{(1-yy)(1-k^{2}yy)}=R$, there results:
\[
\begin{gathered}
\frac{dy}{dx}=\frac{2K}{\pi}.R\\[1.5ex]
\frac{d^{2}y}{dx^{2}}=-\left(\frac{2K}{\pi}\right)^{2}y(1+kk-2k^{2}y^{2})
\end{gathered}
\]

\origpage{114}
\[
\begin{gathered}
\frac{d^{3}y}{dx^{3}}=-\left(\frac{2K}{\pi}\right)^{3}(1+kk-6k^{2}y^{2})R\\[1.5ex]
\frac{d^{4}y}{dx^{4}}=\left(\frac{2K}{\pi}\right)^{4}y\left(1+14kk+k^{4}-20k^{2}(1+k^{2})y^{2}+24k^{4}y^{4}\right)R\\[1.5ex]
\frac{d^{5}y}{dx^{5}}=\left(\frac{2K}{\pi}\right)^{5}\left(1+14kk+k^{4}-60k^{2}(1+k^{2})y^{2}+120k^{4}y^{4}\right)R\\[1.5ex]
\text{etc.}\qquad\qquad\text{etc.},
\end{gathered}
\]
whence:
\[
y=\sin\operatorname{am}\frac{2Kx}{\pi}=\frac{2Kx}{\pi}-\frac{1+k^{2}}{2.3}\left(\frac{2Kx}{\pi}\right)^{3}+\frac{1+14k^{2}+k^{4}}{2.3.4.5}\left(\frac{2Kx}{\pi}\right)^{5}-.\,.,
\]
and therefore:
\[
\frac{\left(\dfrac{2K}{\pi}\right)^{2}}{\sin^{2}\operatorname{am}\dfrac{2Kx}{\pi}}=\frac{1}{xx}+\left(\frac{1+kk}{3}\right)\left(\frac{2K}{\pi}\right)^{2}+\frac{1-kk+k^{4}}{15}\left(\frac{2K}{\pi}\right)^{4}xx+\ldots,
\]
and this formula being compared with 2), we derive:
\[
\left(\frac{1+kk}{3}\right)\left(\frac{2K}{\pi}\right)^{2}=\frac{1}{3}+\left(\frac{2K}{\pi}\right)^{2}-\frac{2K}{\pi}.\,\frac{2E^{\mathrm{I}}}{\pi}-8\left\{\frac{q^{2}}{1-q^{2}}+\frac{2q^{4}}{1-q^{4}}+\frac{3q^{6}}{1-q^{6}}+.\,.\right\},
\]
or:
\[
\text{6)}\quad\frac{q^{2}}{1-q^{2}}+\frac{2q^{4}}{1-q^{4}}+\frac{3q^{6}}{1-q^{6}}+\frac{4q^{8}}{1-q^{8}}+.\,.=\frac{1+\left(\dfrac{2K}{\pi}\right)^{2}(2-kk)-\dfrac{2K}{\pi}.\,\dfrac{2E^{\mathrm{I}}}{\pi}}{2.3.4}.
\]\ednote{The numerator is printed without the factor 3 before $\frac{2K}{\pi}.\,\frac{2E^{\mathrm{I}}}{\pi}$; the author's Corrigenda (p. 114, line 12) ask for $-3\,.\frac{2K}{\pi}.\,\frac{2E^{\mathrm{I}}}{\pi}$, and the scan carries that 3 as a hand-inked correction.}
Further there results:
\[
\left(\frac{1-k^{2}+k^{4}}{15}\right)\left(\frac{2K}{\pi}\right)^{4}=\frac{1}{15}+16\left\{\frac{q^{2}}{1-q^{2}}+\frac{2^{3}q^{4}}{1-q^{4}}+\frac{3^{3}q^{6}}{1-q^{6}}+\frac{4^{3}q^{8}}{1-q^{8}}+.\,.\right\},
\]
or, since $15=2.2^{3}-1$:
\[
\begin{gathered}
(1-k^{2}+k^{4})\left(\frac{2K}{\pi}\right)^{4}=1+2.16\left\{\frac{2^{3}q^{2}}{1-q^{2}}+\frac{4^{3}q^{4}}{1-q^{4}}+\frac{6^{3}q^{6}}{1-q^{6}}+\frac{8^{3}q^{8}}{1-q^{8}}+.\,.\right\}\\[1.5ex]
-16\left\{\frac{q^{2}}{1-q^{2}}+\frac{2^{3}q^{4}}{1-q^{4}}+\frac{3^{3}q^{6}}{1-q^{6}}+\frac{4^{3}q^{8}}{1-q^{8}}+.\,.\right\}.
\end{gathered}
\]
From this formula let the following, §.\ 41.\ 3), be subtracted:
\[
k^{2}\left(\frac{2K}{\pi}\right)^{4}=16\left\{\frac{q}{1-q^{2}}+\frac{2^{3}q^{2}}{1-q^{4}}+\frac{3^{3}q^{3}}{1-q^{6}}+\frac{4^{3}q^{4}}{1-q^{8}}+.\,.\right\},
\]

\origpage{115}
the remainder is:
\[
\text{7)}\quad\left(\frac{2k'K}{\pi}\right)^{4}=1-16\left\{\frac{q}{1-q}-\frac{2^{3}q^{2}}{1-q^{2}}+\frac{3^{3}q^{3}}{1-q^{3}}-\frac{4^{3}q^{4}}{1-q^{4}}+.\,.\right\},
\]
whence also, on changing $q$ into $-q$:
\[
\text{8)}\quad\left(\frac{2K}{\pi}\right)^{4}=1+16\left\{\frac{q}{1+q}+\frac{2^{3}q^{2}}{1-q^{2}}+\frac{3^{3}q^{3}}{1+q^{3}}+\frac{4^{3}q^{4}}{1-q^{4}}+.\,.\right\},
\]
which were the formulae more difficult to investigate. If you join these to those which we found above, you now have the first four powers of $\dfrac{2K}{\pi}$, $\dfrac{2kK}{\pi}$ expanded in fairly neat series.

\subsection*{GENERAL FORMULAE FOR EXPANDING THE FUNCTIONS $\sin^{n}\operatorname{am}\dfrac{2Kx}{\pi}$, $\dfrac{1}{\sin^{n}\operatorname{am}\dfrac{2Kx}{\pi}}$ IN SERIES, PROCEEDING ACCORDING TO THE SINES OR COSINES OF MULTIPLES OF $x$ ITSELF.}

\subsection*{43.}

The expansions of the functions:
\[
\sin\operatorname{am}\frac{2Kx}{\pi},\quad\sin^{2}\operatorname{am}\frac{2Kx}{\pi},\quad\frac{1}{\sin\operatorname{am}\dfrac{2Kx}{\pi}},\quad\frac{1}{\sin^{2}\operatorname{am}\dfrac{2Kx}{\pi}}
\]
having been found, the question now presents itself of carrying out the expansions of the higher powers of
\[
\sin\operatorname{am}\frac{2Kx}{\pi},\quad\frac{1}{\sin\operatorname{am}\dfrac{2Kx}{\pi}}
\]
In Analytic Trigonometry, indeed, an easy way is known, by which, the expansion of $\sin x$, $\cos x$ having been found, you may proceed to the expansion of the expressions $\sin^{n}x$, $\cos^{n}x$; namely, this succeeds by the help of known formulae, by which $\sin^{n}x$, $\cos^{n}x$ are exhibited linearly through the sines or cosines of multiples of $x$. But in the theory of Elliptic Functions

\origpage{116}
that aid fails; recourse will have to be had to another, which we shall set forth in what follows.

The formula, which is evident from the elements:
\[
\frac{d\sin^{n}\operatorname{am}u}{du}=n\sin^{n-1}\operatorname{am}u\sqrt{1-(1+kk)\sin^{2}\operatorname{am}u+kk\sin^{2}\operatorname{am}u},
\]\ednote{The last term under the root is printed $kk\sin^{2}\operatorname{am}u$ where $kk\sin^{4}\operatorname{am}u$ is expected (compare the equations 1) below); an apparent misprint.}
differentiated again, gives:
\[
\text{1)}\quad\frac{d^{2}\sin^{n}\operatorname{am}u}{du^{2}}=n(n-1)\sin^{n-2}\operatorname{am}u-nn(1+k^{2})\sin^{n}\operatorname{am}u+n(n+1)k^{2}\sin^{n+2}\operatorname{am}u.
\]
Putting successively $n=1$, 3, 5, 7 . ., $n=2$, 4, 6, 8 . ., a double series of equations may be formed from this:

\emph{I.}
\[
\begin{gathered}
\frac{d^{2}\sin\operatorname{am}u}{du^{2}}=-(1+k^{2})\sin\operatorname{am}u+2k^{2}\sin^{3}\operatorname{am}u\\[1.5ex]
\frac{d^{2}\sin^{3}\operatorname{am}u}{du^{2}}=6\sin\operatorname{am}u-9(1+k^{2})\sin^{3}\operatorname{am}u+12k^{2}\sin^{5}\operatorname{am}u\\[1.5ex]
\frac{d^{2}\sin^{5}\operatorname{am}u}{du^{2}}=20\sin^{3}\operatorname{am}u-25(1+k^{2})\sin^{5}\operatorname{am}u+30k^{2}\sin^{7}\operatorname{am}u\\[1.5ex]
\frac{d^{2}\sin^{7}\operatorname{am}u}{du^{2}}=42\sin^{5}\operatorname{am}u-49(1+k^{2})\sin^{7}\operatorname{am}u+56k^{2}\sin^{9}\operatorname{am}u\\[1.5ex]
\text{etc.}\qquad\qquad\text{etc.}
\end{gathered}
\]
\emph{II.}
\[
\begin{gathered}
\frac{d^{2}\sin^{2}\operatorname{am}u}{du^{2}}=2-4(1+k^{2})\sin^{2}\operatorname{am}u+6k^{2}\sin^{4}\operatorname{am}u\\[1.5ex]
\frac{d^{2}\sin^{4}\operatorname{am}u}{du^{2}}=12\sin^{2}\operatorname{am}u-16(1+k^{2})\sin^{4}\operatorname{am}u+20k^{2}\sin^{6}\operatorname{am}u\\[1.5ex]
\frac{d^{2}\sin^{6}\operatorname{am}u}{du^{2}}=30\sin^{4}\operatorname{am}u-36(1+k^{2})\sin^{6}\operatorname{am}u+42k^{2}\sin^{8}\operatorname{am}u\\[1.5ex]
\frac{d^{2}\sin^{8}\operatorname{am}u}{du^{2}}=56\sin^{6}\operatorname{am}u-64(1+k^{2})\sin^{8}\operatorname{am}u+72k^{2}\sin^{10}\operatorname{am}u\\[1.5ex]
\text{etc.}\qquad\qquad\text{etc.}
\end{gathered}
\]

\origpage{117}

From the equations I. you derive successively, putting $\Pi n=1.2.3.\,.\,n$:

\emph{I.\ a.}
\[
\begin{gathered}
\Pi 2.k^{2}\sin^{3}\operatorname{am}u=\frac{d^{2}.\sin\operatorname{am}u}{du^{2}}+(1+k^{2})\sin\operatorname{am}u\\[1.5ex]
\Pi 4.k^{4}\sin^{5}\operatorname{am}u=\frac{d^{4}.\sin\operatorname{am}u}{du^{4}}+10(1+k^{2})\frac{d^{2}.\sin\operatorname{am}u}{du^{2}}+3(3+2k^{2}+3k^{4})\sin\operatorname{am}u\\[1.5ex]
\Pi 6.k^{6}\sin^{7}\operatorname{am}u=\frac{d^{6}.\sin\operatorname{am}u}{du^{6}}+35(1+k^{2})\frac{d^{4}.\sin\operatorname{am}u}{du^{4}}+7(37+38k^{2}+37k^{4})\frac{d^{2}.\sin\operatorname{am}u}{du^{2}}\\[1.5ex]
+45(5+3k^{2}+3k^{4}+5k^{6})\sin\operatorname{am}u\\[1.5ex]
\Pi 8.k^{8}\sin^{9}\operatorname{am}u=\frac{d^{8}.\sin\operatorname{am}u}{du^{8}}+84(1+k^{2})\frac{d^{6}.\sin\operatorname{am}u}{du^{6}}+42(47+58k^{2}+47k^{4})\frac{d^{4}.\sin\operatorname{am}u}{du^{4}}\\[1.5ex]
+4(3229+3315k^{2}+3315k^{4}+3329k^{6})\frac{d^{2}.\sin\operatorname{am}u}{du^{2}}\\[1.5ex]
+315(35+20k^{2}+18k^{4}+20k^{6}+35k^{8})\sin\operatorname{am}u\\[1.5ex]
\text{etc.}\qquad\qquad\text{etc.}
\end{gathered}
\]\ednote{In the factor of the term with $d^{2}$ the last coefficient is printed $3329$ where $3229$ is expected by the symmetry of the polynomial (the author's Corrigenda give $3229$); an apparent misprint, corrected by hand in ink in the scan.}

\emph{II.\ a.}
\[
\begin{gathered}
\Pi 3.k^{2}\sin^{4}\operatorname{am}u=\frac{d^{2}.\sin^{2}\operatorname{am}u}{du^{2}}+4(1+k^{2})\sin^{2}\operatorname{am}u-2\\[1.5ex]
\Pi 5.k^{4}\sin^{6}\operatorname{am}u=\frac{d^{4}.\sin^{2}\operatorname{am}u}{du^{4}}+20(1+k^{2})\frac{d^{2}.\sin^{2}\operatorname{am}u}{du^{2}}+8(8+7k^{2}+8k^{4})\sin^{2}\operatorname{am}u-32(1+k^{2})\\[1.5ex]
\Pi 7.k^{6}\sin^{8}\operatorname{am}u=\frac{d^{6}.\sin^{2}\operatorname{am}u}{du^{6}}+56(1+k^{2})\frac{d^{4}.\sin^{2}\operatorname{am}u}{du^{4}}+112(7+8k^{2}+7k^{4})\frac{d^{2}.\sin^{2}\operatorname{am}u}{du^{2}}\\[1.5ex]
+128(18+15k^{2}+15k^{4}+18k^{6})\sin^{2}\operatorname{am}u-48(24+32k^{2}+24k^{4})\\[1.5ex]
\text{etc.}\qquad\qquad\text{etc.}
\end{gathered}
\]\ednote{In the last term of the $\Pi 7$ formula the middle coefficient is printed $32$ where $23$ is expected (the author's Corrigenda give $23$); an apparent misprint, corrected by hand in ink in the scan.}

Thus we see that it may be generally put:
\[
\begin{gathered}
\text{2)}\quad\Pi 2n.k^{2n}\sin^{2n+1}\operatorname{am}u=\\[1.5ex]
\frac{d^{2n}.\sin\operatorname{am}u}{du^{2n}}+A_{n}^{(1)}\frac{d^{2n-2}.\sin\operatorname{am}u}{du^{2n-2}}+A_{n}^{(2)}\frac{d^{2n-4}.\sin\operatorname{am}u}{du^{2n-4}}+.\,.+A_{n}^{(n)}\sin\operatorname{am}u\\[1.5ex]
\text{3)}\quad\Pi(2n-2).k^{2n-2}\sin^{2n}\operatorname{am}u=\\[1.5ex]
\frac{d^{2n-2}.\sin^{2}\operatorname{am}u}{du^{2n-2}}+B_{n}^{(1)}\frac{d^{2n-4}.\sin^{2}\operatorname{am}u}{du^{2n-4}}+B_{n}^{(2)}\frac{d^{2n-6}.\sin^{2}\operatorname{am}u}{du^{2n-6}}+.\,.+B_{n}^{(n-1)}\sin^{2}\operatorname{am}u+B_{n}^{(n)},
\end{gathered}
\]\ednote{The left side of formula 3) is printed $\Pi(2n-2)$ where $\Pi(2n-1)$ is expected (compare $\Pi 3$, $\Pi 5$, $\Pi 7$ above); an apparent misprint, corrected by hand in ink in the scan.}
where $A_{n}^{(m)}$, $B_{n}^{(m)}$ denote rational integral functions of $kk$ of the $m^{\mathrm{ti}}$ order, except $B_{n}^{(n)}$, which is of the $(n-2)^{\mathrm{ti}}$. Further, from the general formula from which we set out:
\[
\frac{d^{2}\sin^{n}\operatorname{am}u}{du^{2}}=n(n-1)\sin^{n-2}\operatorname{am}u-nn(1+k^{2})\sin^{n}\operatorname{am}u+n(n+1)k^{2}\sin^{n+2}\operatorname{am}u
\]

\origpage{118}
it is evident that there will be:
\[
\begin{gathered}
\text{4)}\quad A_{n}^{(m)}=A_{n-1}^{(m)}+(2n-1)^{2}(1+k^{2})A_{n-1}^{(m-1)}-(2n-2)^{2}(2n-1)(2n-3)k^{2}A_{n-2}^{(m-2)}\\[1.5ex]
\text{5)}\quad B_{n}^{(m)}=B_{n-1}^{(m)}+(2n-2)^{2}(1+k^{2})B_{n-1}^{(m-1)}-(2n-3)^{2}(2n-2)(2n-4)k^{2}B_{n-2}^{(m-2)},
\end{gathered}
\]
in which formulae, whenever $m>n$, one must put $A_{n}^{(m)}=0$, $B_{n}^{(m)}=0$.

If $u$ is changed into $u+iK'$, since $\sin\operatorname{am}u$ passes into $\dfrac{1}{k\sin\operatorname{am}u}$, in the proposed formulae $\dfrac{1}{k\sin\operatorname{am}u}$ may be put in place of $\sin\operatorname{am}u$, whence the following arise:
\[
\begin{gathered}
\frac{\Pi 2}{\sin^{3}\operatorname{am}u}=\frac{d^{2}.\dfrac{1}{\sin\operatorname{am}u}}{du^{2}}+(1+k^{2})\frac{1}{\sin\operatorname{am}u}\\[1.5ex]
\frac{\Pi 3}{\sin^{4}\operatorname{am}u}=\frac{d^{2}.\dfrac{1}{\sin^{2}\operatorname{am}u}}{du^{2}}+4(1+k^{2}).\frac{1}{\sin^{2}\operatorname{am}u}-2\\[1.5ex]
\frac{\Pi 4}{\sin^{5}\operatorname{am}u}=\frac{d^{4}.\dfrac{1}{\sin\operatorname{am}u}}{du^{4}}+10(1+k^{2})\frac{d^{2}.\dfrac{1}{\sin\operatorname{am}u}}{du^{2}}+\frac{3(3+2k^{2}+3k^{4})}{\sin\operatorname{am}u}\\[1.5ex]
\frac{\Pi 5}{\sin^{6}\operatorname{am}u}=\frac{d^{4}.\dfrac{1}{\sin^{2}\operatorname{am}u}}{du^{4}}+20(1+k^{2})\frac{d^{2}.\dfrac{1}{\sin^{2}\operatorname{am}u}}{du^{2}}+\frac{8(8+7k^{2}+8k^{4})}{\sin^{2}\operatorname{am}u}-32(1+k^{2}).\\[1.5ex]
\text{etc.}\qquad\qquad\text{etc.},
\end{gathered}
\]\ednote{In the second formula the last term is printed $-2$ where $-2k^{2}$ is expected, and in the fourth formula the last term is printed $-32(1+k^{2})$ where $-32k^{2}(1+k^{2})$ is expected (the author's Corrigenda); apparent omissions of the factor $k^{2}$ (in the fourth formula a period is left standing after the bracket), supplied by hand in ink in the scan.}
and generally:
\[
\begin{gathered}
\text{6)}\quad\frac{\Pi 2n}{\sin^{2n+1}\operatorname{am}u}=\\[1.5ex]
\frac{d^{2n}.\dfrac{1}{\sin\operatorname{am}u}}{du^{2n}}+A_{n}^{(1)}\frac{d^{2n-2}.\dfrac{1}{\sin\operatorname{am}u}}{du^{2n-2}}+A_{n}^{(2)}\frac{d^{2n-4}.\dfrac{1}{\sin\operatorname{am}u}}{du^{2n-4}}+.\,.+\frac{A_{n}^{(n)}}{\sin\operatorname{am}u}\\[1.5ex]
\text{7)}\quad\frac{\Pi(2n-1)}{\sin^{2n}\operatorname{am}u}=\\[1.5ex]
\frac{d^{2n-2}.\dfrac{1}{\sin^{2}\operatorname{am}u}}{du^{2n-2}}+B_{n}^{(1)}\frac{d^{2n-4}.\dfrac{1}{\sin^{2}\operatorname{am}u}}{du^{2n-4}}+B_{n}^{(2)}\frac{d^{2n-6}.\dfrac{1}{\sin^{2}\operatorname{am}u}}{du^{2n-6}}+.\,.+\frac{B_{n}^{(n-1)}}{\sin^{2}\operatorname{am}u}+k^{2}B_{n}^{(n)}.
\end{gathered}
\]\ednote{In the third term of 7) the exponent of $d$ is printed with a damaged character in place of $n$ ($d^{2:-6}$); read here as $d^{2n-6}$.}

\origpage{119}
\subsection*{44.}

Since it has been found in what precedes, if one puts $u=\dfrac{2Kx}{\pi}$, that the expressions
\[
\sin^{n}\operatorname{am}\frac{2Kx}{\pi},\quad\frac{1}{\sin^{n}\operatorname{am}\dfrac{2Kx}{\pi}}
\]
can be expressed linearly through these:
\[
\sin\operatorname{am}\frac{2Kx}{\pi},\quad\sin^{2}\operatorname{am}\frac{2Kx}{\pi},\quad\frac{1}{\sin\operatorname{am}\dfrac{2Kx}{\pi}},\quad\frac{1}{\sin^{2}\operatorname{am}\dfrac{2Kx}{\pi}},
\]
and their differentials, taken with respect to the argument $u$ or $x$, the former follow of themselves from the expansions of the latter, which proceed according to the sines or cosines of multiples of $x$.

Thus we obtain:

\emph{I.}

from the formula:
\[
\frac{2kK}{\pi}.\sin\operatorname{am}\frac{2Kx}{\pi}=4\left\{\frac{\sqrt{q}\sin x}{1-q}+\frac{\sqrt{q^{3}}\sin 3x}{1-q^{3}}+\frac{\sqrt{q^{5}}\sin 5x}{1-q^{5}}+.\,.\right\}
\]
the following:
\[
\begin{gathered}
2\left(\frac{2kK}{\pi}\right)^{3}\sin^{3}\operatorname{am}\frac{2Kx}{\pi}=\\[1.5ex]
4\left\{(1+k^{2})\left(\frac{2K}{\pi}\right)^{2}-1\right\}\frac{\sqrt{q}\sin x}{1-q}+\\[1.5ex]
4\left\{(1+k^{2})\left(\frac{2K}{\pi}\right)^{2}-3^{2}\right\}\frac{\sqrt{q^{3}}\sin 3x}{1-q^{3}}+\\[1.5ex]
4\left\{(1+k^{2})\left(\frac{2K}{\pi}\right)^{2}-5^{2}\right\}\frac{\sqrt{q^{5}}\sin 5x}{1-q^{5}}+\\[1.5ex]
\vdots\\[1.5ex]
2.3.4\left(\frac{2kK}{\pi}\right)^{5}\sin^{5}\operatorname{am}\frac{2Kx}{\pi}=\\[1.5ex]
4\left\{3(3+2k^{2}+3k^{4})\left(\frac{2K}{\pi}\right)^{4}-10(1+k^{2})\left(\frac{2K}{\pi}\right)^{2}+1\right\}\frac{\sqrt{q}\sin x}{1-q}+
\end{gathered}
\]

\origpage{120}
\[
\begin{gathered}
4\left\{3(3+2k^{2}+3k^{4})\left(\frac{2K}{\pi}\right)^{4}-3^{2}.10(1+k^{2})\left(\frac{2K}{\pi}\right)^{2}+3^{4}\right\}\frac{\sqrt{q^{3}}\sin 3x}{1-q^{3}}+\\[1.5ex]
4\left\{3(3+2k^{2}+3k^{4})\left(\frac{2K}{\pi}\right)^{4}-5^{2}.10(1+k^{2})\left(\frac{2K}{\pi}\right)^{2}+5^{4}\right\}\frac{\sqrt{q^{5}}\sin 5x}{1-q^{5}}+\\[1.5ex]
\vdots\\[1.5ex]
\text{etc.}\qquad\qquad\text{etc.};
\end{gathered}
\]

\emph{II.}

from the formula:
\[
\left(\frac{2kK}{\pi}\right)^{2}\sin^{2}\operatorname{am}\frac{2Kx}{\pi}=\frac{2K}{\pi}.\,\frac{2K}{\pi}-\frac{2K}{\pi}.\,\frac{2E^{\mathrm{I}}}{\pi}-4\left\{\frac{2q\cos 2x}{1-q^{2}}+\frac{4q^{2}\cos 4x}{1-q^{4}}+\frac{6q^{3}\cos 6x}{1-q^{6}}+.\,.\right\}
\]
the following:
\[
\begin{gathered}
2.3\left(\frac{2kK}{\pi}\right)^{4}\sin^{4}\operatorname{am}\frac{2Kx}{\pi}=\\[1.5ex]
4(1+k^{2})\left(\frac{2K}{\pi}\right)^{3}\left(\frac{2K}{\pi}-\frac{2E^{\mathrm{I}}}{\pi}\right)-2k^{2}\left(\frac{2K}{\pi}\right)^{4}\\[1.5ex]
-4\left\{2.4(1+k^{2})\left(\frac{2K}{\pi}\right)^{2}-2^{3}\right\}\frac{q\cos 2x}{1-q^{2}}\\[1.5ex]
-4\left\{4.4(1+k^{2})\left(\frac{2K}{\pi}\right)^{2}-4^{3}\right\}\frac{q^{2}\cos 4x}{1-q^{4}}\\[1.5ex]
-4\left\{6.4(1+k^{2})\left(\frac{2K}{\pi}\right)^{2}-6^{3}\right\}\frac{q^{3}\cos 6x}{1-q^{6}}\\[1.5ex]
\vdots\\[1.5ex]
2.3.4.5\left(\frac{2kK}{\pi}\right)^{6}\sin^{6}\operatorname{am}\frac{2Kx}{\pi}=\\[1.5ex]
8(8+7k^{2}+8k^{4})\left(\frac{2K}{\pi}\right)^{5}\left(\frac{2K}{\pi}.\,\frac{2K}{\pi}-\frac{2K}{\pi}.\,\frac{2E^{\mathrm{I}}}{\pi}\right)-32k^{2}(1+k^{2})\left(\frac{2K}{\pi}\right)^{6}\\[1.5ex]
-4\left\{2.8(8+7k^{2}+8k^{4})\left(\frac{2K}{\pi}\right)^{4}-2^{3}.20(1+k^{2})\left(\frac{2K}{\pi}\right)^{2}+2^{5}\right\}\frac{q\cos 2x}{1-q^{2}}\\[1.5ex]
-4\left\{4.8(8+7k^{2}+8k^{4})\left(\frac{2K}{\pi}\right)^{4}-4^{3}.20(1+k^{2})\left(\frac{2K}{\pi}\right)^{2}+4^{5}\right\}\frac{q^{2}\cos 4x}{1-q^{4}}\\[1.5ex]
-4\left\{6.8(8+7k^{2}+8k^{4})\left(\frac{2K}{\pi}\right)^{4}-6^{3}.20(1+k^{2})\left(\frac{2K}{\pi}\right)^{2}+6^{5}\right\}\frac{q^{3}\cos 6x}{1-q^{6}}\\[1.5ex]
\vdots\\[1.5ex]
\text{etc.}\qquad\qquad\text{etc.};
\end{gathered}
\]\ednote{In the first line of the $\sin^{6}$ formula the bracket is printed with the factors $\frac{2K}{\pi}.\,\frac{2K}{\pi}-\frac{2K}{\pi}.\,\frac{2E^{\mathrm{I}}}{\pi}$, where the analogous $\sin^{4}$ formula has $\frac{2K}{\pi}-\frac{2E^{\mathrm{I}}}{\pi}$; the repeated factors are struck through by hand in ink in the scan, and the print is reproduced as it stands.}

\origpage{121}
\emph{III.}

from the formula:
\[
\frac{\dfrac{2K}{\pi}}{\sin\operatorname{am}\dfrac{2Kx}{\pi}}=\frac{1}{\sin x}+\frac{4q\sin x}{1-q}+\frac{4q^{3}\sin 3x}{1-q^{3}}+\frac{4q^{5}\sin 5x}{1-q^{5}}+\,.\,.
\]

the following:
\[
\begin{gathered}
\frac{2\left(\dfrac{2K}{\pi}\right)^{3}}{\sin^{2}\operatorname{am}\dfrac{2Kx}{\pi}}=\\[1.5ex]
(1+k^{2})\left(\frac{2K}{\pi}\right)^{2}\frac{1}{\sin x}+\frac{d^{2}.\dfrac{1}{\sin x}}{dx^{2}}+\\[1.5ex]
4\left\{(1+k^{2})\left(\frac{2K}{\pi}\right)^{2}-1\right\}\frac{q\sin x}{1-q}+\\[1.5ex]
4\left\{(1+k^{2})\left(\frac{2K}{\pi}\right)^{2}-3^{2}\right\}\frac{q^{3}\sin 3x}{1-q^{3}}+\\[1.5ex]
4\left\{(1+k^{2})\left(\frac{2K}{\pi}\right)^{2}-5^{2}\right\}\frac{q^{5}\sin 5x}{1-q^{5}}+\\[1.5ex]
\vdots\\[1.5ex]
\frac{2.3.4\left(\dfrac{2K}{\pi}\right)^{5}}{\sin^{5}\operatorname{am}\dfrac{2Kx}{\pi}}=\\[1.5ex]
\frac{3(3+2k^{2}+3k^{4})\left(\dfrac{2K}{\pi}\right)^{4}}{\sin x}+10(1+k^{2})\left(\frac{2K}{\pi}\right)^{2}\frac{d^{2}.\dfrac{1}{\sin x}}{dx^{2}}+\frac{d^{4}.\dfrac{1}{\sin x}}{dx^{4}}+\\[1.5ex]
4\left\{3(3+2k^{2}+3k^{4})\left(\frac{2K}{\pi}\right)^{4}-10(1+k^{2})\left(\frac{2K}{\pi}\right)^{2}+1\right\}\frac{q\sin x}{1-q}+\\[1.5ex]
4\left\{3(3+2k^{2}+3k^{4})\left(\frac{2K}{\pi}\right)^{4}-3^{2}.10(1+k^{2})\left(\frac{2K}{\pi}\right)^{2}+3^{4}\right\}\frac{q^{3}\sin 3x}{1-q^{3}}+\\[1.5ex]
4\left\{3(3+2k^{2}+3k^{4})\left(\frac{2K}{\pi}\right)^{4}-5^{2}.10(1+k^{2})\left(\frac{2K}{\pi}\right)^{2}+5^{4}\right\}\frac{q^{5}\sin 5x}{1-q^{5}}+\\[1.5ex]
\vdots\\[1.5ex]
\text{etc.}\qquad\qquad\text{etc.};
\end{gathered}
\]

\origpage{122}
\emph{IV.}

From the formula:
\[
\frac{\left(\dfrac{2K}{\pi}\right)^{2}}{\sin^{2}\operatorname{am}\dfrac{2Kx}{\pi}}=\frac{2K}{\pi}\left(\frac{2K}{\pi}-\frac{2E^{\mathrm{I}}}{\pi}\right)+\frac{1}{\sin^{2}x}-4\left\{\frac{2q^{2}\cos 2x}{1-q^{2}}+\frac{4q^{4}\cos 4x}{1-q^{4}}+\frac{6q^{6}\cos 6x}{1-q^{6}}+\,.\,.\right\}
\]

the following:
\[
\begin{gathered}
\frac{2.3\left(\dfrac{2K}{\pi}\right)^{4}}{\sin^{4}\operatorname{am}\dfrac{2Kx}{\pi}}=\\[1.5ex]
4(1+k^{2})\left(\frac{2K}{\pi}\right)^{3}\left(\frac{2K}{\pi}-\frac{2E^{\mathrm{I}}}{\pi}\right)-2k^{2}\left(\frac{2K}{\pi}\right)^{4}+\\[1.5ex]
\frac{4(1+k^{2})\left(\dfrac{2K}{\pi}\right)^{2}}{\sin^{2}x}+\frac{d^{2}.\dfrac{1}{\sin^{2}x}}{dx^{2}}\\[1.5ex]
-4\left\{2.4(1+k^{2})\left(\frac{2K}{\pi}\right)^{2}-2^{3}\right\}\frac{q^{2}\cos 2x}{1-q^{2}}\\[1.5ex]
-4\left\{4.4(1+k^{2})\left(\frac{2K}{\pi}\right)^{2}-4^{3}\right\}\frac{q^{4}\cos 4x}{1-q^{4}}\\[1.5ex]
-4\left\{4.6(1+k^{2})\left(\frac{2K}{\pi}\right)^{2}-6^{3}\right\}\frac{q^{6}\cos 6x}{1-q^{6}}\\[1.5ex]
\vdots\\[1.5ex]
\frac{2.3.4.5\left(\dfrac{2K}{\pi}\right)^{6}}{\sin^{6}\operatorname{am}\dfrac{2Kx}{\pi}}=\\[1.5ex]
8(8+7k^{2}+8k^{4})\left(\frac{2K}{\pi}\right)^{5}\left(\frac{2K}{\pi}-\frac{2E^{\mathrm{I}}}{\pi}\right)-32k^{2}(1+k^{2})\left(\frac{2K}{\pi}\right)^{6}\\[1.5ex]
+\frac{8(8+7k^{2}+8k^{4})\left(\dfrac{2K}{\pi}\right)^{4}}{\sin^{2}x}+20(1+k^{2})\left(\frac{2K}{\pi}\right)^{2}\frac{d^{2}.\dfrac{1}{\sin^{2}x}}{dx^{2}}+\frac{d^{4}.\dfrac{1}{\sin^{2}x}}{dx^{4}}\\[1.5ex]
-4\left\{2.8(8+7k^{2}+8k^{4})\left(\frac{2K}{\pi}\right)^{4}-2^{3}.20(1+k^{2})\left(\frac{2K}{\pi}\right)^{2}+2^{5}\right\}\frac{q^{2}\cos 2x}{1-q^{2}}
\end{gathered}
\]

\origpage{123}
\[
\begin{gathered}
-4\left\{4.8(8+7k^{2}+8k)\left(\frac{2K}{\pi}\right)^{4}-4^{3}.20(1+k^{2})\left(\frac{2K}{\pi}\right)^{2}+4^{5}\right\}\frac{q^{4}\cos 4x}{1-q^{4}}\\[1.5ex]
-4\left\{6.8(8+7k^{2}+8k^{4})\left(\frac{2K}{\pi}\right)^{4}-6^{3}.20(1+k^{2})\left(\frac{2K}{\pi}\right)^{2}+6^{5}\right\}\frac{q^{6}\cos 6x}{1-q^{6}}\\[1.5ex]
\vdots\\[1.5ex]
\text{cet.}\qquad\qquad\text{cet.}
\end{gathered}
\]
\ednote{In the first row the exponent 4 of $k$ in $8k^{4}$ is not printed (apparently dropped ink); it is printed $8k$ and the expected reading is $8k^{4}$.}

\subsection*{45.}
The examples proposed above show how the expansions of the functions $\sin^{n}\operatorname{am}\dfrac{2Kx}{\pi}$, $\dfrac{1}{\sin^{n}\operatorname{am}\dfrac{2Kx}{\pi}}$ are found from the formulae 2), 3), 6), 7) §. 43. The quantities $A_{n}^{(m)}$, $B_{n}^{(m)}$, on which those depend, may be drawn out successively by means of the formulae 4), 5) of the same place. But the question of investigating their general expressions, since these turn out too complicated to be attained by induction, must be taken up from somewhat further back. To this end we premise the following.

The elementary formula is known:
\[
\sin\operatorname{am}(u+v)-\sin\operatorname{am}(u-v)=\frac{2\sin\operatorname{am}v.\cos\operatorname{am}u\,\Delta\operatorname{am}u}{1-k^{2}\sin^{2}\operatorname{am}u\sin^{2}\operatorname{am}v},
\]
and on integrating this with respect to $u$, there results:
\[
\text{1)}\quad\int_{0}^{u}du\left\{\sin\operatorname{am}(u+v)-\sin\operatorname{am}(u-v)\right\}=\frac{1}{k}\log\left(\frac{1+k\sin\operatorname{am}u\sin\operatorname{am}v}{1-k\sin\operatorname{am}u\sin\operatorname{am}v}\right).
\]
From Taylor's theorem it comes about that:
\[
\begin{gathered}
\sin\operatorname{am}(u+v)-\sin\operatorname{am}(u-v)=\\[1.5ex]
2\left\{\frac{d.\sin\operatorname{am}u}{du}.v+\frac{d^{3}.\sin\operatorname{am}u}{du^{3}}.\frac{v^{3}}{\Pi 3}+\frac{d^{5}.\sin\operatorname{am}u}{du^{5}}.\frac{v^{5}}{\Pi 5}+\ldots\right\},
\end{gathered}
\]
whence:
\[
\begin{gathered}
\int_{0}^{u}du\left\{\sin\operatorname{am}(u+v)-\sin\operatorname{am}(u-v)\right\}=\\[1.5ex]
2\left\{\sin\operatorname{am}u.v+\frac{d^{2}.\sin\operatorname{am}u}{du^{2}}.\frac{v^{3}}{\Pi 3}+\frac{d^{4}.\sin\operatorname{am}u}{du^{4}}.\frac{v^{5}}{\Pi 5}+\ldots\right\}.
\end{gathered}
\]

\origpage{124}
For it is readily established that, on putting $u=0$, both $\sin\operatorname{am}u$ and, generally, $\dfrac{d^{2m}\sin\operatorname{am}u}{du^{2m}}$ vanish. Hence equation 1), with its other side also expanded, passes into this:
\[
\begin{gathered}
\text{2)}\quad\sin\operatorname{am}u.v+\frac{d^{2}.\sin\operatorname{am}u}{du^{2}}.\frac{v^{3}}{\Pi 3}+\frac{d^{4}.\sin\operatorname{am}u}{du^{4}}.\frac{v^{5}}{\Pi 5}+\ldots=\\[1.5ex]
\sin\operatorname{am}u\sin\operatorname{am}v+\frac{k^{2}}{3}.\sin^{3}\operatorname{am}u\sin^{3}\operatorname{am}v+\frac{k^{4}}{5}\sin^{5}\operatorname{am}u\sin^{5}\operatorname{am}v+\,.\,.
\end{gathered}
\]

Further, by multiplying into one another the known equations:
\[
\begin{gathered}
\sin\operatorname{am}(u+v)+\sin\operatorname{am}(u-v)=\frac{2\sin\operatorname{am}u\,.\cos\operatorname{am}v\,\Delta\operatorname{am}v}{1-k^{2}\sin^{2}\operatorname{am}u\sin^{2}\operatorname{am}v}\\[1.5ex]
\sin\operatorname{am}(u+v)-\sin\operatorname{am}(u-v)=\frac{2\sin\operatorname{am}v\,.\cos\operatorname{am}u\,\Delta\operatorname{am}u}{1-k^{2}\sin^{2}\operatorname{am}u\sin^{2}\operatorname{am}v}
\end{gathered}
\]
we obtain:
\[
\begin{gathered}
\text{3)}\quad\sin^{2}\operatorname{am}(u+v)-\sin^{2}\operatorname{am}(u-v)=\\[1.5ex]
\frac{4\sin\operatorname{am}u\cos\operatorname{am}u\,\Delta\operatorname{am}u\,.\sin\operatorname{am}v\cos\operatorname{am}v\,\Delta\operatorname{am}v}{\left\{1-k^{2}\sin^{2}\operatorname{am}u\sin^{2}\operatorname{am}v\right\}^{2}}=\\[1.5ex]
\frac{d.\sin^{2}\operatorname{am}u\,.\,d.\sin^{2}\operatorname{am}v}{\left\{1-k^{2}\sin^{2}\operatorname{am}u\sin^{2}\operatorname{am}v\right\}^{2}du\,dv}.
\end{gathered}
\]
When the integration has been made with respect to $v$, there results:
\[
\begin{gathered}
\int_{0}^{v}dv\left\{\sin^{2}\operatorname{am}(u+v)-\sin^{2}\operatorname{am}(u-v)\right\}=\\[1.5ex]
\frac{2\sin\operatorname{am}u\cos\operatorname{am}u\,\Delta\operatorname{am}u\,.\sin^{2}\operatorname{am}v}{1-k^{2}\sin^{2}\operatorname{am}u\sin^{2}\operatorname{am}v}=\frac{\sin^{2}\operatorname{am}v\,.\,d.\sin^{2}\operatorname{am}u}{\left(1-k^{2}\sin^{2}\operatorname{am}u\sin^{2}\operatorname{am}v\right)du}
\end{gathered}
\]
On integrating this again with respect to the other element $u$, we obtain:
\[
\begin{gathered}
\text{4)}\quad\int_{0}^{u}du\int_{0}^{v}dv\left\{\sin^{2}\operatorname{am}(u+v)-\sin^{2}\operatorname{am}(u-v)\right\}=\\[1.5ex]
-\frac{1}{k^{2}}\log\left(1-k^{2}\sin^{2}\operatorname{am}u\sin^{2}\operatorname{am}v\right).
\end{gathered}
\]
From Taylor's theorem it comes about that:
\[
\begin{gathered}
\sin^{2}\operatorname{am}(u+v)-\sin^{2}\operatorname{am}(u-v)=\\[1.5ex]
2\left\{\frac{d.\sin^{2}\operatorname{am}u}{du}.v+\frac{d^{3}.\sin^{2}\operatorname{am}u}{du^{3}}.\frac{v^{3}}{\Pi 3}+\frac{d^{5}.\sin^{2}\operatorname{am}u}{du^{5}}.\frac{v^{5}}{\Pi 5}+\,.\,.\right\},
\end{gathered}
\]

\origpage{125}
whence:
\[
\begin{gathered}
\int_{0}^{v}dv\left\{\sin^{2}\operatorname{am}(u+v)-\sin^{2}\operatorname{am}(u-v)\right\}=\\[1.5ex]
2\left\{\frac{d.\sin^{2}\operatorname{am}u}{du}.\frac{v^{2}}{\Pi 2}+\frac{d^{3}.\sin^{2}\operatorname{am}u}{du^{3}}.\frac{v^{4}}{\Pi 4}+\frac{d^{5}.\sin^{2}\operatorname{am}u}{du^{5}}.\frac{v^{6}}{\Pi 6}+\,.\,.\right\}\\[1.5ex]
\int_{0}^{u}du\int_{0}^{v}dv\left\{\sin^{2}\operatorname{am}(u+v)-\sin^{2}\operatorname{am}(u-v)\right\}=\\[1.5ex]
2\left\{\sin^{2}\operatorname{am}u.\frac{v^{2}}{\Pi 2}+\frac{d^{2}.\sin^{2}\operatorname{am}u}{du^{2}}.\frac{v^{4}}{\Pi 4}+\frac{d^{4}.\sin^{2}\operatorname{am}u}{du^{4}}.\frac{v^{6}}{\Pi 6}+\,.\,.\right\}\\[1.5ex]
-2\left\{U^{(2)}\frac{v^{4}}{\Pi 4}+U^{(4)}\frac{v^{6}}{\Pi 6}+\ldots\right\},
\end{gathered}
\]
provided that by the character $U^{(2m)}$ we denote the value of the expression $\dfrac{d^{2m}.\sin^{2}\operatorname{am}u}{du^{2m}}$ which it obtains on putting $u=0$. Hence equation 4), with its other side also expanded, passes into this:
\[
\begin{gathered}
\text{5)}\quad\sin^{2}\operatorname{am}u.\frac{v^{2}}{\Pi 2}+\frac{d^{2}.\sin^{2}\operatorname{am}u}{du^{2}}.\frac{v^{4}}{\Pi 4}+\frac{d^{4}.\sin^{2}\operatorname{am}u}{du^{4}}.\frac{v^{6}}{\Pi 6}+\ldots\\[1.5ex]
-2\left\{U^{(2)}\frac{v^{4}}{\Pi 4}+U^{(4)}\frac{v^{6}}{\Pi 6}+\ldots\right\}\\[1.5ex]
=\\[1.5ex]
\frac{1}{2}.\sin^{2}\operatorname{am}u\sin^{2}\operatorname{am}v+\frac{k^{2}}{4}.\sin^{4}\operatorname{am}u\sin^{4}\operatorname{am}v+\frac{k^{4}}{6}\sin^{6}\operatorname{am}u\sin^{6}\operatorname{am}v+\ldots
\end{gathered}
\]\ednote{In the second row of formula 5) the print has $-2$ before the brace where $-$ is expected (the author's Corrigenda: for $-2$ read $-$); reproduced as printed.}

These things being duly prepared, let there be put:
\[
u=\sin\operatorname{am}u+R_{1}\sin^{3}\operatorname{am}u+R_{2}\sin^{5}\operatorname{am}u+R_{3}\sin^{7}\operatorname{am}u+\ldots,
\]
and generally:
\[
\begin{gathered}
u^{n}=\left\{\sin\operatorname{am}u+R_{1}\sin^{3}\operatorname{am}u+R_{2}\sin^{5}\operatorname{am}u+R_{3}\sin^{7}\operatorname{am}u+\ldots\right\}^{n}=\\[1.5ex]
\sin^{n}\operatorname{am}u+R_{1}^{(n)}\sin^{n+2}\operatorname{am}u+R_{2}^{(n)}\sin^{n+4}\operatorname{am}u+R_{3}^{(n)}\sin^{n+6}\operatorname{am}u+\ldots:
\end{gathered}
\]
further, from the reversion of the series:
\[
u=\sin\operatorname{am}u+R_{1}\sin^{3}\operatorname{am}u+R_{2}\sin^{5}\operatorname{am}u+R_{3}\sin^{7}\operatorname{am}u+\ldots
\]
let this arise:
\[
\sin\operatorname{am}u=u+S_{1}u^{3}+S_{2}u^{5}+S_{3}u^{7}+\,.\,.,
\]

\origpage{126}
and let there again be:
\[
\begin{gathered}
\sin^{n}\operatorname{am}u=\left\{u+S_{1}u^{3}+S_{2}u^{5}+S_{3}u^{7}+\,.\,.\right\}^{n}=\\[1.5ex]
u^{n}+S_{1}^{(n)}u^{n+2}+S_{2}^{(n)}u^{n+4}+S_{3}^{(n)}u^{n+6}+\ldots
\end{gathered}
\]

Now from equation 2):
\[
\begin{gathered}
\sin\operatorname{am}u.v+\frac{d^{2}.\sin\operatorname{am}u}{du^{2}}.\frac{v^{3}}{\Pi 3}+\frac{d^{4}.\sin\operatorname{am}u}{du^{4}}.\frac{v^{5}}{\Pi 5}+\ldots=\\[1.5ex]
\sin\operatorname{am}u\sin\operatorname{am}v+\frac{k^{2}}{3}\sin^{3}\operatorname{am}u\sin^{3}\operatorname{am}v+\frac{k^{4}}{5}.\sin^{5}\operatorname{am}u\sin^{5}\operatorname{am}v+\,.\,.,
\end{gathered}
\]
when $v$, $v^{3}$, $v^{5}$, etc. have been expanded in series proceeding according to the powers of $\sin\operatorname{am}v$, and the Coefficients of the same power $\sin^{2n+1}\operatorname{am}v$ have been compared with each other in both sides of the equation, there results:
\[
\begin{gathered}
\text{6)}\quad\frac{k^{2n}\sin^{2n+1}\operatorname{am}u}{2n+1}=\\[1.5ex]
R_{n}^{(1)}\sin\operatorname{am}u+R_{n-1}^{(3)}\frac{d^{2}.\sin\operatorname{am}u}{\Pi 3.du^{2}}+R_{n-2}^{(5)}\frac{d^{4}.\sin\operatorname{am}u}{\Pi 5.du^{4}}+\,.\,.+\frac{d^{2n}.\sin\operatorname{am}u}{\Pi(2n+1)du^{2n}}.
\end{gathered}
\]
In the same way from formula 5) there results:
\[
\begin{gathered}
\text{7)}\quad\frac{k^{2n-2}\sin^{2n}\operatorname{am}u}{2n}=\\[1.5ex]
R_{n-1}^{(2)}\frac{\sin^{2}\operatorname{am}u}{\Pi 2}+R_{n-2}^{(4)}\frac{d^{2}.\sin^{2}\operatorname{am}u}{\Pi 4.du^{2}}+R_{n-3}^{(6)}\frac{d^{4}.\sin^{2}\operatorname{am}u}{\Pi 6.du^{4}}+\,.\,.+\frac{d^{2n-2}\sin^{2}\operatorname{am}u}{\Pi 2n.du^{2n-2}}\\[1.5ex]
-\left\{\frac{1}{3.4}+\frac{S_{1}^{(2)}}{5.6}+\frac{S_{2}^{(2)}}{7.8}+\,.\,.+\frac{S_{n-2}^{(2)}}{(2n-1)2n}\right\}{}^{*)}.
\end{gathered}
\]
From 6), 7), on changing $u$ into $u+iK'$, it follows:
\[
\begin{gathered}
\text{8)}\quad\frac{1}{(2n+1)\sin^{2n+1}\operatorname{am}u}=\\[1.5ex]
\frac{R_{n}^{(1)}}{\sin\operatorname{am}u}+R_{n-1}^{(3)}\frac{d^{2}.\dfrac{1}{\sin\operatorname{am}u}}{\Pi 3.du^{2}}+R_{n-2}^{(5)}\frac{d^{4}.\dfrac{1}{\sin\operatorname{am}u}}{\Pi 5.du^{4}}+\,.\,.+\frac{d^{2n}.\dfrac{1}{\sin\operatorname{am}u}}{du^{2n}}
\end{gathered}
\]\ednote{The last term of formula 8) is printed without the factor $\Pi(2n+1)$ in the denominator, which the pattern of the preceding terms and of formula 6) would lead one to expect; reproduced as printed.}

\textbf{*)} For from the notation proposed it comes about that: $\sin^{2}\operatorname{am}u=u^{2}+S_{1}^{(2)}u^{4}+S_{2}^{(2)}u^{6}+S_{3}^{(2)}u^{9}+\,.\,.$, whence, since $U^{(m)}=\dfrac{d^{2m}\sin^{2}\operatorname{am}u}{du^{2m}}$ for the value $u=0$, $U^{(m)}=\Pi 2m.S_{m-1}^{(2)}$.\ednote{The exponent of the last term shown is printed $u^{9}$ where $u^{8}$ is expected; and the index of $U$ is printed $(m)$ in both places where $(2m)$ is used on p. 125. Both reproduced as printed.}

\origpage{127}
\[
\begin{gathered}
\text{9)}\quad\frac{1}{2n\sin^{2n}\operatorname{am}u}=\\[1.5ex]
\frac{R_{n-1}^{(2)}}{\Pi 2.\sin^{2}\operatorname{am}u}+R_{n-2}^{(4)}\frac{d^{2}.\dfrac{1}{\sin^{2}\operatorname{am}u}}{\Pi 4.du^{2}}+R_{n-3}^{(6)}\frac{d^{4}.\dfrac{1}{\sin^{2}\operatorname{am}u}}{\Pi 6.du^{4}}+\,.\,.+\frac{d^{2n-2}.\dfrac{1}{\sin^{2}\operatorname{am}u}}{\Pi 2n.du^{2n-2}}\\[1.5ex]
-k^{2}\left\{\frac{1}{3.4}+\frac{S_{1}^{(2)}}{5.6}+\frac{S_{2}^{(2)}}{7.8}+\,.\,.+\frac{S_{n-2}^{(2)}}{(2n-1)2n}\right\}.
\end{gathered}
\]
These are the general formulae which we sought, by whose aid $\sin^{n}\operatorname{am}u$, $\dfrac{1}{\sin^{n}\operatorname{am}u}$ are found from $\sin\operatorname{am}u$, $\sin^{2}\operatorname{am}u$, $\dfrac{1}{\sin\operatorname{am}u}$, $\dfrac{1}{\sin^{2}\operatorname{am}u}$ and their differentials.

I shall remark on this occasion that, where conversely you expand $\sin\operatorname{am}v$, $\sin^{2}\operatorname{am}v$, $\sin^{3}\operatorname{am}v$, etc. according to the powers of $x$, there are drawn from formulae 2), 5):
\[
\begin{gathered}
\text{10)}\quad\frac{d^{2n}.\sin\operatorname{am}u}{\Pi(2n+1).du^{2n}}=\\[1.5ex]
S_{n}^{(1)}\sin\operatorname{am}u+\frac{k^{2}}{3}S_{n-1}^{(3)}\sin^{3}\operatorname{am}u+\frac{k^{4}}{5}S_{n-2}^{(5)}\sin^{5}\operatorname{am}u+\,.\,.+\frac{k^{2n}}{2n+1}\sin^{2n+1}\operatorname{am}u\\[1.5ex]
\text{11)}\quad\frac{d^{2n}.\sin^{2}\operatorname{am}u}{\Pi(2n+2)du^{2n}}-\frac{S_{n-1}^{(2)}}{(2n+1)(2n+2)}=\\[1.5ex]
\frac{1}{2}S_{n}^{(2)}\sin^{2}\operatorname{am}u+\frac{k^{2}}{4}S_{n-1}^{(4)}\sin^{4}\operatorname{am}u+\frac{k^{4}}{6}S_{n-2}^{(6)}\sin^{6}\operatorname{am}u+\,.\,.+\frac{k^{2n}}{2n+2}\sin^{2n+2}\operatorname{am}u.
\end{gathered}
\]

A few things remain to be added concerning the finding of the quantities $R_{m}^{(n)}$, $S_{m}^{(n)}$. On putting $\sin\operatorname{am}u=y$, it comes about from the definition proposed:
\[
u=\int_{0}^{y}\frac{dy}{\sqrt{(1-y^{2})(1-k^{2}y^{2})}}=y+R_{1}y^{3}+R_{2}y^{5}+R_{3}y^{7}+\ldots,
\]
or:
\[
\frac{1}{\sqrt{(1-y^{2})(1-k^{2}y^{2})}}=1+3R_{1}y^{2}+5R_{2}y^{4}+7R_{3}y^{6}+\ldots;
\]
whence:
\[
\begin{gathered}
3R_{1}=\frac{1+k^{2}}{2},\qquad 5R_{2}=\frac{1.3}{2.4}+\frac{1}{2}.\frac{1}{2}k^{2}+\frac{1.3}{2.4}k^{4}\\[1.5ex]
7R_{3}=\frac{1.3.5}{2.4.6}+\frac{1.3}{2.4}.\frac{1}{2}k^{2}+\frac{1}{2}.\frac{1.3}{2.4}k^{4}+\frac{1.3.5}{2.4.6}k^{6}\\[1.5ex]
9R_{4}=\frac{1.3.5.7}{2.4.6.8}+\frac{1.3.5}{2.4.6}.\frac{1}{2}k^{2}+\frac{1.3}{2.4}.\frac{1.3}{2.4}k^{4}+\frac{1}{2}.\frac{1.3.5}{2.4.6}k^{6}+\frac{1.3.5.7}{2.4.6.8}k^{8}\\[1.5ex]
\text{cet.}\qquad\qquad\text{cet.}
\end{gathered}
\]

\origpage{128}
or also:
\[
\begin{gathered}
3R_{1}=\frac{1}{2}(1+k^{2})\\[1.5ex]
5.R_{2}=\frac{1.3}{2.4}(1+k^{2})^{2}-\frac{1}{2}k^{2}\\[1.5ex]
7R_{3}=\frac{1.3.5}{2.4.6}(1+k^{2})^{3}-\frac{1.3}{2.2}k^{2}(1+k^{2})\\[1.5ex]
9R_{4}=\frac{1.3.5.7}{2.4.6.8}(1+k^{2})^{4}-\frac{1.3.5}{2.4.2}k^{2}(1+k^{2})^{2}+\frac{1.3}{2.4}k^{4}\\[1.5ex]
11R_{5}=\frac{1.3.5.7.9}{2.4.6.8.10}(1+k^{2})^{5}-\frac{1.3.5.7}{2.4.6.2}k^{2}(1+k^{2})^{3}+\frac{1.3.5}{2.2.4}k^{4}(1+k^{2})\\[1.5ex]
13R_{6}=\frac{1.3.\,.\,11}{2.4.\,.\,12}(1+k^{2})^{6}-\frac{1.3.5.7.9}{2.4.6.8.2}k^{2}(1+k^{2})^{4}+\frac{1.3.5.7}{2.4.2.4}k^{4}(1+k^{2})^{2}-\frac{1.3.5}{2.4.6}k^{6}\\[1.5ex]
\text{cet.}\qquad\qquad\text{cet.}
\end{gathered}
\]
or also:
\[
\begin{gathered}
3R_{1}=1-\frac{1}{2}.\,k'k'\\[1.5ex]
5R_{2}=1-\frac{1}{2}.2k'k'+\frac{1.3}{2.4}.\,k'^{4}\\[1.5ex]
7R_{3}=1-\frac{1}{2}.3k'k'+\frac{1.3}{2.4}.3k'^{4}-\frac{1.3.5}{2.4.6}.\,k'^{6}\\[1.5ex]
9R_{4}=1-\frac{1}{2}.4k'k'+\frac{1.3}{2.4}.6k'^{4}-\frac{1.3.5}{2.4.6}.4k'^{6}+\frac{1.3.5.7}{2.4.6.8}k'^{8}\\[1.5ex]
\text{cet.}\qquad\qquad\text{cet.}
\end{gathered}
\]
or finally:
\[
\begin{gathered}
3R_{1}=kk+\frac{1}{2}.\,k'k'\\[1.5ex]
5R_{2}=k^{4}+\frac{1}{2}.2k^{2}k'^{2}+\frac{1.3}{2.4}.\,k'^{4}\\[1.5ex]
7R_{3}=k^{6}+\frac{1}{2}.3k^{4}k'^{2}+\frac{1.3}{2.4}.3k^{2}k'^{4}+\frac{1.3.5}{2.4.6}.\,k'^{6}\\[1.5ex]
9R_{4}=k^{8}+\frac{1}{2}.4k^{6}k'^{2}+\frac{1.3}{2.4}.6k^{4}k'^{4}+\frac{1.3.5}{2.4.6}.4k^{2}k'^{6}+\frac{1.3.5.7}{2.4.6.7}k'^{8}\\[1.5ex]
\text{cet.}\qquad\qquad\text{cet.}
\end{gathered}
\]\ednote{In the last term of the final formula for $9R_{4}$ the denominator is printed $2.4.6.7$ where $2.4.6.8$ is expected (the author's Corrigenda: for $2.4.6.7$ read $2.4.6.8$); reproduced as printed.}

\origpage{129}
Of these four ways of expressing the quantities $R_{m}$, the second way affords a representation of them that is quite remarkable and neat, provided that the quantity is introduced:
\[
r=\frac{1+kk}{2k}.
\]
Thus, e.g., it comes about that
\[
\begin{gathered}
\frac{13R_{6}}{k^{6}}=\\[1.5ex]
\frac{1.3.\,.\,11}{1.2.\,.\,6}r^{6}-\frac{1.3.5.7.9}{1.2.3.4.2}r^{4}+\frac{1.3.5.7}{1.2.2.4}r^{2}-\frac{1.3.5}{2.4.6},
\end{gathered}
\]
and on integrating this expression six times with respect to $r$, we obtain:
\[
\begin{gathered}
13\int^{6}\frac{R_{6}\,dr^{6}}{k^{6}}=\\[1.5ex]
\frac{r^{12}}{2.4.\,.\,.\,12}-\frac{r^{10}}{2.4.6.8.10.2}+\frac{r^{8}}{2.4.6.8.2.4}-\frac{r^{6}}{2.4.6.2.4.6}+C'r^{4}-C''r^{2}+C''',
\end{gathered}
\]
where $C'$, $C''$, $C'''$ denote Arbitrary Constants. When these have been conveniently determined, there results:
\[
13\int^{6}\frac{R_{6}\,dr^{6}}{k^{6}}=\frac{(rr-1)^{6}}{2^{6}.\Pi 6},
\]
whence conversely:
\[
13R_{6}=\frac{k^{6}d^{6}(rr-1)^{6}}{2^{6}.\Pi 6.dr^{6}};
\]
and in the same way there is obtained generally:
\[
\text{12)}\quad(2m+1)R_{m}=\frac{k^{m}d^{m}(rr-1)^{m}}{2^{m}\Pi m.dr^{m}}.
\]
Let the short Commentary (\emph{Crelle Iournal V.\ II.\ p.\ 223}) be compared, entitled:

„Ueber eine besondere Gattung algebraischer Functionen, die aus der Entwicklung der Function $(1-2xz+z^{2})^{-\frac{1}{2}}$ entstehn."

The quantities $R_{m}$ having been found, by known Algorithms one arrives at deriving the quantities $R_{m}^{(n)}$, $S_{m}^{(n)}$, those, namely, such that:
\[
\left\{1+R_{1}x+R_{2}x^{2}+R_{3}x^{3}+\,.\,.\right\}^{n}=1+R_{1}^{(n)}x+R_{2}^{(n)}x^{2}+R_{3}^{(n)}x^{3}+\,.\,.,
\]

\origpage{130}
and further, where there is put:
\[
y=x\left\{1+R_{1}x+R_{2}x^{2}+R_{3}x^{3}+\,.\,.\right\},
\]
let there be made:
\[
x^{n}=y^{n}\left\{1+S_{1}^{(n)}y+S_{2}^{(n)}y^{2}+S_{3}^{(n)}y^{3}+\,.\,.\right\};
\]
which agree with the definition of the quantities $R_{m}^{(n)}$, $S_{m}^{(n)}$ proposed above. But it comes about, on putting:
\[
\phi(x)=1+R_{1}x+R_{2}x^{2}+R_{3}x^{3}+\ldots,
\]
from the theorems found by the Cll.\ \emph{Maclaurin} and \emph{Lagrange}:
\[
\begin{gathered}
R_{m}^{(n)}=\frac{d^{m}.\left\{\phi x\right\}^{n}}{\Pi m.dx^{m}}\\[1.5ex]
S_{m}^{(n)}=\frac{n}{m+n}.\frac{d^{m}\left\{\phi x\right\}^{-(m+n)}}{\Pi m.dx^{m}};
\end{gathered}
\]
provided that, the differentiations having been carried out, $x=0$ is put.

\subsection*{46.}
By the benefit of the formulae 6), 7), 8), 9), §. 45 we acquire the general expansions:
\[
\begin{gathered}
\text{1)}\quad\frac{\left(\dfrac{2kK}{\pi}\right)^{2n+1}\sin^{2n+1}\operatorname{am}\dfrac{2Kx}{\pi}}{2n+1}=\\[1.5ex]
4\left\{R_{n}^{(1)}\left(\frac{2K}{\pi}\right)^{2n}-\frac{R_{n-1}^{(3)}}{\Pi 3}\left(\frac{2K}{\pi}\right)^{2n-2}+\frac{R_{n-2}^{(5)}}{\Pi 5}\left(\frac{2K}{\pi}\right)^{2n-4}-\,.\,.+\frac{(-1)^{n}}{\Pi(2n+1)}\right\}\frac{\sqrt{q}\sin x}{1-q}+\\[1.5ex]
4\left\{R_{n}^{(1)}\left(\frac{2K}{\pi}\right)^{2n}-\frac{3^{2}.R_{n-1}^{(3)}}{\Pi 3}\left(\frac{2K}{\pi}\right)^{2n-2}+\frac{3^{4}R_{n-2}^{(5)}}{\Pi 5}\left(\frac{2K}{\pi}\right)^{2n-4}-\,.\,.+\frac{(-1)^{n}3^{2n}}{\Pi(2n+1)}\right\}\frac{\sqrt{q^{3}}\sin 3x}{1-q^{3}}+\\[1.5ex]
4\left\{R_{n}^{(1)}\left(\frac{2K}{\pi}\right)^{2n}-\frac{5^{2}R_{n-1}^{(3)}}{\Pi 3}\left(\frac{2K}{\pi}\right)^{2n-2}+\frac{5^{4}R_{n-2}^{(5)}}{\Pi 5}\left(\frac{2K}{\pi}\right)^{2n-4}-\,.\,.+\frac{(-1)^{n}5^{2n}}{\Pi(2n+1)}\right\}\frac{\sqrt{q^{5}}\sin 5x}{1-q^{5}}+\\[1.5ex]
\vdots
\end{gathered}
\]
\[
\begin{gathered}
\text{2)}\quad\frac{\left(\dfrac{2kK}{\pi}\right)^{2n}\sin^{2n}\operatorname{am}\dfrac{2Kx}{\pi}}{2n}=\\[1.5ex]
\frac{R_{n-1}^{(2)}}{\Pi 2}\left(\frac{2K}{\pi}\right)^{2n-1}\left(\frac{2K}{\pi}-\frac{2E^{\mathrm{I}}}{\pi}\right)-k^{2}\left(\frac{2K}{\pi}\right)^{2n}\left\{\frac{1}{3.4}+\frac{S_{1}^{(2)}}{5.6}+\frac{S_{2}^{(2)}}{7.8}+\,.\,.+\frac{S_{n-2}^{(2)}}{(2n-1)2n}\right\}
\end{gathered}
\]

\origpage{131}
\[
\begin{gathered}
-4\left\{\frac{2R_{n-1}^{(2)}\left(\dfrac{2K}{\pi}\right)^{2n-2}}{\Pi 2}-\frac{2^{3}R_{n-2}^{(4)}\left(\dfrac{2K}{\pi}\right)^{2n-4}}{\Pi 4}+\,.\,.+\frac{(-1)^{n}2^{n-1}}{\Pi 2n}\right\}\frac{q\cos 2x}{1-q^{2}}\\[1.5ex]
-4\left\{\frac{4R_{n-1}^{(2)}\left(\dfrac{2K}{\pi}\right)^{2n-2}}{\Pi 2}-\frac{4^{3}R_{n-2}^{(4)}\left(\dfrac{2K}{\pi}\right)^{2n-4}}{\Pi 4}+\,.\,.+\frac{(-1)^{n}4^{n-1}}{\Pi 2n}\right\}\frac{q^{2}\cos 4x}{1-q^{4}}\\[1.5ex]
-4\left\{\frac{6R_{n-1}^{(2)}\left(\dfrac{2K}{\pi}\right)^{2n-2}}{\Pi 2}-\frac{6^{3}R_{n-2}^{(4)}\left(\dfrac{2K}{\pi}\right)^{2n-4}}{\Pi 4}+\,.\,.+\frac{(-1)^{n}6^{n-1}}{\Pi 2n}\right\}\frac{q^{3}\cos 6x}{1-q^{6}}\\[1.5ex]
\vdots
\end{gathered}
\]
\ednote{In the numerators of the last terms of the three rows the printed exponents are $(-1)^{n}$ and $2^{n-1}$, $4^{n-1}$, $6^{n-1}$; a hand correction in ink in the margin above gives $(-1)^{n-1}$ and $2^{2n-1}$, $4^{2n-1}$, $6^{2n-1}$, as in the corresponding rows of formula 4) and on p. 132, and the author's Corrigenda (p. 189) list the same correction. The printed form is reproduced; the ink is not transcribed.}
\[
\begin{gathered}
\text{3)}\quad\frac{\left(\dfrac{2K}{\pi}\right)^{2n+1}}{(2n+1)\sin^{2n+1}\operatorname{am}\dfrac{2Kx}{\pi}}=\\[1.5ex]
\frac{R_{n}^{(1)}\left(\dfrac{2K}{\pi}\right)^{2n}}{\sin x}+\frac{R_{n-1}^{(3)}\left(\dfrac{2K}{\pi}\right)^{2n-2}}{\Pi 3.dx^{2}}d^{2}.\frac{1}{\sin x}+\,.\,.+\frac{d^{2n-2}.\dfrac{1}{\sin x}}{\Pi 2n.dx^{2n-2}}+\\[1.5ex]
4\left\{R_{n}^{(1)}\left(\frac{2K}{\pi}\right)^{2n}-\frac{R_{n-1}^{(3)}\left(\dfrac{2K}{\pi}\right)^{2n-2}}{\Pi 3}+\,.\,.+\frac{(-1)^{n}}{\Pi(2n+1)}\right\}\frac{q\sin x}{1-q}+\\[1.5ex]
4\left\{R_{n}^{(1)}\left(\frac{2K}{\pi}\right)^{2n}-\frac{3^{2}R_{n-1}^{(3)}\left(\dfrac{2K}{\pi}\right)^{2n-2}}{\Pi 3}+\,.\,.+\frac{(-1)^{n}3^{2n}}{\Pi(2n+1)}\right\}\frac{q^{3}\sin 3x}{1-q^{3}}+\\[1.5ex]
4\left\{R_{n}^{(1)}\left(\frac{2K}{\pi}\right)^{2n}-\frac{5^{2}R_{n-1}^{(3)}\left(\dfrac{2K}{\pi}\right)^{2n-2}}{\Pi 3}+\,.\,.+\frac{(-1)^{n}5^{2n}}{\Pi(2n+1)}\right\}\frac{q^{5}\sin 5x}{1-q^{5}}+\\[1.5ex]
\vdots
\end{gathered}
\]
\[
\begin{gathered}
\text{4)}\quad\frac{\left(\dfrac{2K}{\pi}\right)^{2n}}{2n.\sin^{2n}\operatorname{am}\dfrac{2Kx}{\pi}}=\\[1.5ex]
\frac{1}{2}R_{n-1}^{(2)}\left(\frac{2K}{\pi}\right)^{2n-1}\left(\frac{2K}{\pi}-\frac{2E^{\mathrm{I}}}{\pi}\right)-k^{2}\left(\frac{2K}{\pi}\right)^{2n}\left\{\frac{1}{3.4}+\frac{S_{1}^{(2)}}{5.6}+\frac{S_{2}^{(2)}}{7.8}+\,.\,.+\frac{S_{n-2}^{(2)}}{(2n-1)2n}\right\}+\\[1.5ex]
\frac{R_{n-1}^{(2)}\left(\dfrac{2K}{\pi}\right)^{2n-2}}{\Pi 2.\sin^{2}x}+\frac{R_{n-2}^{(4)}\left(\dfrac{2K}{\pi}\right)^{2n-4}d^{2}.\dfrac{1}{\sin^{2}x}}{\Pi 4.dx^{2}}+\,.\,.+\frac{d^{2n-2}.\dfrac{1}{\sin^{2}x}}{\Pi 2n,dx^{2n-2}}\\[1.5ex]
-4\left\{\frac{2R_{n-1}^{(2)}\left(\dfrac{2K}{\pi}\right)^{2n-2}}{\Pi 2}-\frac{2^{3}R_{n-2}^{(4)}\left(\dfrac{2K}{\pi}\right)^{2n-4}}{\Pi 4}+\,.\,.+\frac{(-1)^{n-1}2^{2n-1}}{\Pi(2n)}\right\}\frac{q^{2}\cos 2x}{1-q^{2}}
\end{gathered}
\]

\origpage{132}
\[
\begin{gathered}
-4\left\{\frac{4R_{n-1}^{(2)}\left(\dfrac{2K}{\pi}\right)^{2n-2}}{\Pi 2}-\frac{4^{3}R_{n-2}^{(4)}\left(\dfrac{2K}{\pi}\right)^{2n-4}}{\Pi 4}+\,.\,.+\frac{(-1)^{n-1}4^{2n-1}}{\Pi 2n}\right\}\frac{q^{4}\cos 4x}{1-q^{4}}\\[1.5ex]
-4\left\{\frac{6R_{n-1}^{(2)}\left(\dfrac{2K}{\pi}\right)^{2n-2}}{\Pi 2}-\frac{6^{3}R_{n-2}^{(4)}\left(\dfrac{2K}{\pi}\right)^{2n-4}}{\Pi 4}+\,.\,.+\frac{(-1)^{n-1}6^{2n-1}}{\Pi 2n}\right\}\frac{q^{6}\cos 6x}{1-q^{6}}\\[1.5ex]
\vdots
\end{gathered}
\]

From the formulae 6), 7), 8), 9) §. 45 others can be deduced, which with respect to the functions $\cos\operatorname{am}u$, $\operatorname{tang}\operatorname{am}u$, $\Delta\operatorname{am}u$ play the same parts as those do with respect to the function $\sin\operatorname{am}u$. For from the formula:
\[
\sin\operatorname{am}\left(k'u,\frac{ik}{k'}\right)=\cos\operatorname{coam}u,
\]
whence also:
\[
\sin\operatorname{am}\left(k'(K-u),\frac{ik}{k'}\right)=\cos\operatorname{am}u,
\]
we see that in the formulae proposed, where $\dfrac{ik}{k'}$ is put in place of $k$ and $k'(K-u)$ in place of $u$, $\sin\operatorname{am}u$ passes into $\cos\operatorname{am}u$, whence similar formulae are found, which correspond to $\cos\operatorname{am}u$ itself. Further, from the equation:
\[
\sin\operatorname{am}iu=i\operatorname{tang}\operatorname{am}(u,k')
\]
it is evident that at the same time $u$ can be changed into $iu$, $k$ into $k'$, $\sin\operatorname{am}u$ into $\operatorname{tang}\operatorname{am}u$; whence we derive formulae for $\operatorname{tang}\operatorname{am}u$. From these then, because
\[
\operatorname{cotang}\operatorname{am}(u+iK')=i\Delta\operatorname{am}(u),
\]
it is permitted to derive formulae for $\Delta\operatorname{am}u$, which correspond to the formulae 6) - 9) §. 45. These having been found, by a quite similar method, from the expansions of the functions:
\[
\begin{gathered}
\cos\operatorname{am}\frac{2Kx}{\pi},\quad\cos^{2}\operatorname{am}\frac{2Kx}{\pi},\quad\Delta\operatorname{am}\frac{2Kx}{\pi},\quad\Delta^{2}\operatorname{am}\frac{2Kx}{\pi}\\[1.5ex]
\frac{1}{\cos\operatorname{am}\dfrac{2Kx}{\pi}},\quad\frac{1}{\cos^{2}\operatorname{am}\dfrac{2Kx}{\pi}},\quad\frac{1}{\Delta\operatorname{am}\dfrac{2Kx}{\pi}},\quad\frac{1}{\Delta^{2}\operatorname{am}\dfrac{2Kx}{\pi}},
\end{gathered}
\]
proposed by us, you deduce the general expansions of the functions:
\[
\cos^{n}\operatorname{am}\frac{2Kx}{\pi},\quad\Delta^{n}\operatorname{am}\frac{2Kx}{\pi}.
\]
Let it suffice to have indicated these.

\origpage{133}

Notable transformations of the series into which we expand the Elliptic Functions we acquire, on putting $i\,x$ in place of $x$ and employing the formulae which we gave in the first Fundamenta on the reduction of an imaginary argument to a real argument. But since these are ready to hand, we do not wish to dwell longer on them in this place.

\subsection*{THE SECOND SPECIES OF ELLIPTIC INTEGRALS IS EXPANDED INTO SERIES.}

\subsection*{47.}

The formula exhibited above §.\ 41.\ 1) having been integrated:

\[
\left(\frac{2kK}{\pi}\right)^{2}\sin^{2}\operatorname{am}\frac{2Kx}{\pi}=\frac{2K}{\pi}\frac{2K}{\pi}-\frac{2K}{\pi}\frac{2E^{\mathrm{I}}}{\pi}-4\left\{\frac{2q\cos2x}{1-q^{2}}+\frac{4q^{2}\cos4x}{1-q^{4}}+\frac{6q^{3}\cos6x}{1-q^{6}}+.\,.\right\},
\]
thence from $x=0$ up to $x=x$, there results:

\[
\left(\frac{2kK}{\pi}\right)^{2}\int_{0}^{x}\sin^{2}\operatorname{am}\frac{2Kx}{\pi}.dx=
\]
\[
\left\{\frac{2K}{\pi}\frac{2K}{\pi}-\frac{2K}{\pi}\frac{2E^{\mathrm{I}}}{\pi}\right\}x-4\left\{\frac{q\sin2x}{1-q^{2}}+\frac{q^{2}\sin4x}{1-q^{4}}+\frac{q^{3}\sin6x}{1-q^{6}}+\frac{q^{4}\sin8x}{1-q^{8}}+\ldots\right\}.
\]
In what follows let us designate by the character: $\dfrac{2K}{\pi}.Z\left(\dfrac{2Kx}{\pi}\right)$ the expression:

\[
\begin{gathered}
\text{1)}\quad\frac{2K}{\pi}.Z\left(\frac{2Kx}{\pi}\right)=\frac{2Kx}{\pi}\left(\frac{2K}{\pi}-\frac{2E^{\mathrm{I}}}{\pi}\right)-\left(\frac{2kK}{\pi}\right)^{2}\int_{0}^{x}\sin^{2}\operatorname{am}\frac{2Kx}{\pi}.dx=\\[1ex]
4\left\{\frac{q\sin2x}{1-q^{2}}+\frac{q^{2}\sin4x}{1-q^{4}}+\frac{q^{3}\sin6x}{1-q^{6}}+\frac{q^{4}\sin8x}{1-q^{8}}+\ldots\right\}.
\end{gathered}
\]
$E$ the notation of the Cl$^{\mathrm{i}}$ Legendre it will be, on putting $\dfrac{2Kx}{\pi}=u$, $\phi=\operatorname{am}u$:

\[
\text{2)}\quad Z(u)=\frac{F^{\mathrm{I}}E(\phi)-E^{\mathrm{I}}F(\phi)}{K}.
\]
It is convenient to introduce the Function $Z(u)$ in place of $E(\phi)$ into the Analysis of Elliptic Functions; moreover it is ready to hand to recall it by means of formula 2) to the terms customary with the Cl$^{\mathrm{o}}$ Legendre. Let us sketch in a few words how from the expansion itself of the function

\origpage{134}
$Z$, which formula 1) affords, several of its properties, although known, may be derived.

Let $x$ be changed in 1) into $x+\dfrac{\pi}{2}$, there results:

\[
\frac{2K}{\pi}Z\left(\frac{2Kx}{\pi}+K\right)=-4\left\{\frac{q\sin2x}{1-q^{2}}-\frac{q^{2}\sin4x}{1-q^{4}}+\frac{q^{3}\sin6x}{1-q^{6}}+.\,.\right\},
\]
whence:

\[
\frac{2K}{\pi}Z\left(\frac{2Kx}{\pi}\right)-\frac{2K}{\pi}Z\left(\frac{2Kx}{\pi}+K\right)=8\left\{\frac{q\sin2x}{1-q^{2}}+\frac{q^{3}\sin6x}{1-q^{6}}+\frac{q^{5}\sin10x}{1-q^{10}}+.\,.\right\}.
\]
Further let $x$ be changed in 1) into $2x$, $q$ into $q^{2}$, and at the same time $k$ into $k^{(2)}$, $K$ into $K^{(2)}$, there results:

\[
\frac{2K^{(2)}}{\pi}Z\left(\frac{4K^{(2)}x}{\pi},k^{(2)}\right)=4\left\{\frac{q^{2}\sin4x}{1-q^{4}}+\frac{q^{4}\sin8x}{1-q^{8}}+\frac{q^{6}\sin12x}{1-q^{12}}+.\,.\right\},
\]
whence:

\[
\frac{2K}{\pi}Z\left(\frac{2Kx}{\pi}\right)-\frac{2K^{(2)}}{\pi}Z\left(\frac{4K^{(2)}x}{\pi},k^{(2)}\right)=4\left\{\frac{q\sin2x}{1-q^{2}}+\frac{q^{3}\sin6x}{1-q^{6}}+\frac{q^{5}\sin10x}{1-q^{10}}+.\,.\right\}.
\]
But above we found:

\[
\frac{2kK}{\pi}\sin\operatorname{am}\frac{2Kx}{\pi}=4\left\{\frac{\sqrt{q}\sin x}{1-q}+\frac{\sqrt{q^{3}}\sin3x}{1-q^{3}}+\frac{\sqrt{q^{5}}\sin5x}{1-q^{5}}+.\,.\right\}
\]
whence, on changing $q$ into $q^{2}$, $x$ into $2x$:

\[
\frac{2k^{(2)}K^{(2)}}{\pi}\sin\operatorname{am}\left(\frac{4K^{(2)}x}{\pi},k^{(2)}\right)=4\left\{\frac{q\sin2x}{1-q^{2}}+\frac{q^{3}\sin6x}{1-q^{6}}+\frac{q^{5}\sin10x}{1-q^{10}}+.\,.\right\}.
\]
Hence it follows:

\[
\begin{gathered}
\text{3)}\quad\frac{2K}{\pi}\left\{Z\left(\frac{2Kx}{\pi}\right)-Z\left(\frac{2Kx}{\pi}+K\right)\right\}=\frac{4k^{(2)}K^{(2)}}{\pi}\sin\operatorname{am}\left(\frac{4K^{(2)}x}{\pi},k^{(2)}\right)\\[1.5ex]
\text{4)}\quad\frac{2K}{\pi}Z\left(\frac{2Kx}{\pi}\right)-\frac{2K^{(2)}}{\pi}Z\left(\frac{4K^{(2)}x}{\pi},k^{(2)}\right)=\frac{2k^{(2)}K^{(2)}}{\pi}\sin\operatorname{am}\left(\frac{4K^{(2)}x}{\pi},k^{(2)}\right)\\[1.5ex]
\text{5)}\quad\frac{2K}{\pi}Z\left(\frac{2Kx}{\pi}\right)+\frac{2K}{\pi}Z\left(\frac{2Kx}{\pi}+K\right)-\frac{4K^{(2)}}{\pi}Z\left(\frac{4K^{(2)}x}{\pi},k^{(2)}\right)=0.
\end{gathered}
\]
In which formulae, of which 4), 5) afford the transformation of the function $Z$ of the second order, it is:

\[
k^{(2)}=\frac{1-k'}{1+k'},\ K^{(2)}=\frac{1+k'}{2}.K,\ \sin\operatorname{am}\left(\frac{4K^{(2)}x}{\pi},k^{(2)}\right)=(1+k')\sin\operatorname{am}\frac{2Kx}{\pi}.\sin\operatorname{coam}\frac{2Kx}{\pi},
\]

\origpage{135}
as is established concerning the transformation of the second order, proposed by the Cl.\ Legendre. Whence formula 3) may also be represented thus, on putting $u=\dfrac{2Kx}{\pi}$:

\[
\text{6)}\quad Z(u)-Z(u+K)=k^{2}\sin\operatorname{am}u.\sin\operatorname{coam}u.
\]
Let us put for the sake of brevity: $\operatorname{am}\left(\dfrac{2mK^{(m)}x}{\pi},k^{(m)}\right)=\phi^{(m)}$; from formula 4), on putting successively $k^{(2)}$, $k^{(4)}$, $k^{(8)}$, $.\,.$ in place of $k$; $2x$, $4x$, $8x$, $.\,.$ in place of $x$, there results:

\[
\text{7)}\quad K.Z(u)=F^{\mathrm{I}}E(\phi)-E^{\mathrm{I}}F(\phi)=k^{(2)}K^{(2)}\sin\phi^{(2)}+k^{(4)}K^{(4)}\sin\phi^{(4)}+k^{(8)}K^{(8)}\sin\phi^{(8)}+.\,.,
\]
which is the formula that the Cl.\ Legendre gave.

In a similar manner, from the formula of §$^{\mathrm{i}}$ 41:

\[
\frac{2K}{\pi}\frac{2K}{\pi}-\frac{2K}{\pi}\frac{2E^{\mathrm{I}}}{\pi}=8\left\{\frac{q}{(1-q)^{2}}+\frac{q^{3}}{(1-q^{3})^{2}}+\frac{q^{5}}{(1-q^{5})^{2}}+\frac{q^{7}}{(1-q^{7})^{2}}+.\,.\right\},
\]
which may also be expanded in this way:

\[
\frac{2K}{\pi}\frac{2K}{\pi}-\frac{2K}{\pi}\frac{2E^{\mathrm{I}}}{\pi}=8\left\{\frac{q}{1-q^{2}}+\frac{2q^{2}}{1-q^{4}}+\frac{3q^{3}}{1-q^{6}}+\frac{4q^{4}}{1-q^{8}}+.\,.\right\},
\]
when compared with this, which we found above:

\[
\left(\frac{2kK}{\pi}\right)^{2}=16\left\{\frac{q}{1-q^{2}}+\frac{3q^{3}}{1-q^{6}}+\frac{5q^{5}}{1-q^{10}}+\frac{7q^{7}}{1-q^{14}}+\ldots\right\},
\]
there results:

\[
\text{8)}\quad2K(K-E^{\mathrm{I}})=(kK)^{2}+2(k^{(2)}K^{(2)})^{2}+4(k^{(4)}K^{(4)})^{2}+8(k^{(8)}K^{(8)})^{2}+.\,.,
\]
which agrees with that which the Cl.\ \emph{Gauss} gave in the Commentary \emph{Determinatio Attractionis} etc.\ §.\ 17.

\subsection*{48.}

By the same method by which in §.\ 41 we derived the expansion of the expression $\left(\dfrac{2kK}{\pi}\right)^{2}\sin^{2}\operatorname{am}\dfrac{2Kx}{\pi}$, let us inquire into the expression $\left\{\dfrac{2K}{\pi}Z\left(\dfrac{2Kx}{\pi}\right)\right\}^{2}$ to be expanded in series. Let us put

\[
\begin{gathered}
\left(\frac{2K}{\pi}\right)^{2}Z\left(\frac{2Kx}{\pi}\right)Z\left(\frac{2Kx}{\pi}\right)=16\left\{\frac{q\sin2x}{1-q^{2}}+\frac{q^{2}\sin4x}{1-q^{4}}+\frac{q^{3}\sin6x}{1-q^{6}}+\ldots\right\}^{2}\\[1ex]
=8\left\{A+A'\cos2x+A''\cos4x+A'''\cos6x+\ldots\right\},
\end{gathered}
\]

\origpage{136}
which proposed expression we see assume the form, when in place of $2\sin2mx\sin2m'x$ there is put everywhere $\cos(m-m')x-\cos(m+m')x$. First there results:

\[
A=\frac{q^{2}}{(1-q^{2})^{2}}+\frac{q^{4}}{(1'-q^{4})^{2}}+\frac{q^{6}}{(1-q^{6})^{2}}+\frac{q^{8}}{(1-q^{8})^{2}}+\ldots.
\]\ednote{The print has a stray prime after the 1 in the second denominator, apparently a speck or misprint for $(1-q^{4})^{2}$.}
Then in general we obtain $A^{(n)}=2B^{(n)}-C^{(n)}$, provided one puts:

\[
\begin{gathered}
B^{(n)}=\frac{q^{n+2}}{(1-q^{2})(1-q^{2n+2})}+\frac{q^{n+4}}{(1-q^{4})(1-q^{2n+4})}+\frac{q^{n+6}}{(1-q^{6})(1-q^{2n+6})}+\ldots\\[1.5ex]
C^{(n)}=\frac{q^{n}}{(1-q^{2})(1-q^{2n-2})}+\frac{q^{n}}{(1-q^{4})(1-q^{2n-4})}+\ldots+\frac{q^{n}}{(1-q^{2n-2})(1-q^{2})}.
\end{gathered}
\]
In the single terms of these expressions let there be substituted respectively:

\[
\begin{gathered}
\frac{q^{n+m}}{(1-q^{m})(1-q^{2n+m})}=\frac{q^{n}}{1-q^{2n}}\left\{\frac{q^{m}}{1-q^{m}}-\frac{q^{2n+m}}{1-q^{2n+m}}\right\}\\[1.5ex]
\frac{q^{n}}{(1-q^{m})(1-q^{2n-m})}=\frac{q^{n}}{1-q^{2n}}\left\{\frac{q^{m}}{1-q^{m}}+\frac{q^{2n-m}}{1-q^{2n-m}}+1\right\},
\end{gathered}
\]
there results:

\[
\begin{gathered}
B^{(n)}=\frac{q^{n}}{1-q^{2n}}\left\{\frac{q^{2}}{1-q^{2}}+\frac{q^{4}}{1-q^{4}}+\frac{q^{6}}{1-q^{6}}+\ldots\right\}\\[1.5ex]
-\frac{q^{n}}{1-q^{2n}}\left\{\frac{q^{2n+2}}{1-q^{2n+2}}+\frac{q^{2n+4}}{1-q^{2n+4}}+\frac{q^{2n+6}}{1-q^{2n+6}}+\ldots\right\}\\[1.5ex]
=\frac{q^{n}}{1-q^{2n}}\left\{\frac{q^{2}}{1-q^{2}}+\frac{q^{4}}{1-q^{4}}+\frac{q^{6}}{1-q^{6}}+.\,.+\frac{q^{2n}}{1-q^{2n}}\right\}\\[1.5ex]
C^{(n)}=\frac{(n-1)q^{n}}{1-q^{2n}}+\frac{2q^{n}}{1-q^{2n}}\left\{\frac{q^{2}}{1-q^{2}}+\frac{q^{4}}{1-q^{4}}+\frac{q^{6}}{1-q^{6}}+.\,.+\frac{q^{2n-2}}{1-q^{2n-2}}\right\};
\end{gathered}
\]
whence:

\[
A^{(n)}=2B^{(n)}-C^{(n)}=-\frac{(n-1)q^{n}}{1-q^{2n}}+\frac{2q^{3n}}{(1-q^{2n})^{2}}=-\frac{nq^{n}}{1-q^{2n}}+\frac{q^{n}(1+q^{2n})}{(1-q^{2n})^{2}}.
\]
These being collected, the sought expansion is found:

\[
\begin{gathered}
\text{1)}\quad\left(\frac{2K}{\pi}\right)^{2}Z\left(\frac{2Kx}{\pi}\right)Z\left(\frac{2Kx}{\pi}\right)=8A-8\left\{\frac{q\cos2x}{1-q^{2}}+\frac{2q^{2}\cos4x}{1-q^{4}}+\frac{3q^{3}\cos6x}{1-q^{6}}+\ldots\right\}\\[1.5ex]
+8\left\{\frac{q(1+q^{2})\cos2x}{(1-q^{2})^{2}}+\frac{q^{2}(1+q^{4})\cos4x}{(1-q^{4})^{2}}+\frac{q^{3}(1+q^{6})\cos6x}{(1-q^{6})^{2}}+\ldots\right\}.
\end{gathered}
\]
Since the quantity $A=\dfrac{q^{2}}{(1-q^{2})^{2}}+\dfrac{q^{4}}{(1-q^{4})^{2}}+\dfrac{q^{6}}{(1-q^{6})^{2}}+\ldots$ itself may also be expanded in this manner:

\[
A=\frac{q^{2}}{1-q^{2}}+\frac{2q^{4}}{1-q^{4}}+\frac{3q^{6}}{1-q^{6}}+\frac{4q^{8}}{1-q^{8}}+\ldots,
\]

\origpage{137}
we find from §.\ 42.\ 6):

\[
\text{2)}\quad8A=\frac{(2-k^{2})\left(\dfrac{2K}{\pi}\right)^{2}-3\left(\dfrac{2K}{\pi}\right)\left(\dfrac{2E^{\mathrm{I}}}{\pi}\right)+1}{3}.
\]
But further it is established that:

\[
8A=\frac{2}{\pi}.\left(\frac{2K}{\pi}\right)^{2}\int_{0}^{\frac{\pi}{2}}Z\left(\frac{2Kx}{\pi}\right)Z\left(\frac{2Kx}{\pi}\right).dx;
\]
for equation 1) being integrated from $x=0$ up to $x=\dfrac{\pi}{2}$, all terms, except the first, vanish; whence, if it is pleasing to use the notations of Cl$^{\mathrm{i}}$ Legendre:

\[
\text{3)}\quad\int_{0}^{\frac{\pi}{2}}\frac{\left\{F^{\mathrm{I}}E(\phi)-E^{\mathrm{I}}F(\phi)\right\}^{2}}{\Delta(\phi)}.d\phi=\frac{(2-k^{2})F^{\mathrm{I}}F^{\mathrm{I}}F^{\mathrm{I}}-3F^{\mathrm{I}}F^{\mathrm{I}}E^{\mathrm{I}}+\dfrac{\pi\pi F^{\mathrm{I}}}{4}}{3},
\]
which is a determination of a quite abstruse definite Integral.

\subsection*{INDEFINITE ELLIPTIC INTEGRALS OF THE THIRD SPECIES ARE REDUCED TO THE DEFINITE CASE, IN WHICH THE AMPLITUDE EQUALS THE PARAMETER.}

\subsection*{49.}

Before we proceed to expand the third species of Elliptic Integrals into a series, we shall set forth separately a few things which it has fallen to us to add to their Theory, and this almost in the very signs used by their distinguished author. Soon the same will be put forward with new denominations employed.

We start from certain known theorems on the second species of Elliptic Integrals. There results:

\[
\begin{gathered}
\sin\operatorname{am}(u+a)+\sin\operatorname{am}(u-a)=\frac{2\sin\operatorname{am}u.\cos\operatorname{am}a.\Delta\operatorname{am}a}{1-k^{2}\sin^{2}\operatorname{am}a.\sin^{2}\operatorname{am}u}\\[1.5ex]
\sin\operatorname{am}(u+a)-\sin\operatorname{am}(u-a)=\frac{2\sin\operatorname{am}a.\cos\operatorname{am}u.\Delta\operatorname{am}u}{1-k^{2}\sin^{2}\operatorname{am}a.\sin^{2}\operatorname{am}u},
\end{gathered}
\]

\origpage{138}
whence:

\[
\sin^{2}\operatorname{am}(u+a)-\sin^{2}\operatorname{am}(u-a)=\frac{4\sin\operatorname{am}a.\cos\operatorname{am}a.\Delta\operatorname{am}a.\sin\operatorname{am}u.\cos\operatorname{am}u.\Delta\operatorname{am}u}{\left\{1-k^{2}\sin^{2}\operatorname{am}a.\sin^{2}\operatorname{am}u\right\}^{2}},
\]
which formula being integrated with respect to $u$, there results:

\[
\text{1)}\quad\int_{0}^{u}du.\left\{\sin^{2}\operatorname{am}(u+a)-\sin^{2}\operatorname{am}(u-a)\right\}=\frac{2\sin\operatorname{am}a.\cos\operatorname{am}a.\Delta\operatorname{am}a.\sin^{2}\operatorname{am}u}{1-k^{2}\sin^{2}\operatorname{am}a.\sin^{2}\operatorname{am}u},
\]
as we have already found above.

Let there be put: $\operatorname{am}u=\phi$, $\operatorname{am}a=\alpha$, $\operatorname{am}(u+a)=\sigma$, $\operatorname{am}(u-a)=\vartheta$; there will be, from the notation of Cl$^{\mathrm{i}}$ Legendre:

\[
k^{2}\int_{0}^{u}du.\sin^{2}\operatorname{am}u=F(\phi)-E(\phi),
\]
whence also, since $F(\sigma)-F(\alpha)=F(\phi)$, $F(\vartheta)+F(\alpha)=F(\phi)$:

\[
\begin{gathered}
k^{2}\int_{0}^{u}du.\sin^{2}\operatorname{am}(u+a)=F(\phi)-E(\sigma)+E(\alpha)\\[1.5ex]
k^{2}\int_{0}^{u}du.\sin^{2}\operatorname{am}(u-a)=F(\phi)-E(\vartheta)-E(\alpha).
\end{gathered}
\]
Hence equation 1) passes into this:

\[
\text{2)}\quad2E(\alpha)-\left\{E(\sigma)-E(\vartheta)\right\}=\frac{2k^{2}\sin\alpha\cos\alpha\Delta\alpha.\sin^{2}\phi}{1-k^{2}\sin^{2}\alpha.\sin^{2}\phi}.
\]

With $u$ and $a$ interchanged, $\alpha$ passes into $\phi$, $\vartheta$ into $-\vartheta$, $\sigma$ remains unchanged, whence from 1) there results:

\[
2E(\phi)-\left\{E(\sigma)+E(\vartheta)\right\}=\frac{2k^{2}\sin\phi\cos\phi\Delta\phi.\sin^{2}\alpha}{1-k^{2}\sin^{2}\alpha.\sin^{2}\phi},
\]
which being added to equation 1), there comes:

\[
\text{3)}\quad E(\phi)+E(\alpha)-E(\sigma)=k^{2}\sin\alpha.\sin\phi.\sin\sigma,
\]
which is the theorem on the Addition of the function $E$, put forward by Cl.\ Legendre, l.\ c.\ Chap.\ IX.\ p.\ 43.\ c$'$.

\origpage{139}
Integrals of the form:

\[
\int_{0}^{\phi}\frac{\sin^{2}\phi.d\phi}{\left\{1-k^{2}\sin^{2}\alpha.\sin^{2}\phi\right\}\Delta(\phi)}
\]
according to the distribution of Elliptic Integrals into species which Cl.\ Legendre established, constitute the \emph{third} species. The quantity $-k^{2}\sin^{2}\alpha$, which he designates by $n$, he calls the Parameter; we in what follows shall call the angle $\alpha$ itself the \emph{Parameter}. Of which Integrals, equation 2) being multiplied by

\[
\frac{d\phi}{\Delta(\phi)}=\frac{d\sigma}{\Delta(\sigma)}=\frac{d\vartheta}{\Delta(\vartheta)},
\]
and the integration carried out from $\phi=0$ up to $\phi=\phi$, whereupon the limits of $\sigma$ will be: $\sigma=\alpha$, $\sigma=\sigma$, the limits of $\vartheta$: $\vartheta=-\alpha$, $\vartheta=\vartheta$, we derive the following expression:

\[
\int_{0}^{\phi}\frac{2k^{2}\sin\alpha\cos\alpha\Delta\alpha.\sin^{2}\phi.d\phi}{\left\{1-k^{2}\sin^{2}\alpha.\sin^{2}\phi\right\}\Delta(\phi)}=2F(\phi)E(\alpha)-\int_{\alpha}^{\sigma}\frac{E(\sigma).d\sigma}{\Delta(\sigma)}+\int_{-\alpha}^{\vartheta}\frac{E(\vartheta).d\vartheta}{\Delta(\vartheta)}.
\]
It is easily established, since $E(-\phi)=-E(\phi)$, that:

\[
\int_{0}^{\phi}\frac{E(\phi).d\phi}{\Delta(\phi)}=\int_{0}^{-\phi}\frac{E(\phi).d\phi}{\Delta(\phi)},\quad\text{or}\quad\int_{-\phi}^{+\phi}\frac{E(\phi).d\phi}{\Delta(\phi)}=0,
\]
whence, since:

\[
\begin{gathered}
\int_{\alpha}^{\sigma}\frac{E(\sigma).d\sigma}{\Delta(\sigma)}=\int_{0}^{\sigma}\frac{E(\phi).d\phi}{\Delta(\phi)}-\int_{0}^{\alpha}\frac{E(\phi).d\phi}{\Delta(\phi)}\\[1.5ex]
\int_{-\alpha}^{\vartheta}\frac{E(\vartheta).d\vartheta}{\Delta(\vartheta)}=\int_{0}^{\vartheta}\frac{E(\phi).d\phi}{\Delta(\phi)}-\int_{0}^{-\alpha}\frac{E(\phi).d\phi}{\Delta(\phi)}=\int_{0}^{\vartheta}\frac{E(\phi).d\phi}{\Delta(\phi)}-\int_{0}^{\alpha}\frac{E(\phi).d\phi}{\Delta(\phi)},
\end{gathered}
\]
we have obtained the new and memorable

\subsection*{THEOREM I.}

\emph{Let the angles $\vartheta$, $\sigma$ be determined so that:}

\[
F(\phi)+F(\alpha)=F(\sigma);\ F(\phi)-F(\alpha)=F(\vartheta),
\]

\origpage{140}
\emph{there will be:}

\[
\begin{gathered}
\int_{0}^{\phi}\frac{k^{2}\sin\alpha\cos\alpha\Delta\alpha.\sin^{2}\phi.d\phi}{\left\{1-k^{2}\sin^{2}\alpha.\sin^{2}\phi\right\}\Delta(\phi)}=\\[1.5ex]
F(\phi)E(\alpha)-\frac{1}{2}\int_{\vartheta}^{\sigma}\frac{E(\phi).d\phi}{\Delta(\phi)}=F(\phi)E(\alpha)-\frac{1}{2}\int_{0}^{\sigma}\frac{E(\phi).d\phi}{\Delta(\phi)}+\frac{1}{2}\int_{0}^{\vartheta}\frac{E(\phi).d\phi}{\Delta(\phi)},
\end{gathered}
\]
\emph{so that the third species of Elliptic Integrals, which depends on three elements, the Modulus $k$, the Amplitude $\phi$, the Parameter $\alpha$, is reduced to the first and second species, and to a new Transcendent}

\[
\int_{0}^{\phi}\frac{E(\phi).d\phi}{\Delta(\phi)},
\]
\emph{all of which depend only on two elements.}

\subsection*{50.}

Let us put $F(\alpha_{2})=2F(\alpha)$; whenever $\phi=\alpha$, there results $\sigma=\alpha_{2}$, $\vartheta=0$, in which case therefore from the proposed theorem we obtain:

\[
\text{1)}\quad\int_{0}^{\alpha}\frac{k^{2}\sin\alpha\cos\alpha\Delta\alpha.\sin^{2}\phi.d\phi}{\left\{1-k^{2}\sin^{2}\alpha.\sin^{2}\phi\right\}\Delta(\phi)}=F(\alpha)E(\alpha)-\frac{1}{2}\int_{0}^{\alpha_{2}}\frac{E(\phi).d\phi}{\Delta(\phi)}.
\]
Which formula teaches that in place of the new Transcendent there can be substituted also this:

\[
\int_{0}^{\alpha}\frac{\sin^{2}\phi.d\phi}{\left\{1-k^{2}\sin^{2}\alpha.\sin^{2}\phi\right\}\Delta(\phi)},
\]
which is an Integral of the third species that is \emph{definite,} in which the Amplitude equals the Parameter, and which therefore itself also depends only on two elements, on the Modulus $k$ and on that quantity which is at once both Parameter and Amplitude.

Let us put $2F(\mu)=F(\phi)+F(\alpha)=F(\sigma)$, $2F(\partial)=F(\phi)-F(\alpha)=F(\vartheta)$; there will be, from 1):

\[
\frac{1}{2}\int_{0}^{\sigma}\frac{E(\phi).d\phi}{\Delta(\phi)}=F(\mu)E(\mu)-\int_{0}^{\mu}\frac{k^{2}\sin\mu\cos\mu\Delta\mu.\sin^{2}\phi.d\phi}{\left\{1-k^{2}\sin^{2}\mu.\sin^{2}\phi\right\}\Delta(\phi)}
\]

\origpage{141}
\[
\frac{1}{2}\int_{0}^{\vartheta}\frac{E(\phi).d\phi}{\Delta(\phi)}=F(\partial)E(\partial)-\int_{0}^{\partial}\frac{k^{2}\sin\partial\cos\partial\Delta\partial.\sin^{2}\phi.d\phi}{\left\{1-k^{2}\sin^{2}\partial.\sin^{2}\phi\right\}\Delta(\phi)},
\]
which formulae being substituted in the theorem put forward in the preceding §$^{\mathrm{o}}$, we obtain the following

\subsection*{THEOREM II.}

\emph{Let the angles $\mu$, $\partial$ be determined so that:}

\[
F(\mu)=\frac{F(\phi)+F(\alpha)}{2},\ F(\partial)=\frac{F(\phi)-F(\alpha)}{2},
\]
\emph{there will be:}

\[
\begin{gathered}
k^{2}\sin\alpha\cos\alpha\Delta\alpha.\int_{0}^{\phi}\frac{\sin^{2}\phi.d\phi}{\left\{1-k^{2}\sin^{2}\alpha.\sin^{2}\phi\right\}\Delta(\phi)}=\\[1.5ex]
F(\phi)E(\alpha)-F(\mu)E(\mu)+F(\partial)E(\partial)+\\[1.5ex]
k^{2}\sin\mu\cos\mu\Delta\mu.\int_{0}^{\mu}\frac{\sin^{2}\phi.d\phi}{\left\{1-k^{2}\sin^{2}\mu.\sin^{2}\phi\right\}\Delta(\phi)}\\[1.5ex]
-k^{2}\sin\partial\cos\partial\Delta\partial.\int_{0}^{\partial}\frac{\sin^{2}\phi.d\phi}{\left\{1-k^{2}\sin^{2}\partial.\sin^{2}\phi\right\}\Delta(\phi)},
\end{gathered}
\]
\emph{by which formula the indefinite Integrals of the third species are reduced to definite ones, in which the Amplitude equals the Parameter, and hence those which depended on three elements are reduced to other Transcendents, which consist of two only.}

With $\alpha$ and $\phi$ interchanged, $\vartheta$ passes into $-\vartheta$, $\sigma$ remains unchanged, whence since moreover:

\[
\int_{-\vartheta}^{\sigma}\frac{E(\phi).d\phi}{\Delta(\phi)}=\int_{+\vartheta}^{\sigma}\frac{E(\phi).d\phi}{\Delta(\phi)},
\]
from Theorem I:

\[
\int_{0}^{\phi}\frac{k^{2}\sin\alpha\cos\alpha\Delta\alpha.\sin^{2}\phi.d\phi}{\left\{1-k^{2}\sin^{2}\alpha.\sin^{2}\phi\right\}\Delta(\phi)}=F(\phi)E(\alpha)-\frac{1}{2}\int_{\vartheta}^{\sigma}\frac{E(\phi).d\phi}{\Delta(\phi)}
\]

\origpage{142}
we obtain:

\[
\int_{0}^{\alpha}\frac{k^{2}\sin\phi\cos\phi\Delta\phi.\sin^{2}\alpha.d\alpha}{\left\{1-k^{2}\sin^{2}\phi.\sin^{2}\alpha\right\}\Delta(\alpha)}=F(\alpha)E(\phi)-\frac{1}{2}\int_{\vartheta}^{\sigma}\frac{E(\phi).d\phi}{\Delta(\phi)}.
\]
Hence, subtraction being made, there results:

\[
\begin{gathered}
\text{2)}\quad\int_{0}^{\phi}\frac{k^{2}\sin\alpha\cos\alpha\Delta\alpha.\sin^{2}\phi.d\phi}{\left\{1-k^{2}\sin^{2}\alpha.\sin^{2}\phi\right\}\Delta(\phi)}-\int_{0}^{\alpha}\frac{k^{2}\sin\phi\cos\phi\Delta\phi.\sin^{2}\alpha.d\alpha}{\left\{1-k^{2}\sin^{2}\phi.\sin^{2}\alpha\right\}\Delta(\alpha)}=\\[1.5ex]
F(\phi)E(\alpha)-F(\alpha)E(\phi);
\end{gathered}
\]
which formula teaches that \emph{an Integral of the third species can always be reduced to another, in which what was the Parameter becomes the Amplitude, and what was the Amplitude becomes the Parameter.}

Where in formula 2) one puts $\phi=\dfrac{\pi}{2}$, we obtain:

\[
\text{3)}\quad\int_{0}^{\frac{\pi}{2}}\frac{k^{2}\sin\alpha\cos\alpha\Delta\alpha.\sin^{2}\phi.d\phi}{\left\{1-k^{2}\sin^{2}\alpha.\sin^{2}\phi\right\}\Delta(\phi)}=F^{\mathrm{I}}E(\alpha)-E^{\mathrm{I}}F(\alpha).
\]
Formulae 2), 3) agree with those which have been exhibited by Cl.\ Legendre Chap.\ XXIII.\ p.\ 141.\ (n$'$), (p$'$).

\subsection*{ELLIPTIC INTEGRALS OF THE THIRD SPECIES ARE EXPANDED INTO A SERIES. HOW THEY ARE CONVENIENTLY EXPRESSED BY THE NEW TRANSCENDENT $\Theta$.}

\subsection*{51.}

From the formula:

\[
\begin{gathered}
\sin^{2}\operatorname{am}\frac{2K}{\pi}(x+A)-\sin^{2}\operatorname{am}\frac{2K}{\pi}(x-A)=\\[1.5ex]
\frac{4\sin\operatorname{am}\dfrac{2KA}{\pi}\cos\operatorname{am}\dfrac{2KA}{\pi}\Delta\operatorname{am}\dfrac{2KA}{\pi}.\sin\operatorname{am}\dfrac{2Kx}{\pi}\cos\operatorname{am}\dfrac{2Kx}{\pi}\Delta\operatorname{am}\dfrac{2Kx}{\pi}}{\left\{1-k^{2}\sin^{2}\operatorname{am}\dfrac{2KA}{\pi}.\sin^{2}\operatorname{am}\dfrac{2Kx}{\pi}\right\}^{2}},
\end{gathered}
\]

\origpage{143}
which is evident from the elements, we derive by integrating:

\[
\begin{gathered}
\text{1)}\quad\frac{2K}{\pi}\int_{0}^{x}dx.\left\{\sin^{2}\operatorname{am}\frac{2K}{\pi}(x+A)-\sin^{2}\operatorname{am}\frac{2K}{\pi}(x-A)\right\}=\\[1.5ex]
\frac{2\sin\operatorname{am}\dfrac{2KA}{\pi}\cos\operatorname{am}\dfrac{2KA}{\pi}\Delta\operatorname{am}\dfrac{2KA}{\pi}.\sin^{2}\operatorname{am}\dfrac{2Kx}{\pi}}{1-k^{2}\sin^{2}\operatorname{am}\dfrac{2KA}{\pi}.\sin^{2}\operatorname{am}\dfrac{2Kx}{\pi}}.
\end{gathered}
\]
We have already given in §.\ 41 the formula:

\[
\left(\frac{2kK}{\pi}\right)^{2}\sin^{2}\operatorname{am}\frac{2Kx}{\pi}=\frac{2K}{\pi}\frac{2K}{\pi}-\frac{2K}{\pi}\frac{2E^{\mathrm{I}}}{\pi}-4\left\{\frac{2q\cos2x}{1-q^{2}}+\frac{4q^{2}\cos4x}{1-q^{4}}+\frac{6q^{3}\cos6x}{1-q^{6}}+.\,.\right\},
\]
whence:

\[
\begin{gathered}
\left(\frac{2kK}{\pi}\right)^{2}\left\{\sin^{2}\operatorname{am}\frac{2K}{\pi}(x+A-\sin^{2}\operatorname{am}\frac{2K}{\pi}(x-A)\right\}=\\[1.5ex]
4\left\{\frac{2q\cos2(x-A)}{1-q^{2}}+\frac{4q^{2}\cos4(x-A)}{1-q^{4}}+\frac{6q^{3}\cos6(x-A)}{1-q^{6}}+.\,.\right\}\\[1.5ex]
-4\left\{\frac{2q\cos2(x+A)}{1-q^{2}}+\frac{4q^{2}\cos4(x+A)}{1-q^{4}}+\frac{6q^{3}\cos6(x+A)}{1-q^{6}}+.\,.\right\}=\\[1.5ex]
8\left\{\frac{2q\sin2A\sin2x}{1-q^{2}}+\frac{4q^{2}\sin4A\sin4x}{1-q^{4}}+\frac{6q^{3}\sin6A\sin6x}{1-q^{6}}+.\,.\right\}.
\end{gathered}
\]
\ednote{The closing parenthesis after $x+A$ is missing in the print (printed as $(x+A-\sin^{2}\dots$); an apparent misprint for $(x+A)-\sin^{2}\dots$.}
Hence from 1) there results:

\[
\begin{gathered}
\text{2)}\quad\frac{2K}{\pi}.\frac{2k^{2}\sin\operatorname{am}\dfrac{2KA}{\pi}\cos\operatorname{am}\dfrac{2KA}{\pi}\Delta\operatorname{am}\dfrac{2KA}{\pi}.\sin^{2}\operatorname{am}\dfrac{2Kx}{\pi}}{1-k^{2}\sin^{2}\operatorname{am}\dfrac{2KA}{\pi}.\sin^{2}\operatorname{am}\dfrac{2Kx}{\pi}}=\\[1.5ex]
\text{Const.}+4\left\{\frac{q\sin2(x-A)}{1-q^{2}}+\frac{q^{2}\sin4(x-A)}{1-q^{4}}+\frac{q^{3}\sin6(x-A)}{1-q^{6}}+.\,.\right\}\\[1.5ex]
-4\left\{\frac{q\sin2(x+A)}{1-q^{2}}+\frac{q^{2}\sin4(x+A)}{1-q^{4}}+\frac{q^{3}\sin6(x+A)}{1-q^{6}}+.\,.\right\}=\\[1.5ex]
\text{Const.}-8\left\{\frac{q\sin2A\cos2x}{1-q^{2}}+\frac{q^{2}\sin4A\cos4x}{1-q^{4}}+\frac{q^{3}\sin6A\cos6x}{1-q^{6}}+.\,.\right\},
\end{gathered}
\]
where the \emph{Constant} must be so determined that the proposed expression vanishes for $x=0$, whence from §.\ 47.\ 1):

\[
\text{Const.}=8\left\{\frac{q\sin2A}{1-q^{2}}+\frac{q^{2}\sin4A}{1-q^{4}}+\frac{q^{3}\sin6A}{1-q^{6}}+.\,.\right\}=2.\frac{2K}{\pi}Z\left(\frac{2KA}{\pi}\right).
\]

\origpage{144}
Formula 2) being integrated from $x=0$ up to $x=\dfrac{\pi}{2}$, since there comes out $\dfrac{\pi}{2}.\text{Const.}$, the remaining terms vanishing, with $\dfrac{2KA}{\pi}=a$, $\dfrac{2Kx}{\pi}=u$ put, we derive the definite integral:

\[
\int_{0}^{\frac{\pi}{2}}\frac{k^{2}\sin\operatorname{am}a\cos\operatorname{am}a\Delta\operatorname{am}a.\sin^{2}\operatorname{am}u.du}{1-k^{2}\sin^{2}\operatorname{am}a.\sin^{2}\operatorname{am}u}=K.Z(a),
\]
which is the same as 3) §.\ 50.

In what follows we shall designate by the character $\Pi(u,a,k)$ or more briefly by $\Pi(u,a)$ the integral: $\Pi(u,a){}^{*)}=$

\[
\int_{0}^{u}\frac{k^{2}\sin\operatorname{am}a\cos\operatorname{am}a\Delta\operatorname{am}a.\sin^{2}\operatorname{am}u.du}{1-k^{2}\sin^{2}\operatorname{am}u.\sin^{2}\phi}=\int_{0}^{\phi}\frac{k^{2}\sin\alpha\cos\alpha\Delta\alpha.\sin^{2}\phi.d\phi}{\left\{1-k^{2}\sin^{2}\alpha.\sin^{2}\phi\right\}\Delta(\phi)},
\]
\ednote{The denominator of the first integrand is printed $1-k^{2}\sin^{2}\operatorname{am}u.\sin^{2}\phi$, where $1-k^{2}\sin^{2}\operatorname{am}a.\sin^{2}\operatorname{am}u$ is expected (Corrigenda, p. 144, line 7); the copy carries a faint hand correction in ink, which is not transcribed.}
provided $\phi=\operatorname{am}u$, $\alpha=\operatorname{am}a$. These being posited, equation 2) being integrated again from $x=0$ up to $x=x$, there results:

\[
\begin{gathered}
\text{3)}\quad\Pi\left(\frac{2Kx}{\pi},\frac{2KA}{\pi}\right)=\\[1.5ex]
\frac{2Kx}{\pi}Z\left(\frac{2KA}{\pi}\right)-\left\{\frac{q\cos2(x-A)}{1-q^{2}}+\frac{q^{2}\cos4(x-A)}{2(1-q^{4})}+\frac{q^{3}\cos6(x-A)}{3(1-q^{6})}+.\,.\right\}\\[1.5ex]
+\frac{q\cos2(x+A)}{1-q^{2}}+\frac{q^{2}\cos4(x+A)}{2(1-q^{4})}+\frac{q^{3}\cos6(x+A)}{3(1-q^{6})}+.\,.=\\[1.5ex]
\frac{2Kx}{\pi}Z\left(\frac{2KA}{\pi}\right)-2\left\{\frac{q\sin2A\sin2x}{1-q^{2}}+\frac{q^{2}\sin4A\sin4x}{2(1-q^{4})}+\frac{q^{3}\sin6A\sin6x}{3(1-q^{6})}+.\,.\right\},
\end{gathered}
\]
which is the sought expansion of the Elliptic Integral of the third species.

Where we note the known expansion:

\[
-\log(1-2q\cos2x+q^{2})=2\left\{q\cos2x+\frac{q^{2}\cos4x}{2}+\frac{q^{3}\cos6x}{3}+\frac{q^{4}\cos8x}{4}+\ldots\right\},
\]

\textbf{*)} The notation of Cl$^{\mathrm{i}}$ Legendre is somewhat different; for he puts $\Pi(n,\phi)=\displaystyle\int_{0}^{\phi}\frac{d\phi}{\left\{1+n\sin^{2}\phi\right\}\Delta(\phi)}$; so that what for us is $\Pi(u,a)$ will for him be:

\[
\frac{-\cos\alpha\Delta\alpha}{\sin\alpha}F(\phi)+\frac{\cos\alpha\Delta\alpha}{\sin\alpha}\Pi(-k^{2}\sin^{2}\alpha,\phi).
\]
That this sign $\Pi$ should not be confused with the multiplicatory sign $\Pi$, more often employed by us, hardly needs to be mentioned.

\origpage{145}
we see that formula 3), the single denominators $1-q^{2}$, $1-q^{4}$, $1-q^{6}$, etc., being expanded, assumes this form:

\[
\begin{gathered}
\text{4)}\quad \Pi\left(\frac{2Kx}{\pi},\frac{2KA}{\pi}\right)=\\[1.5ex]
\frac{2Kx}{\pi}Z\left(\frac{2KA}{\pi}\right)+\frac{1}{2}\log.\left\{\frac{\left(1-2q\cos2(x-A)+q^{2}\right)\left(1-2q^{3}\cos2(x-A)+q^{6}\right)\ldots}{\left(1-2q\cos2(x+A)+q^{2}\right)\left(1-2q^{3}\cos2(x+A)+q^{6}\right)\ldots}\right\}.
\end{gathered}
\]

\subsection*{52.}

Formula 1) §. 47 being integrated:

\[
\frac{2K}{\pi}Z\left(\frac{2Kx}{\pi}\right)=4\left\{\frac{q\sin2x}{1-q^{2}}+\frac{q^{2}\sin4x}{1-q^{4}}+\frac{q^{3}\sin6x}{1-q^{6}}+\ldots\right\}
\]
from $x=0$ up to $x=x$, there results:

\[
\begin{gathered}
\frac{2K}{\pi}\int_{0}^{x}Z\left(\frac{2Kx}{\pi}\right).dx=-2\left\{\frac{q\cos2x}{1-q^{2}}+\frac{q^{2}\cos4x}{2(1-q^{2})}+\frac{q^{3}\cos6x}{3(1-q^{6})}+.\,.\right\}+\text{Const.}\\[1.5ex]
=\log\left\{(1-2q\cos2x+q^{2})(1-2q^{3}\cos2x+q^{6})(1-2q^{5}\cos2x+q^{10})\ldots\right\}+\text{Const.},
\end{gathered}
\]
\ednote{Printed $2(1-q^{2})$ in the denominator of the second term of the series; an apparent misprint for $2(1-q^{4})$, as in the next display.}
where the \emph{Constant} is so determined that it vanishes for $x=0$; there results $=$\ednote{The print sets an equals sign (or a colon damaged into one) after ``fit''; a colon is expected.}

\[
2\left\{\frac{q}{1-q^{2}}+\frac{q^{2}}{2(1-q^{4})}+\frac{q^{3}}{3(1-q^{6})}+.\,.\right\}=-\log\left\{(1-q)(1-q^{3})(1-q^{5}).\,.\right\}^{2},
\]
and hence:

\[
\text{1)}\quad \frac{2K}{\pi}\int_{0}^{x}Z\left(\frac{2Kx}{\pi}\right).dx=\log\left\{\frac{(1-2q\cos2x+q^{2})(1-2q^{3}\cos2x+q^{6})(1-2q^{5}\cos2x+q^{10})\ldots}{\left\{(1-q)(1-q^{3})(1-q^{5})\ldots\right\}^{2}}\right\}.
\]

We shall designate in future by the character $\Theta(u)$ the expression:

\[
\Theta(u)=\Theta(0)\,e^{\int_{0}^{u}Z(u).du},
\]
$\Theta(0)$ designating a Constant, which we leave as yet undetermined, since we shall obtain a convenient determination of it below; there will be, from 1):

\[
\text{2)}\quad \frac{\Theta\left(\frac{2Kx}{\pi}\right)}{\Theta(0)}=\frac{(1-2q\cos2x+q^{2})(1-2q^{3}\cos2x+q^{6})(1-2q^{5}\cos2x+q^{10})\ldots}{\left\{(1-q)(1-q^{3})(1-q^{5})\ldots\right\}^{2}},
\]

\origpage{146}
whence formula 4) §. 51 passes into this:

\[
\Pi\left(\frac{2Kx}{\pi},\frac{2KA}{\pi}\right)=\frac{2Kx}{\pi}Z\left(\frac{2KA}{\pi}\right)+\frac{1}{2}\log.\frac{\Theta\left(\frac{2K}{\pi}(x-A)\right)}{\Theta\left(\frac{2K}{\pi}(x+A)\right)},
\]
or, with $\dfrac{2Kx}{\pi}=u$, $\dfrac{2KA}{\pi}=a$ again put:

\[
\text{3)}\quad \Pi(u,a)=uZ(a)+\frac{1}{2}\log.\frac{\Theta(u-a)}{\Theta(u+a)}=u.\frac{\Theta'(a)}{\Theta(a)}+\frac{1}{2}\log.\frac{\Theta(u-a)}{\Theta(u+a)},
\]
provided it is put: $\dfrac{d\Theta(u)}{du}=\Theta'(u)$. Which is a convenient expression of the Elliptic Integral $\Pi$ by the new Transcendent $\Theta$.

It is easily established that $\Theta(-u)=\Theta(u)$, whence, $a$ and $u$ being interchanged, from 3) there results:

\[
\Pi(a,u)=aZ(u)+\frac{1}{2}\log.\frac{\Theta(u-a)}{\Theta(u+a)},
\]
which being subtracted from 3), there comes:

\[
\text{4)}\quad \Pi(u,a)-\Pi(a,u)=uZ(a)-aZ(u),
\]
which is the same as formula 2) §. 50. Hence, with $\Pi(K,a)=\Pi^{\mathrm{I}}(a)$ put, $\Pi(a,K)$, $Z(K)$ vanishing, there comes:

\[
\Pi^{\mathrm{I}}(a)=KZ(a),
\]
which is the formula of Cl$^{\mathrm{i}}$ Legendre, which we exhibited above as 3) §. 50.

With $u=a$ put, from 3) there comes:

\[
\text{5)}\quad \Pi(a,a)=aZ(a)+\frac{1}{2}\log.\frac{\Theta(0)}{\Theta(2a)}=aZ(a)-\frac{1}{2}\log.\frac{\Theta(2a)}{\Theta(0)}.
\]
We see therefore that the new Transcendent can be defined either by the Integral $\displaystyle\int\frac{E(\phi).d\phi}{\Delta(\phi)}$ by means of the formula:

\[
\text{6)}\quad \frac{\Theta(u)}{\Theta(0)}=e^{\int_{0}^{u}du.Z(u)}=e^{\int_{0}^{\phi}\frac{F^{\mathrm{I}}E(\phi)-E^{\mathrm{I}}F(\phi)}{F^{\mathrm{I}}\Delta(\phi)}.d\phi},
\]
or by the definite Integral of the third species by means of the formula:

\[
\text{7)}\quad \frac{\Theta(2a)}{\Theta(0)}=e^{2aZ(a)-2\Pi(a,a)}.
\]

\origpage{147}
From formula 5) we obtain:

\[
\begin{gathered}
\frac{1}{2}\log.\frac{\Theta(u-a)}{\Theta(u+a)}=\frac{u-a}{2}.Z\left(\frac{u-a}{2}\right)-\Pi\left(\frac{u-a}{2},\frac{u-a}{2}\right)\\[1.5ex]
-\frac{u+a}{2}.Z\left(\frac{u+a}{2}\right)+\Pi\left(\frac{u+a}{2},\frac{u+a}{2}\right),
\end{gathered}
\]
whence 3) passes into this formula:

\[
\begin{gathered}
\text{8)}\quad \Pi(u,a)=uZ(a)+\frac{u-a}{2}.Z\left(\frac{u-a}{2}\right)-\frac{u+a}{2}.Z\left(\frac{u+a}{2}\right)\\[1.5ex]
-\Pi\left(\frac{u-a}{2},\frac{u-a}{2}\right)+\Pi\left(\frac{u+a}{2},\frac{u+a}{2}\right),
\end{gathered}
\]
which is for the reduction of an Integral of the third sp. that is indefinite to definite ones, and agrees with Theor. II. §. 50.

\subsection*{COROLLARY.}

As we have already above deduced from the expansions found Algorithms suited to computation, less in order that new things be brought forth than that their nature be the better perceived: let us do the same again concerning the expansion found for the function

\[
\begin{gathered}
\frac{\Theta\left(\frac{2Kx}{\pi}\right)}{\Theta(0)}=e^{\int_{0}^{\phi}\frac{F^{\mathrm{I}}E(\phi)-E^{\mathrm{I}}F(\phi)}{F^{\mathrm{I}}\Delta(\phi)}.d\phi}=\\[1.5ex]
\frac{(1-2q\cos2x+q^{2})(1-2q^{3}\cos2x+q^{6})(1-2q^{5}\cos2x+q^{10}).\,.\,.\,.}{\left\{(1-q)(1-q^{3})(1-q^{5})\ldots\right\}}.
\end{gathered}
\]
\ednote{The exponent 2 of the braced denominator is not printed here; it is present in the same expression on the preceding and following pages.}
To which end let us premise the following.

Let the infinite product be put:

\[
T=\left(\frac{1-q}{1+q}\right)\left(\frac{1-q^{2}}{1+q^{2}}\right)^{\frac{1}{2}}\left(\frac{1-q^{4}}{1+q^{4}}\right)^{\frac{1}{4}}\left(\frac{1-q^{8}}{1+q^{8}}\right)^{\frac{1}{8}}\ldots,
\]
provided by repeated turns there is substituted:

\[
(1-q^{2})=(1-q)(1+q),\ (1-q^{4})=(1-q^{2})(1+q^{2}),\ (1-q^{8})=(1-q^{4})(1+q^{4}),\ \ldots,
\]
there results:

\[
T=(1-q)\left(\frac{1-q}{1+q}\right)^{\frac{1}{2}}\left(\frac{1-q^{2}}{1+q^{2}}\right)^{\frac{1}{4}}\left(\frac{1-q^{4}}{1+q^{4}}\right)^{\frac{1}{8}}\left(\frac{1-q^{8}}{1+q^{8}}\right)^{\frac{1}{16}}\ldots
\]

\origpage{148}
\[
\begin{gathered}
=(1-q)(1-q)^{\frac{1}{2}}\left(\frac{1-q}{1+q}\right)^{\frac{1}{4}}\left(\frac{1-q^{2}}{1+q^{2}}\right)^{\frac{1}{8}}\left(\frac{1-q^{4}}{1+q^{4}}\right)^{\frac{1}{16}}.\,.\,.\,.\\[1.5ex]
=(1-q)(1-q)^{\frac{1}{2}}(1-q)^{\frac{1}{4}}\left(\frac{1-q}{1+q}\right)^{\frac{1}{8}}\left(\frac{1-q^{2}}{1+q^{2}}\right)^{\frac{1}{16}}.\,.\,.\,.\\[1.5ex]
.\,.\,.\,.\,.\,.\,,
\end{gathered}
\]
whence we see that there will be:

\[
\text{1)}\quad T=(1-q)(1-q)^{\frac{1}{2}}(1-q)^{\frac{1}{4}}(1-q)^{\frac{1}{8}}(1-q)^{\frac{1}{16}}\ldots=(1-q)^{2}.
\]

Or also, since:

\[
\begin{gathered}
T=\left(\frac{1-q}{1+q}\right)\left(\frac{1-q^{2}}{1+q^{2}}\right)^{\frac{1}{2}}\left(\frac{1-q^{4}}{1+q^{4}}\right)^{\frac{1}{4}}\left(\frac{1-q^{8}}{1+q^{8}}\right)^{\frac{1}{8}}\ldots\\[1.5ex]
=(1-q)\left(\frac{1-q}{1+q}\right)^{\frac{1}{2}}\left(\frac{1-q^{2}}{1+q^{2}}\right)^{\frac{1}{4}}\left(\frac{1-q^{4}}{1+q^{4}}\right)^{\frac{1}{8}}\ldots,
\end{gathered}
\]
there results $T=(1-q)\sqrt{T}$, whence $T=(1-q)^{2}$.

From 1) there results:

\[
\text{2)}\quad 1-q=\left(\frac{1-q}{1+q}\right)^{\frac{1}{2}}\left(\frac{1-q^{2}}{1+q^{2}}\right)^{\frac{1}{4}}\left(\frac{1-q^{4}}{1+q^{4}}\right)^{\frac{1}{8}}\ldots
\]
in which formula in place of $q$ let us successively put $q$, $q^{3}$, $q^{5}$, $q^{7}$, $.\,.$, and let us carry out the infinite multiplication. The formula exhibited above being called in aid:

\[
\sqrt[4]{k'}=\left(\frac{1-q}{1+q}\right)\left(\frac{1-q^{3}}{1+q^{3}}\right)\left(\frac{1-q^{5}}{1+q^{5}}\right)\left(\frac{1-q^{7}}{1+q^{7}}\right)\ldots,
\]
there results:

\[
(1-q)(1-q^{3})(1-q^{5})(1-q^{7})\ldots=\left\{k'\right\}^{\frac{1}{8}}\left\{k^{(2)\prime}\right\}^{\frac{1}{16}}\left\{k^{(4)\prime}\right\}^{\frac{1}{32}}\ldots,
\]
provided we designate, as above, by $k^{(n)\prime}$ the quantity which depends in the same manner on $q^{n}$ as $k'$ on $q$, or the Complement of the Modulus derived by the first transformation of the $n^{\mathrm{ti}}$ order.

Further we found §. 36:

\[
\left\{(1-q)(1-q^{3})(1-q^{5})(1-q^{7}).\,.\right\}^{6}=\frac{2\sqrt[4]{q.k'}}{\sqrt{k}},
\]
whence now:

\[
\text{3)}\quad q=e^{-\frac{\pi K'}{K}}=\frac{kk}{16k'}\left\{k^{(2)\prime}\right\}^{\frac{3}{2}}\left\{k^{(4)\prime}\right\}^{\frac{3}{4}}\left\{k^{(8)\prime}\right\}^{\frac{3}{8}}\ldots
\]

\origpage{149}
With $m=1$, $n=k'$ put; $\dfrac{m+n}{2}=m'$, $\sqrt{mn}=n'$; $\dfrac{m'+n'}{2}=m''$, $\sqrt{m'n'}=n''$, etc.; it is known that $k^{(2)\prime}=\dfrac{n'}{m'}$, $k^{(4)\prime}=\dfrac{n''}{m''}$, $k^{(8)\prime}=\dfrac{n'''}{m'''}$, etc., whence:

\[
\text{4)}\quad q=\frac{mm-nn}{16mn}.\left\{\left(\frac{n'}{m'}\right)^{\frac{1}{2}}\left(\frac{n''}{m''}\right)^{\frac{1}{4}}\left(\frac{n'''}{m'''}\right)^{\frac{1}{8}}\ldots\right\}^{3}.
\]
Hence also follows, $\mu=\dfrac{\pi}{2K}$ designating the common limit to which the quantities $m^{(p)}$, $n^{(p)}$ converge:

\[
\text{5)}\quad K'=\frac{1}{2\mu}\left\{\log\frac{16mn}{mn-nn}+\frac{3}{2}\log\frac{m'}{n'}+\frac{3}{4}\log\frac{m''}{n''}+\frac{3}{8}\log\frac{m'''}{n'''}+\ldots\right\},
\]
\ednote{Printed $mn-nn$ in the denominator; an apparent misprint for $mm-nn$, as in formula 4).}
which formulae supply a most expeditious computation. Formula 5) teaches how, from the same series of quantities which you must have calculated for finding the value of the function $K$, the value of $K'$ also at once results.

Let us transform formula 3). There results, as is known:

\[
k'=\frac{1-k^{(2)}}{1+k^{(2)}};\ k=\frac{2\sqrt{k^{(2)}}}{1+k^{(2)}},\ \text{whence}\ \frac{kk}{k'}=\frac{4k^{(2)}}{k^{(2)\prime}k^{(2)\prime}}.
\]
Hence we obtain, provided that by repeated turns we simultaneously substitute $k^{(2)}$ in place of $k$ and extract the square root:

\[
\begin{gathered}
\frac{kk}{16k'}.\left\{k^{(2)\prime}\right\}^{\frac{3}{2}}=\left\{\frac{k^{(2)}k^{(2)}}{16k^{(2)\prime}}\right\}^{\frac{1}{2}}\\[1.5ex]
\left\{\frac{k^{(2)}k^{(2)}}{16k^{(2)\prime}}\right\}^{\frac{1}{2}}\left\{k^{(4)\prime}\right\}^{\frac{3}{4}}=\left\{\frac{k^{(4)}k^{(4)}}{16k^{(4)\prime}}\right\}^{\frac{1}{4}}\\[1.5ex]
\left\{\frac{k^{(4)}k^{(4)}}{16k^{(4)\prime}}\right\}^{\frac{1}{4}}\left\{k^{(8)\prime}\right\}^{\frac{3}{8}}=\left\{\frac{k^{(8)}k^{(8)}}{16k^{(8)\prime}}\right\}^{\frac{1}{8}}\\[1.5ex]
.\,.\,.\,.\,,
\end{gathered}
\]
whence with $p=2^{m}$ put:

\[
\frac{kk}{16k'}.\left\{k^{(2)\prime}\right\}^{\frac{3}{2}}\left\{k^{(4)}\right\}^{\frac{3}{4}}\left\{k^{(8)\prime}\right\}^{\frac{3}{8}}.\,.\left\{k^{(p)}\right\}^{\frac{3}{p}}=\left\{\frac{k^{(p)}k^{(p)}}{16k^{(p)\prime}}\right\}^{\frac{1}{p}}.
\]
\ednote{Printed $k^{(4)}$ and $k^{(p)}$ without the prime in this line; apparently printer's omissions for $k^{(4)\prime}$ and $k^{(p)\prime}$, as in the preceding lines and on the right.}
Hence we see from formula 3) that $q=e^{-\frac{\pi K'}{K}}$ will be the limit of the expression $\left\{\dfrac{k^{(p)}k^{(p)}}{16}\right\}^{\frac{1}{p}}$, as $m$ or $p$ increases to infinity, which is a theorem discovered by Cl$^{\mathrm{o}}$ Legendre.

\origpage{150}
Moreover, even at a mere glance at the formula exhibited by us:

\[
k=4\sqrt{q}\left\{\frac{(1+q^{2})(1+q^{4})(1+q^{6})(1+q^{8})\ldots}{(1+q)(1+q^{3})(1+q^{5})(1+q^{7})\ldots}\right\}^{4}
\]
it is evident that, neglecting quantities of order $q^{p}$, there will be:

\[
q=\sqrt[p]{\frac{k^{(p)}k^{(p)}}{16}},
\]
which agrees with the said theorem.

Now in our formula

\[
1-q=\left\{\frac{1-q}{1+q}\right\}^{\frac{1}{2}}\left\{\frac{1-q^{2}}{1+q^{2}}\right\}^{\frac{1}{4}}\left\{\frac{1-q^{4}}{1+q^{4}}\right\}^{\frac{1}{8}}\ldots
\]
in place of $q$ let us substitute successively a double series of quantities:

\[
\begin{gathered}
qe^{2ix},\ q^{3}e^{2ix},\ q^{5}e^{2ix},\ q^{7}e^{2ix},\ \ldots\\[1ex]
qe^{-2ix},\ q^{3}e^{-2ix},\ q^{5}e^{-2ix},\ q^{7}e^{-2ix},\ \ldots,
\end{gathered}
\]
and let us carry out an infinite multiplication. Let us call upon the formula §$^{\mathrm{i}}$ 36:

\[
\frac{\Delta\operatorname{am}\dfrac{2Kx}{\pi}}{\sqrt{k'}}=\frac{(1-2q\cos2x+q^{2})(1-2q^{3}\cos2x+q^{6})(1-2q^{5}\cos2x+q^{10})\ldots}{(1+2q\cos2x+q^{2})(1+2q^{3}\cos2x+q^{6})(1+2q^{5}\cos2x+q^{10})\ldots},
\]
and let us denote by $\Delta^{(p)}$ the expression

\[
\frac{\Delta\operatorname{am}\dfrac{2pK^{(p)}x}{\pi}}{\sqrt{k^{(p)\prime}}}=\frac{(1-2q^{p}\cos2px+q^{2p})(1-2q^{3p}\cos2px+q^{6p})(1-2q^{5p}\cos2px+q^{10p})\ldots}{(1+2q^{p}\cos2px+q^{2p})(1+2q^{3p}\cos2px+q^{6p})(1+2q^{5p}\cos2px+q^{10p})\ldots},
\]
there results:

\[
\Delta^{\frac{1}{2}}{\Delta^{(2)}}^{\frac{1}{4}}{\Delta^{(4)}}^{\frac{1}{8}}{\Delta^{(8)}}^{\frac{1}{16}}\ldots=\frac{(1-2q\cos2x+q^{2})(1-2q^{3}\cos2x+q^{6})(1-2q^{5}\cos2x+q^{10})\ldots}{\left\{(1-q)(1-q^{3})(1-q^{5})\ldots\right\}^{2}}.
\]
The constant factor which we have added, $\dfrac{1}{(1-q)^{2}(1-q^{3})^{2}(1-q^{5})^{2}.\,.}$, we have determined from what was found above, or by the fact that each expression, on putting $x=0$, becomes equal to unity. Now we have found:

\[
\frac{\Theta\left(\dfrac{2Kx}{\pi}\right)}{\Theta(0)}=\frac{(1-2q\cos2x+q^{2})(1-2q^{3}\cos2x+q^{6})(1-2q^{5}\cos2x+q^{10})\ldots}{\left\{(1-q)(1-q^{3})(1-q^{5})\ldots\right\}^{2}}
\]
whence

\[
\frac{\Theta\left(\dfrac{2Kx}{\pi}\right)}{\Theta(0)}=\Delta^{\frac{1}{2}}.{\Delta^{(2)}}^{\frac{1}{4}}.{\Delta^{(4)}}^{\frac{1}{8}}.{\Delta^{(8)}}^{\frac{1}{16}}\ldots
\]

\origpage{151}
Hence, on putting $\dfrac{2Kx}{\pi}=u$, $\operatorname{am}u=\phi$, and calling upon the formulae which Cl.\ Legendre has set forth on the transformation of the second order, we obtain the following, which furnishes a ready computation of the function $\Theta$,

\subsection*{THEOREM.}

Let $\operatorname{am}(u)=\phi$, $m=1$, $n=k'$, $\Delta(\phi)=\sqrt{mm\cos^{2}\phi+nn\sin^{2}\phi}=\Delta$ be put, and let the series of quantities be calculated:

\[
\begin{gathered}
m'=\frac{m+n}{2},\ m''=\frac{m'+n'}{2},\ m'''=\frac{m''+n''}{2},\ \ldots\\[1.5ex]
n'=\sqrt{mn},\ n''=\sqrt{m'n'},\ n'''=\sqrt{m''n''},\ \ldots\\[1.5ex]
\Delta'=\frac{\Delta\Delta+n'n'}{2\Delta},\ \Delta''=\frac{\Delta'\Delta'+n''n''}{2\Delta'},\ \Delta'''=\frac{\Delta''\Delta''+n'''n'''}{2\Delta''},\ \ldots,
\end{gathered}
\]
there will be:

\[
\frac{\Theta(u)}{\Theta(0)}=e^{\int_{0}^{\phi}\frac{F^{\mathrm{I}}E(\phi)-E^{\mathrm{I}}F(\phi)}{F^{\mathrm{I}}\Delta(\phi)}.d\phi}=\left\{\frac{\Delta}{m}\right\}^{\frac{1}{2}}.\left\{\frac{\Delta'}{m'}\right\}^{\frac{1}{4}}.\left\{\frac{\Delta''}{m''}\right\}^{\frac{1}{8}}.\left\{\frac{\Delta'''}{m'''}\right\}^{\frac{1}{16}}\ldots
\]

The task of demonstrating this theorem without consideration of expansions, by known and finite formulae, since it is ready to hand, we forgo.

\subsection*{ON THE ADDITION OF ARGUMENTS AND OF THE PARAMETER AND OF THE AMPLITUDE IN THE THIRD SPECIES OF ELLIPTIC INTEGRALS.}

\subsection*{53.}

We obtain the formula fundamental in the Analysis of the Function $\Theta$, and of which we shall have very frequent use in what follows, by the following consideration. For since it has been put:

\[
\Pi(u,a)=\int_{0}^{u}\frac{k^{2}\sin\operatorname{am}a\cos\operatorname{am}a\Delta\operatorname{am}a.\sin^{2}\operatorname{am}u.du}{1-k^{2}\sin^{2}\operatorname{am}a.\sin^{2}\operatorname{am}u},
\]
we have:

\[
\frac{d\Pi(u,a)}{du}=\frac{k^{2}\sin\operatorname{am}a\cos\operatorname{am}a\Delta\operatorname{am}a.\sin^{2}\operatorname{am}u}{1-k^{2}\sin^{2}\operatorname{am}a.\sin^{2}\operatorname{am}u}.
\]

\origpage{152}
This formula integrated with respect to $a$ from $a=0$ up to $a=a$ gives:

\[
\text{1)}\quad \int_{0}^{a}da.\frac{d\Pi(u,a)}{du}=-\frac{1}{2}\log\left(1-k^{2}\sin^{2}\operatorname{am}a\sin^{2}\operatorname{am}u\right).
\]
But from 3) §. 52 we have:

\[
\text{2)}\quad \frac{d\Pi(u,a)}{du}=Z(a)+\frac{1}{2}\frac{\Theta'(u-a)}{\Theta(u-a)}-\frac{1}{2}\frac{\Theta'(u+a)}{\Theta(u+a)},
\]
whence:

\[
\int_{0}^{a}da.\frac{d.\Pi(u,a)}{du}=\log.\frac{\Theta(a)}{\Theta(0)}-\frac{1}{2}\log\Theta(u-a)-\frac{1}{2}\log\Theta(u+a)+\log\Theta(u),
\]
these\ednote{Printed quibos for quibus; see the author's Corrigenda.} being substituted, when you pass from logarithms to numbers\ednote{Printed tranis for transis; see the author's Corrigenda.}, from 1) you obtain:

\[
\text{3)}\quad \Theta(u+a)\Theta(u-a)=\left\{\frac{\Theta u.\Theta a}{\Theta0}\right\}^{2}\left(1-k^{2}\sin^{2}\operatorname{am}a.\sin^{2}\operatorname{am}u\right).
\]

We can represent formula 2) thus:

\[
\frac{k^{2}\sin\operatorname{am}a\cos\operatorname{am}a\Delta\operatorname{am}a.\sin^{2}\operatorname{am}u}{1-k^{2}\sin^{2}\operatorname{am}a\sin^{2}\operatorname{am}u}=Z(a)+\frac{1}{2}Z(u-a)-\frac{1}{2}Z(u+a),
\]
whence, on interchanging $a$ and $u$:

\[
\frac{k^{2}\sin\operatorname{am}u\cos\operatorname{am}u\Delta\operatorname{am}u.\sin^{2}\operatorname{am}a}{1-k^{2}\sin^{2}\operatorname{am}a\sin^{2}\operatorname{am}u}=Z(u)-\frac{1}{2}Z(u-a)-\frac{1}{2}Z(u+a),
\]
and adding these formulae there results:

\[
\text{4)}\quad Z(u)+Z(a)-Z(u+a)=k^{2}\sin\operatorname{am}u.\sin\operatorname{am}a.\sin\operatorname{am}(u+a),
\]
which is for the Addition of the function $Z$, and agrees with formula 3) §. 49:

\[
E(\phi)+E(\alpha)-E(\sigma)=k^{2}\sin\phi.\sin\alpha.\sin\sigma.
\]
On putting $a=K$, since it is easily established that $Z(K)=\dfrac{F^{\mathrm{I}}E^{\mathrm{I}}-E^{\mathrm{I}}F^{\mathrm{I}}}{F^{\mathrm{I}}}=0$, there results from 4):

\[
\text{5)}\quad Z(u)-Z(u+K)=k^{2}\sin\operatorname{am}u.\sin\operatorname{coam}u,
\]
which in §. 47 we derived from the expansion of $Z$ itself. Putting $-u$ in place of $u$, $K-u=v$, we obtain from formula 5):

\[
\text{6)}\quad Z(u)+Z(v)=k^{2}\sin\operatorname{am}u.\sin\operatorname{am}v.
\]
On putting $u=v=\dfrac{K}{2}$, we have: $2Z\left(\dfrac{K}{2}\right)=1-k'{}^{*)}$.

\textbf{*)} For there is: $\sin\operatorname{am}\dfrac{K}{2}=\sqrt{\dfrac{1}{1+k'}}$, $\cos\operatorname{am}\dfrac{K}{2}=\sqrt{\dfrac{k'}{1+k'}}$, $\Delta\operatorname{am}\dfrac{K}{2}=\sqrt{k'}$, $\operatorname{tg}\operatorname{am}\dfrac{K}{2}=\dfrac{1}{\sqrt{k'}}$.

\origpage{153}
Let us integrate formula 5) from $u=0$ up to $u=u$. Since $\displaystyle\int_{0}^{u}Z(u).du=\log\frac{\Theta u}{\Theta0}$, there results:

\[
\log\frac{\Theta u}{\Theta0}-\log.\frac{\Theta(u+K)}{\Theta(K)}=-\log\Delta\operatorname{am}u,
\]
or:

\[
\text{7)}\quad \frac{\Theta0}{\Theta K}.\frac{\Theta(u+K)}{\Theta u}=\Delta\operatorname{am}u.
\]

On putting $u=-K$, we derive from 7) the value of

\[
\text{8)}\quad \frac{\Theta K}{\Theta0}=\frac{1}{\sqrt{k'}},
\]
whence 7) takes the form:

\[
\text{9)}\quad \frac{\Theta(u+K)}{\Theta u}=\frac{\Delta\operatorname{am}u}{\sqrt{k'}}.
\]

Formula 9), from the expansion found:

\[
\frac{\Theta\left(\dfrac{2Kx}{\pi}\right)}{\Theta(0)}=\frac{(1-2q\cos2x+q^{2})(1-2q^{3}\cos2x+q^{6})(1-2q^{5}\cos2x+q^{10})\ldots}{\left\{(1-q)(1-q^{3})(1-q^{5})\ldots\right\}^{2}}
\]
we easily confirm. For on changing $x$ into $x+\dfrac{\pi}{2}$ there is:

\[
\frac{\Theta\left(\dfrac{2Kx}{\pi}+K\right)}{\Theta(0)}=\frac{(1+2q\cos2x+q^{2})(1+2q^{3}\cos2x+q^{6})(1+2q^{5}\cos2x+q^{10})\ldots}{\left\{(1-q)(1-q^{3})(1-q^{5})\ldots\right\}^{2}}
\]
whence:

\[
\frac{\Theta\left(\dfrac{2Kx}{\pi}+K\right)}{\Theta\left(\dfrac{2Kx}{\pi}\right)}=\frac{(1+2q\cos2x+q^{2})(1+2q^{3}\cos2x+q^{6})(1+2q^{5}\cos2x+q^{10})\ldots}{(1-2q\cos2x+q^{2})(1-2q^{3}\cos2x+q^{6})(1-2q^{5}\cos2x+q^{10})\ldots},
\]
which very expression we found in §. 35 to be $=\dfrac{\Delta\operatorname{am}\dfrac{2Kx}{\pi}}{\sqrt{k'}}$, as it ought.

From formula 9) we immediately reduce the expressions $\Pi(u+K,a)$, $\Pi(u,a+K)$ to $\Pi(u,a)$ itself. For there is:

\origpage{154}
\[
\begin{gathered}
\text{10)}\quad \Pi(u+K,a)=(u+K)Z(a)+\frac{1}{2}\log.\frac{\Theta(u+K-a)}{\Theta(u+K+a)}\\[1.5ex]
=(u+K)Z(a)+\frac{1}{2}\log.\frac{\Theta(u-a)}{\Theta(u+a)}+\frac{1}{2}\log.\frac{\Delta\operatorname{am}(u-a)}{\Delta\operatorname{am}(u+a)}\\[1.5ex]
=\Pi(u,a)+K.Z(a)+\frac{1}{2}\log.\frac{\Delta\operatorname{am}(u-a)}{\Delta\operatorname{am}(u+a)}.
\end{gathered}
\]

\[
\begin{gathered}
\text{11)}\quad \Pi(u,a+K)=uZ(a+K)+\frac{1}{2}\log\frac{\Theta(u-a-K)}{\Theta(u+a+K)}=\\[1.5ex]
uZ(a)-k^{2}\sin\operatorname{am}a.\sin\operatorname{coam}a.u+\frac{1}{2}\log\frac{\Theta(u-a)}{\Theta(u+a)}+\frac{1}{2}\log.\frac{\Delta\operatorname{am}(u-a)}{\Delta\operatorname{am}(u+a)}=\\[1.5ex]
\Pi(u,a)-k^{2}\sin\operatorname{am}a\sin\operatorname{coam}a.u+\frac{1}{2}\log\frac{\Delta\operatorname{am}(u-a)}{\Delta\operatorname{am}(u+a)}.
\end{gathered}
\]

\subsection*{54.}

From the fundamental formula, by whose aid the function $\Pi$ is defined by the functions $Z$, $\Theta$:

\[
\text{I)}\quad \Pi(u,a)=uZ(a)+\frac{1}{2}\log.\frac{\Theta(u-a)}{\Theta(u+a)},
\]
calling in the following, themselves also fundamental in the Analysis of the functions $Z$, $\Theta$:

\[
\begin{gathered}
\text{II)}\quad Zu-Z(u+a)=k^{2}\sin\operatorname{am}a.\sin\operatorname{am}u.\sin\operatorname{am}(u+a)\\[1.5ex]
\text{III)}\quad \Theta(u+a)\Theta(u-a)=\left\{\frac{\Theta u.\Theta a}{\Theta0}\right\}^{2}\left(1-k^{2}\sin^{2}\operatorname{am}a.\sin^{2}\operatorname{am}u\right),
\end{gathered}
\]
\ednote{In formula II) the term $+Za$ is missing on the left side (the author's Corrigenda read $Z(u)+Z(a)-Z(u+a)$); an apparent misprint, the omission being supplied only in faint hand correction, which is not transcribed.}
you now easily obtain formulae, both for expressing $\Pi(u+v,a)$ by $\Pi(u,a)$, $\Pi(v,a)$, which we shall call the theorem on the \emph{Addition of the Argument of the Amplitude}, and for expressing $\Pi(u,a+b)$ by $\Pi(u,a)$, $\Pi(u,b)$, which we shall call the theorem on the \emph{Addition of the Argument of the Parameter}. To this end we note the following.

From the formulae:

\[
\begin{gathered}
\Pi(u,a)=u.Za+\frac{1}{2}\log.\frac{\Theta(u-a)}{\Theta(u+a)}\\[1.5ex]
\Pi(v,a)=v.Za+\frac{1}{2}\log.\frac{\Theta(v-a)}{\Theta(v+a)}\\[1.5ex]
\Pi(u+v,a)=(u+v)Za+\frac{1}{2}\log.\frac{\Theta(u+v-a)}{\Theta(u+v+a)}
\end{gathered}
\]
it follows:

\[
\text{1)}\quad \Pi(u,a)+\Pi(v,a)-\Pi(u+v,a)=\frac{1}{2}\log.\frac{\Theta(u-a).\Theta(v-a).\Theta(n+v+a)}{\Theta(u+a).\Theta(v+a).\Theta(u+v-a)}
\]
\ednote{Printed $\Theta(n+v+a)$ in the numerator; an apparent misprint for $\Theta(u+v+a)$, as on the next page.}

\origpage{155}
The expression contained under the logarithmic sign:

\[
\frac{\Theta(u-a).\Theta(v-a).\Theta(u+v+a)}{\Theta(u+a).\Theta(v+a).\Theta(u+v-a)}
\]
it is permitted, by the aid of the fundamental theorem III., to reduce to elliptic functions in a twofold way. For first from it there is:

\[
\begin{gathered}
\Theta(u-a).\Theta(v-a)=\left\{\frac{\Theta\left(\frac{u-v}{2}\right).\Theta\left(\frac{u+v}{2}-a\right)}{\Theta0}\right\}^{2}\left(1-k^{2}\sin^{2}\operatorname{am}\left(\frac{u-v}{2}\right).\sin^{2}\operatorname{am}\left(\frac{u+v}{2}-a\right)\right)\\[2ex]
\Theta(u+a).\Theta(v+a)=\left\{\frac{\Theta\left(\frac{u-v}{2}\right).\Theta\left(\frac{u+v}{2}+a\right)}{\Theta0}\right\}^{2}\left(1-k^{2}\sin^{2}\operatorname{am}\left(\frac{u-v}{2}\right).\sin^{2}\operatorname{am}\left(\frac{u+v}{2}+a\right)\right)\\[2ex]
\Theta(u+v-a)\Theta a=\left\{\frac{\Theta\left(\frac{u+v}{2}\right).\Theta\left(\frac{u+v}{2}-a\right)}{\Theta0}\right\}^{2}\left(1-k^{2}\sin^{2}\operatorname{am}\left(\frac{u+v}{2}\right).\sin^{2}\operatorname{am}\left(\frac{u+v}{2}-a\right)\right)\\[2ex]
\Theta(u+v+a)\Theta a=\left\{\frac{\Theta\left(\frac{u+v}{2}\right).\Theta\left(\frac{u+v}{2}+a\right)}{\Theta0}\right\}^{2}\left(1-k^{2}\sin^{2}\operatorname{am}\left(\frac{u+v}{2}\right).\sin^{2}\operatorname{am}\left(\frac{u+v}{2}+a\right)\right),
\end{gathered}
\]
of which formulae the first and fourth, being multiplied together and divided by the second and third, there results:

\[
\begin{gathered}
\text{2)}\quad \frac{\Theta(u-a).\Theta(v-a).\Theta(u+v+a)}{\Theta(u+a).\Theta(v+a).\Theta(u+v-a)}=\\[1.5ex]
\frac{\left\{1-k^{2}\sin^{2}\operatorname{am}\left(\frac{u-v}{2}\right).\sin^{2}\operatorname{am}\left(\frac{u+v}{2}-a\right)\right\}\left\{1-k^{2}\sin^{2}\operatorname{am}\left(\frac{u+v}{2}\right).\sin^{2}\operatorname{am}\left(\frac{u+v}{2}+a\right)\right\}}{\left\{1-k^{2}\sin^{2}\operatorname{am}\left(\frac{u-v}{2}\right).\sin^{2}\operatorname{am}\left(\frac{u+v}{2}+a\right)\right\}\left\{1-k^{2}\sin^{2}\operatorname{am}\left(\frac{u+v}{2}\right).\sin^{2}\operatorname{am}\left(\frac{u+v}{2}-a\right)\right\}}.
\end{gathered}
\]

So also, which is the second way, when you represent the fundamental theorem III. in this manner:

\[
\left\{\frac{\Theta u.\Theta v}{\Theta0}\right\}^{2}=\frac{\Theta(u+v)\Theta(u-v)}{1-k^{2}\sin^{2}\operatorname{am}u.\sin^{2}\operatorname{am}v},
\]
there is:

\[
\begin{gathered}
\left\{\frac{\Theta(u-a)\Theta(v-a)}{\Theta0}\right\}^{2}=\frac{\Theta(u-v).\Theta(u+v-2a)}{1-k^{2}\sin^{2}\operatorname{am}(u-a).\sin^{2}\operatorname{am}(v-a)}\\[2ex]
\left\{\frac{\Theta(u+a)\Theta(v+a)}{\Theta0}\right\}^{2}=\frac{\Theta(u-v).\Theta(u+v+2a)}{1-k^{2}\sin^{2}\operatorname{am}(u+a).\sin^{2}\operatorname{am}(v+a)}\\[2ex]
\left\{\frac{\Theta a.\Theta(u+v-a)}{\Theta0}\right\}^{2}=\frac{\Theta(u+v).\Theta(u+v-2a)}{1-k^{2}\sin^{2}\operatorname{am}a.\sin^{2}\operatorname{am}(u+v-a)}
\end{gathered}
\]

\origpage{156}
\[
\left\{\frac{\Theta a.\Theta(u+v+a)}{\Theta0}\right\}^{2}=\frac{\Theta(u+v).\Theta(u+v+2a)}{1-k^{2}\sin^{2}\operatorname{am}a.\sin^{2}\operatorname{am}(u+v+a)},
\]
of which formulae again the first and fourth, being multiplied together and divided by the second and third, and roots being extracted, there results:

\[
\begin{gathered}
\text{3)}\quad \frac{\Theta(u-a).\Theta(v-a).\Theta(u+v+a)}{\Theta(u+a).\Theta(v+a).\Theta(u+v-a)}=\\[1.5ex]
\sqrt{\frac{\left\{1-k^{2}\sin^{2}\operatorname{am}(u+a).\sin^{2}\operatorname{am}(v+a)\right\}\left\{1-k^{2}\sin^{2}\operatorname{am}a.\sin^{2}\operatorname{am}(u+v-a)\right\}}{\left\{1-k^{2}\operatorname{siu}^{2}\operatorname{am}(u-a).\sin^{2}\operatorname{am}(v-a)\right\}\left\{1-k^{2}\sin^{2}\operatorname{am}a.\sin^{2}\operatorname{am}(u+v+a)\right\}}}.
\end{gathered}
\]
\ednote{Printed siu instead of sin in the denominator; an apparent misprint.}

That it may be recognised from the elements themselves how the expressions 2), 3) can be transformed one into the other, I note the following.

When in the formula, already often employed:

\[
\sin\operatorname{am}(u+v).\sin\operatorname{am}(u-v)=\frac{\sin^{2}\operatorname{am}u-\sin^{2}\operatorname{am}v}{1-k^{2}\sin^{2}\operatorname{am}u.\sin^{2}\operatorname{am}v}
\]
you put respectively $u+v$, $u-v$ in place of $u$, $v$, there results:

\[
\sin\operatorname{am}2u.\sin\operatorname{am}2v=\frac{\sin^{2}\operatorname{am}(u+v)-\sin^{2}\operatorname{am}(u-v)}{1-k^{2}\sin^{2}\operatorname{am}(u+v).\sin^{2}\operatorname{am}(u-v)}.
\]
Further we have given the formula:

\[
\sin^{2}\operatorname{am}(u+v)-\sin^{2}\operatorname{am}(u-v)=\frac{4\sin\operatorname{am}u.\cos\operatorname{am}u.\Delta\operatorname{am}u.\sin\operatorname{am}v.\cos\operatorname{am}v.\Delta\operatorname{am}v}{\left\{1-k^{2}\sin^{2}\operatorname{am}u\sin^{2}\operatorname{am}v\right\}^{2}},
\]
whence, multiplication being performed, we obtain:

\[
\begin{gathered}
\text{4)}\quad 1-k^{2}\sin^{2}\operatorname{am}(u+v).\sin^{2}\operatorname{am}(u-v)=\frac{4\sin\operatorname{am}u.\cos\operatorname{am}u.\Delta\operatorname{am}u.\sin\operatorname{am}v.\cos\operatorname{am}v.\Delta\operatorname{am}v}{\sin\operatorname{am}2u.\sin\operatorname{am}2v\left\{1-k^{2}\sin^{2}\operatorname{am}u.\sin^{2}\operatorname{am}v\right\}^{2}}\\[1.5ex]
=\frac{\left\{1-k^{2}\sin^{4}\operatorname{am}u\right\}\left\{1-k^{2}\sin^{4}\operatorname{am}v\right\}}{\left\{1-k^{2}\sin^{2}\operatorname{am}u.\sin^{2}\operatorname{am}v\right\}^{2}}{}^{*)},
\end{gathered}
\]
by the benefit of which formula the formulae 2), 3) now pass easily one into the other.

From formula 4) there can be deduced further this more general one:

\[
\begin{gathered}
\text{5)}\quad \frac{\left\{1-k^{2}\sin^{2}\operatorname{am}u.\sin^{2}\operatorname{am}v\right\}\left\{1-k^{2}\sin^{2}\operatorname{am}u'.\sin^{2}\operatorname{am}v'\right\}}{\left\{1-k^{2}\sin^{2}\operatorname{am}u.\sin^{2}\operatorname{am}v'\right\}\left\{1-k^{2}\sin^{2}\operatorname{am}u'.\sin^{2}\operatorname{am}v\right\}}=\\[1.5ex]
\sqrt{\frac{\left\{1-k^{2}\sin^{2}\operatorname{am}(u+u').\sin^{2}\operatorname{am}(u-u')\right\}\left\{1-k^{2}\sin^{2}\operatorname{am}(v+v').\sin^{2}\operatorname{am}(v-v')\right\}}{\left\{1-k^{2}\sin^{2}\operatorname{am}(u+v).\sin^{2}\operatorname{am}(u-v)\right\}\left\{1-k^{2}\sin^{2}\operatorname{am}(u'+v').\sin^{2}\operatorname{am}(u'-v')\right\}}}.
\end{gathered}
\]

\textbf{*)} For the formula is known: $\sin\operatorname{am}2u=\dfrac{2\sin\operatorname{am}u.\cos\operatorname{am}u.\Delta\operatorname{am}u}{1-k^{2}\sin^{4}\operatorname{am}u}$.

\origpage{157}
But Cl.\ Legendre, in that place where he treats of the Addition of the Argument of the Amplitude, (Chap.\ XVI \emph{Comparaison des fonctions elliptiques de la troisième espèce}) exhibits the quantity which is found under the logarithmic sign in this form:
\[
\frac{1-k^{2}\sin\operatorname{am}a.\sin\operatorname{am}u.\sin\operatorname{am}v.\sin\operatorname{am}(u+v-a)}{1+k^{2}\sin\operatorname{am}a.\sin\operatorname{am}u.\sin\operatorname{am}v.\sin\operatorname{am}(u+v+a)},
\]
of which it is not evident at first glance how it agrees with the expressions found by us, either 2) or 3). The rather abstruse transformation is carried out in this manner.

From the elementary formula, of which we have already made very frequent application, there is:
\[
\begin{gathered}
\sin\operatorname{am}u.\sin\operatorname{am}v=\frac{\sin^{2}\operatorname{am}\left(\frac{u+v}{2}\right)-\sin^{2}\operatorname{am}\left(\frac{u-v}{2}\right)}{1-k^{2}\sin^{2}\operatorname{am}\left(\frac{u+v}{2}\right)\sin^{2}\operatorname{am}\left(\frac{u-v}{2}\right)}\\[1.5ex]
\sin\operatorname{am}a.\sin\operatorname{am}(u+v-a)=\frac{\sin^{2}\operatorname{am}\left(\frac{u+v}{2}\right)-\sin^{2}\operatorname{am}\left(\frac{u+v}{2}-a\right)}{1-k^{2}\sin^{2}\operatorname{am}\left(\frac{u+v}{2}\right)\sin^{2}\operatorname{am}\left(\frac{u+v}{2}-a\right)};
\end{gathered}
\]
these equations being multiplied together, there results:
\[
\begin{gathered}
\left\{1-k^{2}\sin^{2}\operatorname{am}\left(\frac{u+v}{2}\right)\sin^{2}\operatorname{am}\left(\frac{u-v}{2}\right)\right\}\left\{1-k^{2}\sin^{2}\operatorname{am}\left(\frac{u+v}{2}\right)\sin^{2}\operatorname{am}\left(\frac{u+v}{2}-a\right)\right\}\times\\[1.5ex]
\left\{1-k^{2}\sin\operatorname{am}a.\sin\operatorname{am}u.\sin\operatorname{am}v.\sin\operatorname{am}(u+v-a)\right\}\\[1.5ex]
=\\[1.5ex]
\left\{1-k^{2}\sin^{2}\operatorname{am}\left(\frac{u+v}{2}\right)\sin^{2}\operatorname{am}\left(\frac{u-v}{2}\right)\right\}\left\{1-k^{2}\sin^{2}\operatorname{am}\left(\frac{u+v}{2}\right)\sin^{2}\operatorname{am}\left(\frac{u+v}{2}-a\right)\right\}\\[1.5ex]
-k^{2}\left\{\sin^{2}\operatorname{am}\left(\frac{u+v}{2}\right)-\sin^{2}\operatorname{am}\left(\frac{u-v}{2}\right)\right\}\left\{\sin^{2}\operatorname{am}\left(\frac{u+v}{2}\right)-\sin^{2}\operatorname{am}\left(\frac{u+v}{2}-a\right)\right\}
\end{gathered}
\]
The other side of the equation being expanded, the terms
\[
\begin{gathered}
-k^{2}\sin^{2}\operatorname{am}\left(\frac{u+v}{2}\right)\left\{\sin^{2}\operatorname{am}\left(\frac{u-v}{2}\right)+\sin^{2}\operatorname{am}\left(\frac{u+v}{2}-a\right)\right\}\\[1.5ex]
+k^{2}\sin^{2}\operatorname{am}\left(\frac{u+v}{2}\right)\left\{\sin^{2}\operatorname{am}\left(\frac{u-v}{2}\right)+\sin^{2}\operatorname{am}\left(\frac{u+v}{2}-a\right)\right\}
\end{gathered}
\]
destroying one another mutually, there is:
\[
\begin{gathered}
1+k^{4}\sin^{4}\operatorname{am}\left(\frac{u+v}{2}\right)\sin^{2}\operatorname{am}\left(\frac{u-v}{2}\right)\sin^{2}\operatorname{am}\left(\frac{u+v}{2}-a\right)\\[1.5ex]
-k^{2}\sin^{4}\operatorname{am}\left(\frac{u+v}{2}\right)-k^{2}\sin^{2}\operatorname{am}\left(\frac{u-v}{2}\right)\sin^{2}\operatorname{am}\left(\frac{u+v}{2}-a\right)=\\[1.5ex]
\left\{1-k^{2}\sin^{4}\operatorname{am}\left(\frac{u+v}{2}\right)\right\}\left\{1-k^{2}\sin^{2}\operatorname{am}\left(\frac{u-v}{2}\right)\sin^{2}\operatorname{am}\left(\frac{u+v}{2}-a\right)\right\},
\end{gathered}
\]

\origpage{158}
whence at length results:
\[
\begin{gathered}
\text{6)}\quad \frac{1-k^{2}\sin^{2}\operatorname{am}\left(\frac{u+v}{2}\right)\sin^{2}\operatorname{am}\left(\frac{u-v}{2}\right)}{1-k^{2}\sin^{4}\operatorname{am}\left(\frac{u+v}{2}\right)}.\left\{1-k^{2}\sin\operatorname{am}a.\sin\operatorname{am}u.\sin\operatorname{am}v.\sin\operatorname{am}(u+v-a)\right\}\\[1.5ex]
=\\[1.5ex]
\frac{1-k^{2}\sin^{2}\operatorname{am}\left(\frac{u-v}{2}\right)\sin^{2}\operatorname{am}\left(\frac{u+v}{2}-a\right)}{1-k^{2}\sin^{2}\operatorname{am}\left(\frac{u+v}{2}\right)\sin^{2}\operatorname{am}\left(\frac{u+v}{2}-a\right)}.
\end{gathered}
\]
Hence, changing $a$ into $-a$, we derive:
\[
\begin{gathered}
\frac{1-k^{2}\sin^{2}\operatorname{am}\left(\frac{u+v}{2}\right)\sin^{2}\operatorname{am}\left(\frac{u-v}{2}\right)}{1-k^{2}\sin^{4}\operatorname{am}\left(\frac{u+v}{2}\right)}\left\{1+k^{2}\sin\operatorname{am}a.\sin\operatorname{am}u.\sin\operatorname{am}v.\sin\operatorname{am}(u+v+a)\right\}\\[1.5ex]
=\\[1.5ex]
\frac{1-k^{2}\sin^{2}\operatorname{am}\left(\frac{u-v}{2}\right)\sin^{2}\operatorname{am}\left(\frac{u+v}{2}+a\right)}{1-k^{2}\sin^{2}\operatorname{am}\left(\frac{u+v}{2}\right)\sin^{2}\operatorname{am}\left(\frac{u+v}{2}+a\right)},
\end{gathered}
\]
whence, division being performed:
\[
\begin{gathered}
\text{7)}\quad \frac{1-k^{2}\sin\operatorname{am}a.\sin\operatorname{am}u.\sin\operatorname{am}v.\sin\operatorname{am}(u+v-a)}{1-k^{2}\sin\operatorname{am}a.\sin\operatorname{am}u.\sin\operatorname{am}v.\sin\operatorname{am}(u+v+a)}=\\[1.5ex]
\frac{1-k^{2}\sin^{2}\operatorname{am}\left(\frac{u-v}{2}\right)\sin^{2}\operatorname{am}\left(\frac{u+v}{2}-a\right)}{1-k^{2}\sin^{2}\operatorname{am}\left(\frac{u-v}{2}\right)\sin^{2}\operatorname{am}\left(\frac{u+v}{2}+a\right)}.\frac{1-k^{2}\sin^{2}\operatorname{am}\left(\frac{u+v}{2}\right)\sin^{2}\operatorname{am}\left(\frac{u+v}{2}+a\right)}{1-k^{2}\sin^{2}\operatorname{am}\left(\frac{u+v}{2}\right)\sin^{2}\operatorname{am}\left(\frac{u+v}{2}-a\right)},
\end{gathered}
\]\ednote{In formula 7) the denominator of the first fraction is printed with the sign $1-k^{2}$; the author's Corrigenda (p. 158, line 10) read $1+k^{2}$ instead.}
which is the sought transformation of the expression proposed by Cl.\ Legendre into expression 2).

Formula 6), on putting $u$, $a$, $v$ in place of $\dfrac{u-v}{2}$, $\dfrac{u+v}{2}$, $\dfrac{u+v}{2}-a$, can be represented also thus:
\[
\begin{gathered}
\text{8)}\quad 1-k^{2}\sin\operatorname{am}(a+u).\sin\operatorname{am}(a-u).\sin\operatorname{am}(a+v).\sin\operatorname{am}(a-v)=\\[1.5ex]
\frac{\left\{1-k^{2}\sin^{4}\operatorname{am}a\right\}\left\{1-k^{2}\sin^{2}\operatorname{am}u.\sin^{2}\operatorname{am}v\right\}}{\left\{1-k^{2}\sin^{2}\operatorname{am}a.\sin^{2}\operatorname{am}u\right\}\left\{1-k^{2}\sin^{2}\operatorname{am}a.\sin^{2}\operatorname{am}v\right\}},
\end{gathered}
\]
whence formula 4) flows as a special case, on putting $u=v$.

\origpage{159}
\subsection*{55.}

From the formulae 1), 2), 3), 7) of the preceding §$^{\mathrm{i}}$ it follows:
\[
\begin{gathered}
\text{1)}\quad \Pi(u,a)+\Pi(v,a)-\Pi(u+v,a)=\\[1.5ex]
\frac{1}{2}\log.\frac{\left\{1-k^{2}\sin^{2}\operatorname{am}\left(\frac{u-v}{2}\right).\sin^{2}\operatorname{am}\left(\frac{u+v}{2}-a\right)\right\}\left\{1-k^{2}\sin^{2}\operatorname{am}\left(\frac{u+v}{2}\right).\sin^{2}\operatorname{am}\left(\frac{u+v}{2}+a\right)\right\}}{\left\{1-k^{2}\sin^{2}\operatorname{am}\left(\frac{u-v}{2}\right).\sin^{2}\operatorname{am}\left(\frac{u+v}{2}+a\right)\right\}\left\{1-k^{2}\sin^{2}\operatorname{am}\left(\frac{u+v}{2}\right).\sin^{2}\operatorname{am}\left(\frac{u+v}{2}-a\right)\right\}}=\\[1.5ex]
\frac{1}{4}\log.\frac{\left\{1-k^{2}\sin^{2}\operatorname{am}(u+a).\sin^{2}\operatorname{am}(v+a)\right\}\left\{1-k^{2}\sin^{2}\operatorname{am}a\sin^{2}\operatorname{am}(u+v-a)\right\}}{\left\{1-k^{2}\sin^{2}\operatorname{am}(u-a).\sin^{2}\operatorname{am}(v-a)\right\}\left\{1-k^{2}\sin^{2}\operatorname{am}a\sin^{2}\operatorname{am}(u+v+a)\right\}}=\\[1.5ex]
\frac{1}{2}\log.\frac{1-k^{2}\sin\operatorname{am}a.\sin\operatorname{am}u.\sin\operatorname{am}v.\sin\operatorname{am}(u+v-a)}{1+k^{2}\sin\operatorname{am}a.\sin\operatorname{am}u.\sin\operatorname{am}v.\sin\operatorname{am}(u+v+a)};
\end{gathered}
\]
which is the theorem on the Addition of the Argument of the \emph{Amplitude}. By exactly the same method the other, on the Addition of the Argument of the \emph{Parameter}, can be investigated, but by the aid of the theorem on the reduction of the Parameter to the Amplitude, which formula 4) §.\ 52 furnished us:
\[
\text{IV)}\quad \Pi(u,a)-\Pi(a,u)=u\,Z(a)-a\,Z(u),
\]
the same flows spontaneously from formula 1). For from IV.\ there is:
\[
\begin{gathered}
\Pi(a,u)-\Pi(u,a)=a\,Z(u)-u\,Z(a)\\
\Pi(b,u)-\Pi(u,b)=b\,Z(u)-u\,Z(b)\\
\Pi(a+b,u)-\Pi(u,u+b)=(a+b)Z(u)-u\,Z(a+b),
\end{gathered}
\]\ednote{In the third line the print has $\Pi(u,u+b)$; the author's Corrigenda (p. 159, line 14) read $\Pi(u,a+b)$ instead.}
whence:
\[
\begin{gathered}
\Pi(u,a)+\Pi(u,b)-\Pi(u,a+b)=\\
\Pi(a,u)+\Pi(b,u)-\Pi(a+b,u)+u\left\{Z(a)+Z(b)-Z(a+b)\right\},
\end{gathered}
\]
or since from 1) there is:
\[
\Pi(a,u)+\Pi(b,u)-\Pi(a+b,u)=\frac{1}{2}\log.\frac{1-k^{2}\sin\operatorname{am}u.\sin\operatorname{am}a.\sin\operatorname{am}b.\sin\operatorname{am}(a+b-u)}{1+k^{2}\sin\operatorname{am}u.\sin\operatorname{am}a.\sin\operatorname{am}b.\sin\operatorname{am}(a+b+u)},
\]
further from II.:
\[
Z(a)+Z(b)-Z(a+b)=k^{2}\sin\operatorname{am}a.\sin\operatorname{am}b.\sin\operatorname{am}(a+b),
\]
there is:
\[
\begin{gathered}
\text{2)}\quad \Pi(u,a)+\Pi(u,b)-\Pi(u,a+b)=\\[1.5ex]
k^{2}\sin\operatorname{am}a\sin\operatorname{am}b\sin\operatorname{am}(a+b).u+\frac{1}{2}\log.\frac{1-k^{2}\sin\operatorname{am}u.\sin\operatorname{am}a.\sin\operatorname{am}b.\sin\operatorname{am}(a+b-u)}{1+k^{2}\sin\operatorname{am}u.\sin\operatorname{am}a.\sin\operatorname{am}b.\sin\operatorname{am}(a+b+u)},
\end{gathered}
\]
which is the sought theorem on the Addition of the Argument of the \emph{Parameter}.

\origpage{160}
We derive other, rather memorable formulae by the following consideration. For from theorem III there is:
\[
\begin{gathered}
\left\{\frac{\Theta(u-a).\Theta(v-b)}{\Theta(0)}\right\}^{2}=\frac{\Theta(u+v-a-b).\Theta(u-v-a+b)}{1-k^{2}\sin^{2}\operatorname{am}(u-a).\sin^{2}\operatorname{am}(v-b)}\\[1.5ex]
\left\{\frac{\Theta(u+a).\Theta(v+b)}{\Theta(0)}\right\}^{2}=\frac{\Theta(u+v+a+b).\Theta(u-v+a-b)}{1-k^{2}\sin^{2}\operatorname{am}(u+a).\sin^{2}\operatorname{am}(v+b)}.
\end{gathered}
\]
Now from theorem I there will be:
\[
\begin{gathered}
\Pi(u,a)+\Pi(v,b)=u\,Z(a)+v\,Z(b)+\frac{1}{2}\log.\frac{\Theta(u-a).\Theta(v-b)}{\Theta(u+a).\Theta(v+b)}\\[1.5ex]
\Pi(u+v,a+b)+\Pi(u-v,a-b)=(u+v)Z(a+b)+(u-v)Z(a-b)+\frac{1}{2}\log.\frac{\Theta(u+v-a-b).\Theta(u-v-a+b)}{\Theta(u+v+a+b).\Theta(u-v+a-b)},
\end{gathered}
\]
whence:
\[
\begin{gathered}
\text{3)}\quad \Pi(u+v,a+b)+\Pi(u-v,a-b)-2\Pi(u,a)-2\Pi(v,b)=\\[1.5ex]
(u+v)Z(a+b)+(u-v)Z(a-b)-2u\,Z(a)-2v\,Z(b)+\\[1.5ex]
\frac{1}{2}\log.\frac{1-k^{2}\sin^{2}\operatorname{am}(u-a).\sin^{2}\operatorname{am}(v-b)}{1-k^{2}\sin^{2}\operatorname{am}(u+a).\sin^{2}\operatorname{am}(v+b)},
\end{gathered}
\]
or since there is:
\[
\begin{gathered}
Z(a)+Z(b)-Z(a+b)=k^{2}\sin\operatorname{am}a.\sin\operatorname{am}b.\sin\operatorname{am}(a+b)\\
Z(a)-Z(b)-Z(a-b)=-k^{2}\sin\operatorname{am}a.\sin\operatorname{am}b.\sin\operatorname{am}(a-b),
\end{gathered}
\]
there results from 2), 3):
\[
\begin{gathered}
\text{4)}\quad \Pi(u+v,a+b)+\Pi(u-v,a-b)-2\Pi(u,a)-2\Pi(v,b)=\\[1.5ex]
-k^{2}\sin\operatorname{am}a.\sin\operatorname{am}b\left\{\sin\operatorname{am}(a+b).(u+v)-\sin\operatorname{am}(a-b).(u-v)\right\}\\[1.5ex]
+\frac{1}{2}\log.\frac{1-k^{2}\sin^{2}\operatorname{am}(u-a)\sin^{2}\operatorname{am}(v-b)}{1-k^{2}\sin^{2}\operatorname{am}(u+a)\sin^{2}\operatorname{am}(v+b)}.
\end{gathered}
\]
Interchanging $u$ and $v$, we obtain:
\[
\begin{gathered}
\text{5)}\quad \Pi(u+v,a+b)-\Pi(u-v,a-b)-2\Pi(v,a)-2\Pi(u,b)=\\[1.5ex]
-k^{2}\sin\operatorname{am}a.\sin\operatorname{am}b\left\{\sin\operatorname{am}(a+b).(u+v)+\sin\operatorname{am}(a-b).(u-v)\right\}\\[1.5ex]
+\frac{1}{2}\log.\frac{1-k^{2}\sin^{2}\operatorname{am}(v-a).\sin^{2}\operatorname{am}(u-b)}{1-k^{2}\sin^{2}\operatorname{am}(v+a).\sin^{2}\operatorname{am}(u+b)}.
\end{gathered}
\]
Adding 4) and 5) we obtain:
\[
\begin{gathered}
\text{6)}\quad \Pi(u+v,a+b)-\Pi(u,a)-\Pi(u,b)-\Pi(v,a)-\Pi(v,b)=\\[1.5ex]
-k^{2}\sin\operatorname{am}a\sin\operatorname{am}b\sin\operatorname{am}(a+b).(u+v)\\[1.5ex]
+\frac{1}{4}\log\left\{\frac{1-k^{2}\sin^{2}\operatorname{am}(u-a)\sin^{2}\operatorname{am}(v-b)}{1-k^{2}\sin^{2}\operatorname{am}(u+a)\sin^{2}\operatorname{am}(v+b)}.\frac{1-k^{2}\sin^{2}\operatorname{am}(v-a)\sin^{2}\operatorname{am}(u-b)}{1-k^{2}\sin^{2}\operatorname{am}(v+a)\sin^{2}\operatorname{am}(u+b)}\right\}.
\end{gathered}
\]

\origpage{161}
On putting $v=0$, from 4), 5) there results:
\[
\begin{gathered}
\text{7)}\quad \Pi(u,a+b)+\Pi(u,a-b)-2\Pi(u,a)=\\[1.5ex]
-k^{2}\sin\operatorname{am}a\sin\operatorname{am}b\left\{\sin\operatorname{am}(a+b)-\sin\operatorname{am}(a-b)\right\}u+\frac{1}{2}\log.\frac{1-k^{2}\sin^{2}\operatorname{am}b\sin^{2}\operatorname{am}(u-a)}{1-k^{2}\sin^{2}\operatorname{am}b\sin^{2}\operatorname{am}(u+a)}\\[1.5ex]
\text{8)}\quad \Pi(u,a+b)-\Pi(u,a-b)-2\Pi(u,b)=\\[1.5ex]
-k^{2}\sin\operatorname{am}a\sin\operatorname{am}b\left\{\sin\operatorname{am}(a+b)+\sin\operatorname{am}(a-b)\right\}u+\frac{1}{2}\log.\frac{1-k^{2}\sin^{2}\operatorname{am}a\sin^{2}\operatorname{am}(u-b)}{1-k^{2}\sin^{2}\operatorname{am}a\sin^{2}\operatorname{am}(u+b)}.
\end{gathered}
\]

On putting $b=0$, from 4), 5) there results:
\[
\begin{gathered}
\text{9)}\quad \Pi(u+v,a)+\Pi(u-v,a)-2\Pi(u,a)=\frac{1}{2}\log.\frac{1-k^{2}\sin^{2}\operatorname{am}v\sin^{2}\operatorname{am}(u-a)}{1-k^{2}\sin^{2}\operatorname{am}v\sin^{2}\operatorname{am}(u+a)}\\[1.5ex]
\text{10)}\quad \Pi(u+v,a)-\Pi(u-v,a)-2\Pi(v,a)=\frac{1}{2}\log.\frac{1-k^{2}\sin^{2}\operatorname{am}u.\sin^{2}\operatorname{am}(v-a)}{1-k^{2}\sin^{2}\operatorname{am}u.\sin^{2}\operatorname{am}(v+a)}.
\end{gathered}
\]

\subsection*{REDUCTIONS OF THE EXPRESSIONS $Z(iu)$, $\Theta(iu)$ TO A REAL ARGUMENT. GENERAL REDUCTION OF THE THIRD SPECIES OF ELLIPTIC INTEGRALS, IN WHICH THE ARGUMENTS BOTH OF THE AMPLITUDE AND OF THE PARAMETER ARE IMAGINARY.}

\subsection*{56.}

We return to the Analysis of the functions $Z$, $\Theta$, whose remarkable use in our theory we have proved in the foregoing. Let us inquire into the reduction of the expressions $Z(iu)$, $\Theta(iu)$ to a real argument. We shall first carry this out with the signs in use by Cl$^{\mathrm{o}}$ Legendre, then adapt it to our notations.

We know in the elements §.\ 19.\ p.\ 34 that the equations hold simultaneously:
\[
\sin\phi=i\operatorname{tang}\psi,\quad \frac{d\phi}{\Delta(\phi)}=\frac{i\,d\psi}{\Delta(\psi,k')},\quad F(\phi)=iF(\psi,k').
\]
Hence there is:
\[
d\phi.\Delta(\phi)=\frac{i\,d\psi(1+kk\operatorname{tg}^{2}\psi)}{\Delta(\psi,k')}=\frac{i\,d\psi.\Delta(\psi,k')}{\cos^{2}\psi},
\]
whence, integration being performed:
\[
\int_{0}^{\phi}\Delta(\phi).d\phi=i\left\{\operatorname{tg}\psi\,\Delta(\psi,k')+\int_{0}^{\psi}\frac{k'k'\sin^{2}\psi}{\Delta(\psi,k')}\right\},
\]\ednote{The integral under the brace is printed without the differential $d\psi$; an apparent omission by the printer.}

\origpage{162}
or:
\[
\text{1)}\quad E(\phi)=i\left\{\operatorname{tg}\psi\,\Delta(\psi,k')+F(\psi,k')-E(\psi,k')\right\}.
\]
Multiplying by $\dfrac{d\phi}{\Delta(\phi)}=\dfrac{i\,d\psi}{\Delta(\psi,k')}$ and integrating we obtain:
\[
\text{2)}\quad \int_{0}^{\phi}\frac{E(\phi).d\phi}{\Delta(\phi)}=\log\cos\psi-\frac{1}{2}\left\{F(\psi,k')\right\}^{2}+\int_{0}^{\psi}\frac{E(\psi,k')}{\Delta(\psi,k')}.
\]\ednote{The last integral is printed without the differential $d\psi$; an apparent omission by the printer.}

From equation 1) it follows:
\[
\frac{F^{\mathrm{I}}E(\phi)-E^{\mathrm{I}}F(\phi)}{i}=F^{\mathrm{I}}\operatorname{tg}\psi\,\Delta(\psi,k')-\left\{F^{\mathrm{I}}E(\psi,k')+(E^{\mathrm{I}}-F^{\mathrm{I}})F(\psi,k')\right\}.
\]
Now let the excellent theorem of Cl$^{\mathrm{i}}$ Legendre be noted (p.\ 61):
\[
F^{\mathrm{I}}E^{\mathrm{I}}(k')+F^{\mathrm{I}}(k')E^{\mathrm{I}}-F^{\mathrm{I}}F^{\mathrm{I}}(k')=\frac{\pi}{2},
\]
whence:
\[
F^{\mathrm{I}}E(\psi,k')+(E^{\mathrm{I}}-F^{\mathrm{I}})F(\psi,k')=\frac{F^{\mathrm{I}}}{F^{\mathrm{I}}(k')}\left\{F^{\mathrm{I}}(k')E(\psi,k')-E^{\mathrm{I}}(k')F(\psi,k')\right\}+\frac{\pi F(\psi,k')}{2F^{\mathrm{I}}(k')},
\]
and therefore:
\[
\text{3)}\quad \frac{F^{\mathrm{I}}E(\phi)-E^{\mathrm{I}}F(\phi)}{iF^{\mathrm{I}}}=\operatorname{tg}\psi\,\Delta(\psi,k')-\frac{F^{\mathrm{I}}(k')E(\psi,k')-E^{\mathrm{I}}(k')F(\psi,k')}{F^{\mathrm{I}}(k')}-\frac{\pi F(\psi,k')}{2F^{\mathrm{I}}F^{\mathrm{I}}(k')}
\]

From our notation we had:
\[
\phi=\operatorname{am}(iu),\quad \psi=\operatorname{am}(u,k'),\quad F(\phi)=iu,\quad F(\psi,k')=u;
\]
further:
\[
\frac{F^{\mathrm{I}}E(\phi)-E^{\mathrm{I}}F(\phi)}{F^{\mathrm{I}}}=Z(iu,k);\quad \frac{F^{\mathrm{I}}(k')E(\psi,k')-E^{\mathrm{I}}(k')F(\psi,k')}{F^{\mathrm{I}}(k')}=Z(u,k'),
\]
whence equation 3) is represented thus:
\[
\text{4)}\quad iZ(iu,k)=-\operatorname{tg}\operatorname{am}(u,k')\,\Delta\operatorname{am}(u,k')+\frac{\pi uu}{2KK'}+Z(u,k').
\]
Hence by integrating there results:
\[
\int_{0}^{u}i\,du\,Z(iu,k)=\log\cos\operatorname{am}(u,k')+\frac{\pi uu}{4KK'}+\int_{0}^{u}Z(u,k')\,du,
\]

\origpage{163}
or, since $\displaystyle\int_{0}^{u}du\,Z(u)=\log\frac{\Theta(u)}{\Theta(0)}$:
\[
\text{5)}\quad \frac{\Theta(iu,k)}{\Theta(0,k)}=e^{\frac{\pi uu}{4KK'}}\cos\operatorname{am}(u,k')\frac{\Theta(u,k')}{\Theta(0,k')}.
\]
Formulae 4), 5) reduce the functions $Z(iu)$, $\Theta(iu)$ to a real argument.

\subsection*{57.}

Let $u$ be changed in 5) into $u+2K'$; there results:
\[
\frac{\Theta(iu+2iK')}{\Theta(0)}=-e^{\frac{\pi(u+2K')^{2}}{4KK'}}\cos\operatorname{am}(u,k')\frac{\Theta(u,k')}{\Theta(0,k')}{}^{*)}=-e^{\frac{\pi(K'+u)}{K}}\frac{\Theta(iu)}{\Theta 0},
\]\ednote{The denominator is printed $\Theta 0$ without parentheses.}
or, putting $u$ in place of $iu$:
\[
\text{1)}\quad \Theta(u+2iK')=-e^{\frac{\pi(K'-iu)}{K}}\Theta(u).
\]
Let $u+K'$ be put in 5) in place of $u$: since
\[
\cos\operatorname{am}(u+K',k')=-\frac{k\sin\operatorname{am}(u,k')}{\Delta\operatorname{am}(u,k')}
\]
and this:
\[
\Theta(u+K',k')=\frac{\Delta\operatorname{am}(u,k')}{\sqrt{k}}.\Theta(u,k'),\quad\text{See }\S\text{.\ 53.\ 9)}
\]
there results:
\[
\begin{gathered}
\frac{\Theta(iu+iK')}{\Theta(0)}=-e^{\frac{\pi(u+K')^{2}}{4KK'}}\sqrt{k}\sin\operatorname{am}(u,k')\frac{\Theta(u,k')}{\Theta(0,k')}\\[1.5ex]
=-e^{\frac{\pi(2u+K')}{4K}}\sqrt{k}\operatorname{tg}\operatorname{am}(u,k')\frac{\Theta(iu)}{\Theta(0)},
\end{gathered}
\]
whence, putting again $u$ in place of $iu$:
\[
\text{2)}\quad \Theta(u+iK')=i\,e^{\frac{\pi(K'-2iu)}{4K}}\sqrt{k}\sin\operatorname{am}(u)\,\Theta(u).
\]

\textbf{*)} For $\Theta(u+2K,k)=\Theta(u)$, and therefore also $\Theta(u+2K',k')=\Theta(u,k')$.

\origpage{164}
Taking logarithms and differentiating, from 1), 2) there results:
\[
\begin{gathered}
\text{3)}\quad Z(u+2iK')=\frac{-i\pi}{K}+Z(u)\\[1.5ex]
\text{4)}\quad Z(u+iK')=\frac{-i\pi}{2K}+\operatorname{cotg}\operatorname{am}(u)\,\Delta\operatorname{am}(u)+Z(u).
\end{gathered}
\]
Putting $u=0$, from 1) - 4) there results:
\[
\text{5)}\quad \left\{\begin{gathered}
\Theta(2iK')=-e^{\frac{\pi K'}{K}}\Theta(0),\quad \Theta(iK')=0\\[1ex]
Z(2iK')=\ \ ;\quad Z(iK')=\infty.
\end{gathered}\right.
\]\ednote{The value of $Z(2iK')$ is missing in the print, which leaves a blank after the equals sign; only a hand-written entry stands there, which is not transcribed.}
Formulae 1), 2) find an excellent confirmation from the nature of the infinite product into which we expand the function $\Theta$:
\[
\begin{gathered}
\text{6)}\quad \frac{\Theta\left(\frac{2Kx}{\pi}\right)}{\Theta(0)}=\frac{(1-2q\cos2x+q^{2})(1-2q^{3}\cos2x+q^{6})(1-2q^{5}\cos2x+q^{10})\ldots}{\left\{(1-q)(1-q^{3})(1-q^{5})\ldots\right\}^{2}}=\\[1.5ex]
\frac{\left\{(1-qe^{2ix})(1-q^{3}e^{2ix})(1-q^{5}e^{2ix}).\,.\right\}\left\{(1-qe^{-2ix})(1-q^{3}e^{-2ix})(1-q^{5}e^{-2ix}).\,.\right\}}{\left\{(1-q)(1-q^{3})(1-q^{5})\ldots\right\}^{2}}.
\end{gathered}
\]
For when $x$ is changed into $x+\dfrac{i\pi K'}{K}$, upon which $e^{ix}$ goes over into $q\,e^{ix}$, the product
\[
\left\{(1-qe^{2ix})(1-q^{3}e^{2ix})(1-q^{5}e^{2ix})\ldots\right\}\left\{(1-qe^{-2ix})(1-q^{3}e^{-2ix})(1-q^{5}e^{-2ix})\ldots\right\}
\]
goes over into this:
\[
\frac{1}{qe^{2ix}}\left\{(1-qe^{2ix})(1-q^{3}e^{2ix})(1-q^{5}e^{2ix})\ldots\right\}\left\{(1-qe^{-2ix})(1-q^{3}e^{-2ix})(1-q^{5}e^{-2ix})\ldots\right\},
\]
whence:
\[
\text{7)}\quad \Theta\left(\frac{2Kx}{\pi}+2iK'\right)=\frac{\Theta\left(\frac{2Kx}{\pi}\right)}{qe^{2ix}}.
\]
But when $x$ is changed into $x+\dfrac{i\pi K'}{2K}$, $e^{ix}$ goes over into $\sqrt{q}\,e^{ix}$, whence the product
\[
\left\{(1-qe^{2ix})(1-q^{3}e^{2ix})(1-q^{5}e^{2ix})\ldots\right\}\left\{(1-qe^{-2ix})(1-q^{3}e^{-2ix})(1-q^{5}e^{-2ix})\ldots\right\}
\]
goes over into this:
\[
\begin{gathered}
(i-e^{-2ix})\left\{(1-q^{2}e^{2ix})(1-q^{4}e^{2ix})\ldots\right\}\left\{(1-q^{2}e^{-2ix})(1-q^{4}e^{-2ix})\ldots\right\}=\\[1.5ex]
\frac{i}{e^{ix}}.2\sin x(1-2q^{2}\cos2x+q^{4})(1-2q^{4}\cos2x+q^{8})(1-q^{6}\cos2x+q^{12}).\,.\,.\,.
\end{gathered}
\]\ednote{The first factor is printed $(i-e^{-2ix})$ with the letter i where the numeral 1 is expected. The factor $(1-q^{6}\cos2x+q^{12})$ is printed without the coefficient 2 of the cosine term; an apparent misprint for $(1-2q^{6}\cos2x+q^{12})$.}

\origpage{165}
But in §.\ 36 we gave the formula:
\[
\sin\operatorname{am}\frac{2Kx}{\pi}=\frac{1}{\sqrt{k}}.\frac{2\sqrt[4]{q}\sin x(1-2q^{2}\cos2x+q^{4})(1-2q^{4}\cos2x+q^{8})\ldots}{(1-2q\cos2x+q^{2})(1-2q^{3}\cos2x+q^{6})(1-2q^{5}\cos2x+q^{10})\ldots},
\]
whence we see that there will be:
\[
\text{8)}\quad \Theta\left(\frac{2Kx}{\pi}+iK'\right)=\frac{i\sqrt{k}\sin\operatorname{am}\frac{2Kx}{\pi}\,\Theta\left(\frac{2Kx}{\pi}\right)}{\sqrt[4]{q}\,e^{ix}}.
\]
Formulae 7), 8), however, on putting $\dfrac{2Kx}{\pi}=u$, agree with formulae 1), 2).

From formula 9) §.\ 53:
\[
\Theta(u+K)=\frac{\Delta\operatorname{am}u}{\sqrt{k'}}.\Theta(u),
\]
putting $iu$ in place of $u$, it follows:
\[
\Theta(iu+K)=\frac{\Delta\operatorname{am}(u,k')}{\sqrt{k'}\cos\operatorname{am}(u,k')}.\Theta(iu),
\]
whence from 5) §.\ 56:
\[
\frac{\Theta(iu+K)}{\Theta(0)}=\frac{1}{\sqrt{k'}}e^{\frac{\pi uu}{4KK'}}\Delta\operatorname{am}(u,k').\frac{\Theta(u,k')}{\Theta(0,k')},
\]
or from the cited formula 9) §.\ 53:
\[
\text{9)}\quad \frac{\Theta(iu+K)}{\Theta(0)}=\sqrt{\frac{k}{k'}}\,e^{\frac{\pi uu}{4KK'}}\frac{\Theta(u+K',k')}{\Theta(0,k')}.
\]
Hence, taking logarithms and differentiating, we obtain:
\[
\text{10)}\quad iZ(iu+K)=\frac{\pi u}{2KK'}+Z(u+K',k').
\]

\subsection*{58.}

The formulae found in §§.\ 56.\ 57 are easily applied to the Analysis of the functions $\Pi$ in the cases in which the Arguments either of the Amplitude or of the Parameter or of both are imaginary.

Let us first demonstrate that the expression $\Pi(u,a+iK')$ can be reduced to $\Pi(u,a)$, whence it is evident, on putting $n=-k^{2}\sin^{2}\operatorname{am}a$, that the integrals

\origpage{166}
\[
\int_{0}^{\phi}\frac{d\phi}{\left\{1+n\sin^{2}\phi\right\}\Delta(\phi)},\quad \int_{0}^{\phi}\frac{d\phi}{\left\{1+\frac{k^{2}}{n}\sin^{2}\phi\right\}\Delta(\phi)}
\]
depend one on the other; which is a notable theorem put forward by Cl.\ Legendre, Chap.\ XV.

We found:
\[
\Pi(u,a+iK')=u\,Z(a+iK')+\frac{1}{2}\log\frac{\Theta(a-u+iK')}{\Theta(a+u+iK')}.
\]
Now from 2), 4) §.\ 57 there is:
\[
\begin{gathered}
\frac{\Theta(a-u+iK')}{\Theta(a+u+iK')}=e^{\frac{i\pi u}{K}}\frac{\sin\operatorname{am}(a-u)}{\sin\operatorname{am}(a+u)}.\frac{\Theta(a-u)}{\Theta(a+u)}\\[1.5ex]
u\,Z(a+iK')=\frac{-i\pi u}{2K}+u\operatorname{cotg}\operatorname{am}a\,\Delta\operatorname{am}a+u\,Z(a),
\end{gathered}
\]
whence, the terms $\dfrac{i\pi u}{2K}-\dfrac{i\pi u}{2K}$ cancelling one another:
\[
\text{1)}\quad \Pi(u,a+iK')=\Pi(u,a)+u\operatorname{cotg}\operatorname{am}a\,\Delta\operatorname{am}a+\frac{1}{2}\log\frac{\sin\operatorname{am}(a-u)}{\sin\operatorname{am}(a+u)}.
\]
Let us put in this formula $ia$ in place of $a$; there results:
\[
\begin{gathered}
\operatorname{cotg}\operatorname{am}(ia)\,\Delta\operatorname{am}(ia)=\frac{-i\Delta\operatorname{am}(a,k')}{\sin\operatorname{am}(a,k')\cos\operatorname{am}(a,k')}\\[1.5ex]
\frac{\sin\operatorname{am}(ia-u)}{\sin\operatorname{am}(ia+u)}=\frac{\Delta\operatorname{am}u-\operatorname{cotg}\operatorname{am}(ia)\,\Delta\operatorname{am}(ia)\operatorname{tg}\operatorname{am}u}{\Delta\operatorname{am}u+\operatorname{cotg}\operatorname{am}(ia)\,\Delta\operatorname{am}(ia)\operatorname{tg}\operatorname{am}u},
\end{gathered}
\]
or, putting for brevity's sake:
\[
\frac{\Delta\operatorname{am}(a,k')}{\sin\operatorname{am}(a,k')\cos\operatorname{am}(a,k')}=\sqrt{\alpha},
\]
there results:
\[
\frac{\sin\operatorname{am}(ia-u)}{\sin\operatorname{am}(ia+u)}=\frac{\Delta\operatorname{am}u+i\sqrt{\alpha}\operatorname{tg}\operatorname{am}u}{\Delta\operatorname{am}u-i\sqrt{\alpha}\operatorname{tg}\operatorname{am}u},
\]
whence 1) goes over into
\[
\text{2)}\quad \frac{\Pi(u,ia+iK')-\Pi(u,ia)}{i}=-\sqrt{\alpha}.u+\operatorname{Arc.}\operatorname{tg}.\frac{\sqrt{\alpha}\operatorname{tg}\operatorname{am}u}{\Delta\operatorname{am}u};
\]
which agrees with the formula $f')$ exhibited by Cl.\ Legendre.

\origpage{167}
\subsection*{59.}

Other formulae, fundamental for the reduction of an imaginary Argument to a real one, we obtain from 9), 10) §.\ 57. Of these I first note this one, by which the imaginary Arguments both of the Amplitude and of the Parameter are reduced to real Arguments:
\[
\text{1)}\quad \Pi(iu,ia+K)=\Pi(u,a+K',k'),
\]
which is demonstrated in this manner. For we have:
\[
\Pi(iu,ia+K)=iu\,Z(ia+K)+\frac{1}{2}\log\frac{\Theta(ia-iu+K)}{\Theta(ia+iu+K)};
\]
further from 10) §.\ 57:
\[
iu\,Z(ia+K)=\frac{\pi ua}{2KK'}+u\,Z(a+K',k'),
\]
from 9) §.\ 57:
\[
\begin{gathered}
\frac{\Theta(ia-iu+K)}{\Theta(0,k)}=\sqrt{\frac{k}{k'}}\,e^{\frac{\pi(a-u)^{2}}{4KK'}}\frac{\Theta(a-u+K',k')}{\Theta(0,k')}\\[1.5ex]
\frac{\Theta(ia+iu+K)}{\Theta(0,k)}=\sqrt{\frac{k}{k'}}\,e^{\frac{\pi(a+u)^{2}}{4KK'}}\frac{\Theta(a+u+K',k')}{\Theta(0,k')},
\end{gathered}
\]
whence:
\[
\frac{\Theta(ia-iu+K)}{\Theta(ia+iu+K)}=e^{\frac{-\pi au}{KK'}}\frac{\Theta(a-u+K',k')}{\Theta(a+u+K',k')},
\]
and therefore, the terms $\dfrac{\pi ua}{2KK'}-\dfrac{\pi ua}{2KK'}$ cancelling one another,
\[
\Pi(iu,ia+K)=u\,Z(a+K',k')+\frac{1}{2}\log\frac{\Theta(a-u+K',k')}{\Theta(a+u+K',k')}=\Pi(u,a+K',k'),
\]
which was to be demonstrated.

Changing in 1) $a$ into $-ia$, there results:
\[
\text{2)}\quad \Pi(iu,a+K)=-\Pi(u,ia+K',k').
\]

Formula 1) is also easily proved by considering the integral itself, by which we defined the function $\Pi$:
\[
\Pi(u,a)=\int_{0}^{u}\frac{k^{2}\sin\operatorname{am}a.\cos\operatorname{am}a.\Delta\operatorname{am}a.\sin^{2}\operatorname{am}u.du}{1-k^{2}\sin^{2}\operatorname{am}a.\sin^{2}\operatorname{am}u},
\]

\origpage{168}
whence:
\[
\Pi(iu,ia+K)=\int_{0}^{u}\frac{ik^{2}\sin\operatorname{am}(ia+K).\cos\operatorname{am}(ia+K).\Delta\operatorname{am}(ia+K).\sin^{2}\operatorname{am}(iu).du}{1-k^{2}\sin^{2}\operatorname{am}(ia+K).\sin^{2}\operatorname{am}(iu)}.
\]
For from the formulae of §$^{\mathrm{i}}$ 19 there is:
\[
\begin{gathered}
\sin\operatorname{am}(ia+K)=\sin\operatorname{coam}(ia)=\frac{\Delta\operatorname{coam}(a,k')}{k}=\frac{\Delta\operatorname{am}(a+K',k')}{k}\\[1.5ex]
\cos\operatorname{am}(ia+K)=-\cos\operatorname{coam}(ia)=\frac{-ik'}{k}\cos\operatorname{coam}(a,k')=\frac{ik'}{k}\cos\operatorname{am}(a+K',k')\\[1.5ex]
\Delta\operatorname{am}(ia+K)=\Delta\operatorname{coam}(ia)=k'\sin\operatorname{coam}(a,k')=k'\sin\operatorname{am}(a+K',k'),
\end{gathered}
\]
whence:
\[
\begin{gathered}
ikk\sin\operatorname{am}(ia+K)\cos\operatorname{am}(ia+K)\Delta\operatorname{am}(ia+K)=\\
-k'k'\sin\operatorname{am}(a+K',k')\cos\operatorname{am}(a+K',k')\Delta\operatorname{am}(a+K',k').
\end{gathered}
\]
Further:
\[
\begin{gathered}
\frac{\sin^{2}\operatorname{am}(iu)}{1-k^{2}\sin^{2}\operatorname{am}(ia+K)\sin^{2}\operatorname{am}(iu)}=\frac{-\operatorname{tg}^{2}\operatorname{am}(u,k')}{1+\Delta^{2}\operatorname{am}(a+K',k')\operatorname{tg}^{2}\operatorname{am}(u,k')}=\\[1.5ex]
\frac{-\sin^{2}\operatorname{am}(u,k')}{\cos^{2}\operatorname{am}(u,k')+\Delta^{2}\operatorname{am}(a+K',k')\sin^{2}\operatorname{am}(u,k')}=\frac{-\sin^{2}\operatorname{am}(u,k')}{1-k'k'\sin^{2}\operatorname{am}(a+K',k')\sin^{2}\operatorname{am}(u,k')},
\end{gathered}
\]
whence:
\[
\begin{gathered}
\Pi(iu.\,ia+K)=\\[1.5ex]
\int_{0}^{u}\frac{k'k'\sin\operatorname{am}(a+K',k').\cos\operatorname{am}(a+K',k').\Delta\operatorname{am}(a+K',k').\sin^{2}\operatorname{am}(u,k').du}{1-k'k'\sin^{2}\operatorname{am}(a+K',k')\sin^{2}\operatorname{am}(u,k')},
\end{gathered}
\]\ednote{The first argument and the second are separated by a period, printed $\Pi(iu.\,ia+K)$; an apparent misprint for a comma.}
or:
\[
\Pi(iu,ia+K)=\Pi(u,a+K',k'),
\]
which was to be demonstrated.

From formulae 9), 10) §.\ 57 we can confirm, in a manner similar to 1), the following formula, which teaches that pairs of functions of imaginary Argument of the Parameter, whose Moduli are each the Complement of the other, can be reduced to one another:
\[
\begin{gathered}
\text{3)}\quad i\Pi(u,ia+K)+i\Pi(a,iu+K',k')=\\[1.5ex]
\frac{\pi au}{2KK'}+u\,Z(a+K',k')+a\,Z(u+K,k).
\end{gathered}
\]
For:
\[
\begin{gathered}
i\Pi(u,ia+K)=iu\,Z(ia+K)+\frac{i}{2}\log.\frac{\Theta(ia+K-u)}{\Theta(ia+K+u)}\\[1.5ex]
i\Pi(a,iu+K',k')=ia\,Z(iu+K',k')+\frac{i}{2}\log.\frac{\Theta(iu+K'-a,k')}{\Theta(iu+K'+a,k')}.
\end{gathered}
\]

\origpage{169}
Now:
\[
\begin{gathered}
\frac{\Theta(ia+K-u)}{\Theta(0)}=\frac{\Theta\{i(a+iu)+K\}}{\Theta(0)}=\sqrt{\frac{k}{k'}}\,e^{\frac{\pi(a+iu)^{2}}{4KK'}}\,\frac{\Theta(a+iu+K',k')}{\Theta(0,k')}\\[1.5ex]
\frac{\Theta(ia+K+u)}{\Theta(0)}=\frac{\Theta\{i(a-iu)+K\}}{\Theta(0)}=\sqrt{\frac{k}{k'}}\,e^{\frac{\pi(a-iu)^{2}}{4KK'}}\,\frac{\Theta(a-iu+K',k')}{\Theta(0,k')},
\end{gathered}
\]
whence, since $\Theta(u+K)=\Theta(K-u)$:
\[
\frac{\Theta(ia+K-u)}{\Theta(ia+K+u)}=e^{\frac{\pi au}{KK'}}\,\frac{\Theta(iu+K'+a,k')}{\Theta(iu+K'-a,k')}
\]
and therefore:
\[
\frac{i}{2}\log.\frac{\Theta(ia+K-u)}{\Theta(ia+K+u)}+\frac{i}{2}\log\frac{\Theta(iu+K'-a,k')}{\Theta(iu+K'+a,k')}=-\frac{\pi au}{2KK'}
\]
Further:
\[
\begin{gathered}
iuZ(ia+K)=\frac{\pi au}{2KK'}+uZ(a+K',k')\\[1.5ex]
iuZ(iu+K',k')=\frac{\pi au}{2KK'}+aZ(u+K,k),
\end{gathered}
\]\ednote{The second equation is printed with $iuZ(iu+K',k')$ on the left; the author's Corrigenda give $iaZ$, which the context requires.}
whence:
\[
i\Pi(u,ia+K)+i\Pi(a,iu+K',k')=\frac{\pi au}{2KK'}+uZ(a+K',k')+aZ(u+K,k);
\]
Q.\ E.\ D.

\subsection*{60.}

It is evident from the formulae:
\[
\begin{gathered}
\sin\operatorname{am}(K+iu)=\frac{1}{k}\,\Delta\operatorname{coam}(u,k')\\[1.5ex]
\sin\operatorname{am}(u+iK')=\frac{1}{k}.\,\frac{1}{\sin\operatorname{am}u},
\end{gathered}
\]
that the Argument $u$, which, while $\sin\operatorname{am}u$ increases from 0 up to 1, passes from 0 to $K$, assumes an imaginary value of the form $K+iv$ when $\sin\operatorname{am}u$ continues to increase from 1 up to $\frac{1}{k}$, so that simultaneously $v$ increases from 0 up to $K'$; then, as $\sin\operatorname{am}u$ increases from $\frac{1}{k}$

\origpage{170}
up to $\infty$, $u$ assumes the form $v+iK'$, so that simultaneously $v$ decreases from $K$ up to 0 ${}^{*)}$.

Hence we see, if in the third species of Elliptic Integrals, which is contained in the scheme:
\[
\int_{0}^{\phi}\frac{d\phi}{(1+n\sin^{2}\phi)\Delta(\phi)},
\]
one puts, as we did, $n=-k^{2}\sin^{2}\operatorname{am}a$, whenever $n$ is negative
\[
\begin{gathered}
\text{between }0\text{ and }-kk,\ \text{one must put }n=-k^{2}\sin^{2}\operatorname{am}a\\[1ex]
\text{- }-kk\text{ and }-1,\quad\text{-}\quad\text{-}\quad n=-k^{2}\sin^{2}\operatorname{am}(ia+K)\\[1ex]
\text{- }-1\text{ and }-\infty,\quad\text{-}\quad\text{-}\quad n=-k^{2}\sin^{2}\operatorname{am}(a+iK'),
\end{gathered}
\]
where $a$ denotes a real quantity. Further, since $-kk\sin^{2}\operatorname{am}(ia)=kk\operatorname{tg}^{2}\operatorname{am}(a,k')$, it is evident that, whenever $n$ is any positive quantity, one must put:
\[
n=-kk\sin^{2}\operatorname{am}(ia).
\]
Hence we have obtained four classes of Elliptic Integrals of the third species, which correspond to the schemes that take on the Arguments
\[
\text{1)}\ a,\quad\text{2)}\ ia+K,\quad\text{3)}\ a+iK',\quad\text{4)}\ ia,
\]
of which the first three belong to negative $n$, the fourth to positive.

But by formula 1) §. 58 we see that the function $\Pi(u,a+iK')$ is reduced to $\Pi(u,a)$, or the third class, in which $n$ is between $-1$ and $-\infty$, is reduced to the first, in which $n$ is between 0 and $-kk$. Further, from formula 11) §. 53 ${}^{**)}$, the function $\Pi(u,ia)$ is always reduced

\textbf{*)} There will be obtained simultaneously:
\[
\begin{gathered}
\sin\operatorname{am}u=0,\quad\frac{1}{\sqrt{1+k'}},\quad 1,\quad\frac{1}{\sqrt{k}},\quad\frac{1}{k},\quad\frac{1}{2}\sqrt{1+k'},\quad\infty\\[1.5ex]
u=0,\quad\frac{K}{2},\quad K,\quad K+\frac{iK'}{2},\quad K+iK',\quad\frac{K}{2}+iK',\quad iK'.
\end{gathered}
\]\ednote{The sixth value is printed as $\frac{1}{2}\sqrt{1+k'}$; since $\sin\operatorname{am}(K/2+iK')=\sqrt{1+k'}/k$, the factor $\frac{1}{k}$ is expected instead of $\frac{1}{2}$, an apparent misprint.}

\textbf{**)} This formula, namely, on putting $ia$ in place of $a$, goes over into the following:
\[
\frac{\Pi(u,ia+K)-\Pi(u,ia)}{i}=-\alpha+\operatorname{Arc}\operatorname{tg}\left\{\alpha\sin\operatorname{am}u.\sin\operatorname{coam}u\right\},
\]
if one puts $\alpha=\dfrac{kk\operatorname{tg}\operatorname{am}(a,k')}{\Delta\operatorname{am}(a,k')}$. This transformation succeeds easily by the elementary formulae of §$^{\mathrm{i}}$ 19.

\origpage{171}
to $\Pi(u,ia+K)$, or the fourth class, in which $n$ is positive, to the second, in which $n$ is negative between $-kk$ and $-1$. Whence we have now obtained the theorem, \emph{that the proposed integral}
\[
\int_{0}^{\phi}\frac{d\phi}{\left\{1+n\sin^{2}\phi\right\}\Delta(\phi)},
\]
\emph{whatever real quantity $n$ be, positive or negative, can always be reduced to a similar integral, in which $n$ is negative between $0$ and $-1$.} This is an excellent discovery of Cl$^{\mathrm{i}}$ Legendre.

Now let us consider the general case, in which both the Amplitude and the Parameter have any imaginary form: it is agreed that this case comprises the expression
\[
\Pi(u+iv,a+ib),
\]
where $u$, $v$, $a$, $b$ denote real quantities. But from the formulae of §$^{\mathrm{i}}$ 55 we see that an expression of this kind is reduced to these four:
\[
\text{1)}\ \Pi(u,a),\quad\text{2)}\ \Pi(iv,ib),\quad\text{3)}\ \Pi(u,ib),\quad\text{4)}\ \Pi(iv,a),
\]
or, if you like, to these four:
\[
\begin{gathered}
\text{1)}\ \Pi(u,a-K),\quad\text{2)}\ \Pi(iv,ib+K)\\[1ex]
\text{3)}\ \Pi(u,ib+K),\quad\text{4)}\ \Pi(iv,a-K).
\end{gathered}
\]
For in general the expression $\Pi(u+v,a+b)$ reduces to the expressions $\Pi(u,a)$, $\Pi(v,b)$, $\Pi(u,b)$, $\Pi(v,a)$, from which the four proposed result, if in place of $v$ you put $iv$, and in place of $a$, $b$ however $a-K$ and $K+ib$. Further, from formulae 1), 2) §$^{\mathrm{i}}$ 59 there is:
\[
\begin{gathered}
\Pi(iv,ib+K)=\Pi(v,b+K',k')\\[1ex]
\Pi(iv,a-K)=-\Pi(v,ia+K',k'),
\end{gathered}
\]
whence expressions 1), 2) are reduced to the first class \ednote{Printed without the word in before classem; the author's Corrigenda read in classem.} $\Pi(u,a)$, expressions 3), 4) to the second class $\Pi(u,ia+K)$; which furnishes us with the

\subsection*{THEOREM.}

\emph{The proposed integral} of the form
\[
\int_{0}^{\phi}\frac{d\phi}{(1+n\sin^{2}\phi)\Delta(\phi)},
\]
\emph{whatever $n$ and $\phi$ be, whether real or imaginary, can be reduced to similar integrals, in which both $\phi$ is real and $n$ is real negative between $0$ and $-1$.}

\origpage{172}
And this theorem \ednote{Printed E; the author's Corrigenda read Et.} is due to Cl$^{\mathrm{o}}$ Legendre, except that he contemplated only real Amplitudes.

By formulae 4), 5) §. 56, $\Pi(u+v,a+b)+\Pi(u-v,a-b)$ is reduced to $\Pi(u,a)$ and $\Pi(v,b)$, and $\Pi(u+v,a+b)-\Pi(u-v,a-b)$ to $\Pi(u,b)$ and $\Pi(v,a)$. Hence it is evident, on putting
\[
\begin{gathered}
\Pi(u+i,a+ib)+\Pi(u-iv,a-ib)=L\\[1.5ex]
\frac{\Pi(u+iv,a+ib)-\Pi(u-iv,a-ib)}{i}=M,
\end{gathered}
\]\ednote{The first line is printed with $\Pi(u+i,a+ib)$; the author's Corrigenda give $\Pi(u+iv,a+ib)$.}
that $L$ depends on the functions $\Pi(u,a-K)$, $\Pi(iv,ib+K)$, $M$ on the functions $\Pi(u,ib+K)$, $\Pi(iv,a-K)$, and therefore that $L$ returns to the first class, $M$ to the second class.

These are the foundations of the theory of the third species of Elliptic Integrals, deduced from new principles. Others will be seen below.

\subsection*{ELLIPTIC FUNCTIONS ARE FRACTIONAL FUNCTIONS. ON THE FUNCTIONS $H$, $\Theta$, WHICH TAKE THE PLACE OF NUMERATOR AND DENOMINATOR.}

\subsection*{61.}

The expansions exhibited in §. 35 declare the genuine nature of the Elliptic functions, namely that they are fractional functions, as we already know from the elements, both vanishing and going off to infinity for innumerable mutually different values of the Argument. Now in the foregoing we have been led to a function which constitutes the denominator of the fraction into which we expand $\sin\operatorname{am}\dfrac{2Kx}{\pi}=$
\[
\frac{1}{\sqrt{k}}.\,\frac{2\sqrt[4]{q}\sin x(1-2q^{2}\cos 2x+q^{4})(1-2q^{4}\cos 2x+q^{8})(1-2q^{6}\cos 2x+q^{12})\ldots}{(1-2q\cos 2x+q^{2})(1-2q^{3}\cos 2x+q^{6})(1-2q^{5}\cos 2x+q^{10})\ldots},
\]
I mean the function
\[
\frac{\Theta\left(\frac{2Kx}{\pi}\right)}{\Theta(0)}=\frac{(1-2q\cos 2x+q^{2})(1-2q^{3}\cos 2x+q^{6})(1-2q^{5}\cos 2x+q^{10})\ldots}{\left\{(1-q)(1-q^{3})(1-q^{5})(1-q^{7})\ldots\right\}^{2}}.
\]
Now let us also denote the numerator by a particular character, and put:
\[
\frac{H\left(\frac{2Kx}{\pi}\right)}{\Theta(0)}=\frac{2\sqrt[4]{q}\sin x(1-2q^{2}\cos 2x+q^{4})(1-2q^{4}\cos 2x+q^{8})(1-2q^{6}\cos 2x+q^{12})\ldots}{\left\{(1-q)(1-q^{3})(1-q^{5})(1-q^{7})\ldots\right\}^{2}},
\]

\origpage{173}
there will be:
\[
\sin\operatorname{am}\frac{2Kx}{\pi}=\frac{1}{\sqrt{k}}.\,\frac{H\left(\frac{2Kx}{\pi}\right)}{\Theta\left(\frac{2Kx}{\pi}\right)}.
\]
Calling in the remaining expansions given in §. 35, we find:
\[
\begin{gathered}
\cos\operatorname{am}\frac{2Kx}{\pi}=\sqrt{\frac{k'}{k}}.\,\frac{H\left(\frac{2K}{\pi}\left(x+\frac{\pi}{2}\right)\right)}{\Theta\left(\frac{2Kx}{\pi}\right)},\\[2ex]
\Delta\operatorname{am}\frac{2Kx}{\pi}=\sqrt{k'}.\,\frac{\Theta\left(\frac{2K}{\pi}\left(x+\frac{\pi}{2}\right)\right)}{\Theta\left(\frac{2Kx}{\pi}\right)},
\end{gathered}
\]
whence, on putting $\dfrac{2Kx}{\pi}=u$:
\[
\text{1)}\quad\sin\operatorname{am}u=\frac{1}{\sqrt{k}}.\,\frac{H(u)}{\Theta(u)};\quad\cos\operatorname{am}u=\sqrt{\frac{k'}{k}}.\,\frac{H(u+K)}{\Theta(u)};\quad\Delta\operatorname{am}u=\sqrt{k'}.\,\frac{\Theta(u+K)}{\Theta(u)}.
\]

Hence flow special formulae:
\[
\text{2)}\quad\Theta(K)=\frac{\Theta(0)}{\sqrt{k'}};\quad H(K)=\sqrt{\frac{k}{k'}}\,\Theta(0).
\]

Putting $H'(u)=\dfrac{dH(u)}{du}$, since:
\[
H'(u)=\sqrt{k}\cos\operatorname{am}u\,\Delta\operatorname{am}u\,\Theta(u)+\sqrt{k}\sin\operatorname{am}u\,\Theta'(u),
\]
for the values $u=0$, $u=K$ we obtain:
\[
\text{3)}\quad H'(0)=\sqrt{k}\,\Theta(0)=\frac{H(K)\Theta(0)}{\Theta(K)};\quad H'(K)=\Theta'(K)=0\,{}^{*)}.
\]

From 2) it follows further:
\[
\text{4)}\quad\sqrt{k}=\frac{H(K)}{\Theta(K)};\quad\sqrt{k'}=\frac{\Theta(0)}{\Theta(K)}.
\]
Moreover:
\[
\begin{gathered}
\text{5)}\quad\Theta(u+2K)=\Theta(-u)=\Theta(u)\\[1ex]
\text{6)}\quad H(u+2K)=H(-u)=-H(u);\quad H(u+4K)=H(u).
\end{gathered}
\]

\textbf{*)} For $Z(K)=0$, whence also $\Theta'(K)=\Theta(K)Z(K)=0$.

\origpage{174}
From formula 2) §. 57:
\[
\Theta(u+iK')=i\,e^{\frac{\pi(K'-2iu)}{4K}}\sqrt{k}\sin\operatorname{am}u.
\]\ednote{The printed formula breaks off after the dot following $\sin\operatorname{am}u$; the missing factor $\Theta(u)$ appears only as a hand-written addition in ink, which is not transcribed.}
it follows:
\[
\text{7)}\quad\Theta(u+iK')=i\,e^{\frac{\pi(K'-2iu)}{4K}}H(u).
\]
Changing in this formula $u$ into $u+iK'$, and calling in 1) §. 57:
\[
\text{8)}\quad\Theta(u+2iK')=-e^{\frac{\pi(K'-iu)}{K}}\Theta(u),
\]
there results:
\[
\text{9)}\quad H(u+iK')=i\,e^{\frac{\pi(K'-2iu)}{4K}}\Theta(u),
\]
whence, changing again $u$ into $u+iK'$, from 7):
\[
\text{10)}\quad H(u+2iK')=-e^{\frac{\pi(K'-iu)}{K}}H(u).
\]

From formulae 7)-10) more general ones can be derived:
\[
\begin{gathered}
\text{11)}\quad e^{\frac{\pi uu}{4KK'}}\Theta(u)=(-1)^{m}e^{\frac{\pi(u+2miK')^{2}}{4KK'}}\Theta(u+2miK')\\[1.5ex]
\text{12)}\quad e^{\frac{\pi uu}{4KK'}}H(u)=(-1)^{m}e^{\frac{\pi(u+2miK')^{2}}{4KK'}}H(u+2miK')\\[1.5ex]
\text{13)}\quad e^{\frac{\pi uu}{4KK'}}H(u)=(-i)^{2m+1}e^{\frac{\pi(u+(2m+1)iK')^{2}}{4KK'}}\Theta(u+(2m+1)iK')\\[1.5ex]
\text{14)}\quad e^{\frac{\pi uu}{4KK'}}\Theta(u)=(-i)^{2m+1}e^{\frac{\pi(u+(2m+1)iK')^{2}}{4KK'}}H(u+(2m+1)iK').
\end{gathered}
\]

From 12), 13) there is:
\[
\text{15)}\quad\Theta\left((2m+1)iK'\right)=0;\ H(2miK')=0.
\]

Formulae 5), 6) demonstrate that the functions $\Theta(u)$, $H(u)$ remain unchanged when $u$ is changed into $u+4K$, formulae 11), 12) that the functions
\[
e^{\frac{\pi uu}{4KK'}}\Theta(u),\quad e^{\frac{\pi uu}{4KK'}}H(u)
\]

\origpage{175}
remain unchanged when $u$ is changed into $u+4iK'$; whence the former have with the Elliptic functions one real Period in common, the latter one imaginary Period.

From formula 5) §. 56:
\[
\frac{\Theta(iu,k)}{\Theta(0,k)}=e^{\frac{\pi uu}{4KK'}}\cos\operatorname{am}(u,k)\frac{\Theta(u,k')}{\Theta(0,k')},
\]
it follows:
\[
\frac{H(iu,k)}{\Theta(0,k)}=\sqrt{k}\sin\operatorname{am}(iu,k).\,\frac{\Theta(iu,k)}{\Theta(0,k)}=i\,e^{\frac{\pi uu}{4KK'}}\sqrt{k}\sin\operatorname{am}(u,k').\,\frac{\Theta(u,k')}{\Theta(0,k')},
\]
whence from 1):
\[
\begin{gathered}
\text{16)}\quad\frac{\Theta(iu,k)}{\Theta(0,k)}=\sqrt{\frac{k}{k'}}\,e^{\frac{\pi uu}{4KK'}}.\,\frac{H(u+K',k')}{\Theta(0,k')}\\[1.5ex]
\text{17)}\quad\frac{H(iu,k)}{\Theta(0,k)}=i\sqrt{\frac{k}{k'}}\,e^{\frac{\pi uu}{4KK'}}.\,\frac{H(u,k')}{\Theta(0,k')}.
\end{gathered}
\]
From 16) it follows, changing $u$ into $iu$, and interchanging $k$ and $k'$:
\[
\text{18)}\quad\frac{H(iu+K,k)}{\Theta(0,k)}=\sqrt{\frac{k}{k'}}\,e^{\frac{\pi uu}{4KK'}}.\,\frac{\Theta(u,k')}{\Theta(0,k')}
\]
to which let 9) §. 57 be adjoined:
\[
\text{19)}\quad\frac{\Theta(iu+K,k)}{\Theta(0,k)}=\sqrt{\frac{k}{k'}}\,e^{\frac{\pi uu}{4KK'}}\,\frac{\Theta(u+K',k')}{\Theta(0,k')}.
\]

From the formula found above:
\[
\Theta(u+v)\Theta(u-v)=\frac{\Theta^{2}u\,\Theta^{2}v}{\Theta^{2}0}\left(1-k^{2}\sin^{2}\operatorname{am}u.\sin^{2}\operatorname{am}v\right),
\]
it follows:
\[
\text{20)}\quad\Theta(u+v)\Theta(u-v)=\frac{\Theta^{2}u\,\Theta^{2}v-H^{2}u\,H^{2}v}{\Theta^{2}0}.
\]
This formula being multiplied into
\[
\begin{gathered}
k\sin\operatorname{am}(u+v)\sin\operatorname{am}(u-v)=\frac{k\sin^{2}\operatorname{am}u-k\sin^{2}\operatorname{am}v}{1-k^{2}\sin^{2}\operatorname{am}u\,\sin^{2}\operatorname{am}v}=\\[1.5ex]
\frac{H^{2}u\,\Theta^{2}v-\Theta^{2}u\,H^{2}v}{\Theta^{2}u\,\Theta^{2}v-H^{2}u\,H^{2}v},
\end{gathered}
\]
there results:
\[
\text{21)}\quad H(u+v)H(u-v)=\frac{H^{2}u\,\Theta^{2}v-\Theta^{2}u\,H^{2}v}{\Theta^{2}(0)}.
\]

\origpage{176}
\subsection*{ON THE EXPANSION OF THE FUNCTIONS $H$, $\Theta$ INTO SERIES. THIRD EXPANSION OF THE ELLIPTIC FUNCTIONS.}

\subsection*{62.}

Let us expand the functions
\[
\begin{gathered}
\frac{\Theta\left(\frac{2Kx}{\pi}\right)}{\Theta(0)}=\frac{(1-2q\cos 2x+q^{2})(1-2q^{3}\cos 2x+q^{6})(1-2q^{5}\cos 2x+q^{10})\ldots}{\left\{(1-q)(1-q^{3})(1-q^{5})\ldots\right\}^{2}}\\[2ex]
\frac{H\left(\frac{2Kx}{\pi}\right)}{\Theta(0)}=\frac{2\sqrt[4]{q}\sin x(1-2q^{2}\cos 2x+q^{4})(1-2q^{4}\cos 2x+q^{8})(1-2q^{6}\cos 2x+q^{12})\ldots}{\left\{(1-q)(1-q^{3})(1-q^{5})\ldots\right\}^{2}}
\end{gathered}
\]
into series
\[
\begin{gathered}
\frac{\Theta\left(\frac{2Kx}{\pi}\right)}{\Theta(0)}=A-2A'\cos 2x+2A''\cos 4x-2A'''\cos 6x+2A''''\cos 8x-\ldots\\[2ex]
\frac{H\left(\frac{2Kx}{\pi}\right)}{\Theta(0)}=2\sqrt[4]{q}\left\{B'\sin x-B''\sin 3x+B'''\sin 5x-B''''\sin 7x+\ldots\right\}.
\end{gathered}
\]
We obtain the determination of the quantities $A,A',A'',A''',\,.\,.;\ B',B'',B''',B'''',\,.\,.$ by means of equations 7)-10) §$^{\mathrm{i}}$ of the preceding, which on putting $u=\dfrac{2Kx}{\pi}$, $q=e^{-\frac{\pi K'}{K}}$ go over into the following:
\[
\begin{gathered}
\Theta\left(\frac{2Kx}{\pi}\right)=-q\,e^{2ix}\,\Theta\left(\frac{2Kx}{\pi}+2iK'\right)\\[2ex]
H\left(\frac{2Kx}{\pi}\right)=-q\,e^{2ix}\,H\left(\frac{2Kx}{\pi}+2iK'\right)\\[2ex]
i\Theta\left(\frac{2Kx}{\pi}\right)=\sqrt[4]{q}\,e^{ix}\,H\left(\frac{2Kx}{\pi}+iK'\right)\\[2ex]
iH\left(\frac{2Kx}{\pi}\right)=\sqrt[4]{q}\,e^{ix}\,\Theta\left(\frac{2Kx}{\pi}+iK'\right).
\end{gathered}
\]
For\ednote{The author's Corrigenda give the reading Quem here.} this purpose we exhibit the proposed expansions thus:
\[
\begin{gathered}
\frac{\Theta\left(\frac{2Kx}{\pi}\right)}{\Theta(0)}=A-A'e^{2ix}+A''e^{4ix}-A'''e^{6ix}+A''''e^{8ix}-\ldots\\[1.5ex]
-A'e^{-2ix}+A''e^{-4ix}-A'''e^{-6ix}+A''''e^{-8ix}-\ldots
\end{gathered}
\]

\origpage{177}
\[
\begin{gathered}
\frac{H\left(\frac{2Kx}{\pi}\right)}{\Theta(0)}=\sqrt[4]{q}\left\{B'e^{ix}-B''e^{3ix}+B'''e^{5ix}-B''''e^{7ix}+\ldots\right\}\\[1.5ex]
-\sqrt[4]{q}\left\{B'e^{-ix}-B''e^{-3ix}+B'''e^{-5ix}-B''''e^{-7ix}+\ldots\right\}.
\end{gathered}
\]
When $x$ is changed into $x-i\log q$, $e^{mix}$ goes over into $q^{m}e^{mix}$, $e^{-mix}$ into $\dfrac{e^{-mix}}{q^{m}}$; further $\Theta\left(\dfrac{2Kx}{\pi}\right)$, $H\left(\dfrac{2Kx}{\pi}\right)$ into $\Theta\left(\dfrac{2Kx}{\pi}+2iK'\right)$, $H\left(\dfrac{2Kx}{\pi}+2iK'\right)$. Hence we obtain:
\[
\begin{gathered}
\frac{\Theta\left(\frac{2Kx}{\pi}\right)}{\Theta(0)}=-q\,e^{2ix}.\,\frac{\Theta\left(\frac{2Kx}{\pi}+2iK'\right)}{\Theta(0)}=\\[1.5ex]
\frac{A'}{q}-Aq\,e^{2ix}+A'q^{3}e^{4ix}-A''q^{5}e^{6ix}+A'''q^{7}e^{8ix}-\ldots\\[1.5ex]
-\frac{A''}{q^{3}}e^{-2ix}+\frac{A'''}{q^{5}}e^{-4ix}-\frac{A''''}{q^{7}}e^{-6ix}+\frac{A'''''}{q^{9}}e^{-8ix}-\ldots
\end{gathered}
\]
\[
\begin{gathered}
\frac{H\left(\frac{2Kx}{\pi}\right)}{\Theta(0)}=-q\,e^{2ix}.\,\frac{H\left(\frac{2Kx}{\pi}+2iK'\right)}{\Theta(0)}=\\[1.5ex]
\sqrt[4]{q}\left\{B'e^{ix}-B'q^{2}e^{3ix}+B''q^{4}e^{5ix}-B'''q^{6}e^{7ix}+\ldots\right\}\\[1.5ex]
-\sqrt[4]{q}\left\{\frac{B''}{q^{2}}e^{-ix}-\frac{B'''}{q^{4}}e^{-3ix}+\frac{B''''}{q^{6}}e^{-5ix}-\frac{B'''''}{q^{8}}e^{-7ix}+\ldots\right\}.
\end{gathered}
\]

Comparing these with the proposed expressions, we derive:
\[
\begin{gathered}
A'=Aq,\quad A''=A'q^{3},\quad A'''=A''q^{5},\quad A''''=A'''q^{7},\quad.\,.\,.\,.\\[1ex]
B''=B'q^{2},\quad B'''=B''q^{4},\quad B''''=B'''q^{6},\quad B'''''=B''''q^{8},\quad.\,.\,.\,.
\end{gathered}
\]
and therefore
\[
\begin{gathered}
A'=Aq,\quad A''=Aq^{4},\quad A'''=Aq^{9},\quad A''''=Aq^{16},\quad\ldots\\[1ex]
B''=B'q^{2},\quad B'''=B'q^{6},\quad B''''=B'q^{12},\quad B'''''=B'q^{20},\quad.\,.\,.\,.
\end{gathered}
\]
whence the sought expansions become:
\[
\begin{gathered}
\frac{\Theta\left(\frac{2Kx}{\pi}\right)}{\Theta(0)}=A\left\{1-2q\cos 2x+2q^{4}\cos 4x-2q^{9}\cos 6x+2q^{16}\cos 8x-\ldots\right\}\\[2ex]
\frac{H\left(\frac{2Kx}{\pi}\right)}{\Theta(0)}=2\sqrt[4]{q}\,B'\left\{\sin x-q^{2}\sin 3x+q^{2.3}\sin 5x-q^{3.4}\sin 7x+q^{4.5}\sin 9x-.\,.\right\}\\[1.5ex]
=B'\left\{2\sqrt[4]{q}\sin x-2\sqrt[4]{q^{9}}\sin 3x+2\sqrt[4]{q^{25}}\sin 5x-2\sqrt[4]{q^{49}}\sin 7x+.\,.\right\}.
\end{gathered}
\]

The expansions found could have been derived one from the other by means of the formula:
\[
iH\left(\frac{2Kx}{\pi}\right)=\sqrt[4]{q}\,e^{ix}\,\Theta\left(\frac{2Kx}{\pi}+iK'\right).
\]

\origpage{178}
For we have found the series:\ednote{The Corrigenda correct the printed word varie to serie.}
\[
\frac{\Theta\left(\frac{2Kx}{\pi}\right)}{\Theta(0)}=A\left\{1-q(e^{2ix}+e^{-2ix})+q^{4}(e^{4ix}+e^{-4ix})+q^{9}(e^{6ix}+e^{-6ix})+.\,.\right\},
\]\ednote{The sign before the term in $q^{9}$ is printed as plus; by the preceding series and the following steps it should be minus, an apparent misprint.}
by changing $x$ into $x-i\log\sqrt{q}$, whereby $e^{2mix}$, $e^{-2mix}$ pass into $q^{m}e^{2mix}$, $\dfrac{e^{-2mix}}{q^{m}}$, $\Theta\left(\dfrac{2Kx}{\pi}\right)$ into $\Theta\left(\dfrac{2Kx}{\pi}+iK'\right)$, and by multiplying by $\sqrt[4]{q}\,e^{ix}$, we obtain:
\[
\begin{gathered}
\frac{iH\left(\frac{2Kx}{\pi}\right)}{\Theta(0)}=\sqrt[4]{q}\,e^{ix}\frac{\Theta\left(\frac{2Kx}{\pi}+iK'\right)}{\Theta(0)}=\\[1.5ex]
A\left\{\sqrt[4]{q}(e^{ix}-e^{-ix})-\sqrt[4]{q^{9}}(e^{3ix}-e^{-3ix})+\sqrt[4]{q^{25}}(e^{5ix}-e^{-5ix})+\ldots\right\},
\end{gathered}
\]
or:
\[
\frac{H\left(\frac{2Kx}{\pi}\right)}{\Theta(0)}=A\left\{2\sqrt[4]{q}\sin x-2\sqrt[4]{q^{9}}\sin 3x+2\sqrt[4]{q^{25}}\sin 5x-2\sqrt[4]{q^{49}}\sin 7x+\ldots\right\}.
\]
By this analysis we have moreover found:
\[
B'=A.
\]

\subsection*{63.}

The determination of $A$ requires particular devices. Let us put, as is permitted from what precedes:
\[
\begin{gathered}
(1-2q\cos 2x+q^{2})(1-2q^{3}\cos 2x+q^{6})(1-2q^{5}\cos 2x+q^{10})\ldots=\\[1ex]
P(q)\left\{1-2q\cos 2x+2q^{4}\cos 4x-2q^{9}\cos 6x+2q^{16}\cos 8x+\ldots\right\}\\[1ex]
\sin x(1-2q^{2}\cos 2x+q^{4})(1-2q^{4}\cos 2x+q^{8})(1-2q^{6}\cos 2x+q^{12})\ldots=\\[1ex]
P(q)\left\{\sin x-q^{1.2}\sin 3x+q^{2.3}\sin 5x-q^{3.4}\sin 7x+q^{4.5}\sin 9x-\ldots\right\};
\end{gathered}
\]
it follows that:
\[
A=\frac{P(q)}{\left\{(1-q)(1-q^{3})(1-q^{5})\ldots\right\}^{2}}.
\]
The second expression remains unchanged when it is multiplied into the first and, the product having been formed, $q^{2}$ is put in place of $q$. Hence we obtain the identical equation:
\[
\begin{gathered}
P(q^{2})P(q^{2})\left\{\sin x-q^{4}\sin 3x+q^{12}\sin 5x-q^{24}\sin 7x+\ldots\right\}\times\\[1ex]
\left\{1-2q^{2}\cos 2x+2q^{8}\cos 4x-2q^{32}\cos 6x+\ldots\right\}=\\[1ex]
P(q)\left\{\sin x-q^{2}\sin 3x+q^{6}\sin 5x-q^{12}\sin 7x+\ldots\right\}.
\end{gathered}
\]

\origpage{179}
Let us now carry out this multiplication itself, so that everywhere $\sin(m+n)x+\sin(m-n)x$ is written in place of $2\sin.mx\cos.nx$: it is easily seen that the coefficient of $\sin x$ in the expanded product will be:
\[
1+q^{2}+q^{6}+q^{12}+q^{20}+\ldots,
\]
so that there results:
\[
\frac{P(q)}{P(q^{2})P(q^{2})}=1+q^{2}+q^{6}+q^{12}+q^{20}+\ldots
\]
But we have found from the second of the proposed formulae, with $x=\dfrac{\pi}{2}$:
\[
\left\{(1+q^{2})(1+q^{4})(1+q^{6})\ldots\right\}^{2}=P(q)\left\{1+q^{2}+q^{6}+q^{12}+q^{20}+\ldots\right\},
\]
whence:
\[
\frac{P(q)P(q)}{P(q^{2})P(q^{2})}=\left\{(1+q^{2})(1+q^{4})(1+q^{6})\ldots\right\}^{2},
\]
or:
\[
\begin{gathered}
\frac{P(q)}{P(q^{2})}=(1+q^{2})(1+q^{4})(1+q^{6})\ldots\\[1.5ex]
=\frac{(1-q^{4})(1-q^{8})(1-q^{12})\ldots}{(1-q^{2})(1-q^{4})(1-q^{6})\ldots}.
\end{gathered}
\]
Hence, by the method already employed more than once ${}^{*)}$, it follows:
\[
P(q)=\frac{1}{(1-q^{2})(1-q^{4})(1-q^{6})(1-q^{8})\ldots}.
\]
Hence at last there results:
\[
\begin{gathered}
A=\frac{1}{(1-q^{2})(1-q^{4})(1-q^{6})\ldots}.\,\frac{1}{\left\{(1-q)(1-q^{3})(1-q^{5})\ldots\right\}^{2}}\\[1.5ex]
=\frac{(1+q)(1+q^{2})(1+q^{3})(1+q^{4})\ldots}{(1-q)(1-q^{2})(1-q^{3})(1-q^{4})\ldots},
\end{gathered}
\]
or, from the expansions which we gave in §. 36:
\[
\frac{1}{A}=\sqrt{\frac{2k'K}{\pi}}.
\]

The quantity which we have so far left undetermined, $\Theta(0)$, let us now put:
\[
\Theta(0)=\frac{1}{A}=\sqrt{\frac{2k'K}{\pi}},
\]

\textbf{*)} Namely, by putting successively $q^{2}$, $q^{4}$, $q^{8}$, $q^{16}$ $.\,.$ in place of $x$ and carrying out the infinite multiplication.\ednote{Printed $x$ in the footnote; the substitution is evidently for $q$, an apparent misprint.}

\origpage{180}
there is found:
\[
\begin{gathered}
\text{1)}\quad\Theta\left(\frac{2Kx}{\pi}\right)=1-2q\cos 2x+2q^{4}\cos 4x-2q^{9}\cos 6x+2q^{16}\cos 8x-\ldots\\[2ex]
\text{2)}\quad H\left(\frac{2Kx}{\pi}\right)=2\sqrt[4]{q}\sin x-2\sqrt[4]{q^{9}}\sin 3x+2\sqrt[4]{q^{25}}\sin 5x-2\sqrt[4]{q^{49}}\sin 7x+\ldots
\end{gathered}
\]

\subsection*{64.}

The identical equation, which we have proved in what precedes:
\[
\begin{gathered}
(1-2q\cos 2x+q^{2})(1-2q^{3}\cos 2x+q^{6})(1-2q^{5}\cos 2x+q^{10}).\,.\,.\,.=\\[1.5ex]
\frac{1-2q\cos 2x+2q^{4}\cos 4x-2q^{9}\cos 6x+2q^{16}\cos 8x-\ldots}{(1-q^{2})(1-q^{4})(1-q^{6})(1-q^{8}).\,.\,.\,.}
\end{gathered}
\]
it pleases us to investigate by yet another way, entirely different from the preceding. To that end\ednote{The Corrigenda correct the printed word Quam to Quem.} let us premise, as lemmas, the two following formulae:
\[
\begin{gathered}
\text{1)}\quad(1+qz)(1+q^{3}z)(1+q^{5}z)(1+q^{7}z)\ldots=\\[1.5ex]
1+\frac{qz}{1-q^{2}}+\frac{q^{4}z^{2}}{(1-q^{2})(1-q^{4})}+\frac{q^{9}z^{3}}{(1-q^{2})(1-q^{4})(1-q^{6})}+\frac{q^{16}z^{4}}{(1-q^{2})(1-q^{4})(1-q^{6})(1-q^{8})}+.\,.\\[2ex]
\text{2)}\quad\frac{1}{(1-qz)(1-q^{2}z)(1-q^{3}z)(1-q^{4}z)\ldots}=\\[1.5ex]
1+\frac{q}{1-q}.\,\frac{z}{1-qz}+\frac{q^{4}}{(1-q)(1-q^{2})}.\,\frac{z^{2}}{(1-qz)(1-q^{2}z)}+\\[1.5ex]
\frac{q^{9}}{(1-q)(1-q^{2})(1-q^{3})}.\,\frac{z^{3}}{(1-qz)(1-q^{2}z)(1-q^{3}z)}+.\,.\,.\,.
\end{gathered}
\]
For the proof of the former I observe that the expression
\[
(1+qz)(1+q^{3}z)(1+q^{5}z)(1+q^{7}z)\ldots
\]
remains unchanged when $q^{2}z$ is put in place of $z$ and multiplication by $(1+qz)$ is carried out; whence, putting:
\[
(1+qz)(1+q^{3}z)(1+q^{5}z)\ldots=1+A'z+A''z^{2}+A'''z^{3}+.\,.,
\]
one obtains:
\[
1+A'z+A''z^{2}+A'''z^{3}+\ldots=(1+qz)(1+A'q^{2}z+A''q^{4}z^{2}+A'''q^{6}z^{3}+\ldots),
\]
and therefore, the expansion having been carried out:
\[
A'=q+q^{2}A',\quad A''=q^{3}A'+q^{4}A'',\quad A'''=q^{5}A''+q^{6}A''',\quad.\,.\,.\,.,
\]
or:
\[
A'=\frac{q}{1-q^{2}},\quad A''=\frac{q^{3}A'}{1-q^{4}},\quad A'''=\frac{q^{5}A''}{1-q^{6}},
\]

\origpage{181}
whence:
\[
A'=\frac{q}{1-q^{2}},\quad A''=\frac{q^{4}}{(1-q^{2})(1-q^{4})},\quad A'''=\frac{q^{9}}{(1-q^{2})(1-q^{4})(1-q^{6})},\quad.\,.\,.\,.,
\]
as was proposed.

For the proof of formula 2) I observe that the expression
\[
\frac{1}{(1-qz)(1-q^{2}z)(1-q^{3}z)(1-q^{4}z)\,.\,.\,.\,.},
\]
remains unchanged when $qz$ is put in place of $z$ and multiplication by $\dfrac{1}{1-qz}$ is carried out; whence, putting:
\[
\begin{gathered}
\frac{1}{(1-qz)(1-q^{2}z)(1-q^{3}z)(1-q^{4}z)\,.\,.\,.\,.}=\\[1ex]
1+\frac{A'z}{1-qz}+\frac{A''z^{2}}{(1-qz)(1-q^{2}z)}+\frac{A'''z^{3}}{(1-qz)(1-q^{2}z)(1-q^{3}z)}+\ldots,
\end{gathered}
\]
we obtain:
\[
\begin{gathered}
1+\frac{A'z}{1-qz}+\frac{A''z^{2}}{(1-qz)(1-q^{2}z)}+\frac{A'''z^{3}}{(1-qz)(1-q^{2}z)(1-q^{3}z)}+\ldots=\\[1.5ex]
\frac{1}{1-qz}+\frac{A'qz}{(1-qz)(1-q^{2}z)}+\frac{A''q^{2}z^{2}}{(1-qz)(1-q^{2}z)(1-q^{3}z)}+\frac{A'''q^{3}z^{3}}{(1-qz)(1-q^{2}z)(1-q^{3}z)(1-q^{4}z)}+\ldots\\[1.5ex]
=1+\frac{(q+A'q)z}{1-qz}+\frac{(q^{3}A'+q^{2}A'')z^{2}}{(1-qz)(1-q^{2}z)}+\frac{(q^{5}A''+q^{3}A''')z^{3}}{(1-qz)(1-q^{2}z)(1-q^{3}z)}+\ldots\,{}^{*)}.
\end{gathered}
\]
Hence it flows:
\[
A'=q+A'q,\quad A''=q^{3}A'+q^{2}A'',\quad A'''=q^{5}A''+q^{3}A''',\quad\ldots,
\]
and therefore:
\[
A'=\frac{q}{1-q},\quad A''=\frac{q^{3}A'}{1-q^{2}},\quad A'''=\frac{q^{5}A''}{1-q^{3}},\quad\ldots,
\]
whence:
\[
A'=\frac{q}{1-q},\quad A''=\frac{q^{4}}{(1-q)(1-q^{2})},\quad A'''=\frac{q^{9}}{(1-q)(1-q^{2})(1-q^{3})},\quad\ldots,
\]
as was proposed.

\textbf{*)} Namely, by substituting in the single terms respectively $\dfrac{1}{1-qz}=1+\dfrac{qz}{1-qz}$, $\dfrac{1}{1-q^{2}z}=1+\dfrac{q^{2}z}{1-q^{2}z}$, $\dfrac{1}{1-q^{3}z}=1+\dfrac{q^{3}z}{1-q^{3}z}$, etc.

\origpage{182}
Now let us form the product:
\[
\begin{gathered}
\left\{(1+qz)(1+q^{3}z)(1+q^{5}z)\,.\,.\right\}\left\{\left(1+\frac{q}{z}\right)\left(1+\frac{q^{3}}{z}\right)\left(1+\frac{q^{5}}{z}\right)\ldots\right\}\\[1.5ex]
=\\[1.5ex]
\left\{1+\frac{q}{1-q^{2}}.\,z+\frac{q^{4}}{(1-q^{2})(1-q^{4})}.\,z^{2}+\frac{q^{9}}{(1-q^{2})(1-q^{4})(1-q^{6})}.\,z^{3}+\ldots\right\}\times\\[1.5ex]
\left\{1+\frac{q}{1-q^{2}}.\,\frac{1}{z}+\frac{q^{4}}{(1-q^{2})(1-q^{4})}.\,\frac{1}{z^{2}}+\frac{q^{9}}{(1-q^{2})(1-q^{4})(1-q^{6})}.\,\frac{1}{z^{3}}+\ldots\right\}.
\end{gathered}
\]
The coefficient of $z^{n}$ or also of $\dfrac{1}{z^{n}}$, which we shall put $B^{(n)}$, we find to be the following:
\[
B^{(n)}=\frac{q^{nn}}{(1-q^{2})(1-q^{4})\ldots(1-q^{2n})}\times\left\{1+\frac{q^{2}}{1-q^{2}}.\,\frac{q^{2n}}{1-q^{2n+2}}+\frac{q^{8}}{(1-q^{2})(1-q^{4})}.\,\frac{q^{4n}}{(1-q^{2n+2})(1-q^{2n+4})}+\frac{q^{18}}{(1-q^{2})(1-q^{4})(1-q^{6})}.\,\frac{q^{6n}}{(1-q^{2n+2})(1-q^{2n+4})(1-q^{2n+6})}+.\,.\right\}
\]
But from formula 2), putting $q^{2}$ in place of $q$ and $z=q^{2.n}$, we find the expression which is seen enclosed in braces to be $=$
\[
\frac{1}{(1-q^{2n+2})(1-q^{2n+4})(1-q^{2n+6})(1-q^{2n+8})\ldots},
\]
whence
\[
B^{(n)}=\frac{q^{nn}}{(1-q^{2})(1-q^{4})(1-q^{6})(1-q^{8})\,.\,.\,.\,.\,.},
\]
and therefore:
\[
\begin{gathered}
\left\{(1+qz)(1+q^{3}z)(1+q^{5}z)\,.\,.\right\}\left\{\left(1+\frac{q}{z}\right)\left(1+\frac{q^{3}}{z}\right)\left(1+\frac{q^{5}}{z}\right)\ldots\right\}=\\[1.5ex]
\frac{1+q\left(z+\dfrac{1}{z}\right)+q^{4}\left(z^{2}+\dfrac{1}{z^{2}}\right)+q^{9}\left(z^{3}+\dfrac{1}{z^{3}}\right)+\ldots}{(1-q^{2})(1-q^{4})(1-q^{6})(1-q^{8})\ldots},
\end{gathered}
\]
or, putting $z=e^{2ix}$, and changing $q$ into $-q$:
\[
\begin{gathered}
(1-2q\cos 2x+q^{2})(1-2q^{3}\cos 2x+q^{6})(1-2q^{5}\cos 2x+q^{10})\ldots=\\[1.5ex]
\frac{1-2q\cos 2x+2q^{4}\cos 4x-2q^{9}\cos 6x+\ldots}{(1-q^{2})(1-q^{4})(1-q^{6})(1-q^{8})\ldots}.
\end{gathered}
\]
Which was to be demonstrated.

\origpage{183}
When $-qz^{2}$ is put in place of $z$ and multiplication by $\sqrt[4]{q}\,z$ is made, there results:
\[
\begin{gathered}
\sqrt[4]{q}\left(z-\frac{1}{z}\right)\left\{(1-q^{2}z^{2})(1-q^{4}z^{2})(1-q^{6}z^{2})\ldots\right\}\times\\[1.5ex]
\left\{\left(1-\frac{q^{2}}{z^{2}}\right)\left(1-\frac{q^{4}}{z^{2}}\right)\left(1-\frac{q^{6}}{z^{2}}\right)\ldots\right\}=\\[1.5ex]
\frac{\sqrt[4]{q}\left(z-\dfrac{1}{z}\right)-\sqrt[4]{q^{9}}\left(z^{3}-\dfrac{1}{z^{3}}\right)+\sqrt[4]{q^{25}}\left(z^{5}-\dfrac{1}{z^{5}}\right)-.\,.\,.\,.}{(1-q^{2})(1-q^{4})(1-q^{6})(1-q^{8})\,.\,.\,.\,.},
\end{gathered}
\]
or, putting $z=e^{ix}$:
\[
\begin{gathered}
2\sqrt[4]{q}\sin x(1-2q^{2}\cos 2x+q^{4})(1-2q^{4}\cos 2x+q^{8})(1-2q^{6}\cos 2x+q^{12})\ldots=\\[1.5ex]
\frac{2\sqrt[4]{q}\sin x-2\sqrt[4]{q^{9}}\sin 3x+2\sqrt[4]{q^{25}}\sin 5x-2\sqrt[4]{q^{49}}\sin 7x+\ldots}{(1-q^{2})(1-q^{4})(1-q^{6})(1-q^{8})\ldots},
\end{gathered}
\]
which is the other expansion found.

\subsection*{65.}

The expansions of the functions
\[
\begin{gathered}
\text{1)}\quad\Theta\left(\frac{2Kx}{\pi}\right)=1-2q\cos 2x+2q^{4}\cos 4x-2q^{9}\cos 6x+2q^{16}\cos 8x-\ldots\\[1.5ex]
\text{2)}\quad H\left(\frac{2Kx}{\pi}\right)=2\sqrt[4]{q}\sin x-2\sqrt[4]{q^{9}}\sin 3x+2\sqrt[4]{q^{25}}\sin 5x-2\sqrt[4]{q^{49}}\sin 7x+\ldots
\end{gathered}
\]
lead of themselves to a new expansion of the Elliptic functions. For from formulae 1) of §. 61, putting $u=\dfrac{2Kx}{\pi}$, we obtain:
\[
\begin{gathered}
\sin\operatorname{am}\frac{2Kx}{\pi}=\frac{1}{\sqrt{k}}.\frac{H\left(\dfrac{2Kx}{\pi}\right)}{\Theta\left(\dfrac{2Kx}{\pi}\right)}\\[2ex]
\cos\operatorname{am}\frac{2Kx}{\pi}=\sqrt{\frac{k'}{k}}\;\frac{H\left(\dfrac{2K}{\pi}\left(x+\dfrac{\pi}{2}\right)\right)}{\Theta\left(\dfrac{2Kx}{\pi}\right)}\\[2ex]
\Delta\operatorname{am}\frac{2Kx}{\pi}=\sqrt{k'}\;\frac{\Theta\left(\dfrac{2K}{\pi}\left(x+\dfrac{\pi}{2}\right)\right)}{\Theta\left(\dfrac{2Kx}{\pi}\right)},
\end{gathered}
\]

\origpage{184}
whence:
\[
\begin{gathered}
\text{3)}\quad\sin\operatorname{am}\frac{2Kx}{\pi}=\frac{1}{\sqrt{k}}.\frac{2\sqrt[4]{q}\sin x-2\sqrt[4]{q^{9}}\sin 3x+2\sqrt[4]{q^{25}}\sin 5x-2\sqrt[4]{q^{49}}\sin 7x+\ldots}{1-2q\cos 2x+2q^{4}\cos 4x-2q^{9}\cos 6x+2q^{16}\cos 8x-\ldots}\\[2ex]
\text{4)}\quad\cos\operatorname{am}\frac{2Kx}{\pi}=\sqrt{\frac{k'}{k}}.\frac{2\sqrt[4]{q}\cos x+2\sqrt[4]{q^{9}}\cos 3x+2\sqrt[4]{q^{25}}\cos 5x+2\sqrt[4]{q^{49}}\cos 7x+\ldots}{1-2q\cos 2x+2q^{4}\cos 4x-2q^{9}\cos 6x+2q^{16}\cos 8x-\ldots}\\[2ex]
\text{5)}\quad\Delta\operatorname{am}\frac{2Kx}{\pi}=\sqrt{k'}.\frac{1+2q\cos 2x+2q^{4}\cos 4x+2q^{9}\cos 6x+2q^{16}\cos 8x+\ldots}{1-2q\cos 2x+2q^{4}\cos 4x-2q^{9}\cos 6x+2q^{16}\cos 8x-\ldots}.
\end{gathered}
\]

Further, from 2), 3) of §. 61, since $\Theta(0)=\sqrt{\dfrac{2k'K}{\pi}}$ has been put, we obtain:
\[
\Theta(K)=\sqrt{\frac{2K}{\pi}},\quad H(K)=\sqrt{\frac{2kK}{\pi}},\quad\Theta(0)=\sqrt{\frac{2k'K}{\pi}},\quad H'(0)=\sqrt{\frac{2kk'K}{\pi}},
\]
whence from 1), 2):
\[
\begin{gathered}
\text{6)}\quad\sqrt{\frac{2K}{\pi}}=1+2q+2q^{4}+2q^{9}+2q^{16}+2q^{25}+\ldots\\[2ex]
\text{7)}\quad\sqrt{\frac{2kK}{\pi}}=2\sqrt[4]{q}+2\sqrt[4]{q^{9}}+2\sqrt[4]{q^{25}}+2\sqrt[4]{q^{49}}+2\sqrt[4]{q^{81}}+\ldots\\[2ex]
\text{8)}\quad\sqrt{\frac{2k'K}{\pi}}=1-2q+2q^{4}-2q^{9}+2q^{16}-2q^{25}+.\,.\\[2ex]
\text{9)}\quad\sqrt{kk'\left(\frac{2K}{\pi}\right)^{3}}=2\sqrt[4]{q}-6\sqrt[4]{q^{9}}+10\sqrt[4]{q^{25}}-14\sqrt[4]{q^{49}}+18\sqrt[4]{q^{81}}-\ldots\,{}^{*)},
\end{gathered}
\]
whence also:
\[
\begin{gathered}
\text{10)}\quad\sqrt{k}=\frac{2\sqrt[4]{q}+2\sqrt[4]{q^{9}}+2\sqrt[4]{q^{25}}+2\sqrt[4]{q^{49}}+2\sqrt[4]{q^{81}}+\ldots}{1+2q+2q^{4}+2q^{9}+2q^{16}+\ldots}\\[2ex]
\text{11)}\quad\sqrt{k'}=\frac{1-2q+2q^{4}-2q^{9}+2q^{16}-2q^{25}+.\,.}{1+2q+2q^{4}+2q^{9}+2q^{16}+2q^{25}+.\,.}
\end{gathered}
\]

It follows further, since $Z(u)=\dfrac{\Theta'(u)}{\Theta(u)}$, $\Pi(u,a)=uZ(a)+\dfrac{1}{2}\log\dfrac{\Theta(u-a)}{\Theta(u+a)}$:
\[
\text{12)}\quad\frac{2K}{\pi}.Z\left(\frac{2Kx}{\pi}\right)=\frac{4q\sin 2x-8q^{4}\sin 4x+12q^{9}\sin 6x-16q^{16}\sin 8x+.\,.}{1-2q\cos 2x+2q^{4}\cos 4x-2q^{9}\cos 6x+2q^{16}\cos 8x-.\,.}
\]

\textbf{*)} For since $\dfrac{dH}{dx}=\dfrac{2K}{\pi}.\dfrac{dH}{du}$, when 2) has been differentiated with respect to $x$ and $x=0$ is then put, there results
\[
\frac{2K}{\pi}H'(0)=\sqrt{kk'\left(\frac{2K}{\pi}\right)^{3}}.
\]

\origpage{185}
\[
\begin{gathered}
\text{13)}\quad\Pi\left(\frac{2Kx}{\pi},\frac{2KA}{\pi}\right)=\frac{2Kx}{\pi}.Z\left(\frac{2KA}{\pi}\right)\\[1.5ex]
+\frac{1}{2}\log\frac{1-2q\cos 2(x-A)+2q^{4}\cos 4(x-A)-2q^{9}\cos 6(x-A)+.\,.}{1-2q\cos 2(x+A)+2q^{4}\cos 4(x+A)-2q^{9}\cos 6(x+A)+.\,.}
\end{gathered}
\]
This is the third expansion of the Elliptic functions.

\subsection*{66.}

From the expansions found:
\[
\begin{gathered}
\text{1)}\quad\left\{(1-q^{2})(1-q^{4})(1-q^{6})\,.\,.\right\}\left\{(1-2q\cos 2x+q^{2})(1-2q^{3}\cos 2x+q^{6})(1-2q^{5}\cos 2x+q^{10})\ldots\right\}\\[1ex]
=\\[1ex]
1-2q\cos 2x+2q^{4}\cos 4x-2q^{9}\cos 6x+2q^{16}\cos 8x-.\,.\\[1.5ex]
\left\{(1-q^{2})(1-q^{4})(1-q^{6})\,.\,.\right\}\sin x(1-2q^{2}\cos 2x+q^{4})(1-2q^{4}\cos 2x+q^{8})\ldots\\[1ex]
=\\[1ex]
\sin x-q^{2}\sin 3x+q^{6}\sin 5x-q^{12}\sin 7x+q^{20}\sin 9x-.\,.
\end{gathered}
\]
the last of which, with $\sqrt{q}$ put in place of $q$, may also be exhibited thus:
\[
\begin{gathered}
\text{2)}\quad\left\{(1-q)(1-q^{2})(1-q^{3})\,.\,.\right\}\sin x(1-2q\cos 2x+q^{2})(1-2q^{2}\cos 2x+q^{4})(1-2q^{3}\cos 2x+q^{6})\ldots\\[1ex]
=\\[1ex]
\sin x-q\sin 3x+q^{3}\sin 5x-q^{6}\sin 7x+q^{10}\sin 9x-q^{15}\sin 11x+\ldots,
\end{gathered}
\]
it follows, putting $x=0$, $x=\dfrac{\pi}{2}$:
\[
\begin{gathered}
\text{3)}\quad\frac{(1-q)(1-q^{2})(1-q^{3})(1-q^{4})\ldots}{(1+q)(1+q^{2})(1+q^{3})(1+q^{4})\ldots}=1-2q+2q^{4}-2q^{9}+2q^{16}-\ldots\\[2ex]
\text{4)}\quad\frac{(1-q^{2})(1-q^{4})(1-q^{6})(1-q^{8})\ldots}{(1-q)(1-q^{3})(1-q^{5})(1-q^{7})\ldots}=1+q+q^{3}+q^{6}+q^{10}+q^{15}+\ldots\\[2ex]
\text{5)}\quad\left\{(1-q)(1-q^{3})(1-q^{5})(1-q^{7})\ldots\right\}^{3}=1-3q+5q^{3}-7q^{6}+9q^{10}-\ldots
\end{gathered}
\]
Let us put in 2) $x=\dfrac{\pi}{3}$, then $\sin x=+\sqrt{\dfrac{3}{4}}$, $\sin 3x=0$, $\sin 5x=-\sqrt{\dfrac{3}{4}}$, $\sin 7x=+\sqrt{\dfrac{3}{4}}$, etc.; further $(1-q)(1-2q\cos 2x+q^{2})=1-q^{3}$, whence 2) passes into this formula:
\[
(1-q^{3})(1-q^{6})(1-q^{9})(1-q^{12})\,.\,.=1-q^{3}-q^{6}+q^{15}+q^{21}-q^{36}-.\,.,
\]
or:
\[
\text{6)}\quad(1-q)(1-q^{2})(1-q^{3})(1-q^{4})\ldots=1-q-q^{2}+q^{5}+q^{7}-q^{12}-.\,.,
\]

\origpage{186}
of which series the general term is:
\[
(-1)^{n}q^{\frac{3nn\pm n}{2}}.
\]
Comparing 5), 6) with one another we obtain:
\[
\text{7)}\quad\left\{1-q-q^{2}+q^{5}+q^{7}-q^{12}-.\,.\right\}^{3}=1-3q+5q^{3}-7q^{6}+9q^{10}-.\,.\,.\,.
\]

Formula 4) was also found by Cl.\ Gauss in the memoir: \emph{Summatio Serierum quarundam singularium.} Comm.\ Gott.\ Vol.\ I.\ a.\ 1808---1811. He deduced it from the following memorable formula:
\[
\begin{gathered}
\text{8)}\quad\frac{(1-qz)(1-q^{3}z)(1-q^{5}z)(1-q^{7}z)\ldots}{(1-q)(1-q^{3})(1-q^{5})(1-q^{7})\ldots}=\\[1.5ex]
1+\frac{q(1-z)}{1-q}+\frac{q^{3}(1-z)(1-qz)}{(1-q)(1-q^{2})}+\frac{q^{6}(1-z)(1-qz)(1-q^{2}z)}{(1-q)(1-q^{2})(1-q^{3})}+.\,.,
\end{gathered}
\]
by putting $z=q$. To this may be added similar formulae, whose demonstration I omit in this place:
\[
\begin{gathered}
\text{9)}\quad\frac{1}{2}\,\frac{(1+z)(1+qz)(1+q^{2}z)\,.\,.}{(1+q)(1+q^{2})(1+q^{3})\,.\,.}+\frac{1}{2}\,\frac{(1-z)(1-qz)(1-q^{2}z)\,.\,.}{(1+q)(1+q^{2})(1+q^{3})\,.\,.}=\\[1.5ex]
1-\frac{q(1-z^{2})}{1-q^{2}}+\frac{q^{4}(1-z^{2})(1-q^{2}z^{2})}{(1-q^{2})(1-q^{4})}-\frac{q^{9}(1-z^{2})(1-q^{2}z^{2})(1-q^{4}z^{2})}{(1-q^{2})(1-q^{4})(1-q^{6})}+\ldots\\[2ex]
\text{10)}\quad\frac{q}{2z}.\,\frac{(1+z)(1+qz)(1+q^{2}z)\,.\,.}{(1+q)(1+q^{2})(1+q^{3})\,.\,.}-\frac{q}{2z}.\,\frac{(1-z)(1-qz)(1-q^{2}z)\,.\,.}{(1+q)(1+q^{2})(1+q^{3})\,.\,.}=\\[1.5ex]
q-\frac{q^{4}(1-z^{2})}{1-q^{2}}+\frac{q^{9}(1-z^{2})(1-q^{2}z^{2})}{(1-q^{2})(1-q^{4})}-\frac{q^{16}(1-z^{2})(1-q^{2}z^{2})(1-q^{4}z^{2})}{(1-q^{2})(1-q^{4})(1-q^{6})}+.\,.
\end{gathered}
\]
of which 9), with $z=q$ put, furnishes:
\[
\frac{1}{2}+\frac{1}{2}.\,\frac{(1-q)(1-q^{2})(1-q^{3})\,.\,.}{(1+q)(1+q^{2})(1+q^{3})\,.\,.}=1-q+q^{4}-q^{9}+.\,.,
\]
or:
\[
\frac{(1-q)(1-q^{2})(1-q^{3})(1-q^{4})\,.\,.}{(1+q)(1+q^{2})(1+q^{3})(1+q^{4})\,.\,.}=1-2q+2q^{4}-2q^{9}+\ldots,
\]
which is formula 3).

Formula 6), which is of the most profound investigation, as one that depends on the trisection of the Elliptic functions, was found and splendidly demonstrated long ago by Cl.\ \emph{Euler}. Of this remarkable demonstration we shall have to treat more fully elsewhere.

\origpage{187}

To these let us add the following expansions:
\[
\begin{gathered}
\text{11)}\quad\frac{\sqrt{kk'\left(\dfrac{2K}{\pi}\right)^{3}}}{\Theta\left(\dfrac{2Kx}{\pi}\right)}=\frac{2\sqrt[4]{q}\left\{(1-q^{2})(1-q^{4})(1-q^{6})(1-q^{8})\ldots\right\}^{2}}{(1-2q\cos 2x+q^{2})(1-2q^{3}\cos 2x+q^{6})(1-2q^{5}\cos 2x+q^{10})\ldots}=\\[2ex]
\frac{2\sqrt[4]{q}(1-q^{2})}{1-2q\cos 2x+q^{3}}-\frac{2\sqrt[4]{q^{9}}(1-q^{6})}{1-2q^{3}\cos 2x+q^{6}}+\frac{2\sqrt[4]{q^{25}}(1-q^{10})}{1-2q^{5}\cos 2x+q^{10}}-.\,.
\end{gathered}
\]\ednote{Printed $q^{3}$ in the denominator of the first term instead of $q^{2}$; an apparent misprint, as the pattern of the following terms and the product above require $1-2q\cos 2x+q^{2}$.}
\[
\begin{gathered}
\text{12)}\quad\frac{\sqrt{kk'\left(\dfrac{2K}{\pi}\right)^{3}}}{H\left(\dfrac{2Kx}{\pi}\right)}=\frac{\left\{(1-q^{2})(1-q^{4})(1-q^{6})(1-q^{8})\,.\,.\right\}^{2}}{\sin x(1-2q^{2}\cos 2x+q^{4})(1-2q^{4}\cos 2x+q^{8})(1-2q^{6}\cos 2x+q^{12})\ldots}=\\[2ex]
\frac{1}{\sin x}-\frac{4q^{2}(1+q^{2})\sin x}{1-2q^{2}\cos 2x+q^{4}}+\frac{4q^{6}(1+q^{4})\sin x}{1-2q^{4}\cos 2x+q^{8}}-\frac{4q^{12}(1+q^{6})\sin x}{1-2q^{6}\cos 2x+q^{12}}+.\,.\\[2ex]
=\frac{1}{\sin x}\left\{\frac{(1-q^{2})(1-q^{4})}{1-2q^{2}\cos 2x+q^{4}}-\frac{q^{2}(1-q^{4})(1-q^{8})}{1-2q^{4}\cos 2x+q^{8}}+\frac{q^{6}(1-q^{6})(1-q^{12})}{1-2q^{6}\cos 2x+q^{12}}-.\,.\right\},
\end{gathered}
\]
which are easily obtained from the well-known theory of the resolution of compound fractions into simple ones.

Hence special expansions are deduced:
\[
\begin{gathered}
\text{13)}\quad\frac{2kK}{\pi}=4\sqrt{q}\left(\frac{1+q}{1-q}\right)-4\sqrt{q^{9}}\left(\frac{1+q^{3}}{1-q^{3}}\right)+4\sqrt{q^{25}}\left(\frac{1+q^{5}}{1-q^{5}}\right)-\ldots\\[2ex]
\text{14)}\quad\frac{2k'K}{\pi}=1-\frac{4q}{1+q}+\frac{4q^{3}}{1+q^{2}}-\frac{4q^{6}}{1+q^{3}}+\frac{4q^{10}}{1+q^{4}}-\ldots
\end{gathered}
\]
When these are compared with the expansions of the expressions $\dfrac{2kK}{\pi}$, $\dfrac{2k'K}{\pi}$ exhibited above, there results:
\[
\begin{gathered}
\frac{\sqrt{q}}{1-q}-\frac{\sqrt{q^{3}}}{1-q^{3}}+\frac{\sqrt{q^{5}}}{1-q^{5}}-\frac{\sqrt{q^{7}}}{1-q^{7}}+.\,.=\\[1.5ex]
\sqrt{q}\left(\frac{1+q}{1-q}\right)-\sqrt{q^{9}}\left(\frac{1+q^{3}}{1-q^{3}}\right)+\sqrt{q^{25}}\left(\frac{1+q^{5}}{1-q^{5}}\right)-\ldots\\[2ex]
1-\frac{4q}{1+q}+\frac{4q^{3}}{1+q^{3}}-\frac{4q^{5}}{1+q^{5}}+\frac{4q^{7}}{1+q^{7}}-\ldots=\\[1.5ex]
1-\frac{4q}{1+q}+\frac{4q^{3}}{1+q^{2}}-\frac{4q^{6}}{1+q^{3}}+\frac{4q^{10}}{1+q^{4}}-\ldots
\end{gathered}
\]
In a similar manner Cl.\ \emph{Clausen} has recently observed ${}^{*)}$, that the series
\[
\frac{q}{1-q}+\frac{q^{2}}{1-q^{2}}+\frac{q^{3}}{1-q^{3}}+\frac{q^{4}}{1-q^{4}}+\ldots
\]

\textbf{*)} \emph{Crelle} Journal etc.\ Vol.\ III.\ p.\ 95.

\origpage{188}
can be transformed into this:
\[
q\left(\frac{1+q}{1-q}\right)+q^{4}\left(\frac{1+q^{2}}{1-q^{2}}\right)+q^{9}\left(\frac{1+q^{3}}{1-q^{3}}\right)+q^{16}\left(\frac{1+q^{4}}{1-q^{4}}\right)+\ldots
\]

We found above the expansions into series of $\dfrac{2K}{\pi}$, $\dfrac{2kK}{\pi}$ and of their second, third, fourth powers. These expansions therefore supply the second, fourth, sixth, eighth power of the expressions
\[
\begin{gathered}
\sqrt{\frac{2K}{\pi}}=1+2q+2q^{4}+2q^{9}+2q^{16}+.\,.\\[2ex]
\sqrt{\frac{2kK}{\pi}}=2\sqrt[4]{q}+2\sqrt[4]{q^{9}}+2\sqrt[4]{q^{25}}+2\sqrt[4]{q^{49}}+.\,.
\end{gathered}
\]
whence various Arithmetical theorems flow. Thus, e.g., from the formula:
\[
\begin{gathered}
\left(\frac{2K}{\pi}\right)^{2}=\left\{1+2q+2q^{4}+2q^{9}+2q^{16}+.\,.\right\}^{4}=\\[1.5ex]
1+8\left\{\frac{q}{1-q}+\frac{q^{2}}{1+q^{2}}+\frac{q^{3}}{1-q^{3}}+\frac{q^{4}}{1+q^{4}}+.\,.\right\}=\\[1.5ex]
1+8\Sigma\phi(p)\left\{q^{p}+3q^{2p}+3q^{4p}+3q^{8p}+.\,.\right\},
\end{gathered}
\]
where $p$ is any odd number, $\phi(p)$ the sum of the factors of $p$, there flows as a Corollary the celebrated theorem of Fermat, that every number is the sum of four squares.

\emph{HALLE, PRINTED BY THE GEBAUERS.}

\origpage{189}

\section*{CORRIGENDA.}

I wish to beg the kind reader to correct before reading the following graver errors which have crept in:

P.\ 3.\ line 11.\ read $M$ for $U$, twice.

--- 7. --- 7. for $\dfrac{A}{B}.\dfrac{1-k^{2}}{(1+mx)^{2}}$ read $\dfrac{A}{C}.\dfrac{1-x^{2}}{(1+mx)^{2}}$.

--- 8. for $k^{2}$ read $x^{2}$.

--- 8. --- 6. for $\sqrt{k.x}$ read $\sqrt{k},x$, twice.

--- 9. $U$ and $V$ must be interchanged everywhere.

--- 10. --- 2. 4. 6. for $-\sqrt{k}$ read $+\sqrt{k}$.

--- 17. --- 8. for $a''y$, $b''y$ read $a''y^{2}$, $b''y^{2}$.

--- 22. --- 12. for $\sqrt{\dfrac{x^{2m+1}}{\lambda}}$ read $\sqrt{\dfrac{k^{2m+1}}{\lambda}}$.

--- 13. for $\dfrac{b^{(m-1)}}{k^{m-4}}$ read $\dfrac{b^{(m-1)}}{k^{m-2}}$.

--- 23. --- 17. for $\dfrac{u(2v+v^{3})}{v^{4}}$ read $\dfrac{u(2v+u^{3})}{v^{4}}$.

--- 25. --- 17. for $(1+2\alpha)^{2}$ read $(1+2\alpha)^{3}$.

--- 29. --- 5. for $+v^{4}$ read $-v^{4}$.

--- 7. for: \emph{v loco u}, read: \emph{u loco v}.

--- 18. for $(1-u^{2}v^{2})$ read $(1-4u^{2}v^{2})$.

last line, for: $\dfrac{u(u+v^{5})y+.\,.}{u^{2}(1+u^{3}v)+.\,.}$ read $\dfrac{u(u+v^{5})y-.\,.}{u^{2}(1+u^{3}v)-.\,.}$.

--- 31. --- 2. for $\dfrac{dx}{\sqrt{1-k^{2}\sin\phi^{2}}}$ read $\dfrac{d\phi}{\sqrt{1-k^{2}\sin\phi^{2}}}$.

--- 35. --- 10. 11. for $k$ read $k'$.

last line, for \emph{profecti} read \emph{perfecti}.

--- 39. --- 16. for $M$ read $(-1)^{\frac{n-1}{2}}M$.

--- 47. --- 10. for $\left\{\ldots\right\}^{2}$ read $\left\{\ldots\right\}^{4}$.

--- 51. --- 8. delete $\sqrt{\dfrac{\lambda'k^{n}}{\lambda k'^{n}}}$, add:
\[
=\sqrt{\frac{\lambda'k^{n}}{\lambda k'^{n}}}.\cos\operatorname{am}(u)\cos\operatorname{am}\left(u+\frac{4K}{n}\right)\cos\operatorname{am}\left(u+\frac{8K}{n}\right).\,.\cos\operatorname{am}\left(u+\frac{4(n-1)K}{n}\right).
\]

--- 9. delete $\sqrt{\dfrac{\lambda'}{k'^{n}}}$, add:
\[
=\sqrt{\frac{\lambda'}{k'^{n}}}.\Delta\operatorname{am}(u)\Delta\operatorname{am}\left(u+\frac{4K}{n}\right)\Delta\operatorname{am}\left(u+\frac{8K}{n}\right).\,.\Delta\operatorname{am}\left(u+\frac{4(n-1)K}{n}\right).
\]

\origpage{190}

p.\ 60. line 19. for $\sqrt{1-\lambda^{2}z^{2}}$ read $\sqrt{1-\lambda^{2}y^{2}}$.

--- 63. --- 6. for $\sin^{2}\operatorname{am}u$ read $yy$.

--- 12. for \emph{earum} read \emph{earumque}.

--- 64. --- 21. for $\dfrac{2m'iK'}{n}$ read $\dfrac{2m'iK'}{nM}$.

--- 67. --- 17. for $-192\lambda^{4}$ read $+192\lambda^{4}$.

--- 76. --- 4. for $\dfrac{1}{n}.\dfrac{\lambda'^{2}.\,.}{k^{2}.\,.}$ read $\dfrac{1}{n}.\dfrac{\lambda'^{2}.\,.}{k'^{2}.\,.}$.

--- 79. --- 10. for $\left(\dfrac{1+\lambda^{2}}{\lambda-\lambda^{3}}\right)$ read $\left(\dfrac{1+\lambda^{2}}{\lambda-\lambda^{3}}\right)^{2}$.

--- 80. --- 9. for $\dfrac{3k'}{k^{5}}$ read $-\dfrac{3k'}{k^{5}}$.

--- 12. for $k^{4}k'^{2}$ read $4k^{4}k'^{2}$, for $k'^{2}(1-k')^{2}$ read $4k'^{2}(1-k')^{2}$.

--- 14. for $k^{5}(1-k')^{2}$ read $k^{5}(1+k')^{2}$.

--- 88. --- 11. for $-q^{5}$ read $-2q^{5}$.

--- 94. --- 11. for $k^{(6)\prime}$ read $k^{(8)\prime}$.

--- 98. last line but one, for \emph{opposuisse} read \emph{apposuisse}.

--- 104. --- 4. for $-\dfrac{8q}{(1+q^{2})}$ read $-\dfrac{8q}{(1+q)^{2}}$.

--- 105. --- 1. for $-\dfrac{24}{10}q^{10}$ read $-\dfrac{24}{5}q^{10}$.

--- 107. --- 8. for \emph{formula} read \emph{unde formula}.

--- 18. for \emph{evolutas, --- --- nanciscimur} read \emph{evolutae --- nanciscuntur.}

--- 108. --- 1. for $\left(\dfrac{2K}{\pi}\right)^{5}$ read $\left(\dfrac{2K}{\pi}\right)^{3}$.

--- 111. last line, for \emph{quem} read \emph{quam}.

--- 114. --- 12. for $-\dfrac{2K}{\pi}.\,\dfrac{2E^{\mathrm{I}}}{\pi}$ read $-3\,.\dfrac{2K}{\pi}.\,\dfrac{2E^{\mathrm{I}}}{\pi}$.

--- 117. --- 8. for 3329 read 3229.

--- 15. for $+32k^{2}$ read $+23k^{2}$.

--- 20. for $\Pi(2n-2)$ read $\Pi(2n-1)$.

last line, for $\dfrac{d^{2}\sin\operatorname{am}u}{du^{2}}$ read $\dfrac{d^{2}\sin^{n}\operatorname{am}u}{du^{2}}$.

--- 118. --- 8. for $-2$ read $-2k^{2}$.

--- 10. for $-32(1+k^{2})$ read $-32k^{2}(1+k^{2})$.

--- 120. --- 17. for $\left(\dfrac{2K}{\pi}.\,\dfrac{2K}{\pi}-\dfrac{2K}{\pi}.\,\dfrac{2E^{\mathrm{I}}}{\pi}\right)$ read $\left(\dfrac{2K}{\pi}-\dfrac{2E^{\mathrm{I}}}{\pi}\right)$.

--- 124. --- 7. for $\sin\operatorname{am}(u-v)-$ read $\sin\operatorname{am}(u+v)-$.

last line but two, \emph{Tayloriana} read \emph{Tayloriano}.

--- 125. --- 10. for $-2$ read $-$.

--- 128. last line but one, for $2.4.6.7$ read $2.4.6.8$.

--- 131. --- 1, 2, 3 for $(-1)^{n}2^{n-1}$, $(-1)^{n}4^{n-1}$, $(-1)^{n}6^{n-1}$ read $(-1)^{n-1}2^{2n-1}$, $(-1)^{n-1}4^{2n-1}$, $(-1)^{n-1}6^{2n-1}$.

--- 7, 8, 9 for $\left(\dfrac{2K}{\pi}\right)^{2n-1}$ read $\left(\dfrac{2K}{\pi}\right)^{2n-2}$.

\origpage{191}

p.\ 132. line 14. for $\operatorname{tang}\operatorname{am}u$ read $i\operatorname{tang}\operatorname{am}u$.

--- 144. --- 7. for $1-k^{2}\sin^{2}\operatorname{am}u.\sin^{2}\phi$ read $1-k^{2}\sin^{2}\operatorname{am}a.\sin^{2}\operatorname{am}u$.

--- 152. --- 7. for \emph{quibos, tranis} read \emph{quibus, transis.}

--- 154. --- 11. for $Z(u)-Z(u+a)$ read $Z(u)+Z(a)-Z(u+a)$.

--- 158. --- 10. in the denominator for $1-k^{2}.\,.$ read $1+k^{2}.\,.$

--- 159. --- 14. for $-\Pi(u,u+b)$ read $-\Pi(u,a+b)$.

--- 160. --- 15: delete 2), 3).

--- 164. --- 10. for $\left\{(1-q)(1-q^{3})(1-q^{5})\,.\,.\right\}$ read $\left\{(1-q)(1-q^{3})(1-q^{5})\,.\,.\right\}^{2}$.

--- 169. --- 10. for $iuZ$ read $iaZ$.

--- 171. --- 21. for \emph{classem} read \emph{in classem.}

--- 172. --- 1. for \emph{E} read \emph{Et.}

--- 6. for $\Pi(u+i,a+ib)$ read $\Pi(u+iv,a+ib)$.

--- 174. --- 2. for $\sin\operatorname{am}u$ read $\sin\operatorname{am}u.\Theta u$.

--- 176. last line but two, for \emph{Quam} read \emph{Quem.}

--- 178. --- 1. for \emph{varie} read \emph{serie.}

--- 180. --- 8. for \emph{Quam} read \emph{Quem.}

\section*{REAL TRANSFORMATIONS OF ELLIPTIC FUNCTIONS AND THEIR COMPLEMENTARY AND SUPPLEMENTARY ONES. (§. 25.)}
\ednote{The folded leaf is printed in four panels: A and B side by side above, and the two complementary panels below. They are transcribed here in the logical order A, complementary to A, B, complementary to B. The braces at the right margin are printed as $\{M.\,k\}$ and so on.}

\subsection*{A. FIRST TRANSFORMATION WITH THE SUPPLEMENTARY.}

\[
\begin{gathered}
\text{a)}\quad \lambda=k^{n}\sin^{4}\operatorname{coam}\dfrac{2K}{n}\sin^{4}\operatorname{coam}\dfrac{4K}{n}.\,.\,.\,.\sin^{4}\operatorname{coam}\left(\dfrac{n-1}{n}\right)K\qquad\{M.\,k\}\\[1.5ex]
\text{aa)}\quad k=\lambda^{n}\sin^{4}\operatorname{coam}\dfrac{2i\Lambda'}{n}\sin^{4}\operatorname{coam}\dfrac{4i\Lambda'}{n}.\,.\,.\,.\sin^{4}\operatorname{coam}\left(\dfrac{n-1}{n}\right)i\Lambda'\qquad\{M.\,\lambda\}\\[1.5ex]
{}=\dfrac{\lambda^{n}}{\Delta^{4}\operatorname{am}\dfrac{2\Lambda'}{n}\,\Delta^{4}\operatorname{am}\dfrac{4\Lambda'}{n}.\,.\,.\,.\Delta^{4}\operatorname{am}\left(\dfrac{n-1}{n}\right)\Lambda'}\qquad\{M.\,\lambda'\}\\[1.5ex]
\text{b)}\quad M=\dfrac{\sin^{2}\operatorname{coam}\dfrac{2K}{n}\sin^{2}\operatorname{coam}\dfrac{4K}{n}.\,.\,.\,.\sin^{2}\operatorname{coam}\left(\dfrac{n-1}{n}\right)K}{\sin^{2}\operatorname{am}\dfrac{2K}{n}\sin^{2}\operatorname{am}\dfrac{4K}{n}.\,.\,.\,.\sin^{2}\operatorname{am}\left(\dfrac{n-1}{n}\right)K}\qquad\{M.\,k\}\\[1.5ex]
\text{bb)}\quad \dfrac{1}{nM}=\dfrac{\sin^{2}\operatorname{coam}\dfrac{2\Lambda'}{n}\sin^{2}\operatorname{coam}\dfrac{4\Lambda'}{n}.\,.\,.\,.\sin^{2}\operatorname{coam}\left(\dfrac{n-1}{n}\right)\Lambda'}{\sin^{2}\operatorname{am}\dfrac{2\Lambda'}{n}\sin^{2}\operatorname{am}\dfrac{4\Lambda'}{n}.\,.\,.\,.\sin^{2}\operatorname{am}\left(\dfrac{n-1}{n}\right)\Lambda'}\qquad\{M.\,\lambda'\}
\end{gathered}
\]
\ednote{In aa) the exponent of the last sin is badly inked and looks like a prime; it is read as the exponent 4 of the other factors.}

\[
\sin\operatorname{am}(u,k)=x;\quad \sin\operatorname{am}\left(\dfrac{u}{M},\lambda\right)=y;\quad \sin\operatorname{am}(nu,k)=z
\]

\[
\begin{gathered}
\text{c)}\quad y=(-1)^{\frac{n-1}{2}}\sqrt{\dfrac{k^{n}}{\lambda}}\sin\operatorname{am}u\sin\operatorname{am}\left(u+\dfrac{4K}{n}\right)\sin\operatorname{am}\left(u+\dfrac{8K}{n}\right).\,.\,.\,.\sin\operatorname{am}\left(u+4\left(\dfrac{n-1}{n}\right)K\right)\qquad\{M.\,k\}\\[1.5ex]
{}=\dfrac{\dfrac{x}{M}\left(1-\dfrac{xx}{\sin^{2}\operatorname{am}\dfrac{2K}{n}}\right)\left(1-\dfrac{xx}{\sin^{2}\operatorname{am}\dfrac{4K}{n}}\right).\,.\,.\,.\left(1-\dfrac{xx}{\sin^{2}\operatorname{am}\left(\dfrac{n-1}{n}\right)K}\right)}{\left(1-k^{2}\sin^{2}\operatorname{am}\dfrac{2K}{n}xx\right)\left(1-k^{2}\sin^{2}\operatorname{am}\dfrac{4K}{n}xx\right).\,.\,.\,.\left(1-k^{2}\sin^{2}\operatorname{am}\left(\dfrac{n-1}{n}\right)Kxx\right)}\qquad\{M.\,k\}\\[1.5ex]
\text{cc)}\quad z=\sqrt{\dfrac{\lambda^{n}}{k}}\sin\operatorname{am}\dfrac{u}{M}\sin\operatorname{am}\left(\dfrac{u}{M}+\dfrac{4i\Lambda'}{n}\right)\sin\operatorname{am}\left(\dfrac{u}{M}+\dfrac{8i\Lambda'}{n}\right)\ldots\sin\operatorname{am}\left(\dfrac{u}{M}+4\left(\dfrac{n-1}{n}\right)i\Lambda'\right)\qquad\{M.\,\lambda\}\\[1.5ex]
{}=\dfrac{nMy\left(1+\dfrac{yy}{\operatorname{tg}^{2}\operatorname{am}\dfrac{2\Lambda'}{n}}\right)\left(1+\dfrac{yy}{\operatorname{tg}^{2}\operatorname{am}\dfrac{4\Lambda'}{n}}\right).\,.\,.\,.\left(1+\dfrac{yy}{\operatorname{tg}^{2}\operatorname{am}\left(\dfrac{n-1}{n}\right)\Lambda'}\right)}{\left(1+\lambda^{2}\operatorname{tg}^{2}\operatorname{am}\dfrac{2\Lambda'}{n}yy\right)\left(1+\lambda^{2}\operatorname{tg}^{2}\operatorname{am}\dfrac{4\Lambda'}{n}yy\right).\,.\,.\,.\left(1+\lambda^{2}\operatorname{tg}^{2}\operatorname{am}\left(\dfrac{n-1}{n}\right)\Lambda'yy\right)}\qquad\{M.\,\lambda'\}
\end{gathered}
\]

\subsection*{COMPLEMENTARY TRANSFORMATIONS.}

\[
\begin{gathered}
\text{a)}\quad \lambda'=k'^{n}\sin^{4}\operatorname{coam}\dfrac{2iK}{n}\sin^{4}\operatorname{coam}\dfrac{4iK}{n}.\,.\,.\,.\sin^{4}\operatorname{coam}\left(\dfrac{n-1}{n}\right)iK\qquad\{M.\,k'\}\\[1.5ex]
{}=\dfrac{k'^{n}}{\Delta^{4}\operatorname{am}\dfrac{2K}{n}\,\Delta^{4}\operatorname{am}\dfrac{4K}{n}.\,.\,.\,.\Delta^{4}\operatorname{am}\left(\dfrac{n-1}{n}\right)K}\qquad\{M.\,k\}\\[1.5ex]
\text{aa)}\quad k'=\lambda'^{n}\sin^{4}\operatorname{coam}\dfrac{2\Lambda'}{n}\sin^{4}\operatorname{coam}\dfrac{4\Lambda'}{n}.\,.\,.\,.\sin^{4}\operatorname{coam}\left(\dfrac{n-1}{n}\right)\Lambda'\qquad\{M.\,\lambda'\}
\end{gathered}
\]

\[\text{b) and bb) the same as above.}\]

\[
\sin\operatorname{am}(u,k')=x;\quad \sin\operatorname{am}\left(\dfrac{u}{M},\lambda'\right)=y;\quad \sin\operatorname{am}(nu,k')=z
\]

\[
\begin{gathered}
\text{c)}\quad y=\sqrt{\dfrac{k'^{n}}{\lambda'}}\sin\operatorname{am}u\sin\operatorname{am}(u+4iK)\sin\operatorname{am}(u+8iK).\,.\,.\,.\sin\operatorname{am}\left(u+\dfrac{4(n-1)}{n}iK\right)\qquad\{M.\,k'\}\\[1.5ex]
{}=\dfrac{\dfrac{x}{M}\left(1+\dfrac{xx}{\operatorname{tg}^{2}\operatorname{am}\dfrac{2K}{n}}\right)\left(1+\dfrac{xx}{\operatorname{tg}^{2}\operatorname{am}\dfrac{4K}{n}}\right).\,.\,.\,.\left(1+\dfrac{xx}{\operatorname{tg}^{2}\operatorname{am}\left(\dfrac{n-1}{n}\right)K}\right)}{\left(1+k'^{2}\operatorname{tg}^{2}\operatorname{am}\dfrac{2K}{n}xx\right)\left(1+k'^{2}\operatorname{tg}^{2}\operatorname{am}\dfrac{4K}{n}xx\right).\,.\,.\,.\left(1+k'^{2}\operatorname{tg}^{2}\operatorname{am}\left(\dfrac{n-1}{n}\right)Kxx\right)}\qquad\{M.\,k\}\\[1.5ex]
\text{cc)}\quad z=(-1)^{\frac{n-1}{2}}\sqrt{\dfrac{\lambda}{k'^{n}}}\sin\operatorname{am}\dfrac{u}{M}\sin\operatorname{am}\left(\dfrac{u}{M}+\dfrac{4\Lambda'}{n}\right)\sin\operatorname{am}\left(\dfrac{u}{M}+\dfrac{8\Lambda'}{n}\right)\ldots\sin\operatorname{am}\left(\dfrac{u}{M}+\dfrac{4(n-1)}{n}\Lambda'\right)\qquad\{M.\,\lambda'\}\\[1.5ex]
{}=\dfrac{nMy\left(1-\dfrac{yy}{\sin^{2}\operatorname{am}\dfrac{2\Lambda'}{n}}\right)\left(1-\dfrac{yy}{\sin^{2}\operatorname{am}\dfrac{4\Lambda'}{n}}\right).\,.\,.\,.\left(1-\dfrac{yy}{\sin^{2}\operatorname{am}\left(\dfrac{n-1}{n}\right)\Lambda'}\right)}{\left(1-\lambda'^{2}\sin^{2}\operatorname{am}\dfrac{2\Lambda'}{n}yy\right)\left(1-\lambda'^{2}\sin^{2}\operatorname{am}\dfrac{4\Lambda'}{n}yy\right).\,.\,.\,.\left(1-\lambda'^{2}\sin^{2}\operatorname{am}\left(\dfrac{n-1}{n}\right)\Lambda'yy\right)}\qquad\{M.\,\lambda'\}
\end{gathered}
\]
\ednote{In the first line of c) the factor $(-1)^{\frac{n-1}{2}}$ is absent in the print, whereas it is printed in the first line of cc) and in the right-hand panel.}

\[
\Lambda=\dfrac{K}{nM};\qquad \Lambda'=\dfrac{K'}{M}
\]

\subsection*{B. SECOND TRANSFORMATION WITH THE SUPPLEMENTARY.}

\[
\begin{gathered}
\text{a)}\quad \lambda=k^{n}\sin^{4}\operatorname{coam}\dfrac{2iK'}{n}\sin^{4}\operatorname{coam}\dfrac{4iK'}{n}.\,.\,.\,.\sin^{4}\operatorname{coam}\left(\dfrac{n-1}{n}\right)iK'\qquad\{M.\,k\}\\[1.5ex]
{}=\dfrac{k^{n}}{\Delta^{4}\operatorname{am}\dfrac{2K'}{n}\,\Delta^{4}\operatorname{am}\dfrac{4K'}{n}.\,.\,.\,.\Delta^{4}\operatorname{am}\left(\dfrac{n-1}{n}\right)K'}\qquad\{M.\,k\}\\[1.5ex]
\text{aa)}\quad k={\lambda_{,}}^{n}\sin^{4}\operatorname{coam}\dfrac{2\Lambda_{,}}{n}\sin^{4}\operatorname{coam}\dfrac{4\Lambda_{,}}{n}.\,.\,.\,.\sin^{4}\operatorname{coam}\left(\dfrac{n-1}{n}\right)\Lambda_{,}\qquad\{M.\,\lambda_{,}\}\\[1.5ex]
\text{b)}\quad M_{,}=\dfrac{\sin^{2}\operatorname{coam}\dfrac{2K'}{n}\sin^{2}\operatorname{coam}\dfrac{4K'}{n}.\,.\,.\,.\sin^{2}\operatorname{coam}\left(\dfrac{n-1}{n}\right)K'}{\sin^{2}\operatorname{am}\dfrac{2K'}{n}\sin^{2}\operatorname{am}\dfrac{4K'}{n}.\,.\,.\,.\sin^{2}\operatorname{am}\left(\dfrac{n-1}{n}\right)K'}\qquad\{M.\,k'\}\\[1.5ex]
\text{bb)}\quad \dfrac{1}{nM_{,}}=\dfrac{\sin^{2}\operatorname{coam}\dfrac{2\Lambda_{,}}{n}\sin^{2}\operatorname{coam}\dfrac{4\Lambda_{,}}{n}.\,.\,.\,.\sin^{2}\operatorname{coam}\left(\dfrac{n-1}{n}\right)\Lambda_{,}}{\sin^{2}\operatorname{am}\dfrac{2\Lambda_{,}}{n}\sin^{2}\operatorname{am}\dfrac{4\Lambda_{,}}{n}.\,.\,.\,.\sin^{2}\operatorname{am}\left(\dfrac{n-1}{n}\right)\Lambda_{,}}\qquad\{M.\,\lambda'\}
\end{gathered}
\]

\[
\sin\operatorname{am}(u,k)=x;\quad \sin\operatorname{am}\left(\dfrac{u}{M_{,}},\lambda_{,}\right)=y;\quad \sin\operatorname{am}(nu,k)=z
\]

\[
\begin{gathered}
\text{c)}\quad y=\sqrt{\dfrac{k^{n}}{\lambda_{,}}}\sin\operatorname{am}u\sin\operatorname{am}\left(u+\dfrac{4iK'}{n}\right)\sin\operatorname{am}\left(u+\dfrac{8iK'}{n}\right)\ldots\sin\operatorname{am}\left(u+\dfrac{4(n-1)}{n}iK'\right)\qquad\{M.\,k\}\\[1.5ex]
{}=\dfrac{\dfrac{x}{M_{,}}\left(1+\dfrac{xx}{\operatorname{tg}^{2}\operatorname{am}\dfrac{2K'}{n}}\right)\left(1+\dfrac{xx}{\operatorname{tg}^{2}\operatorname{am}\dfrac{4K'}{n}}\right).\,.\,.\,.\left(1+\dfrac{xx}{\operatorname{tg}^{2}\operatorname{am}\left(\dfrac{n-1}{n}\right)K'}\right)}{\left(1+k^{2}\operatorname{tg}^{2}\operatorname{am}\dfrac{2K'}{n}xx\right)\left(1+k^{2}\operatorname{tg}^{2}\operatorname{am}\dfrac{4K'}{n}xx\right).\,.\,.\,.\left(1+k^{2}\operatorname{tg}^{2}\operatorname{am}\left(\dfrac{n-1}{n}\right)K'xx\right)}\qquad\{M.\,k'\}\\[1.5ex]
\text{cc)}\quad z=(-1)^{\frac{n-1}{2}}\sqrt{\dfrac{\lambda_{,}^{n}}{k}}\sin\operatorname{am}\dfrac{u}{M_{,}}\sin\operatorname{am}\left(\dfrac{u}{M_{,}}+\dfrac{4\Lambda_{,}}{n}\right)\sin\left(\dfrac{u}{M_{,}}+\dfrac{8\Lambda_{,}}{n}\right)\ldots\sin\operatorname{am}\left(\dfrac{u}{M_{,}}+\dfrac{4(n-1)}{n}\Lambda_{,}\right)\qquad\{M.\,\lambda_{,}\}\\[1.5ex]
{}=\dfrac{\dfrac{y}{M_{,}}\left(1-\dfrac{yy}{\sin^{2}\operatorname{am}\dfrac{2\Lambda_{,}}{n}}\right)\left(1-\dfrac{yy}{\sin^{2}\operatorname{am}\dfrac{4\Lambda_{,}}{n}}\right).\,.\,.\,.\left(1-\dfrac{yy}{\sin^{2}\operatorname{am}\left(\dfrac{n-1}{n}\right)\Lambda_{,}}\right)}{\left(1-\lambda_{,}^{2}\sin^{2}\operatorname{am}\dfrac{2\Lambda_{,}}{n}yy\right)\left(1-\lambda'^{2}\sin^{2}\operatorname{am}\dfrac{4\Lambda_{,}}{n}yy\right).\,.\,.\,.\left(1-\lambda_{,}^{2}\sin^{2}\operatorname{am}\left(\dfrac{n-1}{n}\right)\Lambda_{,}yy\right)}\qquad\{M.\,\lambda_{,}'\}
\end{gathered}
\]
\ednote{In cc) the third factor of the first line is printed $\sin\left(\dots\right)$ without am; the numerator of the second line is printed $\dfrac{y}{M_{,}}$ where the complementary panel has $nM_{,}y$; and the second factor of the denominator is printed with $\lambda'^{2}$ without the sub-comma. All are kept as printed; apparent misprints.}

\subsection*{COMPLEMENTARY TRANSFORMATIONS.}

\[
\begin{gathered}
\text{a)}\quad \lambda_{,}'=k'^{n}\sin^{4}\operatorname{coam}\dfrac{2K'}{n}\sin^{4}\operatorname{coam}\dfrac{4K'}{n}.\,.\,.\,.\sin^{4}\operatorname{coam}\left(\dfrac{n-1}{n}\right)K'\qquad\{M.\,k'\}\\[1.5ex]
\text{aa)}\quad k'={\lambda_{,}'}^{n}\sin^{4}\operatorname{coam}\dfrac{2i\Lambda_{,}}{n}\sin^{4}\operatorname{coam}\dfrac{4i\Lambda_{,}}{n}.\,.\,.\,.\sin^{4}\operatorname{coam}\left(\dfrac{n-1}{n}\right)i\Lambda_{,}\qquad\{M.\,\lambda_{,}'\}\\[1.5ex]
{}=\dfrac{{\lambda_{,}'}^{n}}{\Delta^{4}\operatorname{am}\dfrac{2\Lambda_{,}}{n}\,\Delta^{4}\operatorname{am}\dfrac{4\Lambda_{,}}{n}.\,.\,.\,.\Delta^{4}\operatorname{am}\left(\dfrac{n-1}{n}\right)\Lambda_{,}}\qquad\{M.\,\lambda_{,}\}
\end{gathered}
\]

\[\text{b) and bb) the same as above.}\]

\[
\sin\operatorname{am}(u,k')=x;\quad \sin\operatorname{am}\left(\dfrac{u}{M_{,}},\lambda_{,}'\right)=y;\quad \sin\operatorname{am}(nu,k')=z
\]

\[
\begin{gathered}
\text{c)}\quad y=(-1)^{\frac{n-1}{2}}\sqrt{\dfrac{k'^{n}}{\lambda_{,}'}}\sin\operatorname{am}u\sin\operatorname{am}\left(u+\dfrac{4K'}{n}\right)\sin\operatorname{am}\left(u+\dfrac{8K'}{n}\right).\,.\,.\,.\sin\operatorname{am}\left(u+\dfrac{4(n-1)}{n}K'\right)\qquad\{M.\,k'\}\\[1.5ex]
{}=\dfrac{\dfrac{x}{M_{,}}\left(1-\dfrac{xx}{\sin^{2}\operatorname{am}\dfrac{2K'}{n}}\right)\left(1-\dfrac{xx}{\sin^{2}\operatorname{am}\dfrac{4K'}{n}}\right).\,.\,.\,.\left(1-\dfrac{xx}{\sin^{2}\operatorname{am}\left(\dfrac{n-1}{n}\right)K'}\right)}{\left(1-k'^{2}\sin^{2}\operatorname{am}\dfrac{2K'}{n}xx\right)\left(1-k'^{2}\sin^{2}\operatorname{am}\dfrac{4K'}{n}xx\right).\,.\,.\,.\left(1-k'^{2}\sin^{2}\operatorname{am}\left(\dfrac{n-1}{n}\right)K'xx\right)}\qquad\{M.\,k'\}\\[1.5ex]
\text{cc)}\quad z=\sqrt{\dfrac{{\lambda_{,}'}^{n}}{k'}}\sin\operatorname{am}\left(\dfrac{u}{M_{,}}\right)\sin\operatorname{am}\left(\dfrac{u}{M_{,}}+\dfrac{4i\Lambda_{,}}{n}\right)\sin\operatorname{am}\left(\dfrac{u}{M_{,}}+\dfrac{8i\Lambda_{,}}{n}\right).\,.\,.\,.\sin\operatorname{am}\left(\dfrac{u}{M_{,}}+\dfrac{4(n-1)}{n}i\Lambda_{,}\right)\qquad\{M.\,\lambda_{,}'\}\\[1.5ex]
{}=\dfrac{nM_{,}y\left(1+\dfrac{yy}{\operatorname{tg}^{2}\operatorname{am}\dfrac{2\Lambda_{,}}{n}}\right)\left(1+\dfrac{yy}{\operatorname{tg}^{2}\operatorname{am}\dfrac{4\Lambda_{,}}{n}}\right).\,.\,.\,.\left(1+\dfrac{yy}{\operatorname{tg}^{2}\operatorname{am}\left(\dfrac{n-1}{u}\right)\Lambda_{,}}\right)}{\left(1+{\lambda_{,}'}^{2}\operatorname{tg}^{2}\operatorname{am}\dfrac{2\Lambda_{,}}{n}yy\right)\left(1+{\lambda_{,}'}^{2}\operatorname{tg}^{2}\operatorname{am}\dfrac{4\Lambda_{,}}{n}yy\right).\,.\,.\,.\left(1+{\lambda_{,}'}^{2}\operatorname{tg}^{2}\operatorname{am}\left(\dfrac{n-1}{n}\right)\Lambda_{,}yy\right)}\qquad\{M.\,\lambda'\}
\end{gathered}
\]
\ednote{In the numerator of the last line of cc) the fraction is printed $\dfrac{n-1}{u}$ instead of $\dfrac{n-1}{n}$; an apparent misprint. The marker of that line is printed without the sub-comma.}

\[
\Lambda_{,}=\dfrac{K}{M_{,}};\qquad \Lambda_{,}'=\dfrac{K'}{nM_{,}}.
\]


\end{document}
